% =====================================================================
%  ToScaNA -- User manual
%  Compile:  latexmk -pdf ToScaNA_manual.tex      (or pdflatex three times)
%  The appendices (audit, tests, module reference) are in
%  ToScaNA_appendices.tex, included at the end.
% =====================================================================
\documentclass[11pt,a4paper]{report}
\usepackage[utf8]{inputenc}
\usepackage[T1]{fontenc}
\usepackage{lmodern}
\usepackage{textcomp}
\usepackage[margin=2.3cm]{geometry}
\usepackage{amsmath,amssymb}
\usepackage{booktabs,longtable,array,tabularx}
\usepackage{enumitem}
\usepackage[dvipsnames]{xcolor}
\usepackage{listings}
% keep spaces inside \code{...} and allow line breaks there
\PassOptionsToPackage{obeyspaces,spaces}{url}
\usepackage[hidelinks]{hyperref}
\urlstyle{tt}

% Shared macros (also used by the appendices)
\DeclareRobustCommand{\code}[1]{\path{#1}}
\makeatletter\g@addto@macro\UrlBreaks{\do\_\do\(\do\,\do\=\do\[}\makeatother
\newcommand{\st}[1]{\textsf{\footnotesize\colorbox{black!8}{#1}}}
\newcommand{\fixed}[1]{\textsf{\footnotesize\colorbox{green!15}{fixed #1}}}
\newcommand{\newbeh}{\textsf{\footnotesize\colorbox{orange!20}{results change}}}
\newcommand{\legacy}{\textsf{\footnotesize\colorbox{black!8}{legacy}}}
\newcommand{\warn}[1]{\par\smallskip\noindent\colorbox{orange!12}{\parbox{\dimexpr\linewidth-2\fboxsep}{\textbf{Caution.} #1}}\par\smallskip}
\newcommand{\tip}[1]{\par\smallskip\noindent\colorbox{green!8}{\parbox{\dimexpr\linewidth-2\fboxsep}{\textbf{Tip.} #1}}\par\smallskip}
\lstset{basicstyle=\ttfamily\small,columns=fullflexible,keepspaces=true,
        frame=single,framerule=0.3pt,backgroundcolor=\color{black!3},
        breaklines=true,showstringspaces=false,aboveskip=6pt,belowskip=6pt,
        upquote=true}
\lstdefinelanguage{toscanapar}{morecomment=[l]{\#},morecomment=[l]{!}}
\setlist{itemsep=2pt,topsep=3pt}
\setlength{\emergencystretch}{3em}

\def\ToScaNAManual{}   % tells ToScaNA_appendices.tex it is part of this manual

\title{\textbf{ToScaNA}\\[4pt]\Large Total Scattering Neutron Analysis\\[2pt]
\large User manual}
\author{Version 0.1.0 (branch \code{fixes/group-a})}
\date{November 2026}

\begin{document}
\maketitle
\tableofcontents

% =====================================================================
\chapter{Introduction}
% =====================================================================

\section{What ToScaNA does}

ToScaNA is a Python package for the reduction and analysis of neutron
total-scattering data measured on the diffractometer D4 at the Institut
Laue-Langevin. It covers the path from the raw detector counts to the
real-space correlation functions:

\begin{enumerate}
  \item \textbf{Regrouping} (\code{d4creg}): reading the raw numors,
    correcting for dead time and detector efficiency, normalising to the
    monitor, computing the angle of every detector cell, combining all
    numors and cells on a regular $2\theta$ grid, converting to momentum
    transfer $Q$, and writing the diffractograms in the formats of D4 and of
    the program CORRECT (Chapter~\ref{ch:d4creg}).
  \item \textbf{Calibration} of the incident wavelength $\lambda$ and of the
    zero angle $2\theta_0$ with a nickel powder (Chapter~\ref{ch:nickel}).
  \item \textbf{Description of the experiment}: sample composition,
    geometry and density, container, environment and vanadium standard,
    with the nuclear data needed for the corrections
    (Chapter~\ref{ch:sample}).
  \item \textbf{Models} used when normalising the data: vanadium $Q$
    dependence, inelasticity, self scattering, cross sections
    (Chapter~\ref{ch:models}).
  \item \textbf{Fourier transformation} of the structure factor $S(Q)$ into
    $G(r)$, $g(r)$, the radial distribution function, $T(r)$ and the running
    coordination number $N(r)$, with optional window function and exact
    statistical error bars (Chapter~\ref{ch:realspace}).
\end{enumerate}

\noindent Since November 2026 ToScaNA also contains:
\begin{itemize}
  \item \textbf{pycorrect}, a Python translation of CORRECT 2.24
    (Paalman \& Pings attenuation, Blech \& Averbach multiple scattering,
    Placzek, vanadium normalisation), which runs the usual \code{.com} files;
  \item \textbf{MCAMS}, a Monte Carlo code for attenuation and multiple
    scattering in cylinders, annuli, slabs, discs and spheres, used to check
    and complement the analytical corrections
    (Chapter~\ref{ch:corrections});
  \item \textbf{d4numor}, a survey tool for sets of numors (variables,
    averages, detector movement, preset calculator), and the NeXus reader
    that ToScaNA uses for both the old and the current D4 formats
    (Section~\ref{sec:d4numor}).
\end{itemize}
These three have their own manuals in \code{docs/pycorrect},
\code{docs/mcams} and \code{docs/d4numor}.

\section{How this manual is organised}

Chapters~\ref{ch:install}--\ref{ch:maintenance} explain how to install and
use the package, with worked examples and the formats of all input and
output files. Chapter~\ref{ch:troubleshooting} lists common error messages.
The appendices contain the technical details:
\begin{itemize}
  \item Appendix~\ref{app:audit}: the audit of the code made in November
    2026, every correction with its physical justification, and the open
    items;
  \item Appendix~\ref{app:tests}: the automatic test suite;
  \item Appendix~\ref{app:reference}: a reference of every module, function
    and class.
\end{itemize}
The same information as in Appendix~\ref{app:reference}, in more detail,
is available inside Python: every module, function and class has a
docstring, displayed by \code{help()}:
\begin{lstlisting}[language=Python]
import toscana.numerics.fourier as fourier
help(fourier.sineFT)        # or, in Jupyter:  fourier.sineFT?
\end{lstlisting}

\section{Conventions}

Angles are in degrees, $Q$ in \AA$^{-1}$, $r$ and $\lambda$ in \AA, sample
dimensions in mm, macroscopic densities in g/cm$^3$, atomic number densities
in atoms/\AA$^3$, scattering lengths in fm and cross sections in barns. The
definitions of the real-space functions are given in
Section~\ref{sec:conventions}. Throughout this manual, labels such as \st{C2}
refer to the findings of the audit (Appendix~\ref{app:audit}), and
\newbeh{} marks behaviour that differs from the original version of the
code.

% =====================================================================
\chapter{Installation}\label{ch:install}
% =====================================================================

\section{Requirements}

\begin{longtable}{@{}l l p{0.55\textwidth}@{}}
\toprule
Package & Needed for & Notes \\
\midrule\endhead
Python $\ge$ 3.10 & everything & tested with 3.11.15 \\
NumPy & everything & tested with 2.4.3 \\
SciPy & everything & integration, peak finding, special functions; tested with 1.17.1 \\
matplotlib & plots & imported only by plotting functions \\
lmfit & Ni calibration & tested with 1.3.4 \\
h5py & NeXus raw data & imported only by \code{load_nxs} and \code{d4numor}; tested with 3.15.1 \\
pandas & \code{d4numor} tables & optional, only \code{NumorSet.to_dataframe} \\
periodictable & --- & no longer needed: MCAMS uses ToScaNA's isotope table \\
\bottomrule
\end{longtable}

\section{Making the package importable}

The code imports itself under the name \code{toscana} (for example
\code{from toscana.numerics.fourier import sineFT}), but the project
directory is called \code{ToScaNA-pilot}. There is no installer yet
(\code{pyproject.toml} is an open item), so Python must be told where to
find a directory named \code{toscana}. The simplest way is a symbolic link
in a folder of your choice, added to the search path:
\begin{lstlisting}
mkdir -p ~/python_packages
ln -s /path/to/ToScaNA ~/python_packages/toscana   # your clone
export PYTHONPATH=~/python_packages:$PYTHONPATH     # e.g. in ~/.zshrc
\end{lstlisting}
In a Jupyter notebook, instead of \code{PYTHONPATH}:
\begin{lstlisting}[language=Python]
import os, sys
sys.path.insert(0, os.path.expanduser("~/python_packages"))
import toscana
print(toscana.__version__)          # 0.1.0
\end{lstlisting}

\warn{Do not add the project directory itself (\code{ToScaNA-pilot}) to
\code{sys.path}; add its \emph{parent} folder (or the folder containing the
link). The package must be reached as \code{toscana.xxx}; otherwise its
subpackages (\code{data}, \code{models}, \code{utils}, \ldots) could be
confused with other modules of the same name.}

\section{Checking the installation}

Run the test suite (Appendix~\ref{app:tests}); it needs no installation
and takes about three minutes:
\begin{lstlisting}
python /path/to/ToScaNA/tests/run_tests.py
...
Ran 136 tests in 180 s
OK
\end{lstlisting}

\section{Old import names}

Two subpackages were renamed because their names clashed with modules of
the Python standard library (\st{C21}): \code{toscana.math} is now
\code{toscana.numerics} and \code{toscana.io} is now \code{toscana.dataio}.
The isotope class \code{Elemento} is now \code{Element}. The old names still
work but print a \code{DeprecationWarning}; update notebooks when convenient:
\begin{lstlisting}[language=Python]
from toscana.math.fourier import sineFT        # old: works, with a warning
from toscana.numerics.fourier import sineFT    # new
\end{lstlisting}

\section{Working directory of an experiment}

Several functions read and write files at paths relative to the
\emph{current working directory}. The expected layout is:
\begin{lstlisting}
experiment/
  logfiles/            d4creg.log  (created automatically: ../logfiles from work/)
  regdata/             figures saved by DataXYE.plot  (../regdata from work/)
  work/                <- start Python / the notebook here
    rawdata/           raw numors (path set by <rdp>)
    effd4c.eff         detector efficiencies
    dec.dec            bank angular shifts
    vanadium.par       d4creg parameter file(s)
    *.adat *.reg *.qdat  outputs of d4creg
\end{lstlisting}

% =====================================================================
\chapter{Overview of an analysis}\label{ch:overview}
% =====================================================================

A typical D4 experiment is analysed in this order.
\begin{enumerate}
  \item \textbf{Regroup the Ni powder} measurement with \code{d4creg}, with
    the zero angle set to 0 and the nominal wavelength.
  \item \textbf{Calibrate}: fit the Ni diffractogram to obtain $\lambda$ and
    $2\theta_0$ (Chapter~\ref{ch:nickel}).
  \item \textbf{Regroup every measurement} (sample, empty container,
    environment, vanadium, absorber) with \code{d4creg}, using the
    calibrated $\lambda$ and $2\theta_0$ (Chapter~\ref{ch:d4creg}).
  \item \textbf{Describe the experiment} in a sample-description file and
    build a \code{Measurement} (Chapter~\ref{ch:sample}).
  \item \textbf{Correct} with pycorrect, optionally with MCAMS factors:
    background and container subtraction, attenuation and multiple
    scattering, normalisation to vanadium, Placzek; then
    \textbf{obtain $S(Q)$} (\code{run_correction}, \code{structure_factor},
    Section~\ref{sec:workflow}).
  \item \textbf{Transform to real space} and compute coordination numbers,
    with error bars (Chapter~\ref{ch:realspace}).
\end{enumerate}
The data flow and the function responsible for each step are shown in
Section~\ref{sec:chain}.

% =====================================================================
\chapter{Regrouping raw data: d4creg}\label{ch:d4creg}
% =====================================================================

\section{The parameter file}

The regrouping is controlled by a text file read by
\code{dataio.parameters.readParam}. Example (\code{vanadium.par}):
\begin{lstlisting}[language=toscanapar]
# d4creg parameters for the vanadium rod
<rdp> rawdata/            # directory of the raw numors
<fmt> Nexus               # Nexus (old or current format), anything else: ILL ASCII
<eff> effd4c.eff          # efficiencies ("ones" = no correction)
<dec> dec.dec             # bank shifts ("zeros", "manual")
<tau> 2.4e-6 7.0e-6       # dead times: monitor, detector (s)
<tol> 0.1 1.0             # angle tolerance (cells): warning, error
<nor> monitor 80 1000000  # normalise to 10^6 monitor counts
<wle> 0.49857             # wavelength (A), from the Ni calibration
<zac> -0.1013             # zero angle (deg), from the Ni calibration
<asc> 0.0 140.0 0.125     # angular grid: min max step (deg)
<qsc> 0.0 23.8 0.02       # Q grid: min max step (1/A)
<wri> False               # one file per numor?
<plo> False               # plot the diffractograms?
<out> vanadium            # base name of the output files
<num> 387229-387288       # numors (single or first-last)
<add> 387301              # more numors
<run>                     # do it
\end{lstlisting}
Rules:
\begin{itemize}
  \item Every non-empty line starts with \code{\#} or \code{!} (comment), or
    with a five-character tag such as \code{<wle>}. A line that starts with
    a space is an error.
  \item Values are separated by spaces; anything after \code{\#} on a tag
    line is ignored.
  \item The file is read from top to bottom; \code{<run>} starts the
    regrouping with the values read so far. Several \code{<out>}/\code{<num>}/\code{<run>}
    groups in one file perform several regroupings (e.g.\ sample, container
    and vanadium in one go).
  \item Tags not given keep their defaults. The complete list of tags and
    defaults is in Appendix~\ref{app:reference}, \code{readParam}.
\end{itemize}

\noindent To run it, from the \code{work/} directory:
\begin{lstlisting}[language=Python]
from toscana.dataio.parameters import readParam
runInfo = readParam("vanadium.par")
\end{lstlisting}
\code{readParam} first re-creates \code{../logfiles/d4creg.log}, so the log of
a previous run is overwritten unless it was renamed (see below).

\section{Auxiliary input files}

\paragraph{Efficiencies} (\code{<eff>}): one line per cell,
\code{bank cell efficiency}, both indices starting at 1; lines containing
\code{\#} are ignored. A value $\le 0$ marks a dead cell, which is then
excluded. Cells not listed keep efficiency 1.
\begin{lstlisting}
# bank cell efficiency
1 1 0.9834
1 2 1.0121
...
9 64 0.9977
\end{lstlisting}

\paragraph{Bank shifts} (\code{<dec>}): one line per bank, \code{bank shift}
in degrees, added to the angle of every cell of that bank.

\section{What happens during the regrouping}

For each numor (\code{getOneNumor}):
\begin{enumerate}
  \item The raw counts $N$ of the 9 banks $\times$ 64 cells are read, with
    the monitor $M$, the counting time $t$ and the angle of the first cell.
  \item The read and requested angles are compared; a difference larger than
    \code{<tol>} cell widths gives a warning or stops the program.
  \item Poisson errors $\sigma = \sqrt{N}$ are computed on the raw counts.
  \item Dead time (non-paralysable): $N \to N/(1-N\tau/t)$, cell by cell for
    the detector and once for the monitor; errors scaled likewise.
  \item Efficiency: $N\to N/\varepsilon$, $\sigma\to\sigma/\varepsilon$.
  \item Normalisation: $y = N\,M_0/M$ with $M_0$ = \code{<nor>} monitor counts,
    $\sigma_y^2 = (\sigma M_0/M)^2 + (y/\sqrt{M})^2$.
  \item Angles: $2\theta = 2\theta_\mathrm{read} + \text{shift}_\text{bank} + 15^\circ\,\text{bank} + 0.125^\circ\,\text{cell}$,
    then the zero angle $2\theta_0$ is subtracted.
\end{enumerate}
All numors are then combined (\code{getDiffA}): each cell is a rectangle of
one cell width ($8^\circ/64 = 0.125^\circ$), split between the bins of the
angular grid it overlaps in proportion to the overlap, and weighted by
the monitor share of its numor.
Each bin receives the weighted mean of its contributions; its error is the
error of a weighted mean, $\sqrt{\sum w^2\sigma^2}/\sum w$. This is done for
the sum of the banks and for each bank separately. Finally the regrouped
diffractogram is converted to $Q$ and rebinned onto the $Q$ grid
(\code{rebin}), where each angular point is a rectangle of width
$\Delta Q = \frac{2\pi}{\lambda}\cos\theta\,\Delta(2\theta)$.

\tip{The angular step cannot be smaller than the cell width,
$0.125^\circ$; a smaller value is replaced by $0.125^\circ$ with a warning.
For the $Q$ step, \code{readParam} prints a recommended value,
$\text{step}_{2\theta}\,\Delta Q_\text{range}/\Delta(2\theta)_\text{range}$.}

\section{Output files}

\begin{longtable}{@{}>{\raggedright\arraybackslash}p{0.22\textwidth} p{0.72\textwidth}@{}}
\toprule
File & Content \\
\midrule\endhead
\code{<out>.adat} & CORRECT format, angle scale: a header (\code{\#Block}, instrument, user, numors, title, date, binning, number of points) followed by three columns $2\theta$, intensity, error, for the sum of the banks. Only points with positive, finite intensity are written. \\
\code{<out>.reg} & D4 format: a header with the equivalent \code{d4creg} command line, the efficiency and shift files, sample, user, normalisation, zero angle, wavelength and dead times; then one block per bank with four columns $2\theta$, intensity, error, $Q$. \\
\code{<out>.qdat} & CORRECT format, $Q$ scale: three columns $Q$, intensity, error. \\
\code{<numor>.cdat} & One per numor if \code{<wri> True} (or a single numor): angle, intensity, error, $Q$, cell identifier $1000\times\text{bank} + \text{cell}$. Note: its $Q$ column ignores the zero angle (open item, Section~\ref{sec:open}). \\
\code{../logfiles/d4creg.log} & Count rates, dead-time factors, warnings; renamed at the end to \code{../logfiles/<log>.log}. With the default \code{<log>} the name becomes \code{d4creg.log.log}. \\
\bottomrule
\end{longtable}

\noindent Intensities are counts per \code{<nor>} monitor counts (or per
\code{<nor>} seconds). Errors are one standard deviation, propagated from
counting statistics (they were corrected in this version, \newbeh{}
\st{C2, C5, C29}).

\section{Messages and their meaning}

\begin{longtable}{@{}p{0.42\textwidth} p{0.52\textwidth}@{}}
\toprule
Message & Meaning / action \\
\midrule\endhead
\code{WARNING! The angle differs from the requested angle by more than 0.1 cell} & The detector was not exactly at the requested position; the measured angle is used. Usually harmless. \\
\code{ERROR! ... more than 1.0 cell} (\code{ValueError}) & Probably a motor problem: check the numor, or relax \code{<tol>}. \\
\code{getDTC_dcounts: correction factor < 0} & Count rate times dead time $\ge 1$: the numor cannot be corrected (wrong dead time, or a saturated detector). \\
\code{WARNING: In getDiffA function ... angular step} & The step was smaller than a cell and was enlarged. \\
\code{Invalid numor entry} / \code{Invalid numor range} & \code{<num>} must be 6 digits, or \code{first-last} with first $\le$ last. \\
\code{Wrong input in line} & A line does not start with \code{\#}, \code{!} or \code{<}. \\
\bottomrule
\end{longtable}

% =====================================================================
\chapter{Calibration with nickel powder}\label{ch:nickel}
% =====================================================================

\section{Principle}

The peak positions of a cubic powder depend only on $\lambda/a$ and on the
zero angle: $2\theta_{hkl} = 2\arcsin\bigl(\lambda\sqrt{h^2+k^2+l^2}/2a\bigr) + 2\theta_0$.
Fitting the ten first reflections of fcc nickel, with a known lattice
parameter $a$, gives $\lambda$ and $2\theta_0$. Each reflection is modelled by
an asymmetric Gaussian, on a linear background
(\code{models.ni.niPeaks10}). Because only $\lambda/a$ is measured, $a$ is
\emph{fixed} in the fit, and the fitted $\lambda$ is proportional to it.
The default is $a = 3.52$\,\AA{} (\code{NI_LATTICE_PARAMETER}); use the value
of your powder's certificate if you have it.

\section{Procedure}

First regroup the Ni measurement with \code{d4creg}, with \code{<zac> 0.0} and
the nominal wavelength. Then, in a notebook:
\begin{lstlisting}[language=Python]
from toscana.calibration.d4nifit import (
    setting_model_d4nifit, fitting_model_d4nifit, showing_results_d4nifit)

# 1. read the data, estimate lambda from the (111) peak, build the initial model
axes, limits, xy_res, xy_ref, xy_fit, planes, initial, nickel = \
    setting_model_d4nifit("Ni5.adat", zac_ini=0.0, lattice=3.52)

# (optional) restrict the fitted range or change starting values
limits[0] = 10.0          # lowest angle used in the fit (deg)

# 2. fit (weighted by the errors of the data)
result = fitting_model_d4nifit(nickel, limits, initial, lattice=3.52)

# 3. table of results, plot with residuals, saved as Ni5.png
showing_results_d4nifit(result, nickel, axes, limits,
                        xy_res, xy_ref, xy_fit, planes)
lam = result.params["wavelength"].value
zac = result.params["twotheta0"].value
\end{lstlisting}
Copy $\lambda$ and $2\theta_0$ into \code{<wle>} and \code{<zac>} of the
parameter files of the other measurements.

\paragraph{What the first function does.} It reads the file with
\code{DataXYE}, which finds the peaks between $5^\circ$ and $40^\circ$; it
assumes the first peak is (111) and estimates
$\lambda = 2(a/\sqrt3)\sin\theta_{111}$; it divides the intensities (and
errors) by the first data point above the midpoint between the second and
third peaks (truncated to whole degrees); it
builds the 35 starting parameters (background, $\lambda$, $2\theta_0$, ten
areas proportional to typical Ni intensities, ten asymmetries 0, ten widths)
and plots the starting model. The order of \code{initial} is:
\code{I0, slope, quad, wavelength, twotheta0, A0..A9, G0..G9, S0..S9}.

\paragraph{What is fitted.} Free: background \code{I0} and \code{slope},
\code{wavelength}, \code{twotheta0}, areas \code{A0..A9}, widths
\code{S0..S9}. Fixed: \code{quad}, asymmetries \code{G0..G9}, \code{lattice}.
The fit is weighted by $1/\sigma$ \newbeh{}; the uncertainties reported by
lmfit are scaled by $\sqrt{\chi^2_\nu}$, so they include the effect of model
imperfections. They are statistical uncertainties for a given lattice
parameter: an error $\delta a/a$ adds $\delta\lambda/\lambda = \delta a/a$.

\warn{The automatic peak search needs points between $5^\circ$ and
$40^\circ$ whose height is more than 3.5 times the minimum of that range.
If the (111) peak is not the first one found (e.g.\ because of an impurity),
the initial wavelength is wrong; set \code{initial[3]} by hand before
fitting.}

% =====================================================================
\chapter{Describing the experiment and the sample}\label{ch:sample}
% =====================================================================

\section{The sample-description file}

The data-treatment notebooks read a second kind of parameter file with
\code{dataio.running_params.getRunningParams}, and build a
\code{experiment.Measurement} from it. A complete example, which was
checked to work:
\begin{lstlisting}[language=toscanapar]
# ToScaNA sample-description file (read by getRunningParams)
<ins> D4 # instrument
<pro> 6-01-1234 # proposal
<mpr> A.User # main proposer
<lco> G.Cuello # local contact
<dte> 01/03/2026-05/03/2026 # start-end dates
<cyc> 261 # reactor cycle
<log> 12-34 # notebook-page
<env> Cryofurnace # sample environment
<zac> -0.1013 # zero angle (deg)
<wle> 0.4986 # wavelength (A)
<sli> 3.0 -3.0 # left right slits (mm)
<fla> 25.0 -25.0 # top bottom flags (mm)
<apr> 0.125 # angular precision (deg)
<asc> 1.0 136.0 0.125 # angular grid (deg)
<qsc> 0.3 23.5 0.02 # Q grid (1/A)
<rsc> 0.0 20.0 0.01 # r grid (A)
<Cmater> V # container material
<Cshape> cylinder 5.0 5.2 60.0 # type id od height (mm)
<Nmater> V # normaliser
<Nshape> cylinder 0.0 6.08 50.0 # vanadium rod
<Amater> B # absorber
<Ashape> cylinder 0.0 5.0 60.0 # absorber
<Emater> V # environment
<Eshape> cylinder 40.0 40.2 100.0 # furnace shield
<Sdescr> Silica glass at room temperature # description
<StempK> 295.0 # temperature (K)
<Sato01> Si 1 # element number-of-atoms
<Sato02> O 2 # element number-of-atoms
<Sdegcc> 2.20 # density (g/cm3)
<Smassc> 1.80 # mass in the beam (g)
<Sshape> cylinder 0.0 5.0 60.0 # type id od height (mm)
<Sfulln> 1.0 # fullness
<Stitle> SiO2_RT # title
<CoFile> SiO2_can.qdat # container data
<AbFile> absorber.qdat # absorber data
<NoFile> vanadium.qdat # vanadium data
<EnFile> furnace.qdat # environment data
<SaFile> SiO2_RT.qdat # sample data
\end{lstlisting}

\noindent The syntax is the same as for the d4creg file: values are
separated by spaces (any number), and anything after \code{\#} is a
comment, which is optional. (Before commit \code{ece015d} a value at the
end of a line kept the newline character, so every line needed a comment;
files written that way are read exactly as before.)

\noindent Shapes are given as \code{type inner\_diameter outer\_diameter height}
in mm; only \code{cylinder} is known. Elements are given as
\code{<SatoNN> symbol number} with consecutive \code{NN} = 01, 02, \ldots; use
the isotope names of the nuclear-data table for isotopic substitution
(e.g.\ \code{H2} for deuterium, \code{Li7}). All tags above are required by
\code{Measurement}, which also reads the five data files (3 columns:
$x$, $y$, error).

\paragraph{Extra layers.} \newbeh{} Set-ups with more cylinders than sample,
container and environment (inner wall of an annular can, a second heat
shield, \ldots) are described with optional lines
\begin{lstlisting}[language=toscanapar]
<Lay01> shield environment Nb 30.0 30.1 100.0 # name role material id od height (mm)
<Lay02> wall2 container V 7.0 7.2 60.0 # second container wall
\end{lstlisting}
The role is \code{sample}, \code{container} or \code{environment}: the
layer belongs to that body in the corrections. The material is a name of
\code{corrections.setup.MATERIALS} (V, TiZr, SiO2, Nb, Al, Cu) or
\code{formula@density}, e.g.\ \code{Al2O3@3.95} (g/cm$^3$). They are
available as \code{m.extra_layers}; a line with fewer than six values raises
\code{ValueError}. A non-zero inner diameter in \code{<Sshape>} gives a
hollow sample; the inner diameter of \code{<Cshape>} may be larger than the
sample (gap).

\section{The Measurement object}

\begin{lstlisting}[language=Python]
from toscana.dataio.running_params import getRunningParams
from toscana.experiment.measurement import Measurement

m = Measurement(getRunningParams("SiO2_RT.par"))
m.showTable()                       # summary of the main parameters
print(m.conc)                       # {'Si': 0.333..., 'O': 0.666...}
print(m.AtomicDensity)              # 0.06615 atoms/A^3  (from 2.20 g/cm3)
print(m.Volume, m.PackingFraction)  # 1.178 cm^3, 0.6945
print(m.ssca, m.free, m.sabswl)     # 3.544, 3.170, 0.0158 barn/atom
\end{lstlisting}
The values in the comments are those of the example file. The most useful
attributes are:
\begin{longtable}{@{}>{\raggedright\arraybackslash}p{0.34\textwidth} p{0.60\textwidth}@{}}
\toprule
Attribute & Meaning \\
\midrule\endhead
\code{conc}, \code{atoms}, \code{natoms} & concentrations $c_i$, numbers of atoms in the formula unit, total \\
\code{molarMass} & mean molar mass per atom (g/mol) \\
\code{Density}, \code{AtomicDensity} & bulk density (g/cm$^3$) and atomic density $\rho$ (atoms/\AA$^3$) --- use $\rho$ in the Fourier transform \\
\code{Volume}, \code{EffDensity}, \code{PackingFraction}, \code{EffAtomicDensity} & volume $\pi h(r_o^2-r_i^2)$ (cm$^3$), mass/volume, ratio to the bulk density, effective atomic density (powders) \\
\code{scoh}, \code{sinc}, \code{ssca}, \code{sabs} & mean cross sections per atom (barn); \code{sabs} at 2200\,m/s \\
\code{sabswl} & absorption at the wavelength, $\sigma_\mathrm{abs}\lambda/1.798$ \\
\code{free} & free-atom scattering cross section $\sum c_i\sigma_i[A_i/(A_i+1)]^2$ \newbeh{} \\
\code{b2lowQ}, \code{b2highQ} & self scattering per atom and steradian at low and high $Q$: $\sigma_\mathrm{sca}/4\pi$ and $\sigma_\mathrm{free}/4\pi$ \\
\code{Data}, \code{conData}, \code{norData}, \code{envData}, \code{absData} & the five data sets, arrays $(n, 3)$ \\
\bottomrule
\end{longtable}

\section{Dictionary-based description}

An alternative, used in older notebooks, builds dictionaries with the
functions of \code{experiment.info}: \code{setExpInfo}, \code{setBeamInfo},
\code{setCanInfo}, \code{setBinInfo}, \code{setNumorInfo}, \code{setVanaInfo}
and \code{setSampleInfo}. Each prints a summary and returns a dictionary
(keys listed in their docstrings). \code{setSampleInfo} computes in addition
the sample/vanadium volume and density ratios used to put the data on an
absolute scale; the vanadium defaults are 6.11\,g/cm$^3$,
$b_\mathrm{coh}=-0.443$\,fm, $\sigma_\mathrm{inc}=5.08$\,b.

\section{Nuclear data}

\begin{lstlisting}[language=Python]
from toscana.isotopes import Element
ni = Element("Ni")
print(ni.re_bcoh, ni.sig_coh, ni.sig_inc, ni.sig_abs, ni.weight)
#     10.3 fm     13.3 b      5.2 b        4.49 b       58.6934 u
Element("Ni", table=1)          # print all values with units
Element.getIsotope("Nickel")    # ['Ni58', 'Ni60', 'Ni61', 'Ni62', 'Ni64', 'Ni']
\end{lstlisting}
Values come from \emph{Neutron Scattering Fundamentals}, vol.~44,
Appendix~2. Isotopes are named by symbol and mass number (\code{"Li6"},
\code{"H2"} for deuterium); natural elements by their symbol. Any value can
be overridden for a special case, e.g.\ \code{Element("V", re_bcoh=-0.40)}.
An unknown name raises \code{ValueError}.

\section{Input for CORRECT}

With the corrections inside ToScaNA no command file is needed:
\code{run_correction} builds everything from the \code{Measurement}
(Section~\ref{sec:workflow}). For
the Fortran CORRECT, two functions write command files:
\begin{itemize}
  \item \code{generate_com(template, name, rsample, density, ...)} fills the
    placeholders \texttt{\{pyfname\}}, \texttt{\{pyradius\_cm\}},
    \texttt{\{pydensity\_at\_AA3\}}, \ldots{} of a template file and writes
    \code{<name>.com}. (It then runs a placeholder command; see open items.)
  \item \code{saveCORRECT(expt)} writes \code{<sampleTitle>.com} from an
    object that must provide the attributes \code{instrument},
    \code{sampleTitle}, \code{sampleOuterRadius} (mm), \code{sampleTemperature},
    \code{sampleDensity}, \code{samplePackingFraction}, \code{sampleFullness},
    \code{containerFile}, \code{container} (index 3: outer diameter, mm),
    \code{environmentFile}, \code{vanadiumFile}, \code{vanadium} (index 3:
    outer diameter, mm), \code{wavelength}, \code{zeroAngle},
    \code{beamHeight}, \code{beamWidth} (mm). \code{Measurement} does not
    provide all of them (it uses \code{Title}, \code{OuterRadius},
    \code{TempK}, \ldots), so such an object must be built by hand; check
    the units expected by CORRECT for the density.
\end{itemize}

% =====================================================================
\chapter{Corrections: MCAMS and pycorrect}\label{ch:corrections}
% =====================================================================

\section{The complete workflow}\label{sec:workflow}

\code{toscana.workflows.correction} connects the regrouped data to S(Q):
the experiment is described once, in the sample-description file
(Chapter~\ref{ch:sample}), and no CORRECT command file has to be written.
Two methods are available:
\begin{itemize}
  \item \code{method="mcams"} (default) \newbeh{}: ToScaNA's own Paalman \&
    Pings subtraction with factors computed by MCAMS. Any number of
    concentric cylinders (hollow sample, annular container, container with a
    gap, several environment layers), partial filling, an absorber
    measurement, angle-dependent multiple scattering of every configuration.
  \item \code{method="pycorrect"}: CORRECT~2.24 itself, for the set-ups it can
    describe (solid sample in contact with one container wall, one furnace
    layer, solid vanadium); \code{ValueError} otherwise.
\end{itemize}
Data files may be \code{.adat} (angle) or \code{.qdat} ($Q$, converted to
$2\theta$ with the wavelength) in both methods; the scale is read from the
\code{X caption} header line or, without it, from the extension
(\code{read_regrouped}).
\begin{lstlisting}[language=Python]
import numpy as np
from toscana.dataio.running_params import getRunningParams
from toscana.experiment.measurement import Measurement
from toscana.workflows.correction import (run_correction, structure_factor,
                                          save_structure_factor)
from toscana.numerics.fourier import sineFT
from toscana.dataio.saving import saveRSCF

m = Measurement(getRunningParams("SiO2_RT.par"))   # .adat or .qdat of d4creg
res = run_correction(m)                  # MCAMS factors, ToScaNA's chain
# res = run_correction(m, method="pycorrect")   # CORRECT 2.24
res.save_factors("SiO2_RT", scale="q")   # A_b,C, delta_C, T_C vs Q (or "angle")
res.save_intensities("SiO2_RT_I.dat", scale="q")  # I_s, I_c, I_h, I_v, B_C, MS_C
sq = structure_factor(res, m, q_grid=(0.3, 23.5, 0.02))   # Q, S, error
save_structure_factor("SiO2_RT.soq", sq, ["SiO2 at room temperature"])

r = np.round(np.arange(0.0, 20.0 + 1e-9, 0.01), 10)
fou = sineFT(sq, r, qmax=23.5, density=m.AtomicDensity, window=1, with_errors=True)
saveRSCF("SiO2_RT.rscf", fou, ["SiO2 at room temperature, Lorch, Qmax 23.5"])
\end{lstlisting}

\paragraph{The set-up.} \code{corrections.setup.setups_from_measurement}
turns the \code{Measurement} into two \code{CylinderSetup} objects, one for
the sample measurements and one for the vanadium. Each layer is a cylinder or
an annulus (inner and outer diameter, height, fullness, material) belonging
to a \emph{body}: \code{s} sample, \code{c} container, \code{h} environment,
\code{v} vanadium. The layers come from \code{<Sshape>}, \code{<Cshape>},
\code{<Eshape>} (only with \code{environment="layer"}), \code{<Nshape>} and the
\code{<LayNN>} lines. \textbf{All cross sections come from ToScaNA's isotope
table}; the sample uses its \emph{effective} atomic density (mass/volume, so
the packing fraction is included). The \emph{configurations} are the
measurements, from the outermost body inwards: e.g.\ \code{"h"} (environment
alone), \code{"ch"} (empty can in the environment), \code{"sch"} (sample),
and \code{"v"} or \code{"vh"} for the vanadium.

\paragraph{The factors.} \code{corrections.factors.compute_factors} computes,
for every configuration $C$ and every body $b$ in it, on the angles of the
data:
\begin{itemize}
  \item the Paalman \& Pings factor $A_{b,C}(2\theta)$ (MCAMS direct
    integration over the beam and the scattering points, then a cubic
    spline from a grid of step \code{step_deg}); for a body with several
    layers, the average weighted by the scattering power of the layers;
  \item the multiple-scattering ratio $\delta_C(2\theta)$: total multiple over
    total single scattering of the configuration, the bodies taken as
    isotropic scatterers (MCAMS Monte Carlo, fitted by a polynomial);
  \item the transmission $T_C$ (exact, analytic average of
    $\exp(-\sum\mu\,\ell)$ over the beam);
  \item $n\langle L\rangle_b$ (atoms per unit beam area) and
    $P_b=\sum\rho\,\sigma_s\langle L\rangle$ (scattering power).
\end{itemize}
They are returned as a \code{FactorSet}: \code{res.factors} (sample set-up)
and \code{res.vanadium_factors}. Keys: \code{"A_s,sch"}, \code{"A_c,sch"},
\code{"A_h,sch"}, \code{"A_c,ch"}, \code{"A_h,ch"}, \code{"A_h,h"},
\code{"A_v,v"}, \ldots{} in \code{.A}; \code{"delta_sch"}, \ldots{} in
\code{.delta}; \code{"T_sch"}, \ldots{} in \code{.T}; statistical errors in
\code{.A_err}, \code{.delta_err}. \code{.table("q")} or \code{.table("angle")}
gives the arrays on either scale, \code{.save(file, scale=...)} writes them
with a header (transmissions, $n\langle L\rangle$, $P_b$), and
\code{load_factors(file)} reads them back. With the defaults
(\code{n_points=100\_000}, \code{ms_particles=2\_000\_000}) a three-configuration
set-up takes about one minute; \code{factor_options=dict(...)} changes them.
Precomputed factors (from an earlier run, same angles) are given with
\code{factors=} and \code{vanadium_factors=}.

\paragraph{The subtraction.} With $M_C$ the measurement of configuration
$C$:
\begin{enumerate}
  \item Background. With an absorber measurement (\code{<Amater>} not
    \code{None}), $B_C = T_C\,I_\text{empty} + (1-T_C)\,I_\text{absorber}$;
    otherwise $B_C=I_\text{empty}$. $I_\text{empty}$ is the environment
    measurement with \code{environment="background"} (default), or
    \code{files["empty"]} with \code{environment="layer"}.
  \item Vanadium: $I_v = (M_v - B_v)/(1+\delta_v)/A_{v,v}$ (with an
    environment, its contribution is subtracted first as below), divided by
    $1-P_v$ (Placzek) and smoothed (\code{vanadium_smoothing}). It gives the
    count rate per unit isotropic cross section,
    $K(2\theta) = I_v / (P_v/4\pi)$.
  \item Multiple scattering of each configuration, as in CORRECT (subtracted,
    not divided, since $\delta$ refers to isotropic scatterers):
    $MS_C = \delta_C\,K\sum_{b\in C} A_{b,C}\,P_b/4\pi$.
  \item Nested Paalman \& Pings subtraction from the outermost body inwards:
    \begin{align*}
      I_h &= (M_h-B_h-MS_h)/A_{h,h},\\
      I_c &= (M_{ch}-B_{ch}-MS_{ch}-A_{h,ch}I_h)/A_{c,ch},\\
      I_s &= (M_{sch}-B_{sch}-MS_{sch}-A_{c,sch}I_c-A_{h,sch}I_h)/A_{s,sch}.
    \end{align*}
  \item Normalisation and Placzek correction of the sample:
    $d\sigma/d\Omega = I_s/(K\,n\langle L\rangle_s) + \sum_i c_i\sigma_i P_i/4\pi$.
\end{enumerate}
Errors of the measurements are propagated in quadrature through every step;
without vanadium smoothing (default) the error of the vanadium enters through
$K$, which scales the sample intensity \emph{before} the multiple-scattering
subtraction (checked by a Monte Carlo over noise realisations: the scatter of
the results equals the propagated errors).
\code{res.intensities} holds $I_s$, $I_c$, $I_h$, $I_v$ (with errors),
the smoothed vanadium and $MS_C$; \code{res.backgrounds} the $B_C$; both on the
data grid (\code{res.theta}, \code{res.q}), so that the subtraction can be
redone or modified in a notebook.

\paragraph{Which factors enter.} A factor $A_{b,C}$ is needed for every body
$b$ of every configuration $C$ that was \emph{measured}. With the environment
as a body (E $\equiv$ \code{h}) and the measurements E, CE, SCE, VE, ACE:
$A_{E,E}$; $A_{E,CE}$, $A_{C,CE}$; $A_{E,SCE}$, $A_{C,SCE}$, $A_{S,SCE}$;
$A_{E,VE}$, $A_{V,VE}$; $A_{E,ACE}$, $A_{C,ACE}$, $A_{A,ACE}$. Factors of
configurations that were not measured ($A_{S,S}$, or $A_{S,SC}$, $A_{C,SC}$,
$A_{C,C}$ when the furnace is in the beam) do not enter; the latter are those
of \code{environment="background"}, where the environment is not a body (an
approximation: it neglects the attenuation of the sample and container
scattering by the environment).

\paragraph{The absorber as a configuration.} \newbeh{} With
\code{environment="layer"} and an absorber measurement (e.g.\ $^{10}$B powder
in the container, in the environment; \code{<Amater> [10B]@1.25},
\code{<Ashape>}), the absorber is a body \code{a} and its measurement the
configuration \code{"ach"} (\code{corrections.setup.absorber_setup_from_measurement}).
Every measurement contains the room background $B$, not attenuated by the
set-up; the measurements without sample are then linear, at each angle, in
the unknowns $B$, $I^1_h$, $I^1_c$ and $K$:
\begin{align*}
I_E &= B + A_{h,h}I^1_h + MS_h,\\
I_{CE} &= B + A_{h,ch}I^1_h + A_{c,ch}I^1_c + MS_{ch},\\
I_{ACE} &= B + A_{h,ach}I^1_h + A_{c,ach}I^1_c + A_{a,ach}K P_a/4\pi + MS_{ach},\\
I_{VE} &= B + A_{h,vh}I^1_h + A_{v,vh}K(1-P^{Pl}_v)P_v/4\pi + MS_{vh},
\end{align*}
with $MS_C=\delta_C K\sum_b A_{b,C}P_b/4\pi$. The absorber fixes $B$, the
vanadium $K$. \code{run_correction} solves the system (weighted least squares
when there are more measurements, e.g.\ an empty-instrument run given as
\code{files["empty"]}), then
$I^1_s=(I_{SCE}-B-A_{h,sch}I^1_h-A_{c,sch}I^1_c-MS_{sch})/A_{s,sch}$ and
$d\sigma/d\Omega = I^1_s/(K n\langle L\rangle_s)+\sum_i c_i\sigma_iP_i/4\pi$.
Errors are propagated to first order from all the measurements, with their
correlations. Once $B$ and $K$ are known this is the nested subtraction
above. The result also holds \code{res.absorber_factors},
\code{res.intensities["K"]} and \code{res.backgrounds["B"]}. This is chosen
automatically (\code{absorber_model="auto"}); \code{"transmission"} uses
$B_C=T_C I_\text{empty}+(1-T_C)I_\text{absorber}$ instead.

\paragraph{Options of \code{run_correction}.}
\begin{itemize}
  \item \code{environment="background"} (default): the environment
    measurement is the background of every measurement, as in the CORRECT
    files written by ToScaNA so far; \code{"layer"} (or \code{"furnace"}): the
    environment and the environment \code{<LayNN>} layers are a body $h$,
    corrected for attenuation, and \code{files["empty"]} (empty instrument) is
    the background; \code{"none"}.
  \item \code{vanadium_in_environment=True}: the vanadium was measured inside
    the environment (configurations \code{"h"}, \code{"vh"}).
  \item \code{placzek="series"} (default), \code{"diatomic"} or \code{"none"}.
  \item \code{vanadium_smoothing}: number of polynomial coefficients (0: none;
    negative: CORRECT's Savitzky--Golay smoothing).
  \item \code{files=\{...\}}: data files instead of those of the
    sample-description file; keys \code{sample}, \code{container},
    \code{vanadium}, \code{environment}, \code{empty}, \code{absorber}. Files on
    another angle grid than the sample are interpolated onto it (noted in
    \code{res.notes}).
  \item \code{materials=\{...\}}: extra container or environment materials.
  \item \code{method="pycorrect"}, optionally with \code{mcams="delta"} (MCAMS
    multiple scattering) or \code{"all"} (all factors from MCAMS,
    \code{mcams_options} for the statistics): the CORRECT chain, whose factors
    are also returned as a \code{FactorSet} (keys \code{A_s,sc}, \ldots);
    \code{res.pycorrect} is CORRECT's full result dictionary.
\end{itemize}

\paragraph{Consistency of the two methods.} Given CORRECT's own factors (its
$\delta$ converted to the definition above), ToScaNA's chain reproduces
CORRECT's $d\sigma/d\Omega$ to $10^{-10}$; the differences between the methods
therefore come only from the factors. For the synthetic Ni example the MCAMS
and CORRECT attenuation factors agree within 0.1--0.4\,\%, while the
multiple scattering differs (CORRECT's series underestimates the higher
orders and ignores the container), which changes $d\sigma/d\Omega$ by a few
per cent.

\paragraph{From $d\sigma/d\Omega$ to S(Q).} \code{structure_factor} uses the
Faber--Ziman definition,
\[
  \frac{d\sigma}{d\Omega} = \sum_i c_i\frac{\sigma_i}{4\pi} + \langle b\rangle^2\,[S(Q)-1],
  \qquad
  S(Q) = 1 + \frac{d\sigma/d\Omega - \sum_i c_i\sigma_i/4\pi}{\langle b\rangle^2},
\]
with the bound cross sections $\sigma_i$ and $\langle b\rangle=\sum_i c_i b_i$
of ToScaNA's table, and $\sigma_S = \sigma_{d\sigma/d\Omega}/\langle b\rangle^2$.
With \code{q_grid} the points are rebinned onto a regular grid. The self
scattering is the static value: inelasticity not removed by the Placzek
correction remains in S(Q).

\tip{Both methods were checked on synthetic experiments. pycorrect: Ni powder
in a vanadium can, built with the physics CORRECT inverts, S(Q) recovered
within $3\times10^{-5}$. MCAMS: a hollow Ni sample in a vanadium can with a
gap, a niobium furnace and a heat shield, an absorber, the sample given as
\code{.qdat}; $d\sigma/d\Omega$ and every intermediate intensity are recovered
to $10^{-7}$. To compare the methods on real data, run both and compare
\code{res.factors} and \code{res.dcs}.}

\section{Command line}\label{sec:correct-cli}

The same correction without Python code \newbeh{}:
\begin{lstlisting}
python -m toscana.workflows.correct_cli SAMPLE.par               # MCAMS, files of the .par
python -m toscana.workflows.correct_cli SAMPLE.par --input qdat --scale q --outdir results
python -m toscana.workflows.correct_cli SAMPLE.par --method pycorrect
python -m toscana.workflows.correct_cli --help
\end{lstlisting}
{\sloppy Options: \code{--method mcams|pycorrect}, \code{--input par|adat|qdat} (the
files of the \code{.par} or the same names with that extension),
\code{--scale q|angle|both}, \code{--outdir}, \code{--environment},
\code{--vanadium-in-environment}, \code{--no-absorber}, \code{--placzek},
\code{--vanadium-smoothing}, \code{--n-points}, \code{--ms-particles},
\code{--seed}, \code{--absorber-model}, \code{--background FILE},
\code{--mcams delta|all} (pycorrect with MCAMS factors),
\code{--qgrid}. It writes the factor files (sample and vanadium set-ups), the
intensities and backgrounds on the requested scales, $d\sigma/d\Omega$
(\code{<title>_<method>_dcs.dat}) and S(Q) (\code{.soq}). From Python,
\code{correct_cli.main([...])} does the same and returns the result.\par}

\section{A complete example}\label{sec:example}

The folder \code{example} contains a synthetic D4 experiment with all the
steps \newbeh{}: Ni powder in a vanadium can inside a furnace, a vanadium rod
and a $^{10}$B absorber, i.e.\ the five measurements \code{I_SCE},
\code{I_CE}, \code{I_E}, \code{I_VE}, \code{I_ACE}, each as \code{.adat} and
\code{.qdat}, with the true values. The notebook
\code{D4_corrections_example.ipynb} computes the MCAMS factors, makes the
subtraction step by step (and checks it against \code{run_correction}),
saves $I_S$, $I_V$, the factors, $d\sigma/d\Omega$ and S(Q) on both scales,
and compares with pycorrect. \code{run_example.py} is the same as a script;
\code{correct.py} and \code{run_pycorrect.py} are the command lines of the
previous section and of pycorrect; \code{data/Ni_powder.com} is the CORRECT
command file of the same experiment. \code{example/README.md} lists all the
commands. The environment is a body and the absorber measurement a
configuration (previous paragraphs): the notebook builds and solves the
equations explicitly. On these data the correction recovers the truth with
an rms deviation of $1.0\sigma$ (for $d\sigma/d\Omega$ and for $B$). The
comparison shows CORRECT's furnace factors 2--8\,\% lower than MCAMS's for this
narrow beam (8\,mm) and wide furnace (40\,mm); a brute-force integration
confirms the MCAMS values. Treating the furnace as a background gives
$d\sigma/d\Omega$ within $-6$ to $+3$\,\% here.

\warn{If another package called \code{toscana} is installed in the Python
environment (e.g.\ an older version), \code{python -m toscana...} uses that
one. The example's \code{toscana_path.py} (imported first by its scripts and
notebook) always loads the ToScaNA of the parent folder.}

\section{pycorrect: CORRECT in Python}

\code{toscana.corrections.pycorrect} is a line-by-line translation of
CORRECT 2.24 (M.A.~Howe, R.L.~McGreevy, P.~Zetterstr\"om). It reads the same
command files, so an existing \code{.com} file runs unchanged. It also
reads the \code{.adat} files written by \code{d4creg} (whose header lines
start with \code{\#}). From the folder
of the data:
\begin{lstlisting}
python -m toscana.corrections.pycorrect run.com     # without a file: interactive prompt
\end{lstlisting}
or from Python:
\begin{lstlisting}[language=Python]
import os
from toscana.corrections.pycorrect import Correct
os.chdir("my_data_folder")              # file names in .com files are relative
c = Correct(verbose=False, ask=False)
c.run_file("run.com")                   # writes the output file of the .com
r = c.results[0]                        # dict: theta, dcs, err, factors ...
\end{lstlisting}
The result dictionary contains, besides the differential cross section
\code{dcs} and its error, every factor used: Paalman \& Pings factors
\code{assch}, \code{acsch}, \code{acch}, \ldots, the multiple-scattering
ratios \code{delta_sa}, \code{delta_va}, the Placzek correction. Its manual
(\code{docs/pycorrect/manual}) describes the physics, the command language
and the correspondence with the Fortran routines.

\warn{pycorrect works in \textbf{cm} and cm$^{-1}$, as CORRECT; ToScaNA and
MCAMS work in mm. In \code{.com} files, keep file names relative to the data
folder: in the command language \code{/} starts a qualifier, so an absolute
path containing for instance \code{/T...} is cut there.}

\section{MCAMS: Monte Carlo attenuation and multiple scattering}

\code{toscana.corrections.mcams} simulates neutron paths in layered samples
(sample, can, furnace shields; cylinder, annulus, slab, disc, sphere;
partially filled containers; beams narrower or wider than the sample) and
gives single-scattering attenuation factors, scattering orders and their
angular dependence. Its manual (\code{docs/mcams/manual}, 146 pages) explains
the geometry, the physics, the validation and all the functions. Layers can
be built from a chemical formula and a density; the cross sections then come
from ToScaNA's isotope table:
\begin{lstlisting}[language=Python]
from toscana.corrections.mcams.layer_builder import compute_layer_properties
fe = compute_layer_properties("Fe plate", 1.0, "Fe", 7.87)   # thickness mm, g/cm3
print(fe.summary())        # sigma_scat 11.62 b, sigma_abs 2.56 b, mu 0.12034/mm
\end{lstlisting}

\section{Comparing the two: the bridge}

\code{toscana.corrections.bridge} computes with MCAMS the same quantities as
a pycorrect run, from the pycorrect state, converting the units:
\begin{lstlisting}[language=Python]
from toscana.corrections import bridge
mc = bridge.mcams_factors_like_pycorrect(c.st, r["theta"],
                                         n_particles=300_000, n_batches=8)
# mc["assch"], mc["delta_sa"], ... with mc["assch_err"], ... (MC errors)
\end{lstlisting}
The systematic comparison is in \code{docs/validation} (scripts and a
28-page report). Its main conclusions:
\begin{itemize}
  \item attenuation factors agree within the Monte Carlo error (0.1--0.5\,\%)
    for a beam wider than the sample;
  \item with a beam narrower than the sample, CORRECT is about 1\,\% too high;
  \item CORRECT interpolates the furnace factors with a parabola through
    $0^\circ$, $90^\circ$, $180^\circ$ and misses dips of up to 9\,\%;
  \item the Blech \& Averbach double-scattering table is confirmed, but the
    empirical series used by CORRECT for the higher orders is far too low
    (for the Ni example $\delta = 0.486$ against 0.58--0.63 with MCAMS).
\end{itemize}
When these effects matter (strongly scattering samples, furnaces, narrow
beams), use the MCAMS factors.

% =====================================================================
\chapter{Normalisation models and data operations}\label{ch:models}
% =====================================================================

\section{Vanadium and self scattering}

At low $Q$ a nucleus scatters as if bound, at high $Q$ as if free; the
self-scattering level therefore falls from $\sigma_\mathrm{bound}/4\pi$ to
$\sigma_\mathrm{free}/4\pi = [A/(A+1)]^2\sigma_\mathrm{bound}/4\pi$. ToScaNA
models this with a logistic step centred at $Q_0$ of width $\mathrm{d}Q$
(defaults 7.0 and 2.4\,\AA$^{-1}$, \code{SIGMOID_Q0}, \code{SIGMOID_DQ}):
\begin{align*}
  \texttt{inelastic}(Q) &= \mathrm{lowQ}\,\frac{1 + [A/(1+A)]^2\,\delta}{1+\delta},
  \qquad \delta = e^{(Q-Q_0)/\mathrm{d}Q},\\
  \texttt{vanaQdep}(Q) &= (a_0 + a_1Q + a_2Q^2)\times\texttt{inelastic}(Q),\\
  \texttt{self024}(Q) &= s_0 + q_2Q^2 + q_4Q^4,\\
  \texttt{sigmoidal}(Q) &= \frac{\mathrm{lowQ} + \mathrm{highQ}\,\delta}{1+\delta},
  \qquad \texttt{siglin}(Q) = \texttt{sigmoidal}(Q) + a_1Q .
\end{align*}
These functions have the signature expected by fitting programs (first
argument $Q$, then the parameters), so they can be fitted directly, e.g.
\begin{lstlisting}[language=Python]
import lmfit
from toscana.models.scattering import vanaQdep
model = lmfit.Model(vanaQdep)          # parameters a0, a1, a2, A, lowQ, Q0, dQ
params = model.make_params(a0=1, a1=0, a2=0, A=51.0, lowQ=0.4, Q0=7.0, dQ=2.4)
for name in ("A", "lowQ"):
    params[name].vary = False
fit = model.fit(y, params, x=Q, weights=1/e)
\end{lstlisting}
\code{models.cross_sections} gives the epithermal total cross section
(\code{XS_model}) and the scattering probability including $1/v$ absorption
(\code{scattering_probability}).

\section{Operations on data sets}

All data sets are arrays of three columns $(x, y, \sigma)$.
\begin{itemize}
  \item \code{binary_sum(w1, d1, w2, d2)}: $w_1y_1 + w_2y_2$ with
    $\sigma = \sqrt{w_1^2\sigma_1^2 + w_2^2\sigma_2^2}$ --- e.g.\ sample minus
    container, \code{binary_sum(1, sample, -1, can)}.
  \item \code{wsum2(w1, d1, w2, d2)}: as above, or with \code{w1 = w2 = 0}
    the inverse-variance weighted mean of two measurements, or with
    \code{w2 = 0} the mixture $w_1y_1 + (1-w_1)y_2$.
  \item \code{ratio(y1, e1, y2, e2)}: ratio with errors (lists).
  \item \code{rebin(xdel, wlength, data, xmin, xmax, dbin)}: rebinning onto
    a regular grid.
  \item \code{fit_range}, \code{fittingRange}: points inside a rectangle,
    for fitting a part of a curve.
\end{itemize}
The two data sets of a sum must be on the same $x$ grid; rebin them first
if necessary.

% =====================================================================
\chapter{From S(Q) to real space}\label{ch:realspace}
% =====================================================================

\section{The transform}

\code{numerics.fourier.sineFT} computes $G(r)$, $g(r)$, RDF$(r)$, $T(r)$
and $N(r)$ from $S(Q)$ (definitions in Section~\ref{sec:conventions}):
\begin{lstlisting}[language=Python]
import numpy as np
from toscana.dataio.loading import read_3col
from toscana.numerics.fourier import sineFT
from toscana.dataio.saving import saveRSCF

SQ = read_3col("SiO2_RT.soq")          # columns: Q (1/A), S(Q), error
rho = m.AtomicDensity                  # atoms/A^3 (0.06615 for silica)
r = np.round(np.arange(0.0, 20.0 + 1e-9, 0.01), 10)   # r grid (A)

fou = sineFT(SQ, r, qmax=23.5, density=rho, window=1, with_errors=True)
saveRSCF("SiO2_RT.rscf", fou,
         ["SiO2 at room temperature", "Qmax = 23.5 1/A, Lorch window"])
\end{lstlisting}

\paragraph{Arguments.}
\begin{itemize}
  \item \code{qmax}: the transform uses only $Q < Q_\mathrm{max}$. Choose
    the largest $Q$ where $S(Q)$ is still reliable. Truncation produces
    ripples of period $2\pi/Q_\mathrm{max}$ around each peak of $G(r)$ and
    broadens the peaks.
  \item \code{window}: \code{0} no window; \code{1} the standard Lorch
    function $M(Q)=\sin(\pi Q/Q_\mathrm{max})/(\pi Q/Q_\mathrm{max})$, which
    suppresses the ripples at the cost of a broadening of about
    $\pi/Q_\mathrm{max}$. Both conserve coordination numbers
    (Section~\ref{sec:fourier}) \newbeh{}.
  \item \code{density}: the atomic number density $\rho$ in atoms/\AA$^3$.
    It sets the baseline of $g(r)$ and of RDF$(r)$; an error in $\rho$ shows
    up as $g(r)\neq0$ at low $r$.
  \item \code{selfsca}: the value subtracted from $S(Q)$ (1 for a properly
    normalised $S(Q)$); \code{constant}: an overall factor.
  \item \code{with_errors}: needs a third column with the errors of $S(Q)$;
    the output then has 11 columns.
\end{itemize}

\paragraph{Output.} An array with one row per $r$ and the columns
\begin{center}
\begin{tabular}{@{}l l@{}}
\toprule
\code{with_errors=False} & $r$, $G$, $g$, RDF, $T$, $N$ \\
\code{with_errors=True} & $r$, $G$, $\sigma_G$, $g$, $\sigma_g$, RDF, $\sigma_\mathrm{RDF}$, $T$, $\sigma_T$, $N$, $\sigma_N$ \\
\bottomrule
\end{tabular}
\end{center}
\code{saveRSCF} writes either array with a header that describes every
column; the file reads back with \code{np.loadtxt}.

\section{Coordination numbers}

The number of neighbours between $r_1$ and $r_2$ is the integral of the RDF.
With the 11-column output, the RDF is column 5 and $N(r)$ column 9:
\begin{lstlisting}[language=Python]
from scipy import integrate
r1, r2 = 1.40, 1.85                    # e.g. Si-O first shell in silica
sel = (r >= r1) & (r <= r2)
N12 = integrate.simpson(fou[sel, 5], x=r[sel])
# equivalently, from the running integral:
i1, i2 = np.searchsorted(r, [r1, r2])
N12_bis = fou[i2, 9] - fou[i1, 9]
\end{lstlisting}

\tip{Place $r_1$ where $g(r)=0$ (below the first peak) and $r_2$ at a
minimum. With the limits in these flat regions, the standard Lorch window
and the plain transform give the same coordination number within about
0.1\,\% (Section~\ref{sec:fourier}). Limits on the side of a peak make the
result depend on the window and on $Q_\mathrm{max}$.}

\warn{Make sure the limits are grid points. With \code{np.arange}, the value
3.3 may be stored as 3.3000000000000003, and \code{r <= 3.3} then excludes
it; since the RDF is large, one point can change $N$ by 1\,\%. Round the grid
(\code{np.round(r, 10)}) as in the example.}

\section{Error bars}

The errors of the real-space functions are propagated exactly from the
errors of $S(Q)$ (Section~\ref{sec:errors}). Keep in mind that:
\begin{itemize}
  \item they are statistical only; the uncertainties of the normalisation,
    of $\rho$ and of the self-scattering subtraction usually dominate and
    are not included;
  \item points of $G(r)$ at different $r$ are strongly correlated; do not
    treat them as independent in a fit;
  \item $\sigma_N$ already accounts for these correlations, so it is the
    right error of a coordination number from the running integral
    \emph{starting at the first point of the grid}. The error of
    $N(r_2)-N(r_1)$ is not $\sqrt{\sigma_N(r_2)^2-\sigma_N(r_1)^2}$; when
    $r_1$ lies in the region $g=0$, $\sigma_N(r_2)$ is a good estimate.
\end{itemize}

\section{Back-transform and low-r cleaning}

\code{backFT(Q, r, g, rho, cut)} sets $g(r)=0$ in chosen regions, transforms
to $S(Q)$ and back to real space. It shows how much of $S(Q)$ comes from the
unphysical oscillations of $g(r)$ below the first peak:
\begin{lstlisting}[language=Python]
from toscana.numerics.fourier import backFT
g = fou[:, 3]                                  # g(r) (11-column output)
Q = SQ[:, 0]
S_new, G_new, g_new, rdf_new, T_new, N_new = backFT(Q, r, g, rho, [1.40])
\end{lstlisting}
\code{cut=[r0]} zeroes $r<r_0$; \code{[r0, a1, b1, a2, b2, ...]} zeroes in
addition every interval $(a_k, b_k)$; \code{[-1]} zeroes $g$ up to its last
negative value.

% =====================================================================
\chapter{Working with data files and plots}\label{ch:files}
% =====================================================================

\begin{itemize}
  \item \code{DataXYE(filename)}: reads a 2- or 3-column file (header lines
    starting with \code{\#} or \code{!}); attributes \code{x}, \code{y},
    \code{e}, statistics, and a peak search in $[5^\circ, 40^\circ]$;
    \code{.plot()} chooses the axis label from the extension
    (\code{adat}, \code{reg}: angle; \code{qdat}, \code{soq}: $Q$; \code{pcf},
    \code{pdf}, \code{rdf}, \code{tor}, \code{run}: $r$).
  \item \code{read_3col(filename)} $\to$ array $(n,3)$;
    \code{read_xye(filename)} $\to$ three arrays. A missing error becomes 0.
  \item \code{saveFile_3col}, \code{saveFile_xye}: write 3 columns with a
    header; \code{saveCorrelations}: six real-space columns from separate
    arrays; \code{saveRSCF}: the output of \code{sineFT}.
  \item \code{plotQ}, \code{plotA}, \code{plotR}: set labels, limits, grid
    and legend of the current matplotlib figure.
  \item \code{physics.conversions.ang2q}, \code{q2ang}: $Q\leftrightarrow2\theta$
    for a given wavelength.
\end{itemize}

\section{Surveying numors: d4numor}\label{sec:d4numor}

\code{toscana.d4numor} lists and averages the variables of a set of numors
(temperatures, monitor, counts, angles, \ldots), shows the detector movement
and computes presets. In Python:
\begin{lstlisting}[language=Python]
from toscana import d4numor as d4
ns = d4.NumorSet.from_directory("rawdata", "454600-454620")
ns.summary()                                   # run overview
ns.print_table(["RegTemp", "Monitor"])         # weighted averages
\end{lstlisting}
or on the command line: \code{python -m toscana.d4numor summary -d rawdata -n 454600-454620}.
The manual is \code{docs/d4numor/MANUAL.md}, with an example notebook.

The NeXus files are read by \code{toscana.dataio.nexus} (formerly
\code{nx_reader_D4.py}, version 2), which finds every quantity by name and so
reads both the old format (until 2025: 14 fixed rows in the scan table,
detector data of shape (1, 576, 1)) and the current one (named rows,
detector data (1, 1, 576), motor \code{masked2theta_s}). \code{d4creg} uses the
same reader. To see where each quantity of a file comes from:
\code{python -m toscana.dataio.nexus FILE.nxs}.

% =====================================================================
\chapter{Tests, version control and maintenance}\label{ch:maintenance}
% =====================================================================

\section{Tests}

\code{python tests/run_tests.py} runs the 121 tests in about 80 seconds and
ends with \code{OK} when everything works. Run it after any change to the
code. Appendix~\ref{app:tests} explains what each test checks and how to add
one.

\section{Version control}

The project is a git repository. The original code is commit
\code{1958d36}; all changes are on the branch \code{fixes/group-a}.
Useful commands:
\begin{lstlisting}
git log --oneline                         # history
git show 67b3663                          # one change, with explanation
git diff 1958d36 -- numerics/fourier.py   # compare with the original
git status                                # what has been modified
git checkout main                         # switch branch
\end{lstlisting}

\section{Changing the code safely}

\begin{enumerate}
  \item Write a test that shows the problem (it fails).
  \item Change the code until the test passes; run the whole suite.
  \item Update the docstring of the function, and this manual if needed.
  \item Commit with a message explaining why.
\end{enumerate}

% =====================================================================
\chapter{Troubleshooting}\label{ch:troubleshooting}
% =====================================================================

\begin{longtable}{@{}p{0.40\textwidth} p{0.54\textwidth}@{}}
\toprule
Symptom & Cause and remedy \\
\midrule\endhead
\code{ModuleNotFoundError: No module named 'toscana'} & The package is not on the search path; see Chapter~\ref{ch:install}. \\
\code{partially initialized module 'math'} (old versions) & Python found the old \code{math} folder of the project; fixed by the renaming (\st{C21}). \\
\code{DeprecationWarning: toscana.math is deprecated} & Use \code{toscana.numerics} (and \code{toscana.dataio}, \code{Element}). \\
\code{ValueError: sineFT: window must be 0 ... or 1} & Only no window or standard Lorch exist; the rescaled Lorch was removed. \\
\code{ValueError: ... needs a third column} & \code{with_errors=True} requires $S(Q)$ errors. \\
$g(r)$ clearly non-zero at low $r$ & Wrong $\rho$, wrong normalisation or self-scattering level of $S(Q)$, or $Q_\mathrm{max}$ too small. \\
\code{ValueError: The isotope 'X' does not exist} & Check the symbol; isotopes as \code{"Li6"}, deuterium \code{"H2"}. \\
``The container V is not available'' & Only with versions before \code{ece015d}: end the line with a comment, or update. With the current version the material is really not in the accepted list. \\
Empty figure saved & Old versions saved after \code{show()}; fixed (\st{C22}). \\
Results differ from older analyses & See the table of changes in Section~\ref{sec:migration}. \\
\code{KeyError} reading an old NeXus file (before the integration) & The old \code{load_nxs} read only the current format; update. \\
\code{ModuleNotFoundError: No module named 'beam'} (or \code{pycorrect}, \code{nx_reader_D4}) & A script for the stand-alone versions: import from \code{toscana.corrections.mcams}, \code{toscana.corrections.pycorrect}, \code{toscana.dataio.nexus}. \\
pycorrect cannot find a file whose path contains \code{/T}, \code{/D}, \ldots & Use names relative to the data folder in the \code{.com} file. \\
\bottomrule
\end{longtable}

% =====================================================================
\appendix
% =====================================================================
%  ToScaNA -- Appendices to the user manual
%    A. Audit of the code and corrections made (November 2026)
%    B. The test suite
%    C. Module reference
%
%  Compile stand-alone:      pdflatex ToScaNA_appendices.tex   (twice)
%                       or:  latexmk -pdf ToScaNA_appendices.tex
%
%  Use inside the manual:  in the manual preamble write
%        \def\ToScaNAManual{}
%  load the packages listed below, then % =====================================================================
%  ToScaNA -- Appendices to the user manual
%    A. Audit of the code and corrections made (November 2026)
%    B. The test suite
%    C. Module reference
%
%  Compile stand-alone:      pdflatex ToScaNA_appendices.tex   (twice)
%                       or:  latexmk -pdf ToScaNA_appendices.tex
%
%  Use inside the manual:  in the manual preamble write
%        \def\ToScaNAManual{}
%  load the packages listed below, then % =====================================================================
%  ToScaNA -- Appendices to the user manual
%    A. Audit of the code and corrections made (November 2026)
%    B. The test suite
%    C. Module reference
%
%  Compile stand-alone:      pdflatex ToScaNA_appendices.tex   (twice)
%                       or:  latexmk -pdf ToScaNA_appendices.tex
%
%  Use inside the manual:  in the manual preamble write
%        \def\ToScaNAManual{}
%  load the packages listed below, then \input{ToScaNA_appendices}
%  after \appendix. The stand-alone preamble is then skipped. In a report
%  or book class the three appendices become chapters, otherwise sections.
%  Packages needed: inputenc(utf8) fontenc(T1) lmodern textcomp amsmath amssymb
%  booktabs longtable array tabularx enumitem xcolor(dvipsnames) listings hyperref
%  (with \PassOptionsToPackage{obeyspaces,spaces}{url} before hyperref)
% =====================================================================
\ifdefined\ToScaNAManual\else
\documentclass[11pt,a4paper]{article}
\usepackage[utf8]{inputenc}
\usepackage[T1]{fontenc}
\usepackage{lmodern}
\usepackage{textcomp}
\usepackage[margin=2.3cm]{geometry}
\usepackage{amsmath,amssymb}
\usepackage{booktabs,longtable,array,tabularx}
\usepackage{enumitem}
\usepackage[dvipsnames]{xcolor}
\usepackage{listings}
% keep spaces inside \code{...} and allow line breaks there
\PassOptionsToPackage{obeyspaces,spaces}{url}
\usepackage[hidelinks]{hyperref}
\urlstyle{tt}
\begin{document}
\fi

% ------------------------------------------------------------- macros
% Heading levels: chapters in report/book manuals, sections in articles.
\ifdefined\chapter
  \providecommand{\tsapp}{\chapter}\providecommand{\tssec}{\section}\providecommand{\tssub}{\subsection}
\else
  \providecommand{\tsapp}{\section}\providecommand{\tssec}{\subsection}\providecommand{\tssub}{\subsubsection}
\fi
% \code{...}: verbatim identifiers, breakable after . _ ( , = / (url package)
\ifdefined\code\else\DeclareRobustCommand{\code}[1]{\path{#1}}\fi
\makeatletter\g@addto@macro\UrlBreaks{\do\_\do\(\do\,\do\=\do\[}\makeatother
\providecommand{\fnx}[1]{\par\medskip\noindent\begingroup\raggedright\hangindent=2em{\color{NavyBlue}\bfseries\ttfamily #1}\par\endgroup\nopagebreak\smallskip}
\providecommand{\tn}[1]{{\raggedright\ttfamily #1}}
\providecommand{\fn}[1]{\par\medskip\noindent{\color{NavyBlue}\textbf{\texttt{\detokenize{#1}}}}\par\nopagebreak\smallskip}
\providecommand{\st}[1]{\textsf{\footnotesize\colorbox{black!8}{#1}}}
\providecommand{\fixed}[1]{\textsf{\footnotesize\colorbox{green!15}{fixed #1}}}
\providecommand{\newbeh}{\textsf{\footnotesize\colorbox{orange!20}{results change}}}
\providecommand{\legacy}{\textsf{\footnotesize\colorbox{black!8}{legacy}}}
\lstset{basicstyle=\ttfamily\small,columns=fullflexible,keepspaces=true,
        frame=single,framerule=0.3pt,backgroundcolor=\color{black!3},
        breaklines=true,showstringspaces=false,aboveskip=6pt,belowskip=6pt}
\setlist{itemsep=2pt,topsep=3pt}
\setlength{\emergencystretch}{3em}

\ifdefined\ToScaNAManual\else
\begin{center}
  {\LARGE\bfseries ToScaNA}\\[4pt]
  {\large Total Scattering Neutron Analysis}\\[10pt]
  {\Large Appendices: code audit, corrections, tests and module reference}\\[10pt]
  {\normalsize Version 0.1.0 (branch \code{fixes/group-a}) --- November 2026}
\end{center}
\bigskip
\noindent\textbf{How to read these appendices.}
Appendix~\ref{app:audit} explains what was examined in the code, what was
found and what was changed, with the physical reasoning and the numbers
that justify each change. Appendix~\ref{app:tests} describes the automatic
tests that now protect the code, and how to run them.
Appendix~\ref{app:reference} is a module-by-module reference of the whole
package: what each function does, its arguments and outputs, the formulas
it implements, which files it reads or writes, and whether it was modified.
Throughout, labels such as \st{C2} or \st{P9} refer to the identifiers of
the audit report (\code{ToScaNA_technical_audit.md}): \st{C} for confirmed
errors, \st{P} for possible issues. The marker \newbeh{} flags changes that
modify numerical results with respect to the original code.

\tableofcontents
\clearpage
\appendix
\fi

% =====================================================================
\tsapp{Audit of the code and corrections}\label{app:audit}
% =====================================================================

\tssec{Context and method}

ToScaNA (\emph{Total Scattering Neutron Analysis}) is a Python package for
the reduction and analysis of neutron total-scattering data from the D4
diffractometer at the ILL. Its present form is a port of a single legacy
script, \code{legacy/lib/toscana.py}, into a package. The port was
written ``parity first'': the docstrings repeatedly state that legacy
behaviour, including known defects, was preserved on purpose so that new
and old results could be compared number by number. The legacy script and
its list of known divergences (\code{KNOWN__DIVERGENCES.md}) were not part
of the project tree that was audited.

The work was done in four stages:
\begin{enumerate}
  \item \textbf{Audit (inspection only).} All 52 Python files were read in
    full (about 6\,000 lines plus the 8\,857-line isotope table). Suspected
    problems were reproduced with small, read-only numerical experiments on
    synthetic data, run from a scratch directory so that nothing in the
    project or in the experimental data was modified. The findings were
    classified as \emph{confirmed errors} (C1--C30), \emph{possible issues}
    (P1--P21) and \emph{recommendations}.
  \item \textbf{Version control and tests.} A git repository was created.
    The first commit is the code exactly as received, so that every later
    change can be inspected with \code{git diff}. A test suite (Appendix~\ref{app:tests})
    was written \emph{before} most of the corrections: each test first
    failed on the original code and passes after the correction.
  \item \textbf{Corrections.} Mechanical bugs were fixed directly; changes
    that modify physical results were made only after a decision by the
    author of the code (density and scattering length of vanadium,
    reference wavelength, lattice parameter of Ni, window function).
  \item \textbf{New functionality.} Exact statistical error bars for
    $G(r)$, $g(r)$, RDF$(r)$, $T(r)$ and $N(r)$.
\end{enumerate}

\tssec{Version history}

All changes are on the git branch \code{fixes/group-a}; the branch
\code{main} contains the first two commits only. Each commit message
explains the change in detail (\code{git log}, \code{git show <id>}).

\begin{longtable}{@{}l p{0.78\textwidth}@{}}
\toprule
Commit & Content \\
\midrule\endhead
\tn{1958d36} & Baseline: the code as received, before any change. \\
\tn{67b3663} & Confirmed errors of the D4 reduction chain (Section~\ref{sec:reduction}). \\
\tn{f98cf26} & Test suite (standard library only). \\
\tn{1a9da55} & Mechanical bugs outside the reduction chain (``group A''). \\
\tn{43354d9} & Vanadium density 6.11\,g/cm$^3$. \\
\tn{6785d7d} & Thermal reference wavelength 1.798\,\AA; vanadium $b_\mathrm{coh}=-0.443$\,fm. \\
\tn{337bbaa} & Isotope class renamed \code{Elemento} $\to$ \code{Element}. \\
\tn{67d188a} & Ni lattice parameter 3.52\,\AA{} by default, user-configurable. \\
\tn{d76b9cd} & Fourier window: none, or the standard Lorch function. \\
\tn{0e2a14f} & Subpackages renamed \code{math}$\to$\code{numerics}, \code{io}$\to$\code{dataio}. \\
\tn{255c508} & Physical choices: annulus volume, free cross sections, sigmoid, $Q_0$, weighted Ni fit. \\
\tn{369e286} & Error bars for $G$, $g$, RDF, $T$ and $N$. \\
\bottomrule
\end{longtable}

\noindent To compare any file with the original: \code{git diff 1958d36 -- numerics/fourier.py}
(note that after commit \code{0e2a14f} the old path is \code{math/fourier.py}).
To return to the original code entirely: \code{git checkout 1958d36}.

% ---------------------------------------------------------------------
\tssec{The data-processing chain}\label{sec:chain}

The figure below summarises what the code does from raw data to
real-space functions, and where each step lives.

\begin{lstlisting}
Raw numors (NeXus .nxs, or ILL ASCII)             dataio.acquisition
  -> counts[9 banks x 64 cells], metadata
Angle check |2theta_read - 2theta_requested|       getOneNumor
Poisson errors on RAW counts  sigma = sqrt(N)      getErrorsNumor
Dead time (non-paralysable, per cell)              dataio.corrections
Relative efficiency  N/eps, sigma/eps              getEff
Normalisation to monitor (or time)                 normalise
Angles: 2theta = 2theta_lue + shift[det]
        + 15 deg*det + 0.125 deg*cell              getAngles
Stack numors, subtract zero angle 2theta0          getNumors
Regroup cells onto a 2theta grid                   dataio.d4creg_outputs.getDiffA
  -> .adat (CORRECT format), .reg (D4 format)      saveDiffAngle
Convert to Q and rebin onto a Q grid               numerics.binning.rebin
  -> .qdat (CORRECT format)                        saveDiffQ
--- external program CORRECT: container, environment, absorption,
    multiple scattering, vanadium normalisation (dataio.correct) ---
S(Q)  ->  sine Fourier transform                   numerics.fourier.sineFT
  -> G(r), g(r), RDF(r), T(r), N(r)  [+ errors]    dataio.saving.saveRSCF
Optional: back-transform with low-r cut            numerics.fourier.backFT
Ni powder calibration of lambda and 2theta0        calibration.d4nifit
\end{lstlisting}

\noindent The whole regrouping chain is driven by a parameter file read by
\code{dataio.parameters.readParam}; the line \code{<run>} in that file
launches \code{workflows.d4creg.d4creg}. Corrections for attenuation,
multiple scattering and container are \emph{not} computed in ToScaNA: they
are delegated to CORRECT, for which ToScaNA writes the input file.

% ---------------------------------------------------------------------
\tssec{Conventions, units and formulas}\label{sec:conventions}

\begin{longtable}{@{}l l p{0.52\textwidth}@{}}
\toprule
Quantity & Unit & Notes \\
\midrule\endhead
Scattering angle $2\theta$ & degree & zero angle $2\theta_0$ subtracted in \code{getNumors} \\
Momentum transfer $Q$ & \AA$^{-1}$ & $Q = \frac{4\pi}{\lambda}\sin\theta$ \\
Distance $r$ & \AA & \\
Wavelength $\lambda$ & \AA & \\
Atomic number density $\rho$ & atoms/\AA$^3$ & $\rho = d\,N_A/\bar M \times 10^{-24} = 0.602214076\,d/\bar M$ \\
Macroscopic density $d$ & g/cm$^3$ & \\
Scattering length $b$ & fm & \\
Cross sections $\sigma$ & barn & $1\,\mathrm{b} = 100\,\mathrm{fm}^2$; $\sigma = 4\pi|b|^2/100$ with $b$ in fm \\
Sample dimensions & mm & volumes converted to cm$^3$ (division by 1000) \\
CORRECT input & cm & radius $=$ diameter$/20$ \\
Counts & per monitor & normalised to \code{norm_mon} monitor counts (default $10^6$) or \code{norm_time} seconds \\
Thermal reference & & $v=2200$\,m/s, $\lambda_\mathrm{th}=1.798$\,\AA, $E_\mathrm{th}=25.30$\,meV \\
\bottomrule
\end{longtable}

\paragraph{Real-space functions.} With $F(Q) = Q[S(Q)-1]$ and a window
(modification) function $M(Q)$ that vanishes for $Q \ge Q_\mathrm{max}$,
\code{sineFT} computes (columns of its output)
\begin{align}
  G(r) &= \frac{2}{\pi}\int_0^{Q_\mathrm{max}} Q\,[S(Q)-1]\,M(Q)\,\sin(Qr)\,\mathrm{d}Q
        \;=\; 4\pi\rho r\,[g(r)-1], \label{eq:G}\\
  g(r) &= 1 + \frac{G(r)}{4\pi\rho r} \quad (g(0) \equiv 0),\\
  \mathrm{RDF}(r) &= 4\pi\rho r^2 g(r) = r\,G(r) + 4\pi\rho r^2,\\
  T(r) &= 4\pi\rho r\, g(r) = G(r) + 4\pi\rho r,\\
  N(r) &= \int_{r_0}^{r} \mathrm{RDF}(r')\,\mathrm{d}r' \qquad (r_0 = \text{first point of the $r$ grid}).
\end{align}
$T(r)$ is defined here without the factor $\langle b\rangle^2$ used by some
authors. The integrals are evaluated with Simpson's rule
(\code{scipy.integrate.simpson}) on the given grids; no FFT is used, so
arbitrary, non-uniform grids are allowed. The back-transform \code{backFT}
uses
\begin{equation}
  S(Q) = 1 + \frac{4\pi\rho}{Q}\int r\,[g(r)-1]\sin(Qr)\,\mathrm{d}r,
  \qquad S(0) = 1 + 4\pi\rho\int r^2[g(r)-1]\,\mathrm{d}r .
\end{equation}
These formulas were checked analytically and numerically during the audit
and are correct.

% ---------------------------------------------------------------------
\tssec{Corrections to the reduction chain}\label{sec:reduction}

All items in this section are in commit \code{67b3663}. Before these
corrections, the \emph{intensities} produced by the chain were correct;
the main defects concerned the \emph{error bars}, and the robustness of
the regrouping. The intensities of regrouped and rebinned data were
verified to be identical to the original code (synthetic data, all
in-range cells).

\paragraph{\fixed{C2} Error of the monitor normalisation.} \newbeh{}
\code{normalise} replaced the counts $y$ by $y_n = f\,y$, $f = N_\mathrm{norm}/N_\mathrm{mon}$,
and only afterwards formed the relative error $e/y_n$ with the
\emph{un-normalised} error $e$. The relative error was therefore multiplied
by $N_\mathrm{mon}/N_\mathrm{norm}$. For a cell with $10^4$ counts, the
relative error should be $\sqrt{10^{-4}+1/N_\mathrm{mon}}\approx 0.010$;
the code gave 0.0033 for $N_\mathrm{mon}=10^5$ (3 times too small) and
0.100 for $N_\mathrm{mon}=10^7$ (10 times too large). It was correct only
when the monitor happened to equal $10^6$. Now
\begin{equation}
  y_n = f y, \qquad \sigma_n^2 = (f\,\sigma)^2 + \Bigl(y_n\,\frac{\sigma_{N_\mathrm{mon}}}{N_\mathrm{mon}}\Bigr)^2 .
\end{equation}
Cells with zero counts are flagged \code{NaN} and are therefore excluded by
the regrouping, which is what the original code did in practice (it
produced $0/0$).

\paragraph{\fixed{C5} Poisson errors.} \newbeh{} The error was computed as
$\sqrt{N}$ \emph{after} dead-time and efficiency corrections, i.e.\
$\sqrt{N_\mathrm{raw}/\varepsilon}$ instead of $\sqrt{N_\mathrm{raw}}/\varepsilon$.
Errors are now taken on raw counts and scaled by exactly the same factors
as the counts (dead-time factor, then $1/\varepsilon$). For
$N_\mathrm{raw}=400$ and $\varepsilon=0.8$ the error per cell is now
$25.0$ (it was $22.4$).

\paragraph{\fixed{C4} Dead-time check.} The test
\code{if (correction.all()) > 0} only checks that no factor is exactly
zero: when the measured rate times the dead time exceeded 1, negative
factors were accepted and produced \emph{negative counts} (reproduced:
$\nu\tau = 1.4$ gave $-5\times10^5$ counts). All factors must now be
positive, otherwise a \code{ValueError} is raised. The corrected total
detector counts were computed with the factor of cell [0,0] only; they now
use $\sum N_\mathrm{corr}/\sum N_\mathrm{raw}$. The log message said
``Monitor DT'' for the detector; corrected.

\paragraph{\fixed{C29} Combination of errors when regrouping and rebinning.} \newbeh{}
Both \code{getDiffA} (cells $\to$ angular grid) and \code{rebin} (angle
$\to$ $Q$ grid) form weighted means $\bar y = \sum w_i y_i/\sum w_i$ (weights
= overlap fraction $\times$ monitor share). The errors were averaged in the
same way, $\sum w_i\sigma_i/\sum w_i$, so that they did not decrease when
points were combined. The error of a weighted mean is
\begin{equation}
  \sigma_{\bar y} = \frac{\sqrt{\sum_i w_i^2\sigma_i^2}}{\sum_i w_i}.
\end{equation}
Example: ten points with $\sigma=10$ gave 10.0; now 3.16. On a synthetic
two-numor regrouping the median error fell from 3.0 to 1.58.

\paragraph{\fixed{C6} Size of the angular grid.} \code{getDiffA} sized its
arrays with \code{int((amax-amin)/step)} but filled them with
\code{np.arange(amin, amax, step)}, which has $\lceil(\cdot)\rceil$
elements. Any step that does not divide the range exactly crashed
(reproduced: \code{(1, 140, 0.3)} gave ``could not broadcast (464,) into
(463,)''). The grid size is now taken from the grid itself.

\paragraph{\fixed{C7} Cells outside the angular range.} Bin indices were
obtained with \code{int()}, which truncates towards zero, and were not
checked. A cell just below \code{amin} was counted in bin 0; a cell further
below gave a \emph{negative} index, which Python interprets as counting
from the end of the array: counts measured at $-5^\circ$ appeared near
$135^\circ$. Cells above \code{amax} raised \code{IndexError}. Indices now
use \code{np.floor} and out-of-range parts of a cell are discarded.

\paragraph{\fixed{C8} Last $Q$ point erased.} After rebinning, points with
$y\le0$ are masked together with the previous point; for the first point
the ``previous'' index $-1$ is the \emph{last} point of the array, which was
erased. Corrected.

\paragraph{\fixed{C9} Width used for the $Q$ rebinning.} The angular width
of each regrouped point was hard-coded to $0.125^\circ$ in \code{d4creg}.
It is now the actual step of the regrouped grid (which \code{getDiffA}
enlarges to the cell width if the requested step is smaller). The width
in $Q$ is $\Delta Q = \frac{2\pi}{\lambda}\cos\theta\,\Delta(2\theta)\,\frac{\pi}{180}$.

\paragraph{\fixed{C19} Number of bins.} \code{get_xlim} computed the number
of bins with \code{int()} of a floating-point ratio:
\code{get_xlim(0, 0.3, 0.1)} gave 3 bins instead of 4 because
$0.4/0.1 = 3.9999\ldots$. A tolerance of $10^{-9}$ is now added before
taking the floor; genuinely non-commensurate ranges keep floor semantics.
The same floor/truncation fix was applied to \code{get_bins}.

\paragraph{\fixed{C25} Parsing of the d4creg parameter file.} Lines were
split with \code{split(" ")}, which kept the newline on the last token and
produced empty tokens for repeated spaces. When a value was the last word
of its line it silently became, e.g., \verb|"Nexus\n"|: the data format
then fell back to the ASCII reader, \verb|"True\n"| disabled writing and
plotting, and numor ranges failed their length check. Lines are now cut
at \texttt{\#} (inline comments allowed) and split on any white space.

\paragraph{\fixed{C27} Last line of input files.} Several readers removed
the last character of each line assuming it was a newline; a final line
without newline lost a digit (an efficiency 1.0234 became 1.023). Lines
are now stripped. Affected: \code{getEff}, \code{getDec}, \code{read_xye}.

\paragraph{\fixed{C3} Weighted sum of two data sets.} \newbeh{}
\code{binary_sum} combined $\sqrt{w_1^2 e_1 + w_2^2 e_2}$ instead of
$\sqrt{w_1^2 e_1^2 + w_2^2 e_2^2}$ (dimensionally wrong; $e_1=e_2=2$ gave
2.0 instead of 2.83).

% ---------------------------------------------------------------------
\tssec{Fourier transform and real-space functions}\label{sec:fourier}

\paragraph{\fixed{C1} Window function.} \newbeh{}
The original code offered a step (truncation) and a ``Lorch'' window that
was rescaled so that $\int_0^{Q_\mathrm{max}} M\,\mathrm{d}Q = Q_\mathrm{max}$.
This rescaling gives $M(0) = \pi/\mathrm{Si}(\pi) = 1.696$ instead of 1
($\mathrm{Si}$ is the sine integral), and multiplies $G(r)$ by 1.70 where
the window is smooth. At low $r$, where $g(r)=0$ and $G(r)=-4\pi\rho r$,
the result was $g(r) = 1 - 1.70 = -0.70$ (reproduced). After discussion, the
rescaled version was removed. \code{sineFT(window=...)} and
\code{getSineFT(w=...)} now accept only
\begin{itemize}
  \item \code{0}: no window (truncation at $Q_\mathrm{max}$);
  \item \code{1}: the standard Lorch function
    $M(Q) = \dfrac{\sin(\pi Q/Q_\mathrm{max})}{\pi Q/Q_\mathrm{max}}$, $M(0)=1$;
\end{itemize}
any other value raises \code{ValueError} (the original code treated any
non-zero value as Lorch). The array function \code{LorchN} was replaced by
\code{LorchArray}.

\paragraph{Why the standard Lorch window preserves coordination numbers.}
$M(Q)$ above is the Fourier transform of a box of half-width
$\Delta = \pi/Q_\mathrm{max}$. Multiplying $F(Q)$ by $M(Q)$ is therefore
exactly equivalent to replacing the (truncated) $G(r)$ by its running
average over $[r-\Delta, r+\Delta]$; this is Lorch's original construction
(1969). A running average with a kernel of unit integral, i.e.\ $M(0)=1$:
\begin{enumerate}
  \item conserves the total area of $G(r)$ (the rescaled version inflated
    all areas by 1.70);
  \item changes an integral between finite limits $[r_1, r_2]$ only through
    edge terms of order $\frac{\Delta^2}{6}[G'(r_2) - G'(r_1)]$, which vanish
    when the limits lie where the curve is flat (the region $g=0$ and the
    first minimum);
  \item has no negative lobes, unlike the kernel $\sin(Q_\mathrm{max} r)/r$
    of plain truncation, so it does not move area by oscillation.
\end{enumerate}
Since $N(r)$ integrates $r\,G(r)$ rather than $G(r)$, an additional term
$\frac{\Delta^2}{3}[G(r_2)-G(r_1)]$ appears; with $Q_\mathrm{max}=23.5$\,\AA$^{-1}$
it is about $0.006\,\Delta G$, negligible. Exact conservation for
\emph{arbitrary} limits is impossible with a finite $Q_\mathrm{max}$: it
would require $M(Q)=1$ up to infinite $Q$. Limits that cut through a peak
always lose part of its area. Numerically, on a model liquid whose exact
coordination number is known (Appendix~\ref{sec:synthetic}):

\begin{center}
\begin{tabular}{@{}c c c c c@{}}
\toprule
$Q_\mathrm{max}$ (\AA$^{-1}$) & Limits (\AA) & Exact $N$ & No window & Standard Lorch \\
\midrule
15.0 & 0.0--3.3 & 10.498 & 10.487 & 10.498 \\
15.0 & 2.3--3.3 & 10.498 & 10.497 & 10.490 \\
23.5 & 0.0--3.3 & 10.498 & 10.499 & 10.498 \\
23.5 & 2.3--3.3 & 10.498 & 10.481 & 10.502 \\
\bottomrule
\end{tabular}
\end{center}
With the rescaled window the same coordination numbers came out 4\,\% too
high from the edge of the first shell, and 30\,\% too low from $r=0$.
\textbf{Practical rule:} integrate from a point where $g(r)=0$ to a
minimum. A related pitfall: on a grid produced by \code{np.arange}, the
value \code{3.3} may be stored as 3.3000000000000003 and a mask
\code{r <= 3.3} then drops the end point; since RDF$(3.3\,\text{\AA})\approx 14$
atoms/\AA, this alone changes $N$ by about 1\,\%. Rounding the grid
(\code{np.round(r, 10)}) avoids it.

\paragraph{\fixed{C18} Running coordination number.} \newbeh{} The slice
end of the running integral was computed as \code{1 + int(r/dr)}, which
assumes that the grid starts at 0 and that the division is exact. Even on
a grid starting at 0 in steps of 0.02\,\AA, rounding (e.g.\
$0.58/0.02 = 28.999\ldots$) stopped the integral one point short at 16 of
400 points, producing small steps in $N(r)$; on a grid not starting at 0,
all values were wrong. The loop index is used now; $N(r_0)=0$.

\paragraph{\fixed{C17} Back-transform.} \newbeh{} \code{backFT} divided by
$Q$ and $r$ without protection, producing \code{NaN}/\code{inf} at
$Q=0$ and $r=0$; the analytic limits are now used ($S(0)$ above;
$g(0)\equiv0$ as in \code{sineFT}; $T = 4\pi\rho r g$ computed directly). Its
running integral excluded the current point, unlike \code{sineFT}; made
consistent. The \code{cut} argument: \code{[-1]} zeroes $g(r)$ up to its
last negative value (previously, when $g$ was never negative, the first two
points were zeroed anyway); otherwise any odd-length list
$[r_0, a_1, b_1, a_2, b_2, \ldots]$ zeroes $r<r_0$ and every interval
$(a_k, b_k)$; even lengths raise (previously lengths above 13 were silently
ignored).

\paragraph{\fixed{C20} Sine integral.} \code{getSineIntegral} integrated
numerically on a grid that stopped at $x(1-10^{-4})$ (relative error
$\sim10^{-4}$: $\mathrm{Si}(1)$ gave 0.945999 instead of 0.946083) and
returned \code{NaN} for $x=0$. It now uses the exact \code{scipy.special.sici}.

\paragraph{\fixed{P4} Truncation in \code{getSineFT}.} With \code{w=0} the
argument \code{vmax} was ignored; the integral now stops at \code{vmax}
(results unchanged when \code{vmax} is the default, the end of the data).

\paragraph{\fixed{P5} Error values used as data.} \code{Lorch} and
\code{step} returned $-999$ for $Q_\mathrm{max}\le 0$ after printing a
message; that value was then silently used as a window. They now raise
\code{ValueError}.

% ---------------------------------------------------------------------
\tssec{Error bars of the real-space functions}\label{sec:errors}

\paragraph{Method.} Commit \code{369e286}. Because the transform is linear
in $S(Q)$, statistical errors can be propagated exactly, without Monte
Carlo noise. Writing the Simpson quadrature of Eq.~\eqref{eq:G} as a sum
over the $Q$ points with weights $w_k$,
\begin{equation}
  G(r_j) = \sum_k W_{jk}\,[S(Q_k)-1], \qquad
  W_{jk} = \frac{2}{\pi}\,c\;Q_k\,M(Q_k)\,\sin(Q_k r_j)\,w_k ,
\end{equation}
($c$ is the argument \code{constant}), and assuming uncorrelated $S(Q)$ points
with errors $\sigma_k$,
\begin{equation}
  \sigma_G^2(r_j) = \sum_k W_{jk}^2\,\sigma_k^2,\qquad
  \sigma_T = \sigma_G,\qquad
  \sigma_\mathrm{RDF} = r\,\sigma_G,\qquad
  \sigma_g = \frac{\sigma_G}{4\pi\rho r}\ (r>0).
\end{equation}
The weights $w_k$ are obtained from SciPy itself (by integrating unit
vectors), so they are exactly those used for the values.

\paragraph{The running integral needs the correlations.} Neighbouring
$G(r)$ points are built from the same $S(Q)$ data and are strongly
correlated; adding their variances as if they were independent would be
badly wrong. Instead $N$ is written directly as a linear function of $S$:
$N(r_i) = \sum_j C_{ij}\, r_j\, G(r_j) + \text{const}$, where $C$ is the
matrix of the running Simpson rule, so that
$\sigma_N^2(r_i) = \sum_k \bigl[(C\,\mathrm{diag}(r)\,W)_{ik}\bigr]^2 \sigma_k^2$.
A first version used trapezoidal weights for $C$ and the Monte Carlo test
showed errors up to 11\,\% wrong without window; the final version uses
exactly the running Simpson rule used for $N(r)$ itself. That matrix is
built in $O(N_r^2)$ operations (0.1\,s for 2000 points): odd-length
prefixes are composite Simpson, extended by one three-point panel at a
time; for even lengths SciPy adds a correction on the last interval that
depends only on the last four points. The construction is compared with
SciPy at run time and rebuilt directly (slowly) if a future SciPy uses a
different rule.

\paragraph{Validation.} $S(Q)$ was perturbed 1500 times with Gaussian noise
of the stated errors, each copy was transformed with the ordinary
\code{sineFT}, and the standard deviation of each output was compared with
the analytic error. Ratio analytic/Monte Carlo, median over $r$ [range]:

\begin{center}\small\setlength{\tabcolsep}{4pt}
\begin{tabular}{@{}l c c c c c@{}}
\toprule
Window & $G$ & $g$ & RDF & $T$ & $N$ \\
\midrule
none    & 1.005 [0.95--1.05] & 1.005 [0.95--1.05] & 1.005 [0.95--1.05] & 1.005 [0.95--1.05] & 1.003 [0.96--1.05] \\
Lorch   & 0.998 [0.96--1.06] & 0.997 [0.96--1.06] & 0.998 [0.96--1.06] & 0.998 [0.96--1.06] & 0.998 [0.95--1.06] \\
\bottomrule
\end{tabular}
\end{center}
The spread of 5--6\,\% is the expected statistical spread of a standard
deviation estimated from 1500 samples.

\paragraph{Usage.}
\begin{lstlisting}
fou = sineFT(SQ, r, qmax=23.5, density=rho, window=1, with_errors=True)
saveRSCF("sample.rscf", fou, ["Sample: ...", "Qmax = 23.5, Lorch"])
\end{lstlisting}
\code{SQ} must have three columns ($Q$, $S$, $\sigma_S$). The output has 11
columns, each error beside its function:
$r, G, \sigma_G, g, \sigma_g, \mathrm{RDF}, \sigma_\mathrm{RDF}, T, \sigma_T, N, \sigma_N$.
The file header describes every column, its formula and units. Without
\code{with_errors} the output keeps its 6 columns.

\paragraph{Limits of these error bars.}
\begin{itemize}
  \item They are \emph{statistical} only. Systematic errors (normalisation,
    density, self-scattering subtraction, truncation) are not included, and
    usually dominate.
  \item They assume uncorrelated $S(Q)$ points. The fractional rebinning
    correlates neighbouring points slightly, so the errors may be somewhat
    underestimated.
  \item Errors at different $r$ are strongly correlated. They are correct
    point by point, but a least-squares fit of $G(r)$ that treats them as
    independent overstates the information contained in the data.
\end{itemize}

% ---------------------------------------------------------------------
\tssec{Line shapes, fitting and the Ni calibration}\label{sec:fitting}

\paragraph{\fixed{C11} \code{Lorentzian_error}.} The $x_0$ term contained
\code{e_x0 * 22} (typing error for \code{e_x0**2}) and lacked one factor of
the normalised peak $p = (\gamma/2)^2/[(x-x_0)^2+(\gamma/2)^2]$; the correct
derivative is $\partial L/\partial x_0 = 2A(x-x_0)\,p^2/(\gamma/2)^2$.
Reproduced error: 2.37 instead of 0.032. Now checked against numerical
derivatives.

\paragraph{\fixed{C12} \code{LorGau_error}.} It evaluated an area-normalised
Gaussian while \code{LorGau} uses a unit-height Gaussian, so the returned
value differed from \code{LorGau} (0.600 instead of 1.000 at $x=0$); the
$\eta$, $\sigma$ and $\gamma$ error terms lacked the factor $f_0$. Rewritten
with explicit derivatives.

\paragraph{\fixed{C13} \code{wsum2}.} With $w_1=0$ the error was $w_2 e_2$,
negative for $w_2<0$; now $|w_2|e_2$. With $w_2=0$ and $w_1\notin(0,1]$ it
printed a message and then crashed; it now raises \code{ValueError}.

\paragraph{\fixed{C10} Initial values of the Ni fit.} The amplitudes A2, A3,
A4 started from \code{initial[6..8]} instead of \code{initial[7..9]}
(shifted by one, so A1 and A2 started equal and \code{initial[9]} was never
used). Only starting values were affected, but a non-linear fit can
converge differently.

\paragraph{\fixed{P11} Weighted Ni fit.} \newbeh{} The fit ignored the data
errors. It now uses weights $1/\sigma$; if no valid errors exist
(two-column files store $\sigma=-1$) it fits without weights and prints a
warning. lmfit keeps \code{scale_covar=True}: the reported uncertainties
are multiplied by $\sqrt{\chi^2_\nu}$ and so also account for model misfit.
\code{setting_model_d4nifit} divides the intensities by the value in the
valley between the second and third peaks; the errors are now divided by
the same value.

\paragraph{\fixed{P7} Ni lattice parameter.} \newbeh{} The value
3.52024\,\AA{} was hard-coded in five places. There is now one default,
\code{physics.crystallography.NI_LATTICE_PARAMETER = 3.52}, and every
function that uses it accepts an argument \code{lattice}:
\begin{lstlisting}
setting_model_d4nifit("ni.adat", lattice=3.5238)
fitting_model_d4nifit(nickel, limits, initial, lattice=3.5238)
\end{lstlisting}
In the fit the lattice parameter is a fixed lmfit parameter: the peak
positions determine only $\lambda/a$, so $\lambda$ and $a$ cannot both be
fitted, and the fitted $\lambda$ scales linearly with $a$ (going from
3.52024 to 3.52 lowers $\lambda$ by 0.007\,\%). \code{showing_results_d4nifit}
reads the value back from the fit result.

\paragraph{\fixed{P20} Fitting ranges.} \code{fittingRange0},
\code{fittingRange} and \code{fit_range} were three copies of the same code;
they now share one helper with identical outputs.

% ---------------------------------------------------------------------
\tssec{Physical constants and sample physics}\label{sec:physics}

\paragraph{\fixed{P6} Vanadium density.} \newbeh{} The defaults used
6.51\,g/cm$^3$ (\code{setVanaInfo}) and the corresponding atomic density
0.0769591\,atoms/\AA$^3$ (\code{setSampleInfo}), typed separately. Both now
derive from one constant, \code{VANADIUM_DENSITY_GCC = 6.11}, giving
0.0722305\,atoms/\AA$^3$. The default sample/vanadium density ratio
\code{DenRatioSV} increases by 6.5\,\%.

\paragraph{\fixed{C16} Vanadium scattering length.} \newbeh{} The isotope
table had $b_\mathrm{coh}=-0.443$\,fm but $\sigma_\mathrm{coh}=0.01838$\,b,
which corresponds to $-0.3824$\,fm; \code{setVanaInfo} used $-0.3824$\,fm.
The value $-0.443$\,fm was adopted everywhere and the cross sections now
follow from it: $\sigma_\mathrm{coh} = 4\pi(0.443)^2/100 = 0.02466$\,b and
$\sigma_\mathrm{sca} = \sigma_\mathrm{coh}+\sigma_\mathrm{inc} = 5.10466$\,b
(were 0.01838 and 5.09838\,b). $\sigma_\mathrm{inc}=5.08$\,b and
$\sigma_\mathrm{abs}=5.08$\,b are unchanged, so the vanadium normalisation
changes very little ($\sigma_\mathrm{sca}$ $+0.12$\,\%). A test checks that
the defaults of \code{setVanaInfo} equal the table values.

\paragraph{\fixed{C30} Thermal reference wavelength.} \newbeh{} Tabulated
absorption cross sections refer to neutrons of 2200\,m/s, i.e.\
$\lambda = h/(m_n v) = 1.798$\,\AA{} and $E = 81.8042/\lambda^2 = 25.30$\,meV.
Three different values were in use: 1.8089\,\AA{} (\code{getAbsXS},
$\sigma_\mathrm{abs}$ 0.6\,\% too small), 1.8\,\AA{} (\code{setSampleInfo})
and 25.3\,meV (\code{scattering_probability}). All now derive from
\code{physics.properties.THERMAL_WAVELENGTH = 1.798} and
\code{THERMAL_ENERGY_MEV}; absorption follows the $1/v$ law
$\sigma_\mathrm{abs}(\lambda) = \sigma_\mathrm{abs}^\mathrm{th}\,\lambda/1.798$.

\paragraph{\fixed{C14} Volume of annular samples.} \newbeh{}
\code{Measurement} computed $\pi h (r_o-r_i)^2$ instead of
$\pi h (r_o^2-r_i^2)$. Solid cylinders ($r_i=0$) are unaffected; for a
4/6\,mm annulus the volume was 5 times too small, and the packing fraction
5 times too large.

\paragraph{\fixed{P9} Free-atom cross sections.} \newbeh{} A bound nucleus
of mass ratio $A = m/m_n$ has a free cross section
$\sigma_\mathrm{free} = \sigma_\mathrm{bound}\,[A/(A+1)]^2$. For a mixture
this factor must be applied element by element,
\begin{equation}
  \sigma_\mathrm{free} = \sum_i c_i\,\sigma_i\left(\frac{A_i}{A_i+1}\right)^2 ,
\end{equation}
whereas the code applied it once with the mean mass. The mean mass hides
how light hydrogen is:

\begin{center}
\begin{tabular}{@{}l r r r@{}}
\toprule
Sample & Per element (b) & Mean mass (old, b) & Change \\
\midrule
H$_2$O & 14.908 & 41.116 & $-63.7$\,\% \\
D$_2$O & 3.508 & 4.907 & $-28.5$\,\% \\
SiO$_2$ & 3.170 & 3.212 & $-1.3$\,\% \\
Ni & 17.880 & 17.880 & 0 \\
\bottomrule
\end{tabular}
\end{center}
Applies to \code{Measurement.free} and to \code{FreeCohXS},
\code{FreeIncXS}, \code{FreeScaXS} of \code{setSampleInfo}. Both use
\code{getMassNumber} (neutron mass 1.0086649\,u); previously the isotope
table value 1.00866\,u was also used.

\paragraph{\fixed{C26} \code{sigmoidal}.} \newbeh{} The function returned
$\delta\,\mathrm{lowQ} + (1-\delta)\,\mathrm{highQ}$ with
$\delta = e^{(Q-Q_0)/\mathrm{d}Q}$: unbounded, with negative weights above
$Q_0$, and running from highQ to lowQ, the opposite of what the parameter
names state. It is now the logistic step
\begin{equation}
  f(Q) = \frac{\mathrm{lowQ} + \mathrm{highQ}\,\delta}{1+\delta},\qquad
  f(0)\approx\mathrm{lowQ},\ f(\infty) = \mathrm{highQ},\ f(Q_0) = \tfrac12(\mathrm{lowQ}+\mathrm{highQ}),
\end{equation}
computed with \code{scipy.special.expit} so that large $|Q-Q_0|/\mathrm{d}Q$
cannot overflow. This is the same shape as \code{models.scattering.inelastic}
(with $\mathrm{highQ} = \mathrm{lowQ}\,[A/(A+1)]^2$). \code{siglin}
ignored its argument \code{a1}; it is now $\texttt{sigmoidal}(Q) + a_1 Q$.
\emph{This interpretation of \code{a1} should be confirmed by the author.}

\paragraph{\fixed{P21} One default transition.} \code{inelastic} used
$Q_0=7.0$\,\AA$^{-1}$ and \code{vanaQdep} 7.4\,\AA$^{-1}$. One pair of
defaults, \code{SIGMOID_Q0 = 7.0} and \code{SIGMOID_DQ = 2.4}
(\code{numerics.line_shapes}), now serves \code{sigmoidal}, \code{siglin},
\code{inelastic} and \code{vanaQdep} (effect on \code{vanaQdep}
$\le 0.16$\,\%).

\paragraph{\fixed{C15} Bound cross sections of \code{Element}.} The
attributes \code{scoh_bound} and \code{sinc_bound} were $4\pi|b|^2$ in fm$^2$
but labelled barns (100 times too large: Ni gave 1333 instead of 13.3).
Divided by 100. They were not used elsewhere in the package.

% ---------------------------------------------------------------------
\tssec{Structure and robustness of the code}\label{sec:structure}

\paragraph{\fixed{C21} Names that hid the standard library.} The
subpackages \code{math/} and \code{io/} have the names of standard Python
modules. When Python was started in the project directory (for example a
notebook opened there), \code{import math} found the project folder
instead; \code{import lmfit} failed with ``partially initialized module
`math' has no attribute `floor'''. They were renamed (with \code{git mv},
keeping their history):
\begin{center}
\code{toscana.math} $\to$ \code{toscana.numerics} \qquad
\code{toscana.io} $\to$ \code{toscana.dataio}
\end{center}
The old dotted names remain importable through a small import hook
(\code{_aliases.py}) and emit a \code{DeprecationWarning}; they return the
\emph{same} module objects (an earlier version of the hook loaded a second
copy of each module, which a test detected). No file or folder of the
project now has the name of a standard-library module, and a test checks
this against \code{sys.stdlib_module_names}.

\paragraph{English naming.} The isotope class \code{Elemento} is now
\code{Element}. \code{Elemento} and \code{elemento} remain available as
deprecated aliases. Other Spanish names in the same module
(\code{NOMBRES}, local variables) were not changed.

\paragraph{Other fixes.}
\begin{itemize}
  \item \fixed{C22} \code{DataXYE} crashed on files without any point in
    $[5,40]$ (the range of its automatic peak search), e.g.\ $r$-space
    files; it saved figures \emph{after} \code{plt.show()}, which gives
    empty images with interactive back-ends.
  \item \fixed{P15} An unknown isotope name now raises \code{ValueError}
    instead of returning a half-built object.
  \item \fixed{P13} \code{Measurement.showTable} replaced \code{sys.stdout}
    with an unclosed file and then restored \code{sys.__stdout__}, which
    broke Jupyter output; it now uses \code{contextlib.redirect_stdout}.
  \item \fixed{C23} Characters such as \AA, $^\circ$ and $\pm$ were stored
    doubly encoded (\AA{} appeared as two wrong characters); repaired. \fixed{C24} The \code{.reg}
    header gave the wavelength in degrees; now \AA.
  \item \st{P19} The audit claimed that \code{smooth_curve} had an
    asymmetric window; this was wrong (even sizes become the next odd
    size). Only its documentation was completed.
\end{itemize}

% ---------------------------------------------------------------------
\tssec{What changes for the user}\label{sec:migration}

\begin{longtable}{@{}p{0.36\textwidth} p{0.58\textwidth}@{}}
\toprule
Before & Now \\
\midrule\endhead
\code{from toscana.math...}, \code{toscana.io...} & \code{toscana.numerics...}, \code{toscana.dataio...} (old names work, with a warning) \\
\code{Elemento("Ni")} & \code{Element("Ni")} (old name works, with a warning) \\
\code{LorchN(Q, Qmax)} & removed; \code{LorchArray(Q, Qmax)} is the standard Lorch \\
\code{sineFT(..., window=1)} rescaled Lorch & standard Lorch, $M(0)=1$; values other than 0, 1 raise \\
\code{sineFT} 6 columns & unchanged; \code{with_errors=True} gives 11 columns \\
\code{saveRSCF} 6 columns, no header & 6 or 11 columns, header describing each column \\
Error bars of \code{.adat}, \code{.reg}, \code{.qdat} & smaller and statistically correct (C2, C5, C29) \\
$N(r)$ & correct at every point (C18) \\
\code{Element.scoh_bound} & in barns (100 times smaller) \\
Vanadium defaults & 6.11\,g/cm$^3$, $b=-0.443$\,fm, $\sigma_\mathrm{coh}=0.02466$\,b \\
$\sigma_\mathrm{abs}(\lambda)$ & $\propto\lambda/1.798$ \\
Free cross sections & per element (large change for H, D) \\
\tn{sigmoidal} & a true logistic step \\
Ni calibration & weighted fit, lattice 3.52\,\AA{} by default, \code{lattice=} argument \\
Silent prints on bad input & \code{ValueError} (unknown isotope, invalid \code{cut}, $Q_\mathrm{max}\le0$, invalid window or weights, negative dead-time factor) \\
\bottomrule
\end{longtable}

% ---------------------------------------------------------------------
\tssec{Open items}\label{sec:open}

\paragraph{Need data or instrument knowledge.}
\begin{itemize}
  \item \st{P1} Angle labels of \code{getDiffA}: the angle written for a bin
    is its lower edge, while a cell centred in the bin contributes fully to
    it, i.e.\ a shift of half a step ($0.0625^\circ$ for $0.125^\circ$). It
    is absorbed by the zero angle if the Ni calibration was regrouped in the
    same way; to be confirmed with a real Ni run.
  \item \st{P2} The dead time is applied to the rate of each \emph{cell};
    for position-sensitive banks it is usually governed by the rate of the
    whole bank or of its electronics.
  \item \st{P3} Data are binned twice (cells $\to 2\theta$ $\to Q$), which
    adds broadening and correlations; a direct cells $\to Q$ rebinning
    would avoid one step.
  \item \st{P8} The reactor power is fixed to 1\,MW for NeXus files (only
    the logged rates per MW are affected); the NeXus field is needed.
  \item \st{P16} The ASCII numor reader assumes 10 values per line and a
    block of 11 lines per detector; a sample file is needed to check the
    case \code{ncell} multiple of 10.
\end{itemize}

\paragraph{Found while writing this appendix (not yet changed).}
\begin{itemize}
  \item In \code{getOneNumor}, the $Q$ column of each numor (used in the
    per-numor \code{.cdat} files) is computed from the angles \emph{before}
    the zero angle is subtracted; the main \code{.qdat} output is computed
    from corrected angles and is not affected.
  \item With the default settings the log file is renamed to
    \code{d4creg.log.log} (the extension is added twice).
  \item \code{generate_com} still runs the placeholder command
    \code{executable_name arg1} and writes \code{output.txt} in the current
    directory.
  \item \code{saveCORRECT} writes the keyword \code{wavelenght}; to be
    checked against what CORRECT accepts.
  \item \code{getRunningParams} (sample-description file) still splits lines
    on single spaces and keeps the newline on a text value that ends its
    line (the defect corrected in \code{readParam}, C25): \code{<Cmater> V}
    gives \code{"V\textbackslash n"} and a false ``not available'' warning.
    Work-around: end every line with a comment.
  \item \sloppy \code{saveCORRECT} needs attributes (\code{sampleTitle},
    \code{sampleOuterRadius}, \code{sampleTemperature}, \code{sampleDensity},
    \code{samplePackingFraction}, \code{sampleFullness}, \code{containerFile},
    \code{environmentFile}, \code{vanadium}, \code{vanadiumFile}) that
    \code{Measurement} does not provide.
\end{itemize}

\paragraph{Found while writing the docstrings (not yet changed).}
\begin{itemize}
  \item \code{getDiffA}: for an angular step larger than the cell width, the
    split fractions of a cell straddling a bin edge are computed in bin
    units (0.25/0.75 instead of 0.5/0.5 for a $0.25^\circ$ step); exact at
    the default step $0.125^\circ$.
  \item \code{get_bins}/\code{rebin}: different treatment of the two ends of
    the grid (see Appendix~\ref{app:reference}); \code{rebin} still prints a
    debug line; $Q\lambda/4\pi>1$ or a \code{NaN} abscissa raises a
    \code{ValueError} from \code{int(nan)}.
  \item \code{readParam}: \code{writeNumor} forced to 1 by a single-numor run
    stays 1 for the following \code{<run>} lines of the same file; unknown
    tags and extra values on \code{<num>}/\code{<add>} lines are ignored
    without a message.
  \item \code{getEff}/\code{getDec}: bank or cell numbers $\le 0$ are written
    to elements counted from the end of the array.
  \item \code{setSampleInfo} treats the sample as a solid cylinder (no inner
    diameter); the annulus correction C14 exists only in \code{Measurement}.
  \item Two meanings of ``A'': \code{setSampleInfo()["A"]} averages the
    integer nucleon numbers, \code{setVanaInfo()["A"]} is the mass ratio
    $M/m_n$; \code{isotopes.core.getFree} uses the integer value.
  \item \code{Measurement}: an element given in two \code{<SatoNN>} lines is
    counted twice in \code{natoms} but once in \code{atoms}, so the
    concentrations do not sum to 1.
  \item \code{Element}: overriding \code{bcoh}/\code{binc} does not update
    \code{re_bcoh} etc.; misspelt override keys are ignored silently;
    \code{sinc_bound} is 0 when the table has no incoherent length (e.g.\
    natural V, although $\sigma_\mathrm{inc}=5.08$\,b).
  \item \code{DataXYE.plot} saves to \code{../regdata/} followed by the full
    path of the file without extension, which fails for files in
    sub-directories.
  \item Ni calibration: \code{setting_model_d4nifit} fails (\code{IndexError},
    \code{UnboundLocalError}) if fewer than three peaks are found or no point
    lies above the valley; the fit drops all weights if a single error is
    not positive; the fit-statistics box is cut from lmfit's report at
    fixed character positions; \code{reflections_fcc} gives fewer than ten
    reflections for $\lambda\ge1.355$\,\AA{} (not reachable on D4).
  \item \code{stop()} raises in Jupyter, where standard input is not a
    terminal.
\end{itemize}

\paragraph{Engineering.} A \code{pyproject.toml} so that the package
installs as \code{toscana}; logging instead of \code{print}; output and log
paths as explicit arguments instead of fixed relative paths
(\code{../logfiles}, \code{../regdata}); vectorised transforms and binning;
Monte Carlo propagation of \emph{systematic} parameters ($\lambda$,
$2\theta_0$, density, normalisation); \st{P10} keeping the sign of
$\langle b\rangle$ in \code{Measurement.bcoh}.
% =====================================================================
\tsapp{The test suite}\label{app:tests}
% =====================================================================

\tssec{Purpose and principles}

The tests are small programs that call ToScaNA functions on inputs whose
correct answer is known, and compare. They serve three purposes: they
\emph{document} what each function is expected to do; they \emph{prove}
that the corrections of Appendix~\ref{app:audit} work; and they
\emph{protect} the code: after any future change, running them shows
immediately whether something that used to work is broken.

The suite follows a few rules.
\begin{itemize}
  \item \textbf{Known answers.} Each test compares with an analytic result
    (Bragg's law, an exact integral, a derivative), with an independent
    calculation (finite differences, a Monte Carlo simulation, a direct
    SciPy call), or with an invariant (constant input gives constant
    output; values must not change when only errors are requested).
  \item \textbf{Synthetic data only.} No experimental file is read or
    written. Functions that write files (logs, \code{.adat}, figures) are
    run inside a temporary directory that is deleted afterwards. Functions
    that read raw numors are given a fake reader.
  \item \textbf{Written before the fix.} For every correction, the test was
    written first and seen to fail on the original code (19 tests failed
    for the ``group A'' bugs, for example), then the code was corrected.
  \item \textbf{Standard library only.} The suite uses Python's
    \code{unittest}, so nothing needs to be installed. It can also be run
    by \code{pytest} if available. One test uses lmfit and is skipped
    automatically when lmfit is not installed.
\end{itemize}

\tssec{Running the tests}

From the project directory (or any other directory):
\begin{lstlisting}
python tests/run_tests.py                  # all tests (about 35 s)
python tests/run_tests.py -v               # one line per test
python tests/run_tests.py test_fourier     # one file
python tests/run_tests.py test_fourier.TestBackFT.test_q_zero_limit
\end{lstlisting}
The last line of the output is \code{OK} when every test passes. Otherwise
each failure is listed with the name of the test, the line of the
comparison that failed and both values, for example
\begin{lstlisting}
FAIL: test_bin_count (test_reduction.TestRebin.test_bin_count)
AssertionError: 3 != 4
\end{lstlisting}
which reads: the test expected 4 bins and the function returned 3. About
30 of the 35 seconds are spent in the Monte Carlo validation of the error
bars.

\tssec{Organisation of the \texttt{tests} folder}

\begin{longtable}{@{}l p{0.64\textwidth}@{}}
\toprule
File & Role \\
\midrule\endhead
\code{run_tests.py} & Runs the suite (discovers every \code{test_*.py}). \\
\code{_pkg.py} & Makes the project importable as \code{toscana} without installing it: the directory is called \code{ToScaNA-pilot}, so the package is registered under the name \code{toscana} with \code{importlib}. It also keeps the project directory off \code{sys.path} and disables the writing of \code{.pyc} files into the project. \\
\code{helpers.py} & Shared tools: \code{quiet(f, ...)} calls a function with its printed output suppressed; \code{in_tempdir()} runs a block of code in a fresh temporary directory \code{<tmp>/work} (so that the relative paths \code{../logfiles} and \code{../regdata} used by the code land in \code{<tmp>}); \code{synthetic_liquid()} builds the model liquid described below. \\
\code{test_physics.py} & Unit conversions, crystallography, densities, cross sections, thermal reference, vanadium defaults, scattering models. \\
\code{test_reduction.py} & The D4 reduction chain, including an end-to-end run. \\
\code{test_fourier.py} & Fourier transforms, windows, coordination numbers, back-transform. \\
\code{test_fourier_errors.py} & Error bars of the real-space functions and the output file. \\
\code{test_line_shapes_ops.py} & Line shapes with errors, weighted sums, fitting ranges. \\
\code{test_physical_choices.py} & Annulus volume, free cross sections, sigmoid, $Q_0$, weighted Ni fit. \\
\code{test_misc.py} & Isotopes, data files, Ni calibration, \code{showTable}, project layout, text encoding. \\
\bottomrule
\end{longtable}

\tssec{The synthetic liquid}\label{sec:synthetic}

Several tests need an $S(Q)$ whose real-space functions are known exactly.
\code{helpers.synthetic_liquid(rho=0.1)} builds one from a model pair
correlation function on $r\in[0,30)$\,\AA{} in steps of 0.005\,\AA:
\begin{equation}
  g(r) = \begin{cases} 0 & r < 2.5\ \text{\AA} \\[2pt]
   1 + 0.8\,\exp\!\bigl[-(r-2.8)^2/0.02\bigr] & r \ge 2.5\ \text{\AA} \end{cases}
\end{equation}
a hard core of 2.5\,\AA{} and one coordination shell centred at 2.8\,\AA{}
(Gaussian of standard deviation 0.1\,\AA). Its structure factor is
computed on $Q\in[0,30)$\,\AA$^{-1}$ in steps of 0.01\,\AA$^{-1}$ by direct
integration, $S(Q) = 1 + \frac{4\pi\rho}{Q}\int r[g(r)-1]\sin(Qr)\,\mathrm{d}r$,
with $S(0)$ set equal to the first point. The function returns the
two-column array $(Q, S)$ and the exact $r$ and $g$, from which exact
coordination numbers are obtained. For the error-bar tests an error column
$\sigma(Q) = 0.01\,(1+Q/10)$ is added, growing with $Q$ as in real data.

\tssec{What each test checks}

In the tables below, the test name is given without its \code{test_}
prefix, grouped by test class.

\tssub*{\texttt{test\_physics.py}}
\begin{longtable}{@{}p{0.36\textwidth} p{0.58\textwidth}@{}}
\toprule Test & What is checked \\ \midrule\endhead
\multicolumn{2}{@{}l}{\textit{TestConversions}} \\
\tn{round\_\allowbreak{}trip} & \code{q2ang(ang2q(x))} returns $x$ for $2\theta = 1\ldots140^\circ$. \\
\tn{known\_\allowbreak{}value} & $2\theta=60^\circ$, $\lambda=0.5$\,\AA{} gives $Q = 4\pi$. \\
\tn{bad\_\allowbreak{}wavelength} & $\lambda=0$ raises \code{ValueError}. \\
\multicolumn{2}{@{}l}{\textit{TestCrystallography}} \\
\tn{ni\_\allowbreak{}reflections\_\allowbreak{}match\_\allowbreak{}bragg} & The ten first fcc reflections equal $2\arcsin(\lambda\sqrt{h^2+k^2+l^2}/2a)$ for (111)\ldots(333). \\
\tn{default\_\allowbreak{}lattice\_\allowbreak{}is\_\allowbreak{}ni\_\allowbreak{}3\_\allowbreak{}52} & Default lattice parameter is 3.52\,\AA. \\
\tn{zero\_\allowbreak{}angle\_\allowbreak{}shift} & A zero angle shifts all reflections by the same amount. \\
\multicolumn{2}{@{}l}{\textit{TestProperties}} \\
\tn{water\_\allowbreak{}atomic\_\allowbreak{}density} & 1\,g/cm$^3$ and 6\,g/mol/atom give $0.602214076/6$ atoms/\AA$^3$. \\
\tn{cylinder\_\allowbreak{}volume} & $\pi r^2 h$. \\
\tn{free\_\allowbreak{}cross\_\allowbreak{}section} & $\sigma[A/(A+1)]^2$. \\
\tn{absorption\_\allowbreak{}reference\_\allowbreak{}wavelength} & $\sigma_\mathrm{abs}(1.798\,\text{\AA})$ equals the tabulated value; $1/v$ scaling. \\
\tn{thermal\_\allowbreak{}references\_\allowbreak{}consistent} & 1.798\,\AA{} $\leftrightarrow$ 25.30\,meV $\leftrightarrow$ 2200\,m/s; the energy route (\code{scattering_probability}) and the wavelength route (\code{getAbsXS}) give the same absorption. \\
\multicolumn{2}{@{}l}{\textit{TestVanadiumStandard}} \\
\tn{default\_\allowbreak{}density\_\allowbreak{}and\_\allowbreak{}consistency} & Vanadium density 6.11\,g/cm$^3$; the atomic density default of \code{setSampleInfo} equals that of \code{setVanaInfo}. \\
\multicolumn{2}{@{}l}{\textit{TestModels}} \\
\tn{inelastic\_\allowbreak{}limits} & \code{inelastic} tends to lowQ at low $Q$ and to lowQ$\,[A/(A+1)]^2$ at high $Q$. \\
\tn{xs\_\allowbreak{}model\_\allowbreak{}high\_\allowbreak{}energy\_\allowbreak{}limit} & Epithermal cross section tends to the free value at high energy. \\
\bottomrule
\end{longtable}

\tssub*{\texttt{test\_reduction.py}}
\begin{longtable}{@{}p{0.36\textwidth} p{0.58\textwidth}@{}}
\toprule Test & What is checked \\ \midrule\endhead
\multicolumn{2}{@{}l}{\textit{TestCorrections}} \\
\tn{binary\_\allowbreak{}sum\_\allowbreak{}variances} & $10\pm2$ minus $5\pm2$ gives $5\pm\sqrt8$ (C3). \\
\tn{normalise\_\allowbreak{}relative\_\allowbreak{}error\_\allowbreak{}independent\_\allowbreak{}of\_\allowbreak{}monitor} & For monitors $10^5$, $10^6$, $10^7$ the relative error is $\sqrt{1/N+1/N_\mathrm{mon}}$ (C2). \\
\tn{normalise\_\allowbreak{}zero\_\allowbreak{}counts\_\allowbreak{}flagged} & A cell with zero counts becomes \code{NaN}. \\
\tn{dead\_\allowbreak{}time} & $\nu\tau>1$ raises; corrected counts $=N/(1-N\tau/t)$; total = sum of corrected cells (C4). \\
\tn{efficiency\_\allowbreak{}file\_\allowbreak{}without\_\allowbreak{}final\_\allowbreak{}newline} & The last value of an efficiency file is read intact (C27). \\
\tn{numor\_\allowbreak{}list\_\allowbreak{}expansion} & \code{"000020-000022"} expands to three numors; a reversed range raises. \\
\multicolumn{2}{@{}l}{\textit{TestOneNumor}} \\
\tn{poisson\_\allowbreak{}errors\_\allowbreak{}from\_\allowbreak{}raw\_\allowbreak{}counts} & With a fake NeXus reader (400 counts per cell, efficiency 0.8): counts 500, error $\sqrt{(20/0.8)^2 + (500\times10^{-3})^2}$, i.e.\ Poisson on raw counts plus monitor statistics (C5). \\
\multicolumn{2}{@{}l}{\textit{TestGetDiffA}} (two synthetic numors of 576 cells) \\
\tn{constant\_\allowbreak{}input\_\allowbreak{}gives\_\allowbreak{}constant\_\allowbreak{}output} & Constant counts give the same constant on the regrouped grid. \\
\tn{errors\_\allowbreak{}are\_\allowbreak{}combined\_\allowbreak{}not\_\allowbreak{}averaged} & Combined errors never exceed the input error and their median is clearly smaller (C29). \\
\tn{non\_\allowbreak{}commensurate\_\allowbreak{}step} & Step $0.3^\circ$ on $[1,140)$ works; the grid has the length of \code{np.arange} (C6). \\
\tn{out\_\allowbreak{}of\_\allowbreak{}range\_\allowbreak{}cells\_\allowbreak{}are\_\allowbreak{}dropped} & Cells at $-5^\circ$ and $200^\circ$ do not appear anywhere (C7). \\
\multicolumn{2}{@{}l}{\textit{TestRebin}} \\
\tn{bin\_\allowbreak{}count} & $[0,0.3]$ step 0.1 gives 4 bins; $[0,0.25]$ gives 3 (C19). \\
\tn{error\_\allowbreak{}of\_\allowbreak{}mean} & Ten points of $100\pm10$ per bin give $100\pm10/\sqrt{10}$ (C29). \\
\multicolumn{2}{@{}l}{\textit{TestWorkflow}} \\
\tn{read\_\allowbreak{}param\_\allowbreak{}and\_\allowbreak{}d4creg\_\allowbreak{}end\_\allowbreak{}to\_\allowbreak{}end} & A parameter file (with an inline comment and values at the end of lines) is read by \code{readParam}; its \code{<run>} line executes the whole d4creg chain on four fake numors; \code{.adat} and \code{.qdat} are written, start at the requested $Q$ and contain only finite values (C25, C8). \\
\bottomrule
\end{longtable}

\tssub*{\texttt{test\_fourier.py}}
\begin{longtable}{@{}p{0.36\textwidth} p{0.58\textwidth}@{}}
\toprule Test & What is checked \\ \midrule\endhead
\multicolumn{2}{@{}l}{\textit{TestSineFT}} (synthetic liquid, $Q_\mathrm{max}=25$, $r$ step 0.02\,\AA) \\
\tn{column\_\allowbreak{}relations} & $g = 1+G/4\pi\rho r$, RDF $=4\pi\rho r^2 g$, $T = 4\pi\rho r g$. \\
\tn{low\_\allowbreak{}r\_\allowbreak{}limit} & No window: $|g(r)| < 0.02$ for $0.5<r<1.5$\,\AA. \\
\tn{lorch\_\allowbreak{}low\_\allowbreak{}r\_\allowbreak{}limit} & Same with the Lorch window (C1; it was $-0.70$). \\
\tn{coordination\_\allowbreak{}number} & First-shell $N$ within 5\,\% of the exact value. \\
\tn{running\_\allowbreak{}integral\_\allowbreak{}from\_\allowbreak{}zero} & $N(r_i)$ equals a direct Simpson integral up to $r_i$ (C18). \\
\tn{running\_\allowbreak{}integral\_\allowbreak{}grid\_\allowbreak{}not\_\allowbreak{}starting\_\allowbreak{}at\_\allowbreak{}zero} & Same on a grid starting at 1\,\AA; $N(r_0)=0$. \\
\tn{invalid\_\allowbreak{}window\_\allowbreak{}raises} & \code{window=2} raises. \\
\multicolumn{2}{@{}l}{\textit{TestCoordinationNumbers}} \\
\tn{lorch\_\allowbreak{}conserves\_\allowbreak{}coordination\_\allowbreak{}number} & For $Q_\mathrm{max}=15$ and 23.5\,\AA$^{-1}$, and limits 0--3.3 and 2.3--3.3\,\AA, $N$ with the Lorch window agrees within 0.5\,\% with the exact value and with the no-window value (the property discussed in Section~\ref{sec:fourier}). \\
\multicolumn{2}{@{}l}{\textit{TestBackFT}} \\
\tn{all\_\allowbreak{}outputs\_\allowbreak{}finite} & No \code{NaN}/\code{inf} although the grids contain $Q=0$ and $r=0$ (C17). \\
\tn{q\_\allowbreak{}zero\_\allowbreak{}limit} & $S(0) = 1+4\pi\rho\int r^2(g-1)\,\mathrm{d}r$, continuous with the next point. \\
\tn{r\_\allowbreak{}relations} & $g(0)=0$; $T$, RDF consistent with $g$; running integral includes the current point. \\
\tn{auto\_\allowbreak{}cut\_\allowbreak{}without\_\allowbreak{}negative\_\allowbreak{}values\_\allowbreak{}cuts\_\allowbreak{}nothing} & \code{cut=[-1]} on a $g(r)$ that is never negative cuts nothing. \\
\tn{cut\_\allowbreak{}intervals} & A 17-value cut list works; an even-length list raises. \\
\multicolumn{2}{@{}l}{\textit{TestSineHelpers}} \\
\tn{sine\_\allowbreak{}integral} & \code{getSineIntegral} equals \code{scipy.special.sici} to $10^{-10}$; $\mathrm{Si}(0)=0$ (C20). \\
\tn{get\_\allowbreak{}sine\_\allowbreak{}ft\_\allowbreak{}vmax\_\allowbreak{}truncates} & \code{vmax} truncates the integral; default unchanged (P4). \\
\tn{windows\_\allowbreak{}reject\_\allowbreak{}bad\_\allowbreak{}qmax} & $Q_\mathrm{max}\le 0$ raises (P5). \\
\tn{lorch\_\allowbreak{}shape} & $M(0)=1$, $M(Q_\mathrm{max}/2)=2/\pi$, $M(Q_\mathrm{max})=0$. \\
\tn{lorch\_\allowbreak{}array\_\allowbreak{}is\_\allowbreak{}standard\_\allowbreak{}lorch} & \code{LorchArray} equals the scalar \code{Lorch} everywhere. \\
\tn{get\_\allowbreak{}sine\_\allowbreak{}ft\_\allowbreak{}lorch} & \code{getSineFT(w=1)} equals a direct Simpson integral with the Lorch window; \code{w=2} raises. \\
\bottomrule
\end{longtable}

\tssub*{\texttt{test\_fourier\_errors.py}}
\begin{longtable}{@{}p{0.36\textwidth} p{0.58\textwidth}@{}}
\toprule Test & What is checked \\ \midrule\endhead
\multicolumn{2}{@{}l}{\textit{TestErrorColumns}} \\
\tn{values\_\allowbreak{}identical\_\allowbreak{}to\_\allowbreak{}plain\_\allowbreak{}transform} & Requesting errors does not change the values (bit-identical), both windows. \\
\tn{error\_\allowbreak{}relations} & $\sigma_g=\sigma_G/4\pi\rho r$, $\sigma_\mathrm{RDF}=r\sigma_G$, $\sigma_T=\sigma_G$; $\sigma_N(r_0)=0$ and $\sigma_N$ increases with $r$. \\
\tn{missing\_\allowbreak{}error\_\allowbreak{}column\_\allowbreak{}raises} & Two-column input with \code{with_errors=True} raises. \\
\tn{errors\_\allowbreak{}scale\_\allowbreak{}linearly} & Doubling $\sigma_S$ doubles every error. \\
\multicolumn{2}{@{}l}{\textit{TestRunningSimpsonMatrix}} \\
\tn{matches\_\allowbreak{}scipy\_\allowbreak{}on\_\allowbreak{}any\_\allowbreak{}grid} & The fast running-Simpson matrix reproduces SciPy on uniform, offset, random non-uniform and two-point grids to $10^{-12}$. \\
\multicolumn{2}{@{}l}{\textit{TestAgainstMonteCarlo}} \\
\tn{analytic\_\allowbreak{}errors\_\allowbreak{}match\_\allowbreak{}monte\_\allowbreak{}carlo} & 1500 noisy copies of $S(Q)$ are transformed; for $G$, $g$, RDF, $T$, $N$ and both windows the analytic error agrees with the Monte Carlo standard deviation within 8\,\% at every $r$ and within 2\,\% on average (the statistical precision of the Monte Carlo is about 1.8\,\%). \\
\multicolumn{2}{@{}l}{\textit{TestSaveRSCF}} \\
\tn{file\_\allowbreak{}with\_\allowbreak{}errors} & The 11-column file reloads with \code{np.loadtxt} to the precision written; the header contains the user heading, all 11 column names in order, and the description of the errors. \\
\tn{file\_\allowbreak{}without\_\allowbreak{}errors\_\allowbreak{}keeps\_\allowbreak{}six\_\allowbreak{}columns} & Without errors the file has 6 columns. \\
\tn{wrong\_\allowbreak{}shape\_\allowbreak{}raises} & An array with 5 columns raises. \\
\bottomrule
\end{longtable}

\tssub*{\texttt{test\_line\_shapes\_ops.py}}
\begin{longtable}{@{}p{0.36\textwidth} p{0.58\textwidth}@{}}
\toprule Test & What is checked \\ \midrule\endhead
\multicolumn{2}{@{}l}{\textit{TestLineShapeErrors}} (errors compared with central finite differences, $10^{-5}$) \\
\tn{lorentzian\_\allowbreak{}error} & Value equals \code{Lorentzian}; error equals the propagated error (C11). \\
\tn{gaussian\_\allowbreak{}error} & Same for \code{Gaussian_error} (correct in the original code). \\
\tn{lorgau\_\allowbreak{}error} & Same for \code{LorGau_error} versus \code{LorGau} (C12). \\
\multicolumn{2}{@{}l}{\textit{TestWsum2}} \\
\tn{general\_\allowbreak{}weights} & $w_1 y_1 + w_2 y_2$ with $\sqrt{w_1^2e_1^2+w_2^2e_2^2}$. \\
\tn{negative\_\allowbreak{}single\_\allowbreak{}weight\_\allowbreak{}gives\_\allowbreak{}positive\_\allowbreak{}error} & $w_1=0$, $w_2=-1$: error $e_2$ (C13). \\
\tn{inverse\_\allowbreak{}variance\_\allowbreak{}weighting} & $w_1=w_2=0$: weights $e_2^2/(e_1^2+e_2^2)$, error $e_1e_2/\sqrt{e_1^2+e_2^2}$. \\
\tn{invalid\_\allowbreak{}mixing\_\allowbreak{}weight\_\allowbreak{}raises} & $w_2=0$, $w_1=2$ raises (C13). \\
\multicolumn{2}{@{}l}{\textit{TestFittingRanges}} \\
\tn{three\_\allowbreak{}variants\_\allowbreak{}agree} & \code{fittingRange0}, \code{fittingRange}, \code{fit_range} select the same points as a direct mask (P20). \\
\bottomrule
\end{longtable}

\tssub*{\texttt{test\_physical\_choices.py}}
\begin{longtable}{@{}p{0.36\textwidth} p{0.58\textwidth}@{}}
\toprule Test & What is checked \\ \midrule\endhead
\multicolumn{2}{@{}l}{\textit{TestMeasurement}} (a complete \code{Measurement} of water built from a minimal parameter set) \\
\tn{annulus\_\allowbreak{}volume} & 4/6\,mm annulus, 50\,mm: $\pi h(r_o^2-r_i^2)$ (C14). \\
\tn{solid\_\allowbreak{}cylinder\_\allowbreak{}volume\_\allowbreak{}unchanged} & Solid cylinder: $\pi h r_o^2$. \\
\tn{free\_\allowbreak{}cross\_\allowbreak{}section\_\allowbreak{}per\_\allowbreak{}element} & $\sigma_\mathrm{free}=\sum c_i\sigma_i[A_i/(A_i+1)]^2$; for water it is less than half the old mean-mass value (P9). \\
\multicolumn{2}{@{}l}{\textit{TestSetSampleInfo}} \\
\tn{free\_\allowbreak{}cross\_\allowbreak{}sections\_\allowbreak{}per\_\allowbreak{}element} & Same formula for the coherent, incoherent and total free cross sections of \code{setSampleInfo}. \\
\multicolumn{2}{@{}l}{\textit{TestSigmoid}} \\
\tn{limits\_\allowbreak{}midpoint\_\allowbreak{}monotonic\_\allowbreak{}no\_\allowbreak{}overflow} & Limits lowQ and highQ, value $(\mathrm{lowQ}+\mathrm{highQ})/2$ at $Q_0$, monotonic, bounded, no overflow up to $Q=10^4$ (C26). \\
\tn{siglin\_\allowbreak{}adds\_\allowbreak{}linear\_\allowbreak{}term} & \code{siglin} $=$ \code{sigmoidal} $+\,a_1Q$. \\
\tn{inelastic\_\allowbreak{}is\_\allowbreak{}the\_\allowbreak{}same\_\allowbreak{}logistic} & \code{inelastic} equals \code{sigmoidal} with highQ $=$ lowQ$\,[A/(A+1)]^2$. \\
\tn{single\_\allowbreak{}q0\_\allowbreak{}default} & The four functions share $Q_0=7.0$ (P21). \\
\multicolumn{2}{@{}l}{\textit{TestWeightedNiFit}} (fake lmfit recording the call) \\
\tn{weights\_\allowbreak{}are\_\allowbreak{}inverse\_\allowbreak{}errors} & The fit receives weights $1/\sigma$ (P11). \\
\tn{no\_\allowbreak{}valid\_\allowbreak{}errors\_\allowbreak{}means\_\allowbreak{}unweighted} & With $\sigma=-1$ no weights are passed. \\
\bottomrule
\end{longtable}

\tssub*{\texttt{test\_misc.py}}
\begin{longtable}{@{}p{0.36\textwidth} p{0.58\textwidth}@{}}
\toprule Test & What is checked \\ \midrule\endhead
\multicolumn{2}{@{}l}{\textit{TestIsotopes}} \\
\tn{bound\_\allowbreak{}cross\_\allowbreak{}sections\_\allowbreak{}in\_\allowbreak{}barns} & Ni: \code{scoh_bound} $=4\pi(10.3)^2/100$ b, equal to the tabulated $\sigma_\mathrm{coh}$ (C15). \\
\tn{deprecated\_\allowbreak{}spanish\_\allowbreak{}names} & \code{Elemento}, \code{elemento} still return \code{Element}, with a warning. \\
\tn{unknown\_\allowbreak{}isotope\_\allowbreak{}raises} & An unknown name raises (P15). \\
\tn{vanadium\_\allowbreak{}b\_\allowbreak{}consistent\_\allowbreak{}with\_\allowbreak{}sigma} & V: $b=-0.443$\,fm and $4\pi b^2/100=\sigma_\mathrm{coh}$; $\sigma_\mathrm{sca}=\sigma_\mathrm{coh}+\sigma_\mathrm{inc}$ (C16). \\
\tn{vanadium\_\allowbreak{}standard\_\allowbreak{}defaults\_\allowbreak{}match\_\allowbreak{}database} & \code{setVanaInfo} defaults equal the table values. \\
\multicolumn{2}{@{}l}{\textit{TestDataXYE}} \\
\tn{file\_\allowbreak{}without\_\allowbreak{}points\_\allowbreak{}in\_\allowbreak{}peak\_\allowbreak{}range} & A file with $x<4$ loads; no peaks found (C22). \\
\tn{plot\_\allowbreak{}saved\_\allowbreak{}before\_\allowbreak{}show} & The saved figure contains the curve even when \code{show()} closes the figure (C22). \\
\multicolumn{2}{@{}l}{\textit{TestNiCalibration}} \\
\tn{initial\_\allowbreak{}amplitudes\_\allowbreak{}mapping} & Each A$_i$, G$_i$, S$_i$ starts from its own entry of \code{initial}; lattice 3.52 by default, or the value given (C10, P7). \\
\tn{real\_\allowbreak{}fit\_\allowbreak{}lattice\_\allowbreak{}fixed\_\allowbreak{}and\_\allowbreak{}wavelength\_\allowbreak{}scales} & A real lmfit fit of a synthetic Ni pattern ($\lambda=0.4986$\,\AA) recovers $\lambda$ to $10^{-6}$ with the lattice fixed at 3.52\,\AA; with 3.524\,\AA{} the fitted $\lambda$ is multiplied by $3.524/3.52$. \\
\multicolumn{2}{@{}l}{\textit{TestMeasurementShowTable}} \\
\tn{stdout\_\allowbreak{}restored\_\allowbreak{}and\_\allowbreak{}file\_\allowbreak{}closed} & \code{showTable} writes its file and restores the previous output stream (P13). \\
\multicolumn{2}{@{}l}{\textit{TestProjectLayout}} \\
\tn{no\_\allowbreak{}module\_\allowbreak{}shadows\_\allowbreak{}the\_\allowbreak{}standard\_\allowbreak{}library} & No file or folder name is in \code{sys.stdlib_module_names} (C21). \\
\tn{stdlib\_\allowbreak{}importable\_\allowbreak{}from\_\allowbreak{}project\_\allowbreak{}root} & \code{import math, io, lmfit, numpy} works with the project directory as current directory. \\
\tn{old\_\allowbreak{}package\_\allowbreak{}names\_\allowbreak{}still\_\allowbreak{}work\_\allowbreak{}with\_\allowbreak{}warning} & \code{toscana.math}/\code{toscana.io} return the same modules as \code{numerics}/\code{dataio}, with warnings. \\
\multicolumn{2}{@{}l}{\textit{TestSourceHygiene}} \\
\tn{no\_\allowbreak{}mojibake} & No doubly-encoded characters in any source file (C23). \\
\bottomrule
\end{longtable}

\tssec{Adding a test}

To protect a new function or a new correction, add a method whose name
starts with \code{test_} to a class deriving from \code{unittest.TestCase}
in one of the \code{test_*.py} files (or in a new file of that name). A
minimal example:
\begin{lstlisting}
import unittest
import helpers                      # makes `toscana` importable
import numpy as np
from toscana.physics.conversions import ang2q

class TestMyCase(unittest.TestCase):
    def test_backscattering(self):
        # 2theta = 180 deg -> Q = 4 pi / lambda
        self.assertAlmostEqual(float(ang2q(180.0, 0.5)), 8 * np.pi)
\end{lstlisting}
Useful comparisons: \code{assertAlmostEqual(a, b, places=6)} or
\code{delta=...}; \code{np.testing.assert_allclose(x, y, rtol=...)} for
arrays; \code{with self.assertRaises(ValueError): ...} to require an error.
Wrap calls that print with \code{helpers.quiet}, and calls that write files
in \code{with helpers.in_tempdir():}.
% =====================================================================
\tsapp{Module reference}\label{app:reference}
% =====================================================================

\tssec{Overview of the package}

\begin{lstlisting}
toscana/                       (the project directory, imported as `toscana`)
  __init__.py      imports the subpackages; installs the deprecated aliases
  _aliases.py      toscana.math -> numerics, toscana.io -> dataio (warnings)
  calibration/     Ni powder calibration of wavelength and zero angle
  chemistry/       compatibility re-exports (composition helpers)
  data/            ASCII x-y-e data containers with plotting
  dataio/          raw data, corrections, regrouping, parameter files, output
  experiment/      experiment, beam, container, vanadium, sample bookkeeping
  isotopes/        nuclear data table and the Element class
  models/          model functions: vanadium, self-scattering, cross sections, Ni
  numerics/        binning, Fourier transforms, windows, fitting, line shapes
  physics/         unit conversions, crystallography, densities, cross sections
  plotting/        matplotlib helpers
  utils/           dates, dictionaries, subprocess, interactive stop
  workflows/       the d4creg regrouping workflow
  tests/           the test suite (Appendix B)
  docs/            this document
\end{lstlisting}

\paragraph{Dependencies between modules.} Low-level packages
(\code{numerics}, \code{physics}, \code{utils}, \code{isotopes}) do not import
higher-level ones. \code{dataio} uses \code{physics}; \code{workflows} uses
\code{dataio}, \code{numerics}, \code{plotting}; \code{calibration} uses
\code{data}, \code{numerics}, \code{models}, \code{physics}; \code{experiment}
uses \code{dataio}, \code{isotopes}, \code{physics}. One cycle exists:
\code{dataio.parameters} imports \code{workflows.d4creg} (so that the
\code{<run>} line can launch it), which imports \code{dataio}; it works
because \code{dataio/__init__.py} imports \code{parameters} last.

\paragraph{External libraries.} NumPy and SciPy are required.
matplotlib (plots), lmfit (Ni fit) and h5py (NeXus files) are imported only
inside the functions that need them.

\paragraph{Reading the entries.} Each function is given with its full
signature and default values. \fixed{C$n$} marks functions corrected during
the audit (Appendix~\ref{app:audit}), \legacy{} functions kept as in the
original code. Arrays called ``3-column'' have columns $(x, y, e)$.
``Logs'' means that the function appends to \code{../logfiles/d4creg.log},
a path relative to the current working directory.

% ---------------------------------------------------------------------
\tssec{Package \texttt{dataio}: raw data, corrections and files}

\tssub{\texttt{dataio.acquisition} --- reading numors}

\fnx{printRates(\allowbreak{}detector,\allowbreak\ counts,\allowbreak\ time,\allowbreak\ reactor\_power=52)}
Logs the counts and the counting rate of \code{detector} (a label such as
\code{"monitor"}), and the rate per MW of reactor power. A non-positive
time is replaced by $10^{-6}$\,s, a non-positive power by 52\,MW. Prints
nothing. \legacy

\fnx{getRefLines(\allowbreak{}Lines)}
Locates the marker lines of an ILL ASCII numor (lines starting with ten
\code{R}, \code{A}, \code{F} or \code{S}) and returns the tuple
\code{(lineR, lineAL, lineAI, lineFD, lineFS, lineFM, lineSI)}: run block,
first two \code{A} blocks, three \code{F} (float parameter) blocks, first
\code{S} (spectrum) block. Raises \code{ValueError} if a marker is missing.

\fnx{readD4numorASCII(\allowbreak{}filename)}
Reads one numor in ILL ASCII format. Returns \code{(data, counts)}:
\code{data} is a dictionary of metadata (numor, label, user, local contact,
title, subtitle, proposal, start and end times, and every parameter of the
three float blocks under its label, e.g.\ \code{"CntTime (sec)"},
\code{"MonitorCnts"}, \code{"TotalCnts"}, \code{"Theta_lue(deg)"},
\code{"Theta_des(deg)"}, plus \code{ndet}, \code{ncell}); \code{counts} is
an array (ndet, ncell). The counts are read 10 per line, in blocks of 11
lines per detector. Prints a summary and logs the rates. \legacy

\fnx{load\_nxs(\allowbreak{}nxs\_path,\allowbreak\ ncell=64)}
Reads one numor in NeXus format with h5py (group \code{entry0}). Returns
\code{(metadata, counts)} with the same keys as the ASCII reader where they
exist: monitor (\code{monitor/monsum}), time, angles
(\code{instrument/masked2theta_s} value and target), scanned variables
(\code{TotalAcquisitionSpy}/\code{TotalCount} renamed \code{TotalCnts}),
and counts reshaped to (9, ncell). The reactor power is set to 1\,MW
(open item P8). \legacy

\fnx{getOneNumor(\allowbreak{}numorfile,\allowbreak\ runInfo)} \fixed{C5}
Reads and corrects one numor; this is the core of the reduction. Steps:
\begin{enumerate}
  \item read with \code{load_nxs} (if \code{runInfo["dataFormat"]=="Nexus"},
    adding the extension \code{runInfo["ext"][4]}) or \code{readD4numorASCII};
  \item compare the read and requested angles: more than
    \code{cellTolerance[0]} cell width gives a warning, more than
    \code{cellTolerance[1]} raises \code{ValueError} (cell width =
    \code{angular_range_bank}/\code{ncell});
  \item Poisson errors $\sqrt{N}$ on the raw counts;
  \item dead-time correction of monitor and detector
    (\code{runInfo["dead_time"]}); errors scaled by the same factor;
  \item division of counts and errors by the efficiencies (\code{efffile});
  \item normalisation to monitor or time (\code{runInfo["normalisation"]});
  \item angles of every cell from \code{getAngles}, with the bank shifts of
    \code{decfile}.
\end{enumerate}
Returns lists \code{x} (angles), \code{y} (counts), \code{e} (errors),
array \code{q}, list \code{c} (cell identifiers) and the metadata. If
\code{runInfo["writeNumor"]==1} the numor is also written by
\code{saveOneNumor}. Note: \code{q} is computed before the zero angle is
subtracted (open item, Section~\ref{sec:open}). Logs.

\fnx{getNumors(\allowbreak{}runInfo)}
Expands \code{runInfo["numorLst"]} with \code{getNumorFiles}, reads every
numor from \code{runInfo["path_raw"]} with \code{getOneNumor}, and stacks
them in an array of shape (n\_numors, 5, n\_points), the five rows being:
angle $-$ \code{twotheta0}, counts, errors, $Q$, cell identifier. Returns
\code{(numor, head)} with \code{head} the list of metadata dictionaries.
If there is only one numor, \code{writeNumor} is forced to 1. Logs.

\tssub{\texttt{dataio.corrections} --- detector corrections}

\fnx{getDTC\_mcounts(\allowbreak{}monitor,\allowbreak\ time,\allowbreak\ dead\_time=2.4e-06)}
Dead-time correction of the monitor (non-paralysable):
$\nu = \nu_m/(1-\nu_m\tau)$, $\nu_m = $ monitor$/$time. Returns the
corrected counts; raises \code{ValueError} if $\nu_m\tau\ge1$. Logs the factor.

\fnx{getDTC\_dcounts(\allowbreak{}counts,\allowbreak\ total\_counts,\allowbreak\ time,\allowbreak\ dead\_time=7e-06)} \fixed{C4}
Same correction cell by cell for the detector array. Returns the corrected
array and the total counts scaled by $\sum N_\mathrm{corr}/\sum N_\mathrm{raw}$.
Raises \code{ValueError} if any factor is not positive. Logs the overall
and maximum factor. (Whether the dead time should apply per cell or per
bank is open item P2.)

\fnx{getEff(\allowbreak{}filename=\textquotesingle{}effd4c.eff\textquotesingle{},\allowbreak\ ndet=9,\allowbreak\ ncells=64)} \fixed{C27}
Reads relative efficiencies, one line \code{det cell efficiency} per cell
(1-based); lines containing \texttt{\#} are ignored. Values $\le0$ become
\code{NaN}, which removes the cell from the regrouping. Cells not listed
keep efficiency 1. \code{filename="ones"} returns all ones.

\fnx{getDec(\allowbreak{}filename=\textquotesingle{}dec.dec\textquotesingle{},\allowbreak\ ndet=9)} \fixed{C27}
Reads the angular shift of each bank, lines \code{det shift} (degrees).
\code{"zeros"} returns zeros; \code{"manual"} returns a built-in list.

\fnx{getErrorsNumor(\allowbreak{}counts)}
Poisson errors $\sqrt{N}$. Must be called on raw counts.

\fnx{getNormalFactors(\allowbreak{}monitor,\allowbreak\ time)}
Returns the 2$\times$2 array $[[M, \sqrt M],[t, 0.01]]$: normaliser value
and its absolute error for monitor and time.

\fnx{normalise(\allowbreak{}counts,\allowbreak\ errors,\allowbreak\ normal,\allowbreak\ norm\_type=\textquotesingle{}monitor\textquotesingle{},\allowbreak\ norm\_time=120,\allowbreak\ norm\_mon=1000000)} \fixed{C2}
Scales counts to \code{norm_mon} monitor counts (or \code{norm_time}
seconds if \code{norm_type="time"}): $y_n = f y$, $f=$ norm$/M$, and
$\sigma_n^2 = (f\sigma)^2 + (y_n\sigma_M/M)^2$. Zero-count cells become
\code{NaN}. Returns \code{(counts, errors)}.

\fnx{getAngles(\allowbreak{}counts,\allowbreak\ ttheta\_ref,\allowbreak\ zeros,\allowbreak\ cell\_step=0.125,\allowbreak\ det\_step=15.0)}
Angle of every cell: $2\theta = 2\theta_\mathrm{ref} + \mathrm{zeros}[d] + 15^\circ d + 0.125^\circ c$
($d$, $c$ 0-based); cell identifier $1000(d+1) + c + 1$. Returns two
arrays of the shape of \code{counts}.

\tssub{\texttt{dataio.d4creg\_outputs} --- regrouping and D4 files}

\fnx{getDiffA(\allowbreak{}numor,\allowbreak\ head,\allowbreak\ runInfo)} \fixed{C6, C7, C29}
Regroups all cells of all numors onto the regular angular grid
\code{np.arange(amin, amax, step)} of \code{runInfo["angular_scale"]}. If
the requested step is smaller than a cell
(\code{angular_range_bank}/\code{ncell}), the cell width is used. Each cell,
of the cell width and centred at its angle, is split between the (at
most two) bins it overlaps with fractions $f$ (exact only when the step
equals the cell width, the default; for a larger step the fractions are
computed in units of the bin, not of the cell: with $0.25^\circ$ a cell
overlapping two bins equally is split 0.25/0.75; open item); its weight is $w = f\times
M_i/\sum M$ (monitor share of its numor). For each bin
$y = \sum w y_i/\sum w$ and $\sigma = \sqrt{\sum w^2\sigma_i^2}/\sum w$.
Cells with \code{NaN} counts or errors, and parts of cells outside the grid,
are ignored; empty bins are \code{NaN}. Returns an array (10, 5, n\_angles):
index 0 is the sum of all banks, 1--9 the individual banks; rows are angle,
counts, error, sum of weights, $Q$. The angle written is the lower edge of
the bin (open item P1).

\fnx{saveDiffAngle(\allowbreak{}diffA,\allowbreak\ head,\allowbreak\ runInfo)}
Writes two files: \code{<file>.adat} (CORRECT format: header block, then
angle, counts, error of the total diffractogram, positive finite points
only) and \code{<file>.reg} (D4 format: header with the equivalent d4creg
command line, efficiency and shift files, normalisation, zero angle,
wavelength, dead times; then one block per bank with angle, counts, error,
$Q$). Extensions from \code{runInfo["ext"]}.

\fnx{saveDiffQ(\allowbreak{}diffQ,\allowbreak\ head,\allowbreak\ runInfo)}
Writes \code{<file>.qdat} (CORRECT format) from a 3-column $(Q, y, e)$
array, positive finite points only.

\tssub{\texttt{dataio.parameters} --- the d4creg parameter file}

\fnx{readParam(\allowbreak{}parametersFile)} \fixed{C25}
Reads a d4creg parameter file and returns the dictionary \code{runInfo}
used by the whole chain. Every non-empty line must start with \texttt{\#} or
\code{!} (comments) or with a five-character tag \code{<xxx>} followed by
values; text after \texttt{\#} on a tag line is ignored. Recognised tags and
defaults:

\begin{longtable}{@{}l l p{0.46\textwidth}@{}}
\toprule Tag & Values & Meaning (default) \\ \midrule\endhead
\code{<tol>} & w e & angle tolerance in cells for warning / error (0.1, 1.0) \\
\code{<tau>} & $\tau_M$ $\tau_D$ & dead times in s ($2.4\times10^{-6}$, $7\times10^{-6}$) \\
\code{<ext>} & 6 strings & extensions (\code{.reg .adat .qdat .cdat .nxs .log}) \\
\code{<eff>} & file & efficiency file (\code{effd4c.eff}; \code{ones}) \\
\code{<dec>} & file & bank shifts (\code{dec.dec}; \code{zeros}, \code{manual}) \\
\code{<log>} & name & log file name (\code{d4creg.log}) \\
\code{<fmt>} & \code{Nexus}/other & raw data format (\code{Nexus}) \\
\code{<asc>} & min max step & angular grid, degrees (0, 140, 0.125) \\
\code{<qsc>} & min max step & $Q$ grid, \AA$^{-1}$ (0, 23.8, 0.02); prints a recommended step \\
\code{<wri>} & \code{True}/\code{False} & write one file per numor (False) \\
\code{<plo>} & \code{True}/\code{False} & plot diffractograms (False) \\
\code{<wle>} & $\lambda$ & wavelength, \AA{} (0.49857) \\
\code{<zac>} & $2\theta_0$ & zero angle, degrees ($-0.1013$) \\
\code{<rdp>} & path & raw data directory (\code{rawdata/}) \\
\code{<nor>} & type t M & \code{monitor} or \code{time}, normalisation time and monitor (monitor, 80, $10^6$) \\
\code{<out>} & name & base name of output files (\code{vanadium}) \\
\code{<num>} & list & numors: \code{123456} or \code{123456-123460} (replaces the list) \\
\code{<add>} & list & appends to the list of numors \\
\code{<run>} & & runs \code{d4creg(runInfo)} with the values read so far \\
\bottomrule
\end{longtable}
\noindent Side effect: (re)creates \code{../logfiles/d4creg.log}. The file
is read sequentially, so several \code{<run>} lines with changes in between
perform several reductions.

\tssub{\texttt{dataio.running\_params} --- the analysis parameter file}

\fnx{getRunningParams(\allowbreak{}parfile)}
Reads the parameter file of the data-treatment notebooks and returns a
dictionary \code{inputPar} keyed by tag; each value is a tuple whose first
element is a label and the following ones the values. Tags: experiment
(\code{<ins> <pro> <mpr> <lco> <dte> <cyc> <log> <env>}; \code{<dte>} gives
\code{<sta>}/\code{<end>}, \code{<log>} gives \code{<nbk>}/\code{<pag>}), beam
(\code{<zac> <wle> <sli> <fla> <apr>}; \code{<sli>} gives \code{<lft>}/\code{<rgt>},
\code{<fla>} gives \code{<top>}/\code{<bot>}), grids (\code{<asc> <qsc> <rsc>}),
container, absorber, normaliser, environment (\code{<Cmater> <Cshape>},
\code{<Amater> <Ashape>}, \code{<Nmater> <Nshape>}, \code{<Emater> <Eshape>};
shapes: type, inner diameter, outer diameter, height), sample
(\code{<Sdescr> <StempK> <Sato01>... <Sdegcc> <Smassc> <Sshape> <Sfulln> <Stitle>};
each \code{<SatoNN>} is \code{symbol number}, counted into \code{<SNelem>}),
and data files (\code{<CoFile> <AbFile> <NoFile> <EnFile> <SaFile>}). Unknown
materials or shapes print a warning. Its output is the input of
\code{experiment.Measurement}. \legacy

\tssub{\texttt{dataio.loading}, \texttt{saving}, \texttt{numors}, \texttt{correct}, \texttt{\_logfiles}}

\fnx{read\_xye(\allowbreak{}filename)} \fixed{C27}
Reads 2 or 3 columns, ignoring lines containing \texttt{\#}; a missing or
non-numeric (\code{i}, \code{o}) error becomes 0. Returns three 1-D arrays.

\fnx{read\_3col(\allowbreak{}filename)}
Same, returned as one (n, 3) array.

\fnx{saveFile\_xye(\allowbreak{}filename,\allowbreak\ x,\allowbreak\ y,\allowbreak\ e,\allowbreak\ heading)}, \code{saveFile_3col(filename, data, heading)}
Write a 3-column file with \texttt{\#} header lines (formats \code{9.3f},
\code{18.6f}, \code{18.6f}).

\fnx{saveCorrelations(\allowbreak{}filename,\allowbreak\ ar,\allowbreak\ pcf,\allowbreak\ pdf,\allowbreak\ rdf,\allowbreak\ tor,\allowbreak\ run,\allowbreak\ heading)}
Writes six columns $r, G, g$, RDF, $T, N$ from separate arrays.

\fnx{saveRSCF(\allowbreak{}filename,\allowbreak\ fou,\allowbreak\ heading)} \textsf{\footnotesize\colorbox{green!15}{new header, 11 columns}}
Writes the output of \code{sineFT}: 6 columns, or 11 with errors. After
the file name and the user \code{heading}, a block describes each column
(name, formula, units), the meaning and limits of the errors, and a line of
column names \code{r G dG g dg RDF dRDF T dT N dN}. Other shapes raise.

\fnx{saveOneNumor(\allowbreak{}x,\allowbreak\ y,\allowbreak\ e,\allowbreak\ q,\allowbreak\ c,\allowbreak\ head,\allowbreak\ runInfo)}
Writes \code{<numor><ext[3]>} (angle, counts, error, $Q$, cell) for one numor.

\fnx{getNumorFiles(\allowbreak{}numorLst)}
Expands entries \code{"123456"} and \code{"123456-123460"} into a list of
six-digit numors; malformed or reversed ranges raise.

\fnx{generate\_com(\allowbreak{}CORRECT\_template=\textquotesingle{}correct\_template.com\textquotesingle{},\allowbreak\ name=\textquotesingle{}prueba\textquotesingle{},\allowbreak\ rsample=0.5,\allowbreak\ density=0.11,\allowbreak\ ...)}
Fills the placeholders \texttt{\{pyfname\}}, \texttt{\{pyradius\_cm\}}, \ldots{} of a
CORRECT template and writes \code{<name>.com}. It then runs a placeholder
command (open item). \legacy

\fnx{saveCORRECT(\allowbreak{}expt)}
Writes a CORRECT command file \code{<sampleTitle>.com} from an experiment
object (sample, container, background, vanadium, wavelength, beam size;
zero angle written as 0 since already subtracted). \legacy

\fnx{d4creg\_log\_path(\allowbreak{})}
Returns \code{../logfiles/d4creg.log}, creating the folder and file if needed.

% ---------------------------------------------------------------------
\tssec{Package \texttt{workflows}}

\fnx{d4creg(\allowbreak{}runInfo)} \fixed{C8, C9}
The complete regrouping: \code{getNumors} $\to$ \code{getDiffA} $\to$
\code{saveDiffAngle} (+ plots if \code{plotDiff}) $\to$ conversion of the
regrouped angles to $Q$ $\to$ \code{rebin} onto \code{q_scale} with the
actual angular step as width $\to$ masking of points with $y\le0$ and of
their predecessor $\to$ \code{saveDiffQ} (+ plot) $\to$ renaming of the log
file to \code{../logfiles/<logfile>.log}. Prints the equivalent d4creg
command line. Adds \code{runDate}, \code{nbrNumors}, \code{d4creg_cl} to
\code{runInfo}.

% ---------------------------------------------------------------------
\tssec{Package \texttt{numerics}: numerical methods}

\tssub{\texttt{numerics.fourier} --- sine Fourier transforms}

\fnx{sineFT(\allowbreak{}StrFactor,\allowbreak\ nr,\allowbreak\ qmax=23.5,\allowbreak\ density=0.1,\allowbreak\ constant=1.0,\allowbreak\ selfsca=1.0,\allowbreak\ window=0,\allowbreak\ with\_errors=False)} \fixed{C1, C18} \newbeh
The main transform, Section~\ref{sec:conventions}. \code{StrFactor} is an
array with columns $Q$, $S$ (and $\sigma_S$); \code{nr} the $r$ grid;
\code{density} $\rho$ in atoms/\AA$^3$; \code{constant} a multiplicative
factor; \code{selfsca} the value subtracted from $S$ (1 for a normalised
$S(Q)$); \code{window} 0 (none) or 1 (standard Lorch). Returns an
$(N_r, 6)$ array $r, G, g$, RDF, $T, N$, or with \code{with_errors=True}
an $(N_r, 11)$ array with each error beside its function
(Section~\ref{sec:errors}). The $Q$ integral runs over the whole $Q$ grid
with $M(Q)=0$ for $Q\ge$\code{qmax}, so the zeroed points still take part
in the Simpson panels next to $Q_\mathrm{max}$ (\code{getSineFT} instead
cuts the grid at \code{vmax}); the running integral starts at \code{nr[0]}.

\fnx{getSineFT(\allowbreak{}vvec,\allowbreak\ yvec,\allowbreak\ uvec,\allowbreak\ vmax=0,\allowbreak\ c=1.0,\allowbreak\ s=1.0,\allowbreak\ w=0)} \fixed{P4}
Generic transform
$\mathrm{stf}(u) = \frac2\pi c\int_0^{v_\mathrm{max}} [y(v)-s]\,v\,\sin(vu)\,W(v)\,\mathrm{d}v$,
$W$ = 1 (\code{w=0}) or Lorch (\code{w=1}). \code{vmax}$\le0$ means the
end of the data. Returns \code{None} if the lengths differ. Used by
\code{backFT} in both directions.

\fnx{backFT(\allowbreak{}qu,\allowbreak\ ar,\allowbreak\ pdf,\allowbreak\ density,\allowbreak\ cut)} \fixed{C17}
Cuts $g(r)$ (\code{pdf} on grid \code{ar}) according to \code{cut}
(Section~\ref{sec:fourier}), transforms to $S(Q)$ on \code{qu} and back.
Returns the six arrays $(S, G, g, \mathrm{RDF}, T, N)$. Typical use: remove
the unphysical oscillations of $g(r)$ below the first peak and see their
effect on $S(Q)$.

\fnx{getSineIntegral(\allowbreak{}x)} \fixed{C20}
$\mathrm{Si}(x)=\int_0^x \sin t/t\,\mathrm{d}t$ (\code{scipy.special.sici}).

\noindent Internal helpers: \code{_simpson_weights(x)} (the weights $w_k$ of
SciPy's Simpson rule), \code{_sine_ft_errors(...)} (the error
propagation), \code{_running_simpson_matrix(x)} (matrix $C$ of the running
integral, with self-check), \code{_apply_cut(ar, pdf, cut)}.

\tssub{\texttt{numerics.windows} --- modification functions}

\fnx{Lorch(\allowbreak{}q,\allowbreak\ qmax=23.5)} \fixed{P5}
Standard Lorch value at one point: $\sin(\pi q/q_\mathrm{max})/(\pi q/q_\mathrm{max})$,
1 at 0, 0 outside $[0, q_\mathrm{max})$. Raises if $q_\mathrm{max}\le0$.

\fnx{LorchArray(\allowbreak{}q,\allowbreak\ qmax)} \textsf{\footnotesize\colorbox{green!15}{new}}
The same for an array (replaces the removed, rescaled \code{LorchN}).

\tssub{\texttt{numerics.binning} --- rebinning}

\fnx{get\_xlim(\allowbreak{}xmin,\allowbreak\ xmax,\allowbreak\ dbin)} \fixed{C19}
Bin edges and centres: centres $x_\mathrm{min} + i\,\delta$, edges at
$\pm\delta/2$; returns \code{(nb, x_lim, x_bin)}.

\fnx{get\_bins(\allowbreak{}x,\allowbreak\ xdel,\allowbreak\ x\_lim)}
Bins covered by a rectangle of width \code{xdel} centred at \code{x}, and
the fraction of the rectangle in each. A rectangle that starts below the
first edge is rejected entirely (empty lists). At the upper end the
treatment differs: a rectangle is rejected only if it ends one full bin
beyond the last edge; otherwise its excess goes to an overflow index that
\code{rebin} discards, so a point lying entirely in the bin-width beyond
the last edge is dropped silently and a straddling point contributes only
its overlap.

\fnx{rebin(\allowbreak{}xdel,\allowbreak\ wlength,\allowbreak\ data,\allowbreak\ xmin,\allowbreak\ xmax,\allowbreak\ dbin)} \fixed{C29}
Rebins a 3-column array onto the grid of \code{get_xlim}. Each point is a
rectangle of width \code{xdel} (in the units of $x$ if \code{wlength}$\le0$;
if \code{wlength}$>0$, $x$ is $Q$ and \code{xdel} an angular width in degrees,
converted to $\Delta Q = \frac{2\pi}{\lambda}\cos\theta\,\Delta(2\theta)\frac{\pi}{180}$).
Weighted mean per bin, error $\sqrt{\sum f^2\sigma^2}/\sum f$. Points with
\code{NaN} are skipped; empty bins are 0. Prints a debug line.

\tssub{\texttt{numerics.operations} --- combining data sets}

\fnx{binary\_sum(\allowbreak{}w1,\allowbreak\ data1,\allowbreak\ w2,\allowbreak\ data2)} \fixed{C3}
$w_1y_1+w_2y_2$, error $\sqrt{w_1^2e_1^2+w_2^2e_2^2}$, for two 3-column
arrays of the same shape (otherwise prints and returns \code{None}).

\fnx{wsum2(\allowbreak{}w1,\allowbreak\ data1,\allowbreak\ w2,\allowbreak\ data2)} \fixed{C13}
Four modes: $w_1=w_2=0$: inverse-variance weighted mean; $w_1=0$: $w_2y_2$
with error $|w_2|e_2$; $w_2=0$: mixture $w_1y_1+(1-w_1)y_2$, requires
$0<w_1\le1$ (raises otherwise); else $w_1y_1+w_2y_2$.

\fnx{ratio(\allowbreak{}y1,\allowbreak\ e1,\allowbreak\ y2,\allowbreak\ e2)}
$y_1/y_2$ with error $\sqrt{(y_2e_1)^2+(y_1e_2)^2}/y_2^2$; returns $0\pm0$
where $y_2=0$. Lists, truncated to the shortest input. \legacy

\tssub{\texttt{numerics.line\_shapes} --- peak and step functions}

Constants \code{SIGMOID_Q0 = 7.0} and \code{SIGMOID_DQ = 2.4}
(\AA$^{-1}$): default centre and width of the low-$Q$/high-$Q$ transition.

\fnx{sigmoidal(\allowbreak{}Q,\allowbreak\ lowQ=0.4,\allowbreak\ highQ=0.2,\allowbreak\ Q0=SIGMOID\_Q0,\allowbreak\ dQ=SIGMOID\_DQ)} \fixed{C26} \newbeh
Logistic step $(\mathrm{lowQ}+\mathrm{highQ}\,\delta)/(1+\delta)$, $\delta = e^{(Q-Q_0)/\mathrm{d}Q}$.

\fnx{siglin(\allowbreak{}Q,\allowbreak\ lowQ=0.4,\allowbreak\ highQ=0.2,\allowbreak\ Q0=SIGMOID\_Q0,\allowbreak\ dQ=SIGMOID\_DQ,\allowbreak\ a1=0.001)} \newbeh
\code{sigmoidal(Q)} $+\;a_1Q$.

\fnx{step(\allowbreak{}q,\allowbreak\ qmax=23.5)} \fixed{P5}
1 for $0\le q<q_\mathrm{max}$, else 0; raises if $q_\mathrm{max}\le0$.

\fnx{Lorentzian(\allowbreak{}x,\allowbreak\ A,\allowbreak\ x0,\allowbreak\ gamma,\allowbreak\ bckg)}
$A(\gamma/2)^2/[(x-x_0)^2+(\gamma/2)^2] + b$: height $A$, FWHM $\gamma$.

\fnx{Lorentzian\_error(\allowbreak{}x,\allowbreak\ A=1,\allowbreak\ x0=0,\allowbreak\ gamma=1,\allowbreak\ bckg=0,\allowbreak\ e\_A=0,\allowbreak\ e\_x0=0,\allowbreak\ e\_gamma=0,\allowbreak\ e\_bckg=0)} \fixed{C11}
Value and propagated error of \code{Lorentzian}.

\fnx{Gaussian(\allowbreak{}x,\allowbreak\ A,\allowbreak\ x0,\allowbreak\ sigma,\allowbreak\ bckg)}
Area-normalised Gaussian $\frac{A}{\sigma\sqrt{2\pi}}e^{-(x-x_0)^2/2\sigma^2}+b$ (area $A$).

\fnx{Gaussian\_error(\allowbreak{}x,\allowbreak\ A=1,\allowbreak\ x0=0,\allowbreak\ sigma=1,\allowbreak\ bckg=0,\allowbreak\ e\_A=0,\allowbreak\ e\_x0=0,\allowbreak\ e\_sigma=0,\allowbreak\ e\_bckg=0)}
Value and propagated error of \code{Gaussian}.

\fnx{GaussianA(\allowbreak{}x,\allowbreak\ A,\allowbreak\ x0,\allowbreak\ sigma,\allowbreak\ asym=0)}
Area-normalised Gaussian (area $A$, no background) multiplied by
$e^{-\mathrm{asym}(x-x_0)}$ (asymmetric peak, used for the Ni reflections).
For $\mathrm{asym}\neq0$ the area becomes $A\,e^{\mathrm{asym}^2\sigma^2/2}$
and the maximum moves to $x_0-\mathrm{asym}\,\sigma^2$.

\fnx{LorGau(\allowbreak{}x,\allowbreak\ f0=1.0,\allowbreak\ eta=0.5,\allowbreak\ sigma=2.0,\allowbreak\ gamma=2.0,\allowbreak\ bckg=0.0)}
Mixture centred at 0: $f_0[\eta e^{-x^2/2\sigma^2} + (1-\eta)L(x)] + b$,
both components of unit height.

\fnx{LorGau\_error(\allowbreak{}x,\allowbreak\ f0=(\allowbreak{}1.0,\allowbreak\ 0.0),\allowbreak\ eta=(\allowbreak{}0.5,\allowbreak\ 0.0),\allowbreak\ sigma=(\allowbreak{}2.0,\allowbreak\ 0.0),\allowbreak\ gamma=(\allowbreak{}2.0,\allowbreak\ 0.0),\allowbreak\ bckg=(\allowbreak{}0.0,\allowbreak\ 0.0))} \fixed{C12}
Value and error of \code{LorGau}; each argument is a pair (value, error).

\tssub{\texttt{numerics.fitting}, \texttt{polynomials}, \texttt{peak\_finding}, \texttt{signal\_processing}}

\fnx{get\_chi(\allowbreak{}y1,\allowbreak\ y2)}
Root-mean-square difference $\sqrt{\langle(y_1-y_2)^2\rangle}$.

\fnx{fit\_and\_find\_extremum(\allowbreak{}x,\allowbreak\ y)}
Fits a parabola and returns the extremum $x=-b/2a$, its value and the
fitted curve.

\fnx{fittingRange0(\allowbreak{}xmin,\allowbreak\ xmax,\allowbreak\ ymin,\allowbreak\ ymax,\allowbreak\ abcissas,\allowbreak\ ordinates,\allowbreak\ errors)}, \code{fittingRange(...)}, \code{fit_range(xmin, xmax, ymin, ymax, data)} \fixed{P20}
Select the points inside the rectangle $[x_\mathrm{min},x_\mathrm{max}]\times[y_\mathrm{min},y_\mathrm{max}]$.
\code{fittingRange} returns three $(n,1)$ arrays, the other two one
$(n,3)$ array. Warn below 10 points.

\fnx{polyQ1(\allowbreak{}x,\allowbreak\ a0,\allowbreak\ a1)}, \code{polyQ2(x, a0, a1, a2)}, \code{polyQ4(x, a0, ..., a4)}
Polynomials $\sum a_i x^i$ (Horner scheme).

\fnx{find\_minimum\_within\_range(\allowbreak{}x,\allowbreak\ y,\allowbreak\ x\_min,\allowbreak\ x\_max)}
$(x, y)$ of the minimum of $y$ in $[x_\mathrm{min}, x_\mathrm{max}]$ (raises
on an empty range).

\fnx{find\_peaks\_in\_range(\allowbreak{}x,\allowbreak\ y,\allowbreak\ x\_min,\allowbreak\ x\_max)}
Local maxima (\code{scipy.signal.find_peaks}) whose height divided by the
minimum of the range exceeds 3.5; returns positions and these ratios.
\legacy

\fnx{smooth\_curve(\allowbreak{}x,\allowbreak\ y,\allowbreak\ smoothing\_factor)}
Centred moving average over $2\lfloor w/2\rfloor+1$ points,
$w = \lfloor\text{factor}\times n\rfloor$, shrinking at the edges; \code{x}
is not used.

% ---------------------------------------------------------------------
\tssec{Package \texttt{physics}}

\tssub{\texttt{physics.conversions}, \texttt{geometry}}

\fnx{ang2q(\allowbreak{}x,\allowbreak\ wlength=0.5)}, \code{q2ang(x, wlength=0.5)}
$Q = \frac{4\pi}{\lambda}\sin(\pi\,2\theta/360)$ and its inverse; raise if
$\lambda\le0$.

\fnx{getCylVolume(\allowbreak{}diameter=5.0,\allowbreak\ height=50.0)}
$\pi(d/2)^2h$ (mm$^3$ for mm input).

\tssub{\texttt{physics.crystallography}}

Constant \code{NI_LATTICE_PARAMETER = 3.52} (\AA).

\fnx{sf\_fcc(\allowbreak{}h,\allowbreak\ k,\allowbreak\ l)}, \code{sf_bcc(h, k, l)}, \code{sf_sc(h, k, l)}
Geometric structure factors $\sum_j e^{-i\pi(\ldots)}$ of the fcc (4 or 0),
bcc (2 or 0) and simple cubic (1) lattices.

\fnx{reflections\_fcc(\allowbreak{}wavelength=0.5,\allowbreak\ twotheta0=0.0,\allowbreak\ lattice=NI\_LATTICE\_PARAMETER)}
Angles $2\theta+2\theta_0$ (degrees) of the allowed fcc reflections in
increasing order, followed by the mirrored set $180^\circ-2\theta-2\cdot2\theta_0$.

\tssub{\texttt{physics.properties}}

Constants \code{THERMAL_WAVELENGTH = 1.798} (\AA) and
\code{THERMAL_ENERGY_MEV} $=81.8042/1.798^2 = 25.30$ (meV).

\fnx{getAtomicDensity(\allowbreak{}density=1.0,\allowbreak\ molarMass=6.0)}
$\rho = 0.602214076\,d/\bar M$ atoms/\AA$^3$ ($\bar M$ per atom). A
non-positive density gives the value of water (0.1004) with a message.

\fnx{getDensity(\allowbreak{}atomic\_density=0.1,\allowbreak\ molarM=6.0)}
Inverse of the above.

\fnx{getMassNumber(\allowbreak{}molarMass=6.0)}
$A = M/m_n$, $m_n = 1.0086649$\,u.

\fnx{getFreeXS(\allowbreak{}BoundXS=5.08,\allowbreak\ A=50.9415)}
$\sigma_\mathrm{bound}[A/(A+1)]^2$ for one nucleus, $A=m/m_n$. Note: the
default \code{A=50.9415} is the molar mass of vanadium, not its mass ratio
50.50 (0.03\,\% effect on the factor).

\fnx{getAbsXS(\allowbreak{}AbsXS=5.08,\allowbreak\ wavelength=THERMAL\_WAVELENGTH)} \fixed{C30}
$\sigma_\mathrm{abs}\lambda/1.798$.

\fnx{getNofAtoms(\allowbreak{}atoms)}, \code{getConcentrations(atoms)}, \code{AtomicAvg(concentration, magnitude)}
Number of atoms of a formula unit given as \texttt{\{"Si": 1, "O": 2\}};
concentrations $c_i = n_i/\sum n$; average $\sum c_i m_i$ of a property
given as a dictionary.

% ---------------------------------------------------------------------
\tssec{Package \texttt{models}}

\fnx{inelastic(\allowbreak{}x,\allowbreak\ A=51.0,\allowbreak\ lowQ=0.4,\allowbreak\ Q0=SIGMOID\_Q0,\allowbreak\ dQ=SIGMOID\_DQ)}
Inelasticity (Placzek-like) decrease of the self scattering:
$\mathrm{lowQ}\,[1+(A/(1+A))^2\delta]/(1+\delta)$, from the bound value at
low $Q$ to the free value at high $Q$.

\fnx{self024(\allowbreak{}x,\allowbreak\ self0,\allowbreak\ q2,\allowbreak\ q4)}
$s_0 + q_2Q^2 + q_4Q^4$, an empirical self-scattering curve to fit.

\fnx{vanaQdep(\allowbreak{}x,\allowbreak\ a0,\allowbreak\ a1,\allowbreak\ a2,\allowbreak\ A=51.0,\allowbreak\ lowQ=0.4,\allowbreak\ Q0=SIGMOID\_Q0,\allowbreak\ dQ=SIGMOID\_DQ)}
Vanadium $Q$ dependence: \code{polyQ2} (resolution, attenuation trends)
$\times$ \code{inelastic}.

\fnx{XS\_model(\allowbreak{}Eval,\allowbreak\ Eeff,\allowbreak\ composition\_vec,\allowbreak\ bound\_xs\_vec,\allowbreak\ A\_vec)}
Epithermal total scattering cross section
$\sum c_i\sigma_i[A_i/(A_i+1)]^2[1+E_\mathrm{eff}/(2A_iE)]$.

\fnx{scattering\_probability(\allowbreak{}E0,\allowbreak\ scat\_xs,\allowbreak\ composition\_vec,\allowbreak\ abs\_xs\_vec)}
$1-\sigma_a/(\sigma_a+\sigma_s)$ with $\sigma_a = \sum c_i\sigma_{a,i}/\sqrt{E_0/E_\mathrm{th}}$ ($E_0$ in meV).

\fnx{niPeaks10(\allowbreak{}x,\allowbreak\ I0,\allowbreak\ slope,\allowbreak\ quad,\allowbreak\ wavelength,\allowbreak\ twotheta0,\allowbreak\ A0..A9,\allowbreak\ G0..G9,\allowbreak\ S0..S9,\allowbreak\ lattice=NI\_LATTICE\_PARAMETER)}
Ni powder model: background $I_0 + \mathrm{slope}\,x + \mathrm{quad}\,x^2$
plus ten \code{GaussianA} peaks (areas A$_i$, widths S$_i$, asymmetries
G$_i$) at the first ten fcc reflections for the given wavelength, zero
angle and lattice parameter. $x$ is $2\theta$ in degrees.

% ---------------------------------------------------------------------
\tssec{Package \texttt{calibration}}

\fnx{setting\_model\_d4nifit(\allowbreak{}nifile,\allowbreak\ zac\_ini=0.0,\allowbreak\ lattice=NI\_LATTICE\_PARAMETER)}
Prepares the Ni calibration. Reads the regrouped Ni diffractogram
(\code{DataXYE}), estimates $\lambda$ from the first peak, assumed to be
(111): $\lambda = 2(a/\sqrt3)\sin\theta_{111}$; normalises intensities and
errors to the first data point above the midpoint between the second and
third peaks, truncated to whole degrees; builds the
35 initial parameters (background, $\lambda$, $2\theta_0$, ten areas
proportional to typical Ni intensities, asymmetries 0, widths 0.2--0.4);
plots data, initial model and residuals. Returns
\code{(axes, limits, xy_res, xy_ref, xy_fit, planes, initial, nickel)}.

\fnx{fitting\_model\_d4nifit(\allowbreak{}nickel,\allowbreak\ limits,\allowbreak\ initial,\allowbreak\ lattice=NI\_LATTICE\_PARAMETER)} \fixed{C10, P11}
Fits \code{niPeaks10} with lmfit to the points inside \code{limits}
(rectangle in $2\theta$ and intensity), weighted by $1/\sigma$. Free: I0,
slope, wavelength, twotheta0, A$_i$, S$_i$; fixed: quad, G$_i$, lattice.
Returns the lmfit \code{ModelResult}.

\fnx{showing\_results\_d4nifit(\allowbreak{}result,\allowbreak\ nickel,\allowbreak\ axes,\allowbreak\ limits,\allowbreak\ xy\_res,\allowbreak\ xy\_ref,\allowbreak\ xy\_fit,\allowbreak\ planes)}
Prints $\lambda$, $2\theta_0$ and background with errors, a table of the
ten reflections (position, area, width), plots the fit with the fit
statistics, and saves \code{<basename>.png}.

% ---------------------------------------------------------------------
\tssec{Package \texttt{experiment}}

\tssub{\texttt{experiment.measurement}}

\fnx{class Measurement(\allowbreak{}inputPar)} \fixed{C14, P9, P13}
Collects everything about a measurement from the dictionary returned by
\code{getRunningParams}, and reads the five data files (sample, environment,
container, normaliser, absorber) with \code{read_3col}. Main attributes:
experiment and beam information; \code{vslits}, \code{hslits},
\code{beamHeight}, \code{beamWidth}; \code{aScale}, \code{qScale},
\code{rScale}; \code{container}, \code{normaliser}, \code{absorber},
\code{environ} (material, shape, dimensions); sample \code{Shape},
\code{InnerDiam}, \code{OuterDiam}, \code{Height}, \code{Volume}
$=\pi h(r_o^2-r_i^2)/1000$ cm$^3$, \code{Mass}, \code{Density},
\code{EffDensity} $=$ mass/volume, \code{PackingFraction}; composition
\code{atoms}, \code{natoms}, \code{conc}, \code{nuclei}; averages
\code{molarMass}, \code{scoh}, \code{sinc}, \code{ssca}, \code{sabs}
(barns/atom), \code{bcoh}, \code{binc} (moduli of the averaged lengths,
P10); \code{free} (per-element free cross section); \code{b2lowQ}
$=\sigma_\mathrm{sca}/4\pi$, \code{b2highQ} $=\sigma_\mathrm{free}/4\pi$;
\code{sabswl} (absorption at the wavelength); \code{AtomicDensity},
\code{EffAtomicDensity}; data arrays \code{Data}, \code{envData},
\code{conData}, \code{norData}, \code{absData}. A warning is printed if
the container inner diameter differs from the sample outer diameter.
\code{showTable(outfile=None)} prints the identification, beam and grid
parameters (not the sample, container or composition), to a file if given.
\code{<Nmater>} is not read (the normaliser is always vanadium) and
\code{Fullness} is stored but not used.

\tssub{\texttt{experiment.info}}

Constants \code{VANADIUM_DENSITY_GCC = 6.11}, \code{VANADIUM_MOLAR_MASS = 50.9415},
\code{VANADIUM_ATOMIC_DENSITY} $=0.0722305$ atoms/\AA$^3$.

\fnx{setExpInfo(\allowbreak{}Proposal,\allowbreak\ mainProposer,\allowbreak\ experimenters,\allowbreak\ LC,\allowbreak\ otherUsers,\allowbreak\ startDate,\allowbreak\ endDate,\allowbreak\ environment,\allowbreak\ logBook,\allowbreak\ logPage,\allowbreak\ instrument)}
\fnx{setBeamInfo(\allowbreak{}zeroAngle=0,\allowbreak\ wavelength=0.5,\allowbreak\ LohenSlit=2.5,\allowbreak\ GamsSlit=-2.5,\allowbreak\ topFlag=25.0,\allowbreak\ bottomFlag=-25.0)}
\fnx{setCanInfo(\allowbreak{}material=\textquotesingle{}Vanadium\textquotesingle{},\allowbreak\ shape=\textquotesingle{}Cylinder\textquotesingle{},\allowbreak\ outerDiam=5,\allowbreak\ innerDiam=4.8,\allowbreak\ height=60.0)}
\fnx{setBinInfo(\allowbreak{}AngularResolution=0.125,\allowbreak\ AMin=0.0,\allowbreak\ AMax=140.0,\allowbreak\ AStep=0.125,\allowbreak\ QMin=0.0,\allowbreak\ QMax=23.5,\allowbreak\ QStep=0.02,\allowbreak\ RMin=0.0,\allowbreak\ RMax=20.0,\allowbreak\ RStep=0.01)}
\fnx{setNumorInfo(\allowbreak{}totalNumors,\allowbreak\ containerNumors,\allowbreak\ environmentNumors,\allowbreak\ nickelNumors,\allowbreak\ vanadiumNumors,\allowbreak\ absorberNumors,\allowbreak\ sampleNumors)}
Each builds and prints a dictionary describing one part of the
experiment; every argument has a default (e.g.\ \code{LC="Cuello"},
numor ranges \code{("Exp", 0, 1)}), see the docstrings. The dictionaries describe one part of the
experiment (beam width = left $-$ right slit, height = top $-$ bottom
flag; can wall thickness; number of points of each grid; numor ranges
given as (label, first, last)). \legacy

\fnx{setVanaInfo(\allowbreak{}IncXS=5.08,\allowbreak\ CohXS=0.02466,\allowbreak\ ScaXS=5.10466,\allowbreak\ AbsXS=5.08,\allowbreak\ CohSL=-0.443,\allowbreak\ molarM=50.9415,\allowbreak\ NAtoms=1.0,\allowbreak\ Diam=6.08,\allowbreak\ Height=50.0,\allowbreak\ density=VANADIUM\_DENSITY\_GCC)} \fixed{C16, P6}
Dictionary of the vanadium standard: cross sections, $b$, self scattering
at $Q=0$ ($\sigma_\mathrm{inc}/4\pi$), atomic density, mass number, free
incoherent cross section, volume in cm$^3$.

\fnx{setSampleInfo(\allowbreak{}comp\_sample,\allowbreak\ atoms\_sample,\allowbreak\ natoms\_sample,\allowbreak\ elements,\allowbreak\ c\_sample,\allowbreak\ wavelength=0.5,\allowbreak\ beamH=50.0,\allowbreak\ vanavol=1.0,\allowbreak\ vanadens=VANADIUM\_ATOMIC\_DENSITY,\allowbreak\ height=60.0,\allowbreak\ diameter=5.0,\allowbreak\ mass=1.0,\allowbreak\ density=1.0,\allowbreak\ title=\textquotesingle{}sample\textquotesingle{})} \fixed{C30, P9}
Dictionary of the sample: height in the beam, volumes, effective density,
packing and filling fractions, average molar mass, atomic and effective
atomic densities, averaged cross sections, absorption at $\lambda$,
average $b_\mathrm{coh}$, per-element free cross sections, self scattering
at $Q=0$, and the sample/vanadium volume and density ratios used for
normalisation. \code{elements} maps symbols to \code{Element} objects,
\code{c_sample} symbols to concentrations.

% ---------------------------------------------------------------------
\tssec{Package \texttt{isotopes}}

\fnx{class Element(\allowbreak{}iso\_name,\allowbreak\ table=0,\allowbreak\ *{}*kwargs)} \fixed{C15, P15}
Nuclear data of an element (\code{"Ni"}) or isotope (\code{"Li6"},
\code{"H2"} for deuterium) from \code{database.ISOTOPE_DICT} (values from
\emph{Neutron Scattering Fundamentals}, vol.~44, App.~2). Attributes:
\code{symbol}, \code{iso}, \code{name}, \code{A}, \code{Z}, \code{weight}
(u), \code{abundance}, \code{spin}, \code{re_bcoh}, \code{im_bcoh},
\code{re_binc}, \code{im_binc} (fm), \code{sig_coh}, \code{sig_inc},
\code{sig_sca}, \code{sig_abs} (b), complex \code{bcoh}, \code{binc},
\code{scoh_bound} $=4\pi|b_\mathrm{coh}|^2/100$, \code{sinc_bound} (b),
\code{neut} $=A-Z$. Keyword arguments override table values (with
plausibility warnings). \code{table=1} prints the table of the isotope,
\code{table=2} of all isotopes of the element. Unknown names raise
\code{ValueError}. \code{Element.getIsotope(name)} lists the isotopes of an
element name, or all names with \code{"All"}. \code{Elemento},
\code{elemento}: deprecated aliases.

\fnx{getFree(\allowbreak{}isot,\allowbreak\ key)}
Free value $[A/(A+1)]^2\times$\code{key} (complex) for an \code{Element}.

\noindent Validators \code{Positive}, \code{Negative}, \code{SemiWhole},
\code{Top}, \code{Int}, \code{String}, \code{Complex}: print a warning when
an override is implausible. \code{database.NOMBRES} and \code{UNITS}: names
and units of the table columns.

% ---------------------------------------------------------------------
\tssec{Packages \texttt{data}, \texttt{plotting}, \texttt{utils}, \texttt{chemistry}}

\fnx{class DataXYE(\allowbreak{}filename)} \fixed{C22}
Reads a 2- or 3-column file (header lines start with \texttt{\#} or \code{!};
a missing error is stored as $-1$); attributes \code{x}, \code{y},
\code{e}, \code{head}, averages, minima and maxima, \code{peaks_x},
\code{peaks_y} and the minimum \code{xminr}, \code{yminr} in
$x\in[5,40]$ (empty / \code{NaN} if no point there). Methods
\code{plot(file_format=0, xmin, xmax, ymin, ymax)} (axis label from the
extension; saves to \code{../regdata/<basename>.<format>} if a format is
given), \code{show()}, \code{header(lines=1000)}. \code{DataXYE_old} is an
alias.

\fnx{class DataXYE0(\allowbreak{}filename,\allowbreak\ symbol,\allowbreak\ save,\allowbreak\ scale,\allowbreak\ superplot,\allowbreak\ ...)}, \code{class DataXYE2(file)}
Simpler containers read with \code{np.loadtxt}, with \code{plot} methods
(scale \code{"a"}, \code{"q"} or \code{"r"} for the axis label;
\code{superplot} to overlay); \code{DataXYE0.save(filename)}. \legacy

\fnx{plotQ(\allowbreak{}plt,\allowbreak\ ...)}, \code{plotA(plt, ...)}, \code{plotR(plt, ...)}
Set title, labels, limits, grid and legend of the current matplotlib plot
for $Q$, angle or $r$ scales.

\fnx{defColors(\allowbreak{})}, \code{makePlotsA(runInfo, diffA, head, numor)}, \code{makePlotsQ(runInfo, diffQbin, head, numor)}
Colour list and the diagnostic plots of d4creg (individual numors,
individual banks, total in angle; total in $Q$).

\fnx{getDate(\allowbreak{})}
Current date as \code{" Dayname DD/MM/YYYY HH:MM:SS"}.

\fnx{extractDict(\allowbreak{}dictionary,\allowbreak\ key)}, \code{extractAttr(Dict, attr)}
Sub-dictionary of a dictionary of dictionaries; dictionary of one
attribute of a dictionary of objects (e.g.\ \code{extractAttr(elements, "sig_sca")}).

\fnx{execute\_command\_and\_write\_to\_file(\allowbreak{}command,\allowbreak\ output\_file)}
Runs a command, writing its output to a file; errors are printed, not raised.

\fnx{stop(\allowbreak{}message)}
Prints a message and asks whether to continue (interactive sessions only;
raises \code{RuntimeError} otherwise).

\noindent \code{chemistry.composition} re-exports \code{AtomicAvg},
\code{getConcentrations}, \code{getNofAtoms} (from \code{physics.properties})
and \code{extractAttr}, \code{extractDict} (from \code{utils.dicts}) for
compatibility.

% ---------------------------------------------------------------------
\tssec{Top level: \texttt{\_\_init\_\_} and \texttt{\_aliases}}

\code{toscana/__init__.py} imports the subpackages \code{data},
\code{dataio}, \code{experiment}, \code{isotopes}, \code{models},
\code{numerics}, \code{physics}, \code{plotting}, \code{utils} (not
\code{calibration}, \code{chemistry}, \code{workflows}, which must be
imported explicitly) and defines \code{__version__ = "0.1.0"}.
\code{_aliases.py} installs, at the front of Python's import machinery, a
finder that maps \code{toscana.math...} to \code{toscana.numerics...} and
\code{toscana.io...} to \code{toscana.dataio...}, returning the existing
module objects and emitting a \code{DeprecationWarning}; attribute access
\code{toscana.math} is handled the same way. Remove both once all
notebooks use the new names.

% When \input into the manual, stop here; stand-alone, close the document.
\ifdefined\ToScaNAManual\expandafter\endinput\fi
\end{document}

%  after \appendix. The stand-alone preamble is then skipped. In a report
%  or book class the three appendices become chapters, otherwise sections.
%  Packages needed: inputenc(utf8) fontenc(T1) lmodern textcomp amsmath amssymb
%  booktabs longtable array tabularx enumitem xcolor(dvipsnames) listings hyperref
%  (with \PassOptionsToPackage{obeyspaces,spaces}{url} before hyperref)
% =====================================================================
\ifdefined\ToScaNAManual\else
\documentclass[11pt,a4paper]{article}
\usepackage[utf8]{inputenc}
\usepackage[T1]{fontenc}
\usepackage{lmodern}
\usepackage{textcomp}
\usepackage[margin=2.3cm]{geometry}
\usepackage{amsmath,amssymb}
\usepackage{booktabs,longtable,array,tabularx}
\usepackage{enumitem}
\usepackage[dvipsnames]{xcolor}
\usepackage{listings}
% keep spaces inside \code{...} and allow line breaks there
\PassOptionsToPackage{obeyspaces,spaces}{url}
\usepackage[hidelinks]{hyperref}
\urlstyle{tt}
\begin{document}
\fi

% ------------------------------------------------------------- macros
% Heading levels: chapters in report/book manuals, sections in articles.
\ifdefined\chapter
  \providecommand{\tsapp}{\chapter}\providecommand{\tssec}{\section}\providecommand{\tssub}{\subsection}
\else
  \providecommand{\tsapp}{\section}\providecommand{\tssec}{\subsection}\providecommand{\tssub}{\subsubsection}
\fi
% \code{...}: verbatim identifiers, breakable after . _ ( , = / (url package)
\ifdefined\code\else\DeclareRobustCommand{\code}[1]{\path{#1}}\fi
\makeatletter\g@addto@macro\UrlBreaks{\do\_\do\(\do\,\do\=\do\[}\makeatother
\providecommand{\fnx}[1]{\par\medskip\noindent\begingroup\raggedright\hangindent=2em{\color{NavyBlue}\bfseries\ttfamily #1}\par\endgroup\nopagebreak\smallskip}
\providecommand{\tn}[1]{{\raggedright\ttfamily #1}}
\providecommand{\fn}[1]{\par\medskip\noindent{\color{NavyBlue}\textbf{\texttt{\detokenize{#1}}}}\par\nopagebreak\smallskip}
\providecommand{\st}[1]{\textsf{\footnotesize\colorbox{black!8}{#1}}}
\providecommand{\fixed}[1]{\textsf{\footnotesize\colorbox{green!15}{fixed #1}}}
\providecommand{\newbeh}{\textsf{\footnotesize\colorbox{orange!20}{results change}}}
\providecommand{\legacy}{\textsf{\footnotesize\colorbox{black!8}{legacy}}}
\lstset{basicstyle=\ttfamily\small,columns=fullflexible,keepspaces=true,
        frame=single,framerule=0.3pt,backgroundcolor=\color{black!3},
        breaklines=true,showstringspaces=false,aboveskip=6pt,belowskip=6pt}
\setlist{itemsep=2pt,topsep=3pt}
\setlength{\emergencystretch}{3em}

\ifdefined\ToScaNAManual\else
\begin{center}
  {\LARGE\bfseries ToScaNA}\\[4pt]
  {\large Total Scattering Neutron Analysis}\\[10pt]
  {\Large Appendices: code audit, corrections, tests and module reference}\\[10pt]
  {\normalsize Version 0.1.0 (branch \code{fixes/group-a}) --- November 2026}
\end{center}
\bigskip
\noindent\textbf{How to read these appendices.}
Appendix~\ref{app:audit} explains what was examined in the code, what was
found and what was changed, with the physical reasoning and the numbers
that justify each change. Appendix~\ref{app:tests} describes the automatic
tests that now protect the code, and how to run them.
Appendix~\ref{app:reference} is a module-by-module reference of the whole
package: what each function does, its arguments and outputs, the formulas
it implements, which files it reads or writes, and whether it was modified.
Throughout, labels such as \st{C2} or \st{P9} refer to the identifiers of
the audit report (\code{ToScaNA_technical_audit.md}): \st{C} for confirmed
errors, \st{P} for possible issues. The marker \newbeh{} flags changes that
modify numerical results with respect to the original code.

\tableofcontents
\clearpage
\appendix
\fi

% =====================================================================
\tsapp{Audit of the code and corrections}\label{app:audit}
% =====================================================================

\tssec{Context and method}

ToScaNA (\emph{Total Scattering Neutron Analysis}) is a Python package for
the reduction and analysis of neutron total-scattering data from the D4
diffractometer at the ILL. Its present form is a port of a single legacy
script, \code{legacy/lib/toscana.py}, into a package. The port was
written ``parity first'': the docstrings repeatedly state that legacy
behaviour, including known defects, was preserved on purpose so that new
and old results could be compared number by number. The legacy script and
its list of known divergences (\code{KNOWN__DIVERGENCES.md}) were not part
of the project tree that was audited.

The work was done in four stages:
\begin{enumerate}
  \item \textbf{Audit (inspection only).} All 52 Python files were read in
    full (about 6\,000 lines plus the 8\,857-line isotope table). Suspected
    problems were reproduced with small, read-only numerical experiments on
    synthetic data, run from a scratch directory so that nothing in the
    project or in the experimental data was modified. The findings were
    classified as \emph{confirmed errors} (C1--C30), \emph{possible issues}
    (P1--P21) and \emph{recommendations}.
  \item \textbf{Version control and tests.} A git repository was created.
    The first commit is the code exactly as received, so that every later
    change can be inspected with \code{git diff}. A test suite (Appendix~\ref{app:tests})
    was written \emph{before} most of the corrections: each test first
    failed on the original code and passes after the correction.
  \item \textbf{Corrections.} Mechanical bugs were fixed directly; changes
    that modify physical results were made only after a decision by the
    author of the code (density and scattering length of vanadium,
    reference wavelength, lattice parameter of Ni, window function).
  \item \textbf{New functionality.} Exact statistical error bars for
    $G(r)$, $g(r)$, RDF$(r)$, $T(r)$ and $N(r)$.
\end{enumerate}

\tssec{Version history}

All changes are on the git branch \code{fixes/group-a}; the branch
\code{main} contains the first two commits only. Each commit message
explains the change in detail (\code{git log}, \code{git show <id>}).

\begin{longtable}{@{}l p{0.78\textwidth}@{}}
\toprule
Commit & Content \\
\midrule\endhead
\tn{1958d36} & Baseline: the code as received, before any change. \\
\tn{67b3663} & Confirmed errors of the D4 reduction chain (Section~\ref{sec:reduction}). \\
\tn{f98cf26} & Test suite (standard library only). \\
\tn{1a9da55} & Mechanical bugs outside the reduction chain (``group A''). \\
\tn{43354d9} & Vanadium density 6.11\,g/cm$^3$. \\
\tn{6785d7d} & Thermal reference wavelength 1.798\,\AA; vanadium $b_\mathrm{coh}=-0.443$\,fm. \\
\tn{337bbaa} & Isotope class renamed \code{Elemento} $\to$ \code{Element}. \\
\tn{67d188a} & Ni lattice parameter 3.52\,\AA{} by default, user-configurable. \\
\tn{d76b9cd} & Fourier window: none, or the standard Lorch function. \\
\tn{0e2a14f} & Subpackages renamed \code{math}$\to$\code{numerics}, \code{io}$\to$\code{dataio}. \\
\tn{255c508} & Physical choices: annulus volume, free cross sections, sigmoid, $Q_0$, weighted Ni fit. \\
\tn{369e286} & Error bars for $G$, $g$, RDF, $T$ and $N$. \\
\bottomrule
\end{longtable}

\noindent To compare any file with the original: \code{git diff 1958d36 -- numerics/fourier.py}
(note that after commit \code{0e2a14f} the old path is \code{math/fourier.py}).
To return to the original code entirely: \code{git checkout 1958d36}.

% ---------------------------------------------------------------------
\tssec{The data-processing chain}\label{sec:chain}

The figure below summarises what the code does from raw data to
real-space functions, and where each step lives.

\begin{lstlisting}
Raw numors (NeXus .nxs, or ILL ASCII)             dataio.acquisition
  -> counts[9 banks x 64 cells], metadata
Angle check |2theta_read - 2theta_requested|       getOneNumor
Poisson errors on RAW counts  sigma = sqrt(N)      getErrorsNumor
Dead time (non-paralysable, per cell)              dataio.corrections
Relative efficiency  N/eps, sigma/eps              getEff
Normalisation to monitor (or time)                 normalise
Angles: 2theta = 2theta_lue + shift[det]
        + 15 deg*det + 0.125 deg*cell              getAngles
Stack numors, subtract zero angle 2theta0          getNumors
Regroup cells onto a 2theta grid                   dataio.d4creg_outputs.getDiffA
  -> .adat (CORRECT format), .reg (D4 format)      saveDiffAngle
Convert to Q and rebin onto a Q grid               numerics.binning.rebin
  -> .qdat (CORRECT format)                        saveDiffQ
--- external program CORRECT: container, environment, absorption,
    multiple scattering, vanadium normalisation (dataio.correct) ---
S(Q)  ->  sine Fourier transform                   numerics.fourier.sineFT
  -> G(r), g(r), RDF(r), T(r), N(r)  [+ errors]    dataio.saving.saveRSCF
Optional: back-transform with low-r cut            numerics.fourier.backFT
Ni powder calibration of lambda and 2theta0        calibration.d4nifit
\end{lstlisting}

\noindent The whole regrouping chain is driven by a parameter file read by
\code{dataio.parameters.readParam}; the line \code{<run>} in that file
launches \code{workflows.d4creg.d4creg}. Corrections for attenuation,
multiple scattering and container are \emph{not} computed in ToScaNA: they
are delegated to CORRECT, for which ToScaNA writes the input file.

% ---------------------------------------------------------------------
\tssec{Conventions, units and formulas}\label{sec:conventions}

\begin{longtable}{@{}l l p{0.52\textwidth}@{}}
\toprule
Quantity & Unit & Notes \\
\midrule\endhead
Scattering angle $2\theta$ & degree & zero angle $2\theta_0$ subtracted in \code{getNumors} \\
Momentum transfer $Q$ & \AA$^{-1}$ & $Q = \frac{4\pi}{\lambda}\sin\theta$ \\
Distance $r$ & \AA & \\
Wavelength $\lambda$ & \AA & \\
Atomic number density $\rho$ & atoms/\AA$^3$ & $\rho = d\,N_A/\bar M \times 10^{-24} = 0.602214076\,d/\bar M$ \\
Macroscopic density $d$ & g/cm$^3$ & \\
Scattering length $b$ & fm & \\
Cross sections $\sigma$ & barn & $1\,\mathrm{b} = 100\,\mathrm{fm}^2$; $\sigma = 4\pi|b|^2/100$ with $b$ in fm \\
Sample dimensions & mm & volumes converted to cm$^3$ (division by 1000) \\
CORRECT input & cm & radius $=$ diameter$/20$ \\
Counts & per monitor & normalised to \code{norm_mon} monitor counts (default $10^6$) or \code{norm_time} seconds \\
Thermal reference & & $v=2200$\,m/s, $\lambda_\mathrm{th}=1.798$\,\AA, $E_\mathrm{th}=25.30$\,meV \\
\bottomrule
\end{longtable}

\paragraph{Real-space functions.} With $F(Q) = Q[S(Q)-1]$ and a window
(modification) function $M(Q)$ that vanishes for $Q \ge Q_\mathrm{max}$,
\code{sineFT} computes (columns of its output)
\begin{align}
  G(r) &= \frac{2}{\pi}\int_0^{Q_\mathrm{max}} Q\,[S(Q)-1]\,M(Q)\,\sin(Qr)\,\mathrm{d}Q
        \;=\; 4\pi\rho r\,[g(r)-1], \label{eq:G}\\
  g(r) &= 1 + \frac{G(r)}{4\pi\rho r} \quad (g(0) \equiv 0),\\
  \mathrm{RDF}(r) &= 4\pi\rho r^2 g(r) = r\,G(r) + 4\pi\rho r^2,\\
  T(r) &= 4\pi\rho r\, g(r) = G(r) + 4\pi\rho r,\\
  N(r) &= \int_{r_0}^{r} \mathrm{RDF}(r')\,\mathrm{d}r' \qquad (r_0 = \text{first point of the $r$ grid}).
\end{align}
$T(r)$ is defined here without the factor $\langle b\rangle^2$ used by some
authors. The integrals are evaluated with Simpson's rule
(\code{scipy.integrate.simpson}) on the given grids; no FFT is used, so
arbitrary, non-uniform grids are allowed. The back-transform \code{backFT}
uses
\begin{equation}
  S(Q) = 1 + \frac{4\pi\rho}{Q}\int r\,[g(r)-1]\sin(Qr)\,\mathrm{d}r,
  \qquad S(0) = 1 + 4\pi\rho\int r^2[g(r)-1]\,\mathrm{d}r .
\end{equation}
These formulas were checked analytically and numerically during the audit
and are correct.

% ---------------------------------------------------------------------
\tssec{Corrections to the reduction chain}\label{sec:reduction}

All items in this section are in commit \code{67b3663}. Before these
corrections, the \emph{intensities} produced by the chain were correct;
the main defects concerned the \emph{error bars}, and the robustness of
the regrouping. The intensities of regrouped and rebinned data were
verified to be identical to the original code (synthetic data, all
in-range cells).

\paragraph{\fixed{C2} Error of the monitor normalisation.} \newbeh{}
\code{normalise} replaced the counts $y$ by $y_n = f\,y$, $f = N_\mathrm{norm}/N_\mathrm{mon}$,
and only afterwards formed the relative error $e/y_n$ with the
\emph{un-normalised} error $e$. The relative error was therefore multiplied
by $N_\mathrm{mon}/N_\mathrm{norm}$. For a cell with $10^4$ counts, the
relative error should be $\sqrt{10^{-4}+1/N_\mathrm{mon}}\approx 0.010$;
the code gave 0.0033 for $N_\mathrm{mon}=10^5$ (3 times too small) and
0.100 for $N_\mathrm{mon}=10^7$ (10 times too large). It was correct only
when the monitor happened to equal $10^6$. Now
\begin{equation}
  y_n = f y, \qquad \sigma_n^2 = (f\,\sigma)^2 + \Bigl(y_n\,\frac{\sigma_{N_\mathrm{mon}}}{N_\mathrm{mon}}\Bigr)^2 .
\end{equation}
Cells with zero counts are flagged \code{NaN} and are therefore excluded by
the regrouping, which is what the original code did in practice (it
produced $0/0$).

\paragraph{\fixed{C5} Poisson errors.} \newbeh{} The error was computed as
$\sqrt{N}$ \emph{after} dead-time and efficiency corrections, i.e.\
$\sqrt{N_\mathrm{raw}/\varepsilon}$ instead of $\sqrt{N_\mathrm{raw}}/\varepsilon$.
Errors are now taken on raw counts and scaled by exactly the same factors
as the counts (dead-time factor, then $1/\varepsilon$). For
$N_\mathrm{raw}=400$ and $\varepsilon=0.8$ the error per cell is now
$25.0$ (it was $22.4$).

\paragraph{\fixed{C4} Dead-time check.} The test
\code{if (correction.all()) > 0} only checks that no factor is exactly
zero: when the measured rate times the dead time exceeded 1, negative
factors were accepted and produced \emph{negative counts} (reproduced:
$\nu\tau = 1.4$ gave $-5\times10^5$ counts). All factors must now be
positive, otherwise a \code{ValueError} is raised. The corrected total
detector counts were computed with the factor of cell [0,0] only; they now
use $\sum N_\mathrm{corr}/\sum N_\mathrm{raw}$. The log message said
``Monitor DT'' for the detector; corrected.

\paragraph{\fixed{C29} Combination of errors when regrouping and rebinning.} \newbeh{}
Both \code{getDiffA} (cells $\to$ angular grid) and \code{rebin} (angle
$\to$ $Q$ grid) form weighted means $\bar y = \sum w_i y_i/\sum w_i$ (weights
= overlap fraction $\times$ monitor share). The errors were averaged in the
same way, $\sum w_i\sigma_i/\sum w_i$, so that they did not decrease when
points were combined. The error of a weighted mean is
\begin{equation}
  \sigma_{\bar y} = \frac{\sqrt{\sum_i w_i^2\sigma_i^2}}{\sum_i w_i}.
\end{equation}
Example: ten points with $\sigma=10$ gave 10.0; now 3.16. On a synthetic
two-numor regrouping the median error fell from 3.0 to 1.58.

\paragraph{\fixed{C6} Size of the angular grid.} \code{getDiffA} sized its
arrays with \code{int((amax-amin)/step)} but filled them with
\code{np.arange(amin, amax, step)}, which has $\lceil(\cdot)\rceil$
elements. Any step that does not divide the range exactly crashed
(reproduced: \code{(1, 140, 0.3)} gave ``could not broadcast (464,) into
(463,)''). The grid size is now taken from the grid itself.

\paragraph{\fixed{C7} Cells outside the angular range.} Bin indices were
obtained with \code{int()}, which truncates towards zero, and were not
checked. A cell just below \code{amin} was counted in bin 0; a cell further
below gave a \emph{negative} index, which Python interprets as counting
from the end of the array: counts measured at $-5^\circ$ appeared near
$135^\circ$. Cells above \code{amax} raised \code{IndexError}. Indices now
use \code{np.floor} and out-of-range parts of a cell are discarded.

\paragraph{\fixed{C8} Last $Q$ point erased.} After rebinning, points with
$y\le0$ are masked together with the previous point; for the first point
the ``previous'' index $-1$ is the \emph{last} point of the array, which was
erased. Corrected.

\paragraph{\fixed{C9} Width used for the $Q$ rebinning.} The angular width
of each regrouped point was hard-coded to $0.125^\circ$ in \code{d4creg}.
It is now the actual step of the regrouped grid (which \code{getDiffA}
enlarges to the cell width if the requested step is smaller). The width
in $Q$ is $\Delta Q = \frac{2\pi}{\lambda}\cos\theta\,\Delta(2\theta)\,\frac{\pi}{180}$.

\paragraph{\fixed{C19} Number of bins.} \code{get_xlim} computed the number
of bins with \code{int()} of a floating-point ratio:
\code{get_xlim(0, 0.3, 0.1)} gave 3 bins instead of 4 because
$0.4/0.1 = 3.9999\ldots$. A tolerance of $10^{-9}$ is now added before
taking the floor; genuinely non-commensurate ranges keep floor semantics.
The same floor/truncation fix was applied to \code{get_bins}.

\paragraph{\fixed{C25} Parsing of the d4creg parameter file.} Lines were
split with \code{split(" ")}, which kept the newline on the last token and
produced empty tokens for repeated spaces. When a value was the last word
of its line it silently became, e.g., \verb|"Nexus\n"|: the data format
then fell back to the ASCII reader, \verb|"True\n"| disabled writing and
plotting, and numor ranges failed their length check. Lines are now cut
at \texttt{\#} (inline comments allowed) and split on any white space.

\paragraph{\fixed{C27} Last line of input files.} Several readers removed
the last character of each line assuming it was a newline; a final line
without newline lost a digit (an efficiency 1.0234 became 1.023). Lines
are now stripped. Affected: \code{getEff}, \code{getDec}, \code{read_xye}.

\paragraph{\fixed{C3} Weighted sum of two data sets.} \newbeh{}
\code{binary_sum} combined $\sqrt{w_1^2 e_1 + w_2^2 e_2}$ instead of
$\sqrt{w_1^2 e_1^2 + w_2^2 e_2^2}$ (dimensionally wrong; $e_1=e_2=2$ gave
2.0 instead of 2.83).

% ---------------------------------------------------------------------
\tssec{Fourier transform and real-space functions}\label{sec:fourier}

\paragraph{\fixed{C1} Window function.} \newbeh{}
The original code offered a step (truncation) and a ``Lorch'' window that
was rescaled so that $\int_0^{Q_\mathrm{max}} M\,\mathrm{d}Q = Q_\mathrm{max}$.
This rescaling gives $M(0) = \pi/\mathrm{Si}(\pi) = 1.696$ instead of 1
($\mathrm{Si}$ is the sine integral), and multiplies $G(r)$ by 1.70 where
the window is smooth. At low $r$, where $g(r)=0$ and $G(r)=-4\pi\rho r$,
the result was $g(r) = 1 - 1.70 = -0.70$ (reproduced). After discussion, the
rescaled version was removed. \code{sineFT(window=...)} and
\code{getSineFT(w=...)} now accept only
\begin{itemize}
  \item \code{0}: no window (truncation at $Q_\mathrm{max}$);
  \item \code{1}: the standard Lorch function
    $M(Q) = \dfrac{\sin(\pi Q/Q_\mathrm{max})}{\pi Q/Q_\mathrm{max}}$, $M(0)=1$;
\end{itemize}
any other value raises \code{ValueError} (the original code treated any
non-zero value as Lorch). The array function \code{LorchN} was replaced by
\code{LorchArray}.

\paragraph{Why the standard Lorch window preserves coordination numbers.}
$M(Q)$ above is the Fourier transform of a box of half-width
$\Delta = \pi/Q_\mathrm{max}$. Multiplying $F(Q)$ by $M(Q)$ is therefore
exactly equivalent to replacing the (truncated) $G(r)$ by its running
average over $[r-\Delta, r+\Delta]$; this is Lorch's original construction
(1969). A running average with a kernel of unit integral, i.e.\ $M(0)=1$:
\begin{enumerate}
  \item conserves the total area of $G(r)$ (the rescaled version inflated
    all areas by 1.70);
  \item changes an integral between finite limits $[r_1, r_2]$ only through
    edge terms of order $\frac{\Delta^2}{6}[G'(r_2) - G'(r_1)]$, which vanish
    when the limits lie where the curve is flat (the region $g=0$ and the
    first minimum);
  \item has no negative lobes, unlike the kernel $\sin(Q_\mathrm{max} r)/r$
    of plain truncation, so it does not move area by oscillation.
\end{enumerate}
Since $N(r)$ integrates $r\,G(r)$ rather than $G(r)$, an additional term
$\frac{\Delta^2}{3}[G(r_2)-G(r_1)]$ appears; with $Q_\mathrm{max}=23.5$\,\AA$^{-1}$
it is about $0.006\,\Delta G$, negligible. Exact conservation for
\emph{arbitrary} limits is impossible with a finite $Q_\mathrm{max}$: it
would require $M(Q)=1$ up to infinite $Q$. Limits that cut through a peak
always lose part of its area. Numerically, on a model liquid whose exact
coordination number is known (Appendix~\ref{sec:synthetic}):

\begin{center}
\begin{tabular}{@{}c c c c c@{}}
\toprule
$Q_\mathrm{max}$ (\AA$^{-1}$) & Limits (\AA) & Exact $N$ & No window & Standard Lorch \\
\midrule
15.0 & 0.0--3.3 & 10.498 & 10.487 & 10.498 \\
15.0 & 2.3--3.3 & 10.498 & 10.497 & 10.490 \\
23.5 & 0.0--3.3 & 10.498 & 10.499 & 10.498 \\
23.5 & 2.3--3.3 & 10.498 & 10.481 & 10.502 \\
\bottomrule
\end{tabular}
\end{center}
With the rescaled window the same coordination numbers came out 4\,\% too
high from the edge of the first shell, and 30\,\% too low from $r=0$.
\textbf{Practical rule:} integrate from a point where $g(r)=0$ to a
minimum. A related pitfall: on a grid produced by \code{np.arange}, the
value \code{3.3} may be stored as 3.3000000000000003 and a mask
\code{r <= 3.3} then drops the end point; since RDF$(3.3\,\text{\AA})\approx 14$
atoms/\AA, this alone changes $N$ by about 1\,\%. Rounding the grid
(\code{np.round(r, 10)}) avoids it.

\paragraph{\fixed{C18} Running coordination number.} \newbeh{} The slice
end of the running integral was computed as \code{1 + int(r/dr)}, which
assumes that the grid starts at 0 and that the division is exact. Even on
a grid starting at 0 in steps of 0.02\,\AA, rounding (e.g.\
$0.58/0.02 = 28.999\ldots$) stopped the integral one point short at 16 of
400 points, producing small steps in $N(r)$; on a grid not starting at 0,
all values were wrong. The loop index is used now; $N(r_0)=0$.

\paragraph{\fixed{C17} Back-transform.} \newbeh{} \code{backFT} divided by
$Q$ and $r$ without protection, producing \code{NaN}/\code{inf} at
$Q=0$ and $r=0$; the analytic limits are now used ($S(0)$ above;
$g(0)\equiv0$ as in \code{sineFT}; $T = 4\pi\rho r g$ computed directly). Its
running integral excluded the current point, unlike \code{sineFT}; made
consistent. The \code{cut} argument: \code{[-1]} zeroes $g(r)$ up to its
last negative value (previously, when $g$ was never negative, the first two
points were zeroed anyway); otherwise any odd-length list
$[r_0, a_1, b_1, a_2, b_2, \ldots]$ zeroes $r<r_0$ and every interval
$(a_k, b_k)$; even lengths raise (previously lengths above 13 were silently
ignored).

\paragraph{\fixed{C20} Sine integral.} \code{getSineIntegral} integrated
numerically on a grid that stopped at $x(1-10^{-4})$ (relative error
$\sim10^{-4}$: $\mathrm{Si}(1)$ gave 0.945999 instead of 0.946083) and
returned \code{NaN} for $x=0$. It now uses the exact \code{scipy.special.sici}.

\paragraph{\fixed{P4} Truncation in \code{getSineFT}.} With \code{w=0} the
argument \code{vmax} was ignored; the integral now stops at \code{vmax}
(results unchanged when \code{vmax} is the default, the end of the data).

\paragraph{\fixed{P5} Error values used as data.} \code{Lorch} and
\code{step} returned $-999$ for $Q_\mathrm{max}\le 0$ after printing a
message; that value was then silently used as a window. They now raise
\code{ValueError}.

% ---------------------------------------------------------------------
\tssec{Error bars of the real-space functions}\label{sec:errors}

\paragraph{Method.} Commit \code{369e286}. Because the transform is linear
in $S(Q)$, statistical errors can be propagated exactly, without Monte
Carlo noise. Writing the Simpson quadrature of Eq.~\eqref{eq:G} as a sum
over the $Q$ points with weights $w_k$,
\begin{equation}
  G(r_j) = \sum_k W_{jk}\,[S(Q_k)-1], \qquad
  W_{jk} = \frac{2}{\pi}\,c\;Q_k\,M(Q_k)\,\sin(Q_k r_j)\,w_k ,
\end{equation}
($c$ is the argument \code{constant}), and assuming uncorrelated $S(Q)$ points
with errors $\sigma_k$,
\begin{equation}
  \sigma_G^2(r_j) = \sum_k W_{jk}^2\,\sigma_k^2,\qquad
  \sigma_T = \sigma_G,\qquad
  \sigma_\mathrm{RDF} = r\,\sigma_G,\qquad
  \sigma_g = \frac{\sigma_G}{4\pi\rho r}\ (r>0).
\end{equation}
The weights $w_k$ are obtained from SciPy itself (by integrating unit
vectors), so they are exactly those used for the values.

\paragraph{The running integral needs the correlations.} Neighbouring
$G(r)$ points are built from the same $S(Q)$ data and are strongly
correlated; adding their variances as if they were independent would be
badly wrong. Instead $N$ is written directly as a linear function of $S$:
$N(r_i) = \sum_j C_{ij}\, r_j\, G(r_j) + \text{const}$, where $C$ is the
matrix of the running Simpson rule, so that
$\sigma_N^2(r_i) = \sum_k \bigl[(C\,\mathrm{diag}(r)\,W)_{ik}\bigr]^2 \sigma_k^2$.
A first version used trapezoidal weights for $C$ and the Monte Carlo test
showed errors up to 11\,\% wrong without window; the final version uses
exactly the running Simpson rule used for $N(r)$ itself. That matrix is
built in $O(N_r^2)$ operations (0.1\,s for 2000 points): odd-length
prefixes are composite Simpson, extended by one three-point panel at a
time; for even lengths SciPy adds a correction on the last interval that
depends only on the last four points. The construction is compared with
SciPy at run time and rebuilt directly (slowly) if a future SciPy uses a
different rule.

\paragraph{Validation.} $S(Q)$ was perturbed 1500 times with Gaussian noise
of the stated errors, each copy was transformed with the ordinary
\code{sineFT}, and the standard deviation of each output was compared with
the analytic error. Ratio analytic/Monte Carlo, median over $r$ [range]:

\begin{center}\small\setlength{\tabcolsep}{4pt}
\begin{tabular}{@{}l c c c c c@{}}
\toprule
Window & $G$ & $g$ & RDF & $T$ & $N$ \\
\midrule
none    & 1.005 [0.95--1.05] & 1.005 [0.95--1.05] & 1.005 [0.95--1.05] & 1.005 [0.95--1.05] & 1.003 [0.96--1.05] \\
Lorch   & 0.998 [0.96--1.06] & 0.997 [0.96--1.06] & 0.998 [0.96--1.06] & 0.998 [0.96--1.06] & 0.998 [0.95--1.06] \\
\bottomrule
\end{tabular}
\end{center}
The spread of 5--6\,\% is the expected statistical spread of a standard
deviation estimated from 1500 samples.

\paragraph{Usage.}
\begin{lstlisting}
fou = sineFT(SQ, r, qmax=23.5, density=rho, window=1, with_errors=True)
saveRSCF("sample.rscf", fou, ["Sample: ...", "Qmax = 23.5, Lorch"])
\end{lstlisting}
\code{SQ} must have three columns ($Q$, $S$, $\sigma_S$). The output has 11
columns, each error beside its function:
$r, G, \sigma_G, g, \sigma_g, \mathrm{RDF}, \sigma_\mathrm{RDF}, T, \sigma_T, N, \sigma_N$.
The file header describes every column, its formula and units. Without
\code{with_errors} the output keeps its 6 columns.

\paragraph{Limits of these error bars.}
\begin{itemize}
  \item They are \emph{statistical} only. Systematic errors (normalisation,
    density, self-scattering subtraction, truncation) are not included, and
    usually dominate.
  \item They assume uncorrelated $S(Q)$ points. The fractional rebinning
    correlates neighbouring points slightly, so the errors may be somewhat
    underestimated.
  \item Errors at different $r$ are strongly correlated. They are correct
    point by point, but a least-squares fit of $G(r)$ that treats them as
    independent overstates the information contained in the data.
\end{itemize}

% ---------------------------------------------------------------------
\tssec{Line shapes, fitting and the Ni calibration}\label{sec:fitting}

\paragraph{\fixed{C11} \code{Lorentzian_error}.} The $x_0$ term contained
\code{e_x0 * 22} (typing error for \code{e_x0**2}) and lacked one factor of
the normalised peak $p = (\gamma/2)^2/[(x-x_0)^2+(\gamma/2)^2]$; the correct
derivative is $\partial L/\partial x_0 = 2A(x-x_0)\,p^2/(\gamma/2)^2$.
Reproduced error: 2.37 instead of 0.032. Now checked against numerical
derivatives.

\paragraph{\fixed{C12} \code{LorGau_error}.} It evaluated an area-normalised
Gaussian while \code{LorGau} uses a unit-height Gaussian, so the returned
value differed from \code{LorGau} (0.600 instead of 1.000 at $x=0$); the
$\eta$, $\sigma$ and $\gamma$ error terms lacked the factor $f_0$. Rewritten
with explicit derivatives.

\paragraph{\fixed{C13} \code{wsum2}.} With $w_1=0$ the error was $w_2 e_2$,
negative for $w_2<0$; now $|w_2|e_2$. With $w_2=0$ and $w_1\notin(0,1]$ it
printed a message and then crashed; it now raises \code{ValueError}.

\paragraph{\fixed{C10} Initial values of the Ni fit.} The amplitudes A2, A3,
A4 started from \code{initial[6..8]} instead of \code{initial[7..9]}
(shifted by one, so A1 and A2 started equal and \code{initial[9]} was never
used). Only starting values were affected, but a non-linear fit can
converge differently.

\paragraph{\fixed{P11} Weighted Ni fit.} \newbeh{} The fit ignored the data
errors. It now uses weights $1/\sigma$; if no valid errors exist
(two-column files store $\sigma=-1$) it fits without weights and prints a
warning. lmfit keeps \code{scale_covar=True}: the reported uncertainties
are multiplied by $\sqrt{\chi^2_\nu}$ and so also account for model misfit.
\code{setting_model_d4nifit} divides the intensities by the value in the
valley between the second and third peaks; the errors are now divided by
the same value.

\paragraph{\fixed{P7} Ni lattice parameter.} \newbeh{} The value
3.52024\,\AA{} was hard-coded in five places. There is now one default,
\code{physics.crystallography.NI_LATTICE_PARAMETER = 3.52}, and every
function that uses it accepts an argument \code{lattice}:
\begin{lstlisting}
setting_model_d4nifit("ni.adat", lattice=3.5238)
fitting_model_d4nifit(nickel, limits, initial, lattice=3.5238)
\end{lstlisting}
In the fit the lattice parameter is a fixed lmfit parameter: the peak
positions determine only $\lambda/a$, so $\lambda$ and $a$ cannot both be
fitted, and the fitted $\lambda$ scales linearly with $a$ (going from
3.52024 to 3.52 lowers $\lambda$ by 0.007\,\%). \code{showing_results_d4nifit}
reads the value back from the fit result.

\paragraph{\fixed{P20} Fitting ranges.} \code{fittingRange0},
\code{fittingRange} and \code{fit_range} were three copies of the same code;
they now share one helper with identical outputs.

% ---------------------------------------------------------------------
\tssec{Physical constants and sample physics}\label{sec:physics}

\paragraph{\fixed{P6} Vanadium density.} \newbeh{} The defaults used
6.51\,g/cm$^3$ (\code{setVanaInfo}) and the corresponding atomic density
0.0769591\,atoms/\AA$^3$ (\code{setSampleInfo}), typed separately. Both now
derive from one constant, \code{VANADIUM_DENSITY_GCC = 6.11}, giving
0.0722305\,atoms/\AA$^3$. The default sample/vanadium density ratio
\code{DenRatioSV} increases by 6.5\,\%.

\paragraph{\fixed{C16} Vanadium scattering length.} \newbeh{} The isotope
table had $b_\mathrm{coh}=-0.443$\,fm but $\sigma_\mathrm{coh}=0.01838$\,b,
which corresponds to $-0.3824$\,fm; \code{setVanaInfo} used $-0.3824$\,fm.
The value $-0.443$\,fm was adopted everywhere and the cross sections now
follow from it: $\sigma_\mathrm{coh} = 4\pi(0.443)^2/100 = 0.02466$\,b and
$\sigma_\mathrm{sca} = \sigma_\mathrm{coh}+\sigma_\mathrm{inc} = 5.10466$\,b
(were 0.01838 and 5.09838\,b). $\sigma_\mathrm{inc}=5.08$\,b and
$\sigma_\mathrm{abs}=5.08$\,b are unchanged, so the vanadium normalisation
changes very little ($\sigma_\mathrm{sca}$ $+0.12$\,\%). A test checks that
the defaults of \code{setVanaInfo} equal the table values.

\paragraph{\fixed{C30} Thermal reference wavelength.} \newbeh{} Tabulated
absorption cross sections refer to neutrons of 2200\,m/s, i.e.\
$\lambda = h/(m_n v) = 1.798$\,\AA{} and $E = 81.8042/\lambda^2 = 25.30$\,meV.
Three different values were in use: 1.8089\,\AA{} (\code{getAbsXS},
$\sigma_\mathrm{abs}$ 0.6\,\% too small), 1.8\,\AA{} (\code{setSampleInfo})
and 25.3\,meV (\code{scattering_probability}). All now derive from
\code{physics.properties.THERMAL_WAVELENGTH = 1.798} and
\code{THERMAL_ENERGY_MEV}; absorption follows the $1/v$ law
$\sigma_\mathrm{abs}(\lambda) = \sigma_\mathrm{abs}^\mathrm{th}\,\lambda/1.798$.

\paragraph{\fixed{C14} Volume of annular samples.} \newbeh{}
\code{Measurement} computed $\pi h (r_o-r_i)^2$ instead of
$\pi h (r_o^2-r_i^2)$. Solid cylinders ($r_i=0$) are unaffected; for a
4/6\,mm annulus the volume was 5 times too small, and the packing fraction
5 times too large.

\paragraph{\fixed{P9} Free-atom cross sections.} \newbeh{} A bound nucleus
of mass ratio $A = m/m_n$ has a free cross section
$\sigma_\mathrm{free} = \sigma_\mathrm{bound}\,[A/(A+1)]^2$. For a mixture
this factor must be applied element by element,
\begin{equation}
  \sigma_\mathrm{free} = \sum_i c_i\,\sigma_i\left(\frac{A_i}{A_i+1}\right)^2 ,
\end{equation}
whereas the code applied it once with the mean mass. The mean mass hides
how light hydrogen is:

\begin{center}
\begin{tabular}{@{}l r r r@{}}
\toprule
Sample & Per element (b) & Mean mass (old, b) & Change \\
\midrule
H$_2$O & 14.908 & 41.116 & $-63.7$\,\% \\
D$_2$O & 3.508 & 4.907 & $-28.5$\,\% \\
SiO$_2$ & 3.170 & 3.212 & $-1.3$\,\% \\
Ni & 17.880 & 17.880 & 0 \\
\bottomrule
\end{tabular}
\end{center}
Applies to \code{Measurement.free} and to \code{FreeCohXS},
\code{FreeIncXS}, \code{FreeScaXS} of \code{setSampleInfo}. Both use
\code{getMassNumber} (neutron mass 1.0086649\,u); previously the isotope
table value 1.00866\,u was also used.

\paragraph{\fixed{C26} \code{sigmoidal}.} \newbeh{} The function returned
$\delta\,\mathrm{lowQ} + (1-\delta)\,\mathrm{highQ}$ with
$\delta = e^{(Q-Q_0)/\mathrm{d}Q}$: unbounded, with negative weights above
$Q_0$, and running from highQ to lowQ, the opposite of what the parameter
names state. It is now the logistic step
\begin{equation}
  f(Q) = \frac{\mathrm{lowQ} + \mathrm{highQ}\,\delta}{1+\delta},\qquad
  f(0)\approx\mathrm{lowQ},\ f(\infty) = \mathrm{highQ},\ f(Q_0) = \tfrac12(\mathrm{lowQ}+\mathrm{highQ}),
\end{equation}
computed with \code{scipy.special.expit} so that large $|Q-Q_0|/\mathrm{d}Q$
cannot overflow. This is the same shape as \code{models.scattering.inelastic}
(with $\mathrm{highQ} = \mathrm{lowQ}\,[A/(A+1)]^2$). \code{siglin}
ignored its argument \code{a1}; it is now $\texttt{sigmoidal}(Q) + a_1 Q$.
\emph{This interpretation of \code{a1} should be confirmed by the author.}

\paragraph{\fixed{P21} One default transition.} \code{inelastic} used
$Q_0=7.0$\,\AA$^{-1}$ and \code{vanaQdep} 7.4\,\AA$^{-1}$. One pair of
defaults, \code{SIGMOID_Q0 = 7.0} and \code{SIGMOID_DQ = 2.4}
(\code{numerics.line_shapes}), now serves \code{sigmoidal}, \code{siglin},
\code{inelastic} and \code{vanaQdep} (effect on \code{vanaQdep}
$\le 0.16$\,\%).

\paragraph{\fixed{C15} Bound cross sections of \code{Element}.} The
attributes \code{scoh_bound} and \code{sinc_bound} were $4\pi|b|^2$ in fm$^2$
but labelled barns (100 times too large: Ni gave 1333 instead of 13.3).
Divided by 100. They were not used elsewhere in the package.

% ---------------------------------------------------------------------
\tssec{Structure and robustness of the code}\label{sec:structure}

\paragraph{\fixed{C21} Names that hid the standard library.} The
subpackages \code{math/} and \code{io/} have the names of standard Python
modules. When Python was started in the project directory (for example a
notebook opened there), \code{import math} found the project folder
instead; \code{import lmfit} failed with ``partially initialized module
`math' has no attribute `floor'''. They were renamed (with \code{git mv},
keeping their history):
\begin{center}
\code{toscana.math} $\to$ \code{toscana.numerics} \qquad
\code{toscana.io} $\to$ \code{toscana.dataio}
\end{center}
The old dotted names remain importable through a small import hook
(\code{_aliases.py}) and emit a \code{DeprecationWarning}; they return the
\emph{same} module objects (an earlier version of the hook loaded a second
copy of each module, which a test detected). No file or folder of the
project now has the name of a standard-library module, and a test checks
this against \code{sys.stdlib_module_names}.

\paragraph{English naming.} The isotope class \code{Elemento} is now
\code{Element}. \code{Elemento} and \code{elemento} remain available as
deprecated aliases. Other Spanish names in the same module
(\code{NOMBRES}, local variables) were not changed.

\paragraph{Other fixes.}
\begin{itemize}
  \item \fixed{C22} \code{DataXYE} crashed on files without any point in
    $[5,40]$ (the range of its automatic peak search), e.g.\ $r$-space
    files; it saved figures \emph{after} \code{plt.show()}, which gives
    empty images with interactive back-ends.
  \item \fixed{P15} An unknown isotope name now raises \code{ValueError}
    instead of returning a half-built object.
  \item \fixed{P13} \code{Measurement.showTable} replaced \code{sys.stdout}
    with an unclosed file and then restored \code{sys.__stdout__}, which
    broke Jupyter output; it now uses \code{contextlib.redirect_stdout}.
  \item \fixed{C23} Characters such as \AA, $^\circ$ and $\pm$ were stored
    doubly encoded (\AA{} appeared as two wrong characters); repaired. \fixed{C24} The \code{.reg}
    header gave the wavelength in degrees; now \AA.
  \item \st{P19} The audit claimed that \code{smooth_curve} had an
    asymmetric window; this was wrong (even sizes become the next odd
    size). Only its documentation was completed.
\end{itemize}

% ---------------------------------------------------------------------
\tssec{What changes for the user}\label{sec:migration}

\begin{longtable}{@{}p{0.36\textwidth} p{0.58\textwidth}@{}}
\toprule
Before & Now \\
\midrule\endhead
\code{from toscana.math...}, \code{toscana.io...} & \code{toscana.numerics...}, \code{toscana.dataio...} (old names work, with a warning) \\
\code{Elemento("Ni")} & \code{Element("Ni")} (old name works, with a warning) \\
\code{LorchN(Q, Qmax)} & removed; \code{LorchArray(Q, Qmax)} is the standard Lorch \\
\code{sineFT(..., window=1)} rescaled Lorch & standard Lorch, $M(0)=1$; values other than 0, 1 raise \\
\code{sineFT} 6 columns & unchanged; \code{with_errors=True} gives 11 columns \\
\code{saveRSCF} 6 columns, no header & 6 or 11 columns, header describing each column \\
Error bars of \code{.adat}, \code{.reg}, \code{.qdat} & smaller and statistically correct (C2, C5, C29) \\
$N(r)$ & correct at every point (C18) \\
\code{Element.scoh_bound} & in barns (100 times smaller) \\
Vanadium defaults & 6.11\,g/cm$^3$, $b=-0.443$\,fm, $\sigma_\mathrm{coh}=0.02466$\,b \\
$\sigma_\mathrm{abs}(\lambda)$ & $\propto\lambda/1.798$ \\
Free cross sections & per element (large change for H, D) \\
\tn{sigmoidal} & a true logistic step \\
Ni calibration & weighted fit, lattice 3.52\,\AA{} by default, \code{lattice=} argument \\
Silent prints on bad input & \code{ValueError} (unknown isotope, invalid \code{cut}, $Q_\mathrm{max}\le0$, invalid window or weights, negative dead-time factor) \\
\bottomrule
\end{longtable}

% ---------------------------------------------------------------------
\tssec{Open items}\label{sec:open}

\paragraph{Need data or instrument knowledge.}
\begin{itemize}
  \item \st{P1} Angle labels of \code{getDiffA}: the angle written for a bin
    is its lower edge, while a cell centred in the bin contributes fully to
    it, i.e.\ a shift of half a step ($0.0625^\circ$ for $0.125^\circ$). It
    is absorbed by the zero angle if the Ni calibration was regrouped in the
    same way; to be confirmed with a real Ni run.
  \item \st{P2} The dead time is applied to the rate of each \emph{cell};
    for position-sensitive banks it is usually governed by the rate of the
    whole bank or of its electronics.
  \item \st{P3} Data are binned twice (cells $\to 2\theta$ $\to Q$), which
    adds broadening and correlations; a direct cells $\to Q$ rebinning
    would avoid one step.
  \item \st{P8} The reactor power is fixed to 1\,MW for NeXus files (only
    the logged rates per MW are affected); the NeXus field is needed.
  \item \st{P16} The ASCII numor reader assumes 10 values per line and a
    block of 11 lines per detector; a sample file is needed to check the
    case \code{ncell} multiple of 10.
\end{itemize}

\paragraph{Found while writing this appendix (not yet changed).}
\begin{itemize}
  \item In \code{getOneNumor}, the $Q$ column of each numor (used in the
    per-numor \code{.cdat} files) is computed from the angles \emph{before}
    the zero angle is subtracted; the main \code{.qdat} output is computed
    from corrected angles and is not affected.
  \item With the default settings the log file is renamed to
    \code{d4creg.log.log} (the extension is added twice).
  \item \code{generate_com} still runs the placeholder command
    \code{executable_name arg1} and writes \code{output.txt} in the current
    directory.
  \item \code{saveCORRECT} writes the keyword \code{wavelenght}; to be
    checked against what CORRECT accepts.
  \item \code{getRunningParams} (sample-description file) still splits lines
    on single spaces and keeps the newline on a text value that ends its
    line (the defect corrected in \code{readParam}, C25): \code{<Cmater> V}
    gives \code{"V\textbackslash n"} and a false ``not available'' warning.
    Work-around: end every line with a comment.
  \item \sloppy \code{saveCORRECT} needs attributes (\code{sampleTitle},
    \code{sampleOuterRadius}, \code{sampleTemperature}, \code{sampleDensity},
    \code{samplePackingFraction}, \code{sampleFullness}, \code{containerFile},
    \code{environmentFile}, \code{vanadium}, \code{vanadiumFile}) that
    \code{Measurement} does not provide.
\end{itemize}

\paragraph{Found while writing the docstrings (not yet changed).}
\begin{itemize}
  \item \code{getDiffA}: for an angular step larger than the cell width, the
    split fractions of a cell straddling a bin edge are computed in bin
    units (0.25/0.75 instead of 0.5/0.5 for a $0.25^\circ$ step); exact at
    the default step $0.125^\circ$.
  \item \code{get_bins}/\code{rebin}: different treatment of the two ends of
    the grid (see Appendix~\ref{app:reference}); \code{rebin} still prints a
    debug line; $Q\lambda/4\pi>1$ or a \code{NaN} abscissa raises a
    \code{ValueError} from \code{int(nan)}.
  \item \code{readParam}: \code{writeNumor} forced to 1 by a single-numor run
    stays 1 for the following \code{<run>} lines of the same file; unknown
    tags and extra values on \code{<num>}/\code{<add>} lines are ignored
    without a message.
  \item \code{getEff}/\code{getDec}: bank or cell numbers $\le 0$ are written
    to elements counted from the end of the array.
  \item \code{setSampleInfo} treats the sample as a solid cylinder (no inner
    diameter); the annulus correction C14 exists only in \code{Measurement}.
  \item Two meanings of ``A'': \code{setSampleInfo()["A"]} averages the
    integer nucleon numbers, \code{setVanaInfo()["A"]} is the mass ratio
    $M/m_n$; \code{isotopes.core.getFree} uses the integer value.
  \item \code{Measurement}: an element given in two \code{<SatoNN>} lines is
    counted twice in \code{natoms} but once in \code{atoms}, so the
    concentrations do not sum to 1.
  \item \code{Element}: overriding \code{bcoh}/\code{binc} does not update
    \code{re_bcoh} etc.; misspelt override keys are ignored silently;
    \code{sinc_bound} is 0 when the table has no incoherent length (e.g.\
    natural V, although $\sigma_\mathrm{inc}=5.08$\,b).
  \item \code{DataXYE.plot} saves to \code{../regdata/} followed by the full
    path of the file without extension, which fails for files in
    sub-directories.
  \item Ni calibration: \code{setting_model_d4nifit} fails (\code{IndexError},
    \code{UnboundLocalError}) if fewer than three peaks are found or no point
    lies above the valley; the fit drops all weights if a single error is
    not positive; the fit-statistics box is cut from lmfit's report at
    fixed character positions; \code{reflections_fcc} gives fewer than ten
    reflections for $\lambda\ge1.355$\,\AA{} (not reachable on D4).
  \item \code{stop()} raises in Jupyter, where standard input is not a
    terminal.
\end{itemize}

\paragraph{Engineering.} A \code{pyproject.toml} so that the package
installs as \code{toscana}; logging instead of \code{print}; output and log
paths as explicit arguments instead of fixed relative paths
(\code{../logfiles}, \code{../regdata}); vectorised transforms and binning;
Monte Carlo propagation of \emph{systematic} parameters ($\lambda$,
$2\theta_0$, density, normalisation); \st{P10} keeping the sign of
$\langle b\rangle$ in \code{Measurement.bcoh}.
% =====================================================================
\tsapp{The test suite}\label{app:tests}
% =====================================================================

\tssec{Purpose and principles}

The tests are small programs that call ToScaNA functions on inputs whose
correct answer is known, and compare. They serve three purposes: they
\emph{document} what each function is expected to do; they \emph{prove}
that the corrections of Appendix~\ref{app:audit} work; and they
\emph{protect} the code: after any future change, running them shows
immediately whether something that used to work is broken.

The suite follows a few rules.
\begin{itemize}
  \item \textbf{Known answers.} Each test compares with an analytic result
    (Bragg's law, an exact integral, a derivative), with an independent
    calculation (finite differences, a Monte Carlo simulation, a direct
    SciPy call), or with an invariant (constant input gives constant
    output; values must not change when only errors are requested).
  \item \textbf{Synthetic data only.} No experimental file is read or
    written. Functions that write files (logs, \code{.adat}, figures) are
    run inside a temporary directory that is deleted afterwards. Functions
    that read raw numors are given a fake reader.
  \item \textbf{Written before the fix.} For every correction, the test was
    written first and seen to fail on the original code (19 tests failed
    for the ``group A'' bugs, for example), then the code was corrected.
  \item \textbf{Standard library only.} The suite uses Python's
    \code{unittest}, so nothing needs to be installed. It can also be run
    by \code{pytest} if available. One test uses lmfit and is skipped
    automatically when lmfit is not installed.
\end{itemize}

\tssec{Running the tests}

From the project directory (or any other directory):
\begin{lstlisting}
python tests/run_tests.py                  # all tests (about 35 s)
python tests/run_tests.py -v               # one line per test
python tests/run_tests.py test_fourier     # one file
python tests/run_tests.py test_fourier.TestBackFT.test_q_zero_limit
\end{lstlisting}
The last line of the output is \code{OK} when every test passes. Otherwise
each failure is listed with the name of the test, the line of the
comparison that failed and both values, for example
\begin{lstlisting}
FAIL: test_bin_count (test_reduction.TestRebin.test_bin_count)
AssertionError: 3 != 4
\end{lstlisting}
which reads: the test expected 4 bins and the function returned 3. About
30 of the 35 seconds are spent in the Monte Carlo validation of the error
bars.

\tssec{Organisation of the \texttt{tests} folder}

\begin{longtable}{@{}l p{0.64\textwidth}@{}}
\toprule
File & Role \\
\midrule\endhead
\code{run_tests.py} & Runs the suite (discovers every \code{test_*.py}). \\
\code{_pkg.py} & Makes the project importable as \code{toscana} without installing it: the directory is called \code{ToScaNA-pilot}, so the package is registered under the name \code{toscana} with \code{importlib}. It also keeps the project directory off \code{sys.path} and disables the writing of \code{.pyc} files into the project. \\
\code{helpers.py} & Shared tools: \code{quiet(f, ...)} calls a function with its printed output suppressed; \code{in_tempdir()} runs a block of code in a fresh temporary directory \code{<tmp>/work} (so that the relative paths \code{../logfiles} and \code{../regdata} used by the code land in \code{<tmp>}); \code{synthetic_liquid()} builds the model liquid described below. \\
\code{test_physics.py} & Unit conversions, crystallography, densities, cross sections, thermal reference, vanadium defaults, scattering models. \\
\code{test_reduction.py} & The D4 reduction chain, including an end-to-end run. \\
\code{test_fourier.py} & Fourier transforms, windows, coordination numbers, back-transform. \\
\code{test_fourier_errors.py} & Error bars of the real-space functions and the output file. \\
\code{test_line_shapes_ops.py} & Line shapes with errors, weighted sums, fitting ranges. \\
\code{test_physical_choices.py} & Annulus volume, free cross sections, sigmoid, $Q_0$, weighted Ni fit. \\
\code{test_misc.py} & Isotopes, data files, Ni calibration, \code{showTable}, project layout, text encoding. \\
\bottomrule
\end{longtable}

\tssec{The synthetic liquid}\label{sec:synthetic}

Several tests need an $S(Q)$ whose real-space functions are known exactly.
\code{helpers.synthetic_liquid(rho=0.1)} builds one from a model pair
correlation function on $r\in[0,30)$\,\AA{} in steps of 0.005\,\AA:
\begin{equation}
  g(r) = \begin{cases} 0 & r < 2.5\ \text{\AA} \\[2pt]
   1 + 0.8\,\exp\!\bigl[-(r-2.8)^2/0.02\bigr] & r \ge 2.5\ \text{\AA} \end{cases}
\end{equation}
a hard core of 2.5\,\AA{} and one coordination shell centred at 2.8\,\AA{}
(Gaussian of standard deviation 0.1\,\AA). Its structure factor is
computed on $Q\in[0,30)$\,\AA$^{-1}$ in steps of 0.01\,\AA$^{-1}$ by direct
integration, $S(Q) = 1 + \frac{4\pi\rho}{Q}\int r[g(r)-1]\sin(Qr)\,\mathrm{d}r$,
with $S(0)$ set equal to the first point. The function returns the
two-column array $(Q, S)$ and the exact $r$ and $g$, from which exact
coordination numbers are obtained. For the error-bar tests an error column
$\sigma(Q) = 0.01\,(1+Q/10)$ is added, growing with $Q$ as in real data.

\tssec{What each test checks}

In the tables below, the test name is given without its \code{test_}
prefix, grouped by test class.

\tssub*{\texttt{test\_physics.py}}
\begin{longtable}{@{}p{0.36\textwidth} p{0.58\textwidth}@{}}
\toprule Test & What is checked \\ \midrule\endhead
\multicolumn{2}{@{}l}{\textit{TestConversions}} \\
\tn{round\_\allowbreak{}trip} & \code{q2ang(ang2q(x))} returns $x$ for $2\theta = 1\ldots140^\circ$. \\
\tn{known\_\allowbreak{}value} & $2\theta=60^\circ$, $\lambda=0.5$\,\AA{} gives $Q = 4\pi$. \\
\tn{bad\_\allowbreak{}wavelength} & $\lambda=0$ raises \code{ValueError}. \\
\multicolumn{2}{@{}l}{\textit{TestCrystallography}} \\
\tn{ni\_\allowbreak{}reflections\_\allowbreak{}match\_\allowbreak{}bragg} & The ten first fcc reflections equal $2\arcsin(\lambda\sqrt{h^2+k^2+l^2}/2a)$ for (111)\ldots(333). \\
\tn{default\_\allowbreak{}lattice\_\allowbreak{}is\_\allowbreak{}ni\_\allowbreak{}3\_\allowbreak{}52} & Default lattice parameter is 3.52\,\AA. \\
\tn{zero\_\allowbreak{}angle\_\allowbreak{}shift} & A zero angle shifts all reflections by the same amount. \\
\multicolumn{2}{@{}l}{\textit{TestProperties}} \\
\tn{water\_\allowbreak{}atomic\_\allowbreak{}density} & 1\,g/cm$^3$ and 6\,g/mol/atom give $0.602214076/6$ atoms/\AA$^3$. \\
\tn{cylinder\_\allowbreak{}volume} & $\pi r^2 h$. \\
\tn{free\_\allowbreak{}cross\_\allowbreak{}section} & $\sigma[A/(A+1)]^2$. \\
\tn{absorption\_\allowbreak{}reference\_\allowbreak{}wavelength} & $\sigma_\mathrm{abs}(1.798\,\text{\AA})$ equals the tabulated value; $1/v$ scaling. \\
\tn{thermal\_\allowbreak{}references\_\allowbreak{}consistent} & 1.798\,\AA{} $\leftrightarrow$ 25.30\,meV $\leftrightarrow$ 2200\,m/s; the energy route (\code{scattering_probability}) and the wavelength route (\code{getAbsXS}) give the same absorption. \\
\multicolumn{2}{@{}l}{\textit{TestVanadiumStandard}} \\
\tn{default\_\allowbreak{}density\_\allowbreak{}and\_\allowbreak{}consistency} & Vanadium density 6.11\,g/cm$^3$; the atomic density default of \code{setSampleInfo} equals that of \code{setVanaInfo}. \\
\multicolumn{2}{@{}l}{\textit{TestModels}} \\
\tn{inelastic\_\allowbreak{}limits} & \code{inelastic} tends to lowQ at low $Q$ and to lowQ$\,[A/(A+1)]^2$ at high $Q$. \\
\tn{xs\_\allowbreak{}model\_\allowbreak{}high\_\allowbreak{}energy\_\allowbreak{}limit} & Epithermal cross section tends to the free value at high energy. \\
\bottomrule
\end{longtable}

\tssub*{\texttt{test\_reduction.py}}
\begin{longtable}{@{}p{0.36\textwidth} p{0.58\textwidth}@{}}
\toprule Test & What is checked \\ \midrule\endhead
\multicolumn{2}{@{}l}{\textit{TestCorrections}} \\
\tn{binary\_\allowbreak{}sum\_\allowbreak{}variances} & $10\pm2$ minus $5\pm2$ gives $5\pm\sqrt8$ (C3). \\
\tn{normalise\_\allowbreak{}relative\_\allowbreak{}error\_\allowbreak{}independent\_\allowbreak{}of\_\allowbreak{}monitor} & For monitors $10^5$, $10^6$, $10^7$ the relative error is $\sqrt{1/N+1/N_\mathrm{mon}}$ (C2). \\
\tn{normalise\_\allowbreak{}zero\_\allowbreak{}counts\_\allowbreak{}flagged} & A cell with zero counts becomes \code{NaN}. \\
\tn{dead\_\allowbreak{}time} & $\nu\tau>1$ raises; corrected counts $=N/(1-N\tau/t)$; total = sum of corrected cells (C4). \\
\tn{efficiency\_\allowbreak{}file\_\allowbreak{}without\_\allowbreak{}final\_\allowbreak{}newline} & The last value of an efficiency file is read intact (C27). \\
\tn{numor\_\allowbreak{}list\_\allowbreak{}expansion} & \code{"000020-000022"} expands to three numors; a reversed range raises. \\
\multicolumn{2}{@{}l}{\textit{TestOneNumor}} \\
\tn{poisson\_\allowbreak{}errors\_\allowbreak{}from\_\allowbreak{}raw\_\allowbreak{}counts} & With a fake NeXus reader (400 counts per cell, efficiency 0.8): counts 500, error $\sqrt{(20/0.8)^2 + (500\times10^{-3})^2}$, i.e.\ Poisson on raw counts plus monitor statistics (C5). \\
\multicolumn{2}{@{}l}{\textit{TestGetDiffA}} (two synthetic numors of 576 cells) \\
\tn{constant\_\allowbreak{}input\_\allowbreak{}gives\_\allowbreak{}constant\_\allowbreak{}output} & Constant counts give the same constant on the regrouped grid. \\
\tn{errors\_\allowbreak{}are\_\allowbreak{}combined\_\allowbreak{}not\_\allowbreak{}averaged} & Combined errors never exceed the input error and their median is clearly smaller (C29). \\
\tn{non\_\allowbreak{}commensurate\_\allowbreak{}step} & Step $0.3^\circ$ on $[1,140)$ works; the grid has the length of \code{np.arange} (C6). \\
\tn{out\_\allowbreak{}of\_\allowbreak{}range\_\allowbreak{}cells\_\allowbreak{}are\_\allowbreak{}dropped} & Cells at $-5^\circ$ and $200^\circ$ do not appear anywhere (C7). \\
\multicolumn{2}{@{}l}{\textit{TestRebin}} \\
\tn{bin\_\allowbreak{}count} & $[0,0.3]$ step 0.1 gives 4 bins; $[0,0.25]$ gives 3 (C19). \\
\tn{error\_\allowbreak{}of\_\allowbreak{}mean} & Ten points of $100\pm10$ per bin give $100\pm10/\sqrt{10}$ (C29). \\
\multicolumn{2}{@{}l}{\textit{TestWorkflow}} \\
\tn{read\_\allowbreak{}param\_\allowbreak{}and\_\allowbreak{}d4creg\_\allowbreak{}end\_\allowbreak{}to\_\allowbreak{}end} & A parameter file (with an inline comment and values at the end of lines) is read by \code{readParam}; its \code{<run>} line executes the whole d4creg chain on four fake numors; \code{.adat} and \code{.qdat} are written, start at the requested $Q$ and contain only finite values (C25, C8). \\
\bottomrule
\end{longtable}

\tssub*{\texttt{test\_fourier.py}}
\begin{longtable}{@{}p{0.36\textwidth} p{0.58\textwidth}@{}}
\toprule Test & What is checked \\ \midrule\endhead
\multicolumn{2}{@{}l}{\textit{TestSineFT}} (synthetic liquid, $Q_\mathrm{max}=25$, $r$ step 0.02\,\AA) \\
\tn{column\_\allowbreak{}relations} & $g = 1+G/4\pi\rho r$, RDF $=4\pi\rho r^2 g$, $T = 4\pi\rho r g$. \\
\tn{low\_\allowbreak{}r\_\allowbreak{}limit} & No window: $|g(r)| < 0.02$ for $0.5<r<1.5$\,\AA. \\
\tn{lorch\_\allowbreak{}low\_\allowbreak{}r\_\allowbreak{}limit} & Same with the Lorch window (C1; it was $-0.70$). \\
\tn{coordination\_\allowbreak{}number} & First-shell $N$ within 5\,\% of the exact value. \\
\tn{running\_\allowbreak{}integral\_\allowbreak{}from\_\allowbreak{}zero} & $N(r_i)$ equals a direct Simpson integral up to $r_i$ (C18). \\
\tn{running\_\allowbreak{}integral\_\allowbreak{}grid\_\allowbreak{}not\_\allowbreak{}starting\_\allowbreak{}at\_\allowbreak{}zero} & Same on a grid starting at 1\,\AA; $N(r_0)=0$. \\
\tn{invalid\_\allowbreak{}window\_\allowbreak{}raises} & \code{window=2} raises. \\
\multicolumn{2}{@{}l}{\textit{TestCoordinationNumbers}} \\
\tn{lorch\_\allowbreak{}conserves\_\allowbreak{}coordination\_\allowbreak{}number} & For $Q_\mathrm{max}=15$ and 23.5\,\AA$^{-1}$, and limits 0--3.3 and 2.3--3.3\,\AA, $N$ with the Lorch window agrees within 0.5\,\% with the exact value and with the no-window value (the property discussed in Section~\ref{sec:fourier}). \\
\multicolumn{2}{@{}l}{\textit{TestBackFT}} \\
\tn{all\_\allowbreak{}outputs\_\allowbreak{}finite} & No \code{NaN}/\code{inf} although the grids contain $Q=0$ and $r=0$ (C17). \\
\tn{q\_\allowbreak{}zero\_\allowbreak{}limit} & $S(0) = 1+4\pi\rho\int r^2(g-1)\,\mathrm{d}r$, continuous with the next point. \\
\tn{r\_\allowbreak{}relations} & $g(0)=0$; $T$, RDF consistent with $g$; running integral includes the current point. \\
\tn{auto\_\allowbreak{}cut\_\allowbreak{}without\_\allowbreak{}negative\_\allowbreak{}values\_\allowbreak{}cuts\_\allowbreak{}nothing} & \code{cut=[-1]} on a $g(r)$ that is never negative cuts nothing. \\
\tn{cut\_\allowbreak{}intervals} & A 17-value cut list works; an even-length list raises. \\
\multicolumn{2}{@{}l}{\textit{TestSineHelpers}} \\
\tn{sine\_\allowbreak{}integral} & \code{getSineIntegral} equals \code{scipy.special.sici} to $10^{-10}$; $\mathrm{Si}(0)=0$ (C20). \\
\tn{get\_\allowbreak{}sine\_\allowbreak{}ft\_\allowbreak{}vmax\_\allowbreak{}truncates} & \code{vmax} truncates the integral; default unchanged (P4). \\
\tn{windows\_\allowbreak{}reject\_\allowbreak{}bad\_\allowbreak{}qmax} & $Q_\mathrm{max}\le 0$ raises (P5). \\
\tn{lorch\_\allowbreak{}shape} & $M(0)=1$, $M(Q_\mathrm{max}/2)=2/\pi$, $M(Q_\mathrm{max})=0$. \\
\tn{lorch\_\allowbreak{}array\_\allowbreak{}is\_\allowbreak{}standard\_\allowbreak{}lorch} & \code{LorchArray} equals the scalar \code{Lorch} everywhere. \\
\tn{get\_\allowbreak{}sine\_\allowbreak{}ft\_\allowbreak{}lorch} & \code{getSineFT(w=1)} equals a direct Simpson integral with the Lorch window; \code{w=2} raises. \\
\bottomrule
\end{longtable}

\tssub*{\texttt{test\_fourier\_errors.py}}
\begin{longtable}{@{}p{0.36\textwidth} p{0.58\textwidth}@{}}
\toprule Test & What is checked \\ \midrule\endhead
\multicolumn{2}{@{}l}{\textit{TestErrorColumns}} \\
\tn{values\_\allowbreak{}identical\_\allowbreak{}to\_\allowbreak{}plain\_\allowbreak{}transform} & Requesting errors does not change the values (bit-identical), both windows. \\
\tn{error\_\allowbreak{}relations} & $\sigma_g=\sigma_G/4\pi\rho r$, $\sigma_\mathrm{RDF}=r\sigma_G$, $\sigma_T=\sigma_G$; $\sigma_N(r_0)=0$ and $\sigma_N$ increases with $r$. \\
\tn{missing\_\allowbreak{}error\_\allowbreak{}column\_\allowbreak{}raises} & Two-column input with \code{with_errors=True} raises. \\
\tn{errors\_\allowbreak{}scale\_\allowbreak{}linearly} & Doubling $\sigma_S$ doubles every error. \\
\multicolumn{2}{@{}l}{\textit{TestRunningSimpsonMatrix}} \\
\tn{matches\_\allowbreak{}scipy\_\allowbreak{}on\_\allowbreak{}any\_\allowbreak{}grid} & The fast running-Simpson matrix reproduces SciPy on uniform, offset, random non-uniform and two-point grids to $10^{-12}$. \\
\multicolumn{2}{@{}l}{\textit{TestAgainstMonteCarlo}} \\
\tn{analytic\_\allowbreak{}errors\_\allowbreak{}match\_\allowbreak{}monte\_\allowbreak{}carlo} & 1500 noisy copies of $S(Q)$ are transformed; for $G$, $g$, RDF, $T$, $N$ and both windows the analytic error agrees with the Monte Carlo standard deviation within 8\,\% at every $r$ and within 2\,\% on average (the statistical precision of the Monte Carlo is about 1.8\,\%). \\
\multicolumn{2}{@{}l}{\textit{TestSaveRSCF}} \\
\tn{file\_\allowbreak{}with\_\allowbreak{}errors} & The 11-column file reloads with \code{np.loadtxt} to the precision written; the header contains the user heading, all 11 column names in order, and the description of the errors. \\
\tn{file\_\allowbreak{}without\_\allowbreak{}errors\_\allowbreak{}keeps\_\allowbreak{}six\_\allowbreak{}columns} & Without errors the file has 6 columns. \\
\tn{wrong\_\allowbreak{}shape\_\allowbreak{}raises} & An array with 5 columns raises. \\
\bottomrule
\end{longtable}

\tssub*{\texttt{test\_line\_shapes\_ops.py}}
\begin{longtable}{@{}p{0.36\textwidth} p{0.58\textwidth}@{}}
\toprule Test & What is checked \\ \midrule\endhead
\multicolumn{2}{@{}l}{\textit{TestLineShapeErrors}} (errors compared with central finite differences, $10^{-5}$) \\
\tn{lorentzian\_\allowbreak{}error} & Value equals \code{Lorentzian}; error equals the propagated error (C11). \\
\tn{gaussian\_\allowbreak{}error} & Same for \code{Gaussian_error} (correct in the original code). \\
\tn{lorgau\_\allowbreak{}error} & Same for \code{LorGau_error} versus \code{LorGau} (C12). \\
\multicolumn{2}{@{}l}{\textit{TestWsum2}} \\
\tn{general\_\allowbreak{}weights} & $w_1 y_1 + w_2 y_2$ with $\sqrt{w_1^2e_1^2+w_2^2e_2^2}$. \\
\tn{negative\_\allowbreak{}single\_\allowbreak{}weight\_\allowbreak{}gives\_\allowbreak{}positive\_\allowbreak{}error} & $w_1=0$, $w_2=-1$: error $e_2$ (C13). \\
\tn{inverse\_\allowbreak{}variance\_\allowbreak{}weighting} & $w_1=w_2=0$: weights $e_2^2/(e_1^2+e_2^2)$, error $e_1e_2/\sqrt{e_1^2+e_2^2}$. \\
\tn{invalid\_\allowbreak{}mixing\_\allowbreak{}weight\_\allowbreak{}raises} & $w_2=0$, $w_1=2$ raises (C13). \\
\multicolumn{2}{@{}l}{\textit{TestFittingRanges}} \\
\tn{three\_\allowbreak{}variants\_\allowbreak{}agree} & \code{fittingRange0}, \code{fittingRange}, \code{fit_range} select the same points as a direct mask (P20). \\
\bottomrule
\end{longtable}

\tssub*{\texttt{test\_physical\_choices.py}}
\begin{longtable}{@{}p{0.36\textwidth} p{0.58\textwidth}@{}}
\toprule Test & What is checked \\ \midrule\endhead
\multicolumn{2}{@{}l}{\textit{TestMeasurement}} (a complete \code{Measurement} of water built from a minimal parameter set) \\
\tn{annulus\_\allowbreak{}volume} & 4/6\,mm annulus, 50\,mm: $\pi h(r_o^2-r_i^2)$ (C14). \\
\tn{solid\_\allowbreak{}cylinder\_\allowbreak{}volume\_\allowbreak{}unchanged} & Solid cylinder: $\pi h r_o^2$. \\
\tn{free\_\allowbreak{}cross\_\allowbreak{}section\_\allowbreak{}per\_\allowbreak{}element} & $\sigma_\mathrm{free}=\sum c_i\sigma_i[A_i/(A_i+1)]^2$; for water it is less than half the old mean-mass value (P9). \\
\multicolumn{2}{@{}l}{\textit{TestSetSampleInfo}} \\
\tn{free\_\allowbreak{}cross\_\allowbreak{}sections\_\allowbreak{}per\_\allowbreak{}element} & Same formula for the coherent, incoherent and total free cross sections of \code{setSampleInfo}. \\
\multicolumn{2}{@{}l}{\textit{TestSigmoid}} \\
\tn{limits\_\allowbreak{}midpoint\_\allowbreak{}monotonic\_\allowbreak{}no\_\allowbreak{}overflow} & Limits lowQ and highQ, value $(\mathrm{lowQ}+\mathrm{highQ})/2$ at $Q_0$, monotonic, bounded, no overflow up to $Q=10^4$ (C26). \\
\tn{siglin\_\allowbreak{}adds\_\allowbreak{}linear\_\allowbreak{}term} & \code{siglin} $=$ \code{sigmoidal} $+\,a_1Q$. \\
\tn{inelastic\_\allowbreak{}is\_\allowbreak{}the\_\allowbreak{}same\_\allowbreak{}logistic} & \code{inelastic} equals \code{sigmoidal} with highQ $=$ lowQ$\,[A/(A+1)]^2$. \\
\tn{single\_\allowbreak{}q0\_\allowbreak{}default} & The four functions share $Q_0=7.0$ (P21). \\
\multicolumn{2}{@{}l}{\textit{TestWeightedNiFit}} (fake lmfit recording the call) \\
\tn{weights\_\allowbreak{}are\_\allowbreak{}inverse\_\allowbreak{}errors} & The fit receives weights $1/\sigma$ (P11). \\
\tn{no\_\allowbreak{}valid\_\allowbreak{}errors\_\allowbreak{}means\_\allowbreak{}unweighted} & With $\sigma=-1$ no weights are passed. \\
\bottomrule
\end{longtable}

\tssub*{\texttt{test\_misc.py}}
\begin{longtable}{@{}p{0.36\textwidth} p{0.58\textwidth}@{}}
\toprule Test & What is checked \\ \midrule\endhead
\multicolumn{2}{@{}l}{\textit{TestIsotopes}} \\
\tn{bound\_\allowbreak{}cross\_\allowbreak{}sections\_\allowbreak{}in\_\allowbreak{}barns} & Ni: \code{scoh_bound} $=4\pi(10.3)^2/100$ b, equal to the tabulated $\sigma_\mathrm{coh}$ (C15). \\
\tn{deprecated\_\allowbreak{}spanish\_\allowbreak{}names} & \code{Elemento}, \code{elemento} still return \code{Element}, with a warning. \\
\tn{unknown\_\allowbreak{}isotope\_\allowbreak{}raises} & An unknown name raises (P15). \\
\tn{vanadium\_\allowbreak{}b\_\allowbreak{}consistent\_\allowbreak{}with\_\allowbreak{}sigma} & V: $b=-0.443$\,fm and $4\pi b^2/100=\sigma_\mathrm{coh}$; $\sigma_\mathrm{sca}=\sigma_\mathrm{coh}+\sigma_\mathrm{inc}$ (C16). \\
\tn{vanadium\_\allowbreak{}standard\_\allowbreak{}defaults\_\allowbreak{}match\_\allowbreak{}database} & \code{setVanaInfo} defaults equal the table values. \\
\multicolumn{2}{@{}l}{\textit{TestDataXYE}} \\
\tn{file\_\allowbreak{}without\_\allowbreak{}points\_\allowbreak{}in\_\allowbreak{}peak\_\allowbreak{}range} & A file with $x<4$ loads; no peaks found (C22). \\
\tn{plot\_\allowbreak{}saved\_\allowbreak{}before\_\allowbreak{}show} & The saved figure contains the curve even when \code{show()} closes the figure (C22). \\
\multicolumn{2}{@{}l}{\textit{TestNiCalibration}} \\
\tn{initial\_\allowbreak{}amplitudes\_\allowbreak{}mapping} & Each A$_i$, G$_i$, S$_i$ starts from its own entry of \code{initial}; lattice 3.52 by default, or the value given (C10, P7). \\
\tn{real\_\allowbreak{}fit\_\allowbreak{}lattice\_\allowbreak{}fixed\_\allowbreak{}and\_\allowbreak{}wavelength\_\allowbreak{}scales} & A real lmfit fit of a synthetic Ni pattern ($\lambda=0.4986$\,\AA) recovers $\lambda$ to $10^{-6}$ with the lattice fixed at 3.52\,\AA; with 3.524\,\AA{} the fitted $\lambda$ is multiplied by $3.524/3.52$. \\
\multicolumn{2}{@{}l}{\textit{TestMeasurementShowTable}} \\
\tn{stdout\_\allowbreak{}restored\_\allowbreak{}and\_\allowbreak{}file\_\allowbreak{}closed} & \code{showTable} writes its file and restores the previous output stream (P13). \\
\multicolumn{2}{@{}l}{\textit{TestProjectLayout}} \\
\tn{no\_\allowbreak{}module\_\allowbreak{}shadows\_\allowbreak{}the\_\allowbreak{}standard\_\allowbreak{}library} & No file or folder name is in \code{sys.stdlib_module_names} (C21). \\
\tn{stdlib\_\allowbreak{}importable\_\allowbreak{}from\_\allowbreak{}project\_\allowbreak{}root} & \code{import math, io, lmfit, numpy} works with the project directory as current directory. \\
\tn{old\_\allowbreak{}package\_\allowbreak{}names\_\allowbreak{}still\_\allowbreak{}work\_\allowbreak{}with\_\allowbreak{}warning} & \code{toscana.math}/\code{toscana.io} return the same modules as \code{numerics}/\code{dataio}, with warnings. \\
\multicolumn{2}{@{}l}{\textit{TestSourceHygiene}} \\
\tn{no\_\allowbreak{}mojibake} & No doubly-encoded characters in any source file (C23). \\
\bottomrule
\end{longtable}

\tssec{Adding a test}

To protect a new function or a new correction, add a method whose name
starts with \code{test_} to a class deriving from \code{unittest.TestCase}
in one of the \code{test_*.py} files (or in a new file of that name). A
minimal example:
\begin{lstlisting}
import unittest
import helpers                      # makes `toscana` importable
import numpy as np
from toscana.physics.conversions import ang2q

class TestMyCase(unittest.TestCase):
    def test_backscattering(self):
        # 2theta = 180 deg -> Q = 4 pi / lambda
        self.assertAlmostEqual(float(ang2q(180.0, 0.5)), 8 * np.pi)
\end{lstlisting}
Useful comparisons: \code{assertAlmostEqual(a, b, places=6)} or
\code{delta=...}; \code{np.testing.assert_allclose(x, y, rtol=...)} for
arrays; \code{with self.assertRaises(ValueError): ...} to require an error.
Wrap calls that print with \code{helpers.quiet}, and calls that write files
in \code{with helpers.in_tempdir():}.
% =====================================================================
\tsapp{Module reference}\label{app:reference}
% =====================================================================

\tssec{Overview of the package}

\begin{lstlisting}
toscana/                       (the project directory, imported as `toscana`)
  __init__.py      imports the subpackages; installs the deprecated aliases
  _aliases.py      toscana.math -> numerics, toscana.io -> dataio (warnings)
  calibration/     Ni powder calibration of wavelength and zero angle
  chemistry/       compatibility re-exports (composition helpers)
  data/            ASCII x-y-e data containers with plotting
  dataio/          raw data, corrections, regrouping, parameter files, output
  experiment/      experiment, beam, container, vanadium, sample bookkeeping
  isotopes/        nuclear data table and the Element class
  models/          model functions: vanadium, self-scattering, cross sections, Ni
  numerics/        binning, Fourier transforms, windows, fitting, line shapes
  physics/         unit conversions, crystallography, densities, cross sections
  plotting/        matplotlib helpers
  utils/           dates, dictionaries, subprocess, interactive stop
  workflows/       the d4creg regrouping workflow
  tests/           the test suite (Appendix B)
  docs/            this document
\end{lstlisting}

\paragraph{Dependencies between modules.} Low-level packages
(\code{numerics}, \code{physics}, \code{utils}, \code{isotopes}) do not import
higher-level ones. \code{dataio} uses \code{physics}; \code{workflows} uses
\code{dataio}, \code{numerics}, \code{plotting}; \code{calibration} uses
\code{data}, \code{numerics}, \code{models}, \code{physics}; \code{experiment}
uses \code{dataio}, \code{isotopes}, \code{physics}. One cycle exists:
\code{dataio.parameters} imports \code{workflows.d4creg} (so that the
\code{<run>} line can launch it), which imports \code{dataio}; it works
because \code{dataio/__init__.py} imports \code{parameters} last.

\paragraph{External libraries.} NumPy and SciPy are required.
matplotlib (plots), lmfit (Ni fit) and h5py (NeXus files) are imported only
inside the functions that need them.

\paragraph{Reading the entries.} Each function is given with its full
signature and default values. \fixed{C$n$} marks functions corrected during
the audit (Appendix~\ref{app:audit}), \legacy{} functions kept as in the
original code. Arrays called ``3-column'' have columns $(x, y, e)$.
``Logs'' means that the function appends to \code{../logfiles/d4creg.log},
a path relative to the current working directory.

% ---------------------------------------------------------------------
\tssec{Package \texttt{dataio}: raw data, corrections and files}

\tssub{\texttt{dataio.acquisition} --- reading numors}

\fnx{printRates(\allowbreak{}detector,\allowbreak\ counts,\allowbreak\ time,\allowbreak\ reactor\_power=52)}
Logs the counts and the counting rate of \code{detector} (a label such as
\code{"monitor"}), and the rate per MW of reactor power. A non-positive
time is replaced by $10^{-6}$\,s, a non-positive power by 52\,MW. Prints
nothing. \legacy

\fnx{getRefLines(\allowbreak{}Lines)}
Locates the marker lines of an ILL ASCII numor (lines starting with ten
\code{R}, \code{A}, \code{F} or \code{S}) and returns the tuple
\code{(lineR, lineAL, lineAI, lineFD, lineFS, lineFM, lineSI)}: run block,
first two \code{A} blocks, three \code{F} (float parameter) blocks, first
\code{S} (spectrum) block. Raises \code{ValueError} if a marker is missing.

\fnx{readD4numorASCII(\allowbreak{}filename)}
Reads one numor in ILL ASCII format. Returns \code{(data, counts)}:
\code{data} is a dictionary of metadata (numor, label, user, local contact,
title, subtitle, proposal, start and end times, and every parameter of the
three float blocks under its label, e.g.\ \code{"CntTime (sec)"},
\code{"MonitorCnts"}, \code{"TotalCnts"}, \code{"Theta_lue(deg)"},
\code{"Theta_des(deg)"}, plus \code{ndet}, \code{ncell}); \code{counts} is
an array (ndet, ncell). The counts are read 10 per line, in blocks of 11
lines per detector. Prints a summary and logs the rates. \legacy

\fnx{load\_nxs(\allowbreak{}nxs\_path,\allowbreak\ ncell=64)}
Reads one numor in NeXus format with h5py (group \code{entry0}). Returns
\code{(metadata, counts)} with the same keys as the ASCII reader where they
exist: monitor (\code{monitor/monsum}), time, angles
(\code{instrument/masked2theta_s} value and target), scanned variables
(\code{TotalAcquisitionSpy}/\code{TotalCount} renamed \code{TotalCnts}),
and counts reshaped to (9, ncell). The reactor power is set to 1\,MW
(open item P8). \legacy

\fnx{getOneNumor(\allowbreak{}numorfile,\allowbreak\ runInfo)} \fixed{C5}
Reads and corrects one numor; this is the core of the reduction. Steps:
\begin{enumerate}
  \item read with \code{load_nxs} (if \code{runInfo["dataFormat"]=="Nexus"},
    adding the extension \code{runInfo["ext"][4]}) or \code{readD4numorASCII};
  \item compare the read and requested angles: more than
    \code{cellTolerance[0]} cell width gives a warning, more than
    \code{cellTolerance[1]} raises \code{ValueError} (cell width =
    \code{angular_range_bank}/\code{ncell});
  \item Poisson errors $\sqrt{N}$ on the raw counts;
  \item dead-time correction of monitor and detector
    (\code{runInfo["dead_time"]}); errors scaled by the same factor;
  \item division of counts and errors by the efficiencies (\code{efffile});
  \item normalisation to monitor or time (\code{runInfo["normalisation"]});
  \item angles of every cell from \code{getAngles}, with the bank shifts of
    \code{decfile}.
\end{enumerate}
Returns lists \code{x} (angles), \code{y} (counts), \code{e} (errors),
array \code{q}, list \code{c} (cell identifiers) and the metadata. If
\code{runInfo["writeNumor"]==1} the numor is also written by
\code{saveOneNumor}. Note: \code{q} is computed before the zero angle is
subtracted (open item, Section~\ref{sec:open}). Logs.

\fnx{getNumors(\allowbreak{}runInfo)}
Expands \code{runInfo["numorLst"]} with \code{getNumorFiles}, reads every
numor from \code{runInfo["path_raw"]} with \code{getOneNumor}, and stacks
them in an array of shape (n\_numors, 5, n\_points), the five rows being:
angle $-$ \code{twotheta0}, counts, errors, $Q$, cell identifier. Returns
\code{(numor, head)} with \code{head} the list of metadata dictionaries.
If there is only one numor, \code{writeNumor} is forced to 1. Logs.

\tssub{\texttt{dataio.corrections} --- detector corrections}

\fnx{getDTC\_mcounts(\allowbreak{}monitor,\allowbreak\ time,\allowbreak\ dead\_time=2.4e-06)}
Dead-time correction of the monitor (non-paralysable):
$\nu = \nu_m/(1-\nu_m\tau)$, $\nu_m = $ monitor$/$time. Returns the
corrected counts; raises \code{ValueError} if $\nu_m\tau\ge1$. Logs the factor.

\fnx{getDTC\_dcounts(\allowbreak{}counts,\allowbreak\ total\_counts,\allowbreak\ time,\allowbreak\ dead\_time=7e-06)} \fixed{C4}
Same correction cell by cell for the detector array. Returns the corrected
array and the total counts scaled by $\sum N_\mathrm{corr}/\sum N_\mathrm{raw}$.
Raises \code{ValueError} if any factor is not positive. Logs the overall
and maximum factor. (Whether the dead time should apply per cell or per
bank is open item P2.)

\fnx{getEff(\allowbreak{}filename=\textquotesingle{}effd4c.eff\textquotesingle{},\allowbreak\ ndet=9,\allowbreak\ ncells=64)} \fixed{C27}
Reads relative efficiencies, one line \code{det cell efficiency} per cell
(1-based); lines containing \texttt{\#} are ignored. Values $\le0$ become
\code{NaN}, which removes the cell from the regrouping. Cells not listed
keep efficiency 1. \code{filename="ones"} returns all ones.

\fnx{getDec(\allowbreak{}filename=\textquotesingle{}dec.dec\textquotesingle{},\allowbreak\ ndet=9)} \fixed{C27}
Reads the angular shift of each bank, lines \code{det shift} (degrees).
\code{"zeros"} returns zeros; \code{"manual"} returns a built-in list.

\fnx{getErrorsNumor(\allowbreak{}counts)}
Poisson errors $\sqrt{N}$. Must be called on raw counts.

\fnx{getNormalFactors(\allowbreak{}monitor,\allowbreak\ time)}
Returns the 2$\times$2 array $[[M, \sqrt M],[t, 0.01]]$: normaliser value
and its absolute error for monitor and time.

\fnx{normalise(\allowbreak{}counts,\allowbreak\ errors,\allowbreak\ normal,\allowbreak\ norm\_type=\textquotesingle{}monitor\textquotesingle{},\allowbreak\ norm\_time=120,\allowbreak\ norm\_mon=1000000)} \fixed{C2}
Scales counts to \code{norm_mon} monitor counts (or \code{norm_time}
seconds if \code{norm_type="time"}): $y_n = f y$, $f=$ norm$/M$, and
$\sigma_n^2 = (f\sigma)^2 + (y_n\sigma_M/M)^2$. Zero-count cells become
\code{NaN}. Returns \code{(counts, errors)}.

\fnx{getAngles(\allowbreak{}counts,\allowbreak\ ttheta\_ref,\allowbreak\ zeros,\allowbreak\ cell\_step=0.125,\allowbreak\ det\_step=15.0)}
Angle of every cell: $2\theta = 2\theta_\mathrm{ref} + \mathrm{zeros}[d] + 15^\circ d + 0.125^\circ c$
($d$, $c$ 0-based); cell identifier $1000(d+1) + c + 1$. Returns two
arrays of the shape of \code{counts}.

\tssub{\texttt{dataio.d4creg\_outputs} --- regrouping and D4 files}

\fnx{getDiffA(\allowbreak{}numor,\allowbreak\ head,\allowbreak\ runInfo)} \fixed{C6, C7, C29}
Regroups all cells of all numors onto the regular angular grid
\code{np.arange(amin, amax, step)} of \code{runInfo["angular_scale"]}. If
the requested step is smaller than a cell
(\code{angular_range_bank}/\code{ncell}), the cell width is used. Each cell,
of the cell width and centred at its angle, is split between the (at
most two) bins it overlaps with fractions $f$ (exact only when the step
equals the cell width, the default; for a larger step the fractions are
computed in units of the bin, not of the cell: with $0.25^\circ$ a cell
overlapping two bins equally is split 0.25/0.75; open item); its weight is $w = f\times
M_i/\sum M$ (monitor share of its numor). For each bin
$y = \sum w y_i/\sum w$ and $\sigma = \sqrt{\sum w^2\sigma_i^2}/\sum w$.
Cells with \code{NaN} counts or errors, and parts of cells outside the grid,
are ignored; empty bins are \code{NaN}. Returns an array (10, 5, n\_angles):
index 0 is the sum of all banks, 1--9 the individual banks; rows are angle,
counts, error, sum of weights, $Q$. The angle written is the lower edge of
the bin (open item P1).

\fnx{saveDiffAngle(\allowbreak{}diffA,\allowbreak\ head,\allowbreak\ runInfo)}
Writes two files: \code{<file>.adat} (CORRECT format: header block, then
angle, counts, error of the total diffractogram, positive finite points
only) and \code{<file>.reg} (D4 format: header with the equivalent d4creg
command line, efficiency and shift files, normalisation, zero angle,
wavelength, dead times; then one block per bank with angle, counts, error,
$Q$). Extensions from \code{runInfo["ext"]}.

\fnx{saveDiffQ(\allowbreak{}diffQ,\allowbreak\ head,\allowbreak\ runInfo)}
Writes \code{<file>.qdat} (CORRECT format) from a 3-column $(Q, y, e)$
array, positive finite points only.

\tssub{\texttt{dataio.parameters} --- the d4creg parameter file}

\fnx{readParam(\allowbreak{}parametersFile)} \fixed{C25}
Reads a d4creg parameter file and returns the dictionary \code{runInfo}
used by the whole chain. Every non-empty line must start with \texttt{\#} or
\code{!} (comments) or with a five-character tag \code{<xxx>} followed by
values; text after \texttt{\#} on a tag line is ignored. Recognised tags and
defaults:

\begin{longtable}{@{}l l p{0.46\textwidth}@{}}
\toprule Tag & Values & Meaning (default) \\ \midrule\endhead
\code{<tol>} & w e & angle tolerance in cells for warning / error (0.1, 1.0) \\
\code{<tau>} & $\tau_M$ $\tau_D$ & dead times in s ($2.4\times10^{-6}$, $7\times10^{-6}$) \\
\code{<ext>} & 6 strings & extensions (\code{.reg .adat .qdat .cdat .nxs .log}) \\
\code{<eff>} & file & efficiency file (\code{effd4c.eff}; \code{ones}) \\
\code{<dec>} & file & bank shifts (\code{dec.dec}; \code{zeros}, \code{manual}) \\
\code{<log>} & name & log file name (\code{d4creg.log}) \\
\code{<fmt>} & \code{Nexus}/other & raw data format (\code{Nexus}) \\
\code{<asc>} & min max step & angular grid, degrees (0, 140, 0.125) \\
\code{<qsc>} & min max step & $Q$ grid, \AA$^{-1}$ (0, 23.8, 0.02); prints a recommended step \\
\code{<wri>} & \code{True}/\code{False} & write one file per numor (False) \\
\code{<plo>} & \code{True}/\code{False} & plot diffractograms (False) \\
\code{<wle>} & $\lambda$ & wavelength, \AA{} (0.49857) \\
\code{<zac>} & $2\theta_0$ & zero angle, degrees ($-0.1013$) \\
\code{<rdp>} & path & raw data directory (\code{rawdata/}) \\
\code{<nor>} & type t M & \code{monitor} or \code{time}, normalisation time and monitor (monitor, 80, $10^6$) \\
\code{<out>} & name & base name of output files (\code{vanadium}) \\
\code{<num>} & list & numors: \code{123456} or \code{123456-123460} (replaces the list) \\
\code{<add>} & list & appends to the list of numors \\
\code{<run>} & & runs \code{d4creg(runInfo)} with the values read so far \\
\bottomrule
\end{longtable}
\noindent Side effect: (re)creates \code{../logfiles/d4creg.log}. The file
is read sequentially, so several \code{<run>} lines with changes in between
perform several reductions.

\tssub{\texttt{dataio.running\_params} --- the analysis parameter file}

\fnx{getRunningParams(\allowbreak{}parfile)}
Reads the parameter file of the data-treatment notebooks and returns a
dictionary \code{inputPar} keyed by tag; each value is a tuple whose first
element is a label and the following ones the values. Tags: experiment
(\code{<ins> <pro> <mpr> <lco> <dte> <cyc> <log> <env>}; \code{<dte>} gives
\code{<sta>}/\code{<end>}, \code{<log>} gives \code{<nbk>}/\code{<pag>}), beam
(\code{<zac> <wle> <sli> <fla> <apr>}; \code{<sli>} gives \code{<lft>}/\code{<rgt>},
\code{<fla>} gives \code{<top>}/\code{<bot>}), grids (\code{<asc> <qsc> <rsc>}),
container, absorber, normaliser, environment (\code{<Cmater> <Cshape>},
\code{<Amater> <Ashape>}, \code{<Nmater> <Nshape>}, \code{<Emater> <Eshape>};
shapes: type, inner diameter, outer diameter, height), sample
(\code{<Sdescr> <StempK> <Sato01>... <Sdegcc> <Smassc> <Sshape> <Sfulln> <Stitle>};
each \code{<SatoNN>} is \code{symbol number}, counted into \code{<SNelem>}),
and data files (\code{<CoFile> <AbFile> <NoFile> <EnFile> <SaFile>}). Unknown
materials or shapes print a warning. Its output is the input of
\code{experiment.Measurement}. \legacy

\tssub{\texttt{dataio.loading}, \texttt{saving}, \texttt{numors}, \texttt{correct}, \texttt{\_logfiles}}

\fnx{read\_xye(\allowbreak{}filename)} \fixed{C27}
Reads 2 or 3 columns, ignoring lines containing \texttt{\#}; a missing or
non-numeric (\code{i}, \code{o}) error becomes 0. Returns three 1-D arrays.

\fnx{read\_3col(\allowbreak{}filename)}
Same, returned as one (n, 3) array.

\fnx{saveFile\_xye(\allowbreak{}filename,\allowbreak\ x,\allowbreak\ y,\allowbreak\ e,\allowbreak\ heading)}, \code{saveFile_3col(filename, data, heading)}
Write a 3-column file with \texttt{\#} header lines (formats \code{9.3f},
\code{18.6f}, \code{18.6f}).

\fnx{saveCorrelations(\allowbreak{}filename,\allowbreak\ ar,\allowbreak\ pcf,\allowbreak\ pdf,\allowbreak\ rdf,\allowbreak\ tor,\allowbreak\ run,\allowbreak\ heading)}
Writes six columns $r, G, g$, RDF, $T, N$ from separate arrays.

\fnx{saveRSCF(\allowbreak{}filename,\allowbreak\ fou,\allowbreak\ heading)} \textsf{\footnotesize\colorbox{green!15}{new header, 11 columns}}
Writes the output of \code{sineFT}: 6 columns, or 11 with errors. After
the file name and the user \code{heading}, a block describes each column
(name, formula, units), the meaning and limits of the errors, and a line of
column names \code{r G dG g dg RDF dRDF T dT N dN}. Other shapes raise.

\fnx{saveOneNumor(\allowbreak{}x,\allowbreak\ y,\allowbreak\ e,\allowbreak\ q,\allowbreak\ c,\allowbreak\ head,\allowbreak\ runInfo)}
Writes \code{<numor><ext[3]>} (angle, counts, error, $Q$, cell) for one numor.

\fnx{getNumorFiles(\allowbreak{}numorLst)}
Expands entries \code{"123456"} and \code{"123456-123460"} into a list of
six-digit numors; malformed or reversed ranges raise.

\fnx{generate\_com(\allowbreak{}CORRECT\_template=\textquotesingle{}correct\_template.com\textquotesingle{},\allowbreak\ name=\textquotesingle{}prueba\textquotesingle{},\allowbreak\ rsample=0.5,\allowbreak\ density=0.11,\allowbreak\ ...)}
Fills the placeholders \texttt{\{pyfname\}}, \texttt{\{pyradius\_cm\}}, \ldots{} of a
CORRECT template and writes \code{<name>.com}. It then runs a placeholder
command (open item). \legacy

\fnx{saveCORRECT(\allowbreak{}expt)}
Writes a CORRECT command file \code{<sampleTitle>.com} from an experiment
object (sample, container, background, vanadium, wavelength, beam size;
zero angle written as 0 since already subtracted). \legacy

\fnx{d4creg\_log\_path(\allowbreak{})}
Returns \code{../logfiles/d4creg.log}, creating the folder and file if needed.

% ---------------------------------------------------------------------
\tssec{Package \texttt{workflows}}

\fnx{d4creg(\allowbreak{}runInfo)} \fixed{C8, C9}
The complete regrouping: \code{getNumors} $\to$ \code{getDiffA} $\to$
\code{saveDiffAngle} (+ plots if \code{plotDiff}) $\to$ conversion of the
regrouped angles to $Q$ $\to$ \code{rebin} onto \code{q_scale} with the
actual angular step as width $\to$ masking of points with $y\le0$ and of
their predecessor $\to$ \code{saveDiffQ} (+ plot) $\to$ renaming of the log
file to \code{../logfiles/<logfile>.log}. Prints the equivalent d4creg
command line. Adds \code{runDate}, \code{nbrNumors}, \code{d4creg_cl} to
\code{runInfo}.

% ---------------------------------------------------------------------
\tssec{Package \texttt{numerics}: numerical methods}

\tssub{\texttt{numerics.fourier} --- sine Fourier transforms}

\fnx{sineFT(\allowbreak{}StrFactor,\allowbreak\ nr,\allowbreak\ qmax=23.5,\allowbreak\ density=0.1,\allowbreak\ constant=1.0,\allowbreak\ selfsca=1.0,\allowbreak\ window=0,\allowbreak\ with\_errors=False)} \fixed{C1, C18} \newbeh
The main transform, Section~\ref{sec:conventions}. \code{StrFactor} is an
array with columns $Q$, $S$ (and $\sigma_S$); \code{nr} the $r$ grid;
\code{density} $\rho$ in atoms/\AA$^3$; \code{constant} a multiplicative
factor; \code{selfsca} the value subtracted from $S$ (1 for a normalised
$S(Q)$); \code{window} 0 (none) or 1 (standard Lorch). Returns an
$(N_r, 6)$ array $r, G, g$, RDF, $T, N$, or with \code{with_errors=True}
an $(N_r, 11)$ array with each error beside its function
(Section~\ref{sec:errors}). The $Q$ integral runs over the whole $Q$ grid
with $M(Q)=0$ for $Q\ge$\code{qmax}, so the zeroed points still take part
in the Simpson panels next to $Q_\mathrm{max}$ (\code{getSineFT} instead
cuts the grid at \code{vmax}); the running integral starts at \code{nr[0]}.

\fnx{getSineFT(\allowbreak{}vvec,\allowbreak\ yvec,\allowbreak\ uvec,\allowbreak\ vmax=0,\allowbreak\ c=1.0,\allowbreak\ s=1.0,\allowbreak\ w=0)} \fixed{P4}
Generic transform
$\mathrm{stf}(u) = \frac2\pi c\int_0^{v_\mathrm{max}} [y(v)-s]\,v\,\sin(vu)\,W(v)\,\mathrm{d}v$,
$W$ = 1 (\code{w=0}) or Lorch (\code{w=1}). \code{vmax}$\le0$ means the
end of the data. Returns \code{None} if the lengths differ. Used by
\code{backFT} in both directions.

\fnx{backFT(\allowbreak{}qu,\allowbreak\ ar,\allowbreak\ pdf,\allowbreak\ density,\allowbreak\ cut)} \fixed{C17}
Cuts $g(r)$ (\code{pdf} on grid \code{ar}) according to \code{cut}
(Section~\ref{sec:fourier}), transforms to $S(Q)$ on \code{qu} and back.
Returns the six arrays $(S, G, g, \mathrm{RDF}, T, N)$. Typical use: remove
the unphysical oscillations of $g(r)$ below the first peak and see their
effect on $S(Q)$.

\fnx{getSineIntegral(\allowbreak{}x)} \fixed{C20}
$\mathrm{Si}(x)=\int_0^x \sin t/t\,\mathrm{d}t$ (\code{scipy.special.sici}).

\noindent Internal helpers: \code{_simpson_weights(x)} (the weights $w_k$ of
SciPy's Simpson rule), \code{_sine_ft_errors(...)} (the error
propagation), \code{_running_simpson_matrix(x)} (matrix $C$ of the running
integral, with self-check), \code{_apply_cut(ar, pdf, cut)}.

\tssub{\texttt{numerics.windows} --- modification functions}

\fnx{Lorch(\allowbreak{}q,\allowbreak\ qmax=23.5)} \fixed{P5}
Standard Lorch value at one point: $\sin(\pi q/q_\mathrm{max})/(\pi q/q_\mathrm{max})$,
1 at 0, 0 outside $[0, q_\mathrm{max})$. Raises if $q_\mathrm{max}\le0$.

\fnx{LorchArray(\allowbreak{}q,\allowbreak\ qmax)} \textsf{\footnotesize\colorbox{green!15}{new}}
The same for an array (replaces the removed, rescaled \code{LorchN}).

\tssub{\texttt{numerics.binning} --- rebinning}

\fnx{get\_xlim(\allowbreak{}xmin,\allowbreak\ xmax,\allowbreak\ dbin)} \fixed{C19}
Bin edges and centres: centres $x_\mathrm{min} + i\,\delta$, edges at
$\pm\delta/2$; returns \code{(nb, x_lim, x_bin)}.

\fnx{get\_bins(\allowbreak{}x,\allowbreak\ xdel,\allowbreak\ x\_lim)}
Bins covered by a rectangle of width \code{xdel} centred at \code{x}, and
the fraction of the rectangle in each. A rectangle that starts below the
first edge is rejected entirely (empty lists). At the upper end the
treatment differs: a rectangle is rejected only if it ends one full bin
beyond the last edge; otherwise its excess goes to an overflow index that
\code{rebin} discards, so a point lying entirely in the bin-width beyond
the last edge is dropped silently and a straddling point contributes only
its overlap.

\fnx{rebin(\allowbreak{}xdel,\allowbreak\ wlength,\allowbreak\ data,\allowbreak\ xmin,\allowbreak\ xmax,\allowbreak\ dbin)} \fixed{C29}
Rebins a 3-column array onto the grid of \code{get_xlim}. Each point is a
rectangle of width \code{xdel} (in the units of $x$ if \code{wlength}$\le0$;
if \code{wlength}$>0$, $x$ is $Q$ and \code{xdel} an angular width in degrees,
converted to $\Delta Q = \frac{2\pi}{\lambda}\cos\theta\,\Delta(2\theta)\frac{\pi}{180}$).
Weighted mean per bin, error $\sqrt{\sum f^2\sigma^2}/\sum f$. Points with
\code{NaN} are skipped; empty bins are 0. Prints a debug line.

\tssub{\texttt{numerics.operations} --- combining data sets}

\fnx{binary\_sum(\allowbreak{}w1,\allowbreak\ data1,\allowbreak\ w2,\allowbreak\ data2)} \fixed{C3}
$w_1y_1+w_2y_2$, error $\sqrt{w_1^2e_1^2+w_2^2e_2^2}$, for two 3-column
arrays of the same shape (otherwise prints and returns \code{None}).

\fnx{wsum2(\allowbreak{}w1,\allowbreak\ data1,\allowbreak\ w2,\allowbreak\ data2)} \fixed{C13}
Four modes: $w_1=w_2=0$: inverse-variance weighted mean; $w_1=0$: $w_2y_2$
with error $|w_2|e_2$; $w_2=0$: mixture $w_1y_1+(1-w_1)y_2$, requires
$0<w_1\le1$ (raises otherwise); else $w_1y_1+w_2y_2$.

\fnx{ratio(\allowbreak{}y1,\allowbreak\ e1,\allowbreak\ y2,\allowbreak\ e2)}
$y_1/y_2$ with error $\sqrt{(y_2e_1)^2+(y_1e_2)^2}/y_2^2$; returns $0\pm0$
where $y_2=0$. Lists, truncated to the shortest input. \legacy

\tssub{\texttt{numerics.line\_shapes} --- peak and step functions}

Constants \code{SIGMOID_Q0 = 7.0} and \code{SIGMOID_DQ = 2.4}
(\AA$^{-1}$): default centre and width of the low-$Q$/high-$Q$ transition.

\fnx{sigmoidal(\allowbreak{}Q,\allowbreak\ lowQ=0.4,\allowbreak\ highQ=0.2,\allowbreak\ Q0=SIGMOID\_Q0,\allowbreak\ dQ=SIGMOID\_DQ)} \fixed{C26} \newbeh
Logistic step $(\mathrm{lowQ}+\mathrm{highQ}\,\delta)/(1+\delta)$, $\delta = e^{(Q-Q_0)/\mathrm{d}Q}$.

\fnx{siglin(\allowbreak{}Q,\allowbreak\ lowQ=0.4,\allowbreak\ highQ=0.2,\allowbreak\ Q0=SIGMOID\_Q0,\allowbreak\ dQ=SIGMOID\_DQ,\allowbreak\ a1=0.001)} \newbeh
\code{sigmoidal(Q)} $+\;a_1Q$.

\fnx{step(\allowbreak{}q,\allowbreak\ qmax=23.5)} \fixed{P5}
1 for $0\le q<q_\mathrm{max}$, else 0; raises if $q_\mathrm{max}\le0$.

\fnx{Lorentzian(\allowbreak{}x,\allowbreak\ A,\allowbreak\ x0,\allowbreak\ gamma,\allowbreak\ bckg)}
$A(\gamma/2)^2/[(x-x_0)^2+(\gamma/2)^2] + b$: height $A$, FWHM $\gamma$.

\fnx{Lorentzian\_error(\allowbreak{}x,\allowbreak\ A=1,\allowbreak\ x0=0,\allowbreak\ gamma=1,\allowbreak\ bckg=0,\allowbreak\ e\_A=0,\allowbreak\ e\_x0=0,\allowbreak\ e\_gamma=0,\allowbreak\ e\_bckg=0)} \fixed{C11}
Value and propagated error of \code{Lorentzian}.

\fnx{Gaussian(\allowbreak{}x,\allowbreak\ A,\allowbreak\ x0,\allowbreak\ sigma,\allowbreak\ bckg)}
Area-normalised Gaussian $\frac{A}{\sigma\sqrt{2\pi}}e^{-(x-x_0)^2/2\sigma^2}+b$ (area $A$).

\fnx{Gaussian\_error(\allowbreak{}x,\allowbreak\ A=1,\allowbreak\ x0=0,\allowbreak\ sigma=1,\allowbreak\ bckg=0,\allowbreak\ e\_A=0,\allowbreak\ e\_x0=0,\allowbreak\ e\_sigma=0,\allowbreak\ e\_bckg=0)}
Value and propagated error of \code{Gaussian}.

\fnx{GaussianA(\allowbreak{}x,\allowbreak\ A,\allowbreak\ x0,\allowbreak\ sigma,\allowbreak\ asym=0)}
Area-normalised Gaussian (area $A$, no background) multiplied by
$e^{-\mathrm{asym}(x-x_0)}$ (asymmetric peak, used for the Ni reflections).
For $\mathrm{asym}\neq0$ the area becomes $A\,e^{\mathrm{asym}^2\sigma^2/2}$
and the maximum moves to $x_0-\mathrm{asym}\,\sigma^2$.

\fnx{LorGau(\allowbreak{}x,\allowbreak\ f0=1.0,\allowbreak\ eta=0.5,\allowbreak\ sigma=2.0,\allowbreak\ gamma=2.0,\allowbreak\ bckg=0.0)}
Mixture centred at 0: $f_0[\eta e^{-x^2/2\sigma^2} + (1-\eta)L(x)] + b$,
both components of unit height.

\fnx{LorGau\_error(\allowbreak{}x,\allowbreak\ f0=(\allowbreak{}1.0,\allowbreak\ 0.0),\allowbreak\ eta=(\allowbreak{}0.5,\allowbreak\ 0.0),\allowbreak\ sigma=(\allowbreak{}2.0,\allowbreak\ 0.0),\allowbreak\ gamma=(\allowbreak{}2.0,\allowbreak\ 0.0),\allowbreak\ bckg=(\allowbreak{}0.0,\allowbreak\ 0.0))} \fixed{C12}
Value and error of \code{LorGau}; each argument is a pair (value, error).

\tssub{\texttt{numerics.fitting}, \texttt{polynomials}, \texttt{peak\_finding}, \texttt{signal\_processing}}

\fnx{get\_chi(\allowbreak{}y1,\allowbreak\ y2)}
Root-mean-square difference $\sqrt{\langle(y_1-y_2)^2\rangle}$.

\fnx{fit\_and\_find\_extremum(\allowbreak{}x,\allowbreak\ y)}
Fits a parabola and returns the extremum $x=-b/2a$, its value and the
fitted curve.

\fnx{fittingRange0(\allowbreak{}xmin,\allowbreak\ xmax,\allowbreak\ ymin,\allowbreak\ ymax,\allowbreak\ abcissas,\allowbreak\ ordinates,\allowbreak\ errors)}, \code{fittingRange(...)}, \code{fit_range(xmin, xmax, ymin, ymax, data)} \fixed{P20}
Select the points inside the rectangle $[x_\mathrm{min},x_\mathrm{max}]\times[y_\mathrm{min},y_\mathrm{max}]$.
\code{fittingRange} returns three $(n,1)$ arrays, the other two one
$(n,3)$ array. Warn below 10 points.

\fnx{polyQ1(\allowbreak{}x,\allowbreak\ a0,\allowbreak\ a1)}, \code{polyQ2(x, a0, a1, a2)}, \code{polyQ4(x, a0, ..., a4)}
Polynomials $\sum a_i x^i$ (Horner scheme).

\fnx{find\_minimum\_within\_range(\allowbreak{}x,\allowbreak\ y,\allowbreak\ x\_min,\allowbreak\ x\_max)}
$(x, y)$ of the minimum of $y$ in $[x_\mathrm{min}, x_\mathrm{max}]$ (raises
on an empty range).

\fnx{find\_peaks\_in\_range(\allowbreak{}x,\allowbreak\ y,\allowbreak\ x\_min,\allowbreak\ x\_max)}
Local maxima (\code{scipy.signal.find_peaks}) whose height divided by the
minimum of the range exceeds 3.5; returns positions and these ratios.
\legacy

\fnx{smooth\_curve(\allowbreak{}x,\allowbreak\ y,\allowbreak\ smoothing\_factor)}
Centred moving average over $2\lfloor w/2\rfloor+1$ points,
$w = \lfloor\text{factor}\times n\rfloor$, shrinking at the edges; \code{x}
is not used.

% ---------------------------------------------------------------------
\tssec{Package \texttt{physics}}

\tssub{\texttt{physics.conversions}, \texttt{geometry}}

\fnx{ang2q(\allowbreak{}x,\allowbreak\ wlength=0.5)}, \code{q2ang(x, wlength=0.5)}
$Q = \frac{4\pi}{\lambda}\sin(\pi\,2\theta/360)$ and its inverse; raise if
$\lambda\le0$.

\fnx{getCylVolume(\allowbreak{}diameter=5.0,\allowbreak\ height=50.0)}
$\pi(d/2)^2h$ (mm$^3$ for mm input).

\tssub{\texttt{physics.crystallography}}

Constant \code{NI_LATTICE_PARAMETER = 3.52} (\AA).

\fnx{sf\_fcc(\allowbreak{}h,\allowbreak\ k,\allowbreak\ l)}, \code{sf_bcc(h, k, l)}, \code{sf_sc(h, k, l)}
Geometric structure factors $\sum_j e^{-i\pi(\ldots)}$ of the fcc (4 or 0),
bcc (2 or 0) and simple cubic (1) lattices.

\fnx{reflections\_fcc(\allowbreak{}wavelength=0.5,\allowbreak\ twotheta0=0.0,\allowbreak\ lattice=NI\_LATTICE\_PARAMETER)}
Angles $2\theta+2\theta_0$ (degrees) of the allowed fcc reflections in
increasing order, followed by the mirrored set $180^\circ-2\theta-2\cdot2\theta_0$.

\tssub{\texttt{physics.properties}}

Constants \code{THERMAL_WAVELENGTH = 1.798} (\AA) and
\code{THERMAL_ENERGY_MEV} $=81.8042/1.798^2 = 25.30$ (meV).

\fnx{getAtomicDensity(\allowbreak{}density=1.0,\allowbreak\ molarMass=6.0)}
$\rho = 0.602214076\,d/\bar M$ atoms/\AA$^3$ ($\bar M$ per atom). A
non-positive density gives the value of water (0.1004) with a message.

\fnx{getDensity(\allowbreak{}atomic\_density=0.1,\allowbreak\ molarM=6.0)}
Inverse of the above.

\fnx{getMassNumber(\allowbreak{}molarMass=6.0)}
$A = M/m_n$, $m_n = 1.0086649$\,u.

\fnx{getFreeXS(\allowbreak{}BoundXS=5.08,\allowbreak\ A=50.9415)}
$\sigma_\mathrm{bound}[A/(A+1)]^2$ for one nucleus, $A=m/m_n$. Note: the
default \code{A=50.9415} is the molar mass of vanadium, not its mass ratio
50.50 (0.03\,\% effect on the factor).

\fnx{getAbsXS(\allowbreak{}AbsXS=5.08,\allowbreak\ wavelength=THERMAL\_WAVELENGTH)} \fixed{C30}
$\sigma_\mathrm{abs}\lambda/1.798$.

\fnx{getNofAtoms(\allowbreak{}atoms)}, \code{getConcentrations(atoms)}, \code{AtomicAvg(concentration, magnitude)}
Number of atoms of a formula unit given as \texttt{\{"Si": 1, "O": 2\}};
concentrations $c_i = n_i/\sum n$; average $\sum c_i m_i$ of a property
given as a dictionary.

% ---------------------------------------------------------------------
\tssec{Package \texttt{models}}

\fnx{inelastic(\allowbreak{}x,\allowbreak\ A=51.0,\allowbreak\ lowQ=0.4,\allowbreak\ Q0=SIGMOID\_Q0,\allowbreak\ dQ=SIGMOID\_DQ)}
Inelasticity (Placzek-like) decrease of the self scattering:
$\mathrm{lowQ}\,[1+(A/(1+A))^2\delta]/(1+\delta)$, from the bound value at
low $Q$ to the free value at high $Q$.

\fnx{self024(\allowbreak{}x,\allowbreak\ self0,\allowbreak\ q2,\allowbreak\ q4)}
$s_0 + q_2Q^2 + q_4Q^4$, an empirical self-scattering curve to fit.

\fnx{vanaQdep(\allowbreak{}x,\allowbreak\ a0,\allowbreak\ a1,\allowbreak\ a2,\allowbreak\ A=51.0,\allowbreak\ lowQ=0.4,\allowbreak\ Q0=SIGMOID\_Q0,\allowbreak\ dQ=SIGMOID\_DQ)}
Vanadium $Q$ dependence: \code{polyQ2} (resolution, attenuation trends)
$\times$ \code{inelastic}.

\fnx{XS\_model(\allowbreak{}Eval,\allowbreak\ Eeff,\allowbreak\ composition\_vec,\allowbreak\ bound\_xs\_vec,\allowbreak\ A\_vec)}
Epithermal total scattering cross section
$\sum c_i\sigma_i[A_i/(A_i+1)]^2[1+E_\mathrm{eff}/(2A_iE)]$.

\fnx{scattering\_probability(\allowbreak{}E0,\allowbreak\ scat\_xs,\allowbreak\ composition\_vec,\allowbreak\ abs\_xs\_vec)}
$1-\sigma_a/(\sigma_a+\sigma_s)$ with $\sigma_a = \sum c_i\sigma_{a,i}/\sqrt{E_0/E_\mathrm{th}}$ ($E_0$ in meV).

\fnx{niPeaks10(\allowbreak{}x,\allowbreak\ I0,\allowbreak\ slope,\allowbreak\ quad,\allowbreak\ wavelength,\allowbreak\ twotheta0,\allowbreak\ A0..A9,\allowbreak\ G0..G9,\allowbreak\ S0..S9,\allowbreak\ lattice=NI\_LATTICE\_PARAMETER)}
Ni powder model: background $I_0 + \mathrm{slope}\,x + \mathrm{quad}\,x^2$
plus ten \code{GaussianA} peaks (areas A$_i$, widths S$_i$, asymmetries
G$_i$) at the first ten fcc reflections for the given wavelength, zero
angle and lattice parameter. $x$ is $2\theta$ in degrees.

% ---------------------------------------------------------------------
\tssec{Package \texttt{calibration}}

\fnx{setting\_model\_d4nifit(\allowbreak{}nifile,\allowbreak\ zac\_ini=0.0,\allowbreak\ lattice=NI\_LATTICE\_PARAMETER)}
Prepares the Ni calibration. Reads the regrouped Ni diffractogram
(\code{DataXYE}), estimates $\lambda$ from the first peak, assumed to be
(111): $\lambda = 2(a/\sqrt3)\sin\theta_{111}$; normalises intensities and
errors to the first data point above the midpoint between the second and
third peaks, truncated to whole degrees; builds the
35 initial parameters (background, $\lambda$, $2\theta_0$, ten areas
proportional to typical Ni intensities, asymmetries 0, widths 0.2--0.4);
plots data, initial model and residuals. Returns
\code{(axes, limits, xy_res, xy_ref, xy_fit, planes, initial, nickel)}.

\fnx{fitting\_model\_d4nifit(\allowbreak{}nickel,\allowbreak\ limits,\allowbreak\ initial,\allowbreak\ lattice=NI\_LATTICE\_PARAMETER)} \fixed{C10, P11}
Fits \code{niPeaks10} with lmfit to the points inside \code{limits}
(rectangle in $2\theta$ and intensity), weighted by $1/\sigma$. Free: I0,
slope, wavelength, twotheta0, A$_i$, S$_i$; fixed: quad, G$_i$, lattice.
Returns the lmfit \code{ModelResult}.

\fnx{showing\_results\_d4nifit(\allowbreak{}result,\allowbreak\ nickel,\allowbreak\ axes,\allowbreak\ limits,\allowbreak\ xy\_res,\allowbreak\ xy\_ref,\allowbreak\ xy\_fit,\allowbreak\ planes)}
Prints $\lambda$, $2\theta_0$ and background with errors, a table of the
ten reflections (position, area, width), plots the fit with the fit
statistics, and saves \code{<basename>.png}.

% ---------------------------------------------------------------------
\tssec{Package \texttt{experiment}}

\tssub{\texttt{experiment.measurement}}

\fnx{class Measurement(\allowbreak{}inputPar)} \fixed{C14, P9, P13}
Collects everything about a measurement from the dictionary returned by
\code{getRunningParams}, and reads the five data files (sample, environment,
container, normaliser, absorber) with \code{read_3col}. Main attributes:
experiment and beam information; \code{vslits}, \code{hslits},
\code{beamHeight}, \code{beamWidth}; \code{aScale}, \code{qScale},
\code{rScale}; \code{container}, \code{normaliser}, \code{absorber},
\code{environ} (material, shape, dimensions); sample \code{Shape},
\code{InnerDiam}, \code{OuterDiam}, \code{Height}, \code{Volume}
$=\pi h(r_o^2-r_i^2)/1000$ cm$^3$, \code{Mass}, \code{Density},
\code{EffDensity} $=$ mass/volume, \code{PackingFraction}; composition
\code{atoms}, \code{natoms}, \code{conc}, \code{nuclei}; averages
\code{molarMass}, \code{scoh}, \code{sinc}, \code{ssca}, \code{sabs}
(barns/atom), \code{bcoh}, \code{binc} (moduli of the averaged lengths,
P10); \code{free} (per-element free cross section); \code{b2lowQ}
$=\sigma_\mathrm{sca}/4\pi$, \code{b2highQ} $=\sigma_\mathrm{free}/4\pi$;
\code{sabswl} (absorption at the wavelength); \code{AtomicDensity},
\code{EffAtomicDensity}; data arrays \code{Data}, \code{envData},
\code{conData}, \code{norData}, \code{absData}. A warning is printed if
the container inner diameter differs from the sample outer diameter.
\code{showTable(outfile=None)} prints the identification, beam and grid
parameters (not the sample, container or composition), to a file if given.
\code{<Nmater>} is not read (the normaliser is always vanadium) and
\code{Fullness} is stored but not used.

\tssub{\texttt{experiment.info}}

Constants \code{VANADIUM_DENSITY_GCC = 6.11}, \code{VANADIUM_MOLAR_MASS = 50.9415},
\code{VANADIUM_ATOMIC_DENSITY} $=0.0722305$ atoms/\AA$^3$.

\fnx{setExpInfo(\allowbreak{}Proposal,\allowbreak\ mainProposer,\allowbreak\ experimenters,\allowbreak\ LC,\allowbreak\ otherUsers,\allowbreak\ startDate,\allowbreak\ endDate,\allowbreak\ environment,\allowbreak\ logBook,\allowbreak\ logPage,\allowbreak\ instrument)}
\fnx{setBeamInfo(\allowbreak{}zeroAngle=0,\allowbreak\ wavelength=0.5,\allowbreak\ LohenSlit=2.5,\allowbreak\ GamsSlit=-2.5,\allowbreak\ topFlag=25.0,\allowbreak\ bottomFlag=-25.0)}
\fnx{setCanInfo(\allowbreak{}material=\textquotesingle{}Vanadium\textquotesingle{},\allowbreak\ shape=\textquotesingle{}Cylinder\textquotesingle{},\allowbreak\ outerDiam=5,\allowbreak\ innerDiam=4.8,\allowbreak\ height=60.0)}
\fnx{setBinInfo(\allowbreak{}AngularResolution=0.125,\allowbreak\ AMin=0.0,\allowbreak\ AMax=140.0,\allowbreak\ AStep=0.125,\allowbreak\ QMin=0.0,\allowbreak\ QMax=23.5,\allowbreak\ QStep=0.02,\allowbreak\ RMin=0.0,\allowbreak\ RMax=20.0,\allowbreak\ RStep=0.01)}
\fnx{setNumorInfo(\allowbreak{}totalNumors,\allowbreak\ containerNumors,\allowbreak\ environmentNumors,\allowbreak\ nickelNumors,\allowbreak\ vanadiumNumors,\allowbreak\ absorberNumors,\allowbreak\ sampleNumors)}
Each builds and prints a dictionary describing one part of the
experiment; every argument has a default (e.g.\ \code{LC="Cuello"},
numor ranges \code{("Exp", 0, 1)}), see the docstrings. The dictionaries describe one part of the
experiment (beam width = left $-$ right slit, height = top $-$ bottom
flag; can wall thickness; number of points of each grid; numor ranges
given as (label, first, last)). \legacy

\fnx{setVanaInfo(\allowbreak{}IncXS=5.08,\allowbreak\ CohXS=0.02466,\allowbreak\ ScaXS=5.10466,\allowbreak\ AbsXS=5.08,\allowbreak\ CohSL=-0.443,\allowbreak\ molarM=50.9415,\allowbreak\ NAtoms=1.0,\allowbreak\ Diam=6.08,\allowbreak\ Height=50.0,\allowbreak\ density=VANADIUM\_DENSITY\_GCC)} \fixed{C16, P6}
Dictionary of the vanadium standard: cross sections, $b$, self scattering
at $Q=0$ ($\sigma_\mathrm{inc}/4\pi$), atomic density, mass number, free
incoherent cross section, volume in cm$^3$.

\fnx{setSampleInfo(\allowbreak{}comp\_sample,\allowbreak\ atoms\_sample,\allowbreak\ natoms\_sample,\allowbreak\ elements,\allowbreak\ c\_sample,\allowbreak\ wavelength=0.5,\allowbreak\ beamH=50.0,\allowbreak\ vanavol=1.0,\allowbreak\ vanadens=VANADIUM\_ATOMIC\_DENSITY,\allowbreak\ height=60.0,\allowbreak\ diameter=5.0,\allowbreak\ mass=1.0,\allowbreak\ density=1.0,\allowbreak\ title=\textquotesingle{}sample\textquotesingle{})} \fixed{C30, P9}
Dictionary of the sample: height in the beam, volumes, effective density,
packing and filling fractions, average molar mass, atomic and effective
atomic densities, averaged cross sections, absorption at $\lambda$,
average $b_\mathrm{coh}$, per-element free cross sections, self scattering
at $Q=0$, and the sample/vanadium volume and density ratios used for
normalisation. \code{elements} maps symbols to \code{Element} objects,
\code{c_sample} symbols to concentrations.

% ---------------------------------------------------------------------
\tssec{Package \texttt{isotopes}}

\fnx{class Element(\allowbreak{}iso\_name,\allowbreak\ table=0,\allowbreak\ *{}*kwargs)} \fixed{C15, P15}
Nuclear data of an element (\code{"Ni"}) or isotope (\code{"Li6"},
\code{"H2"} for deuterium) from \code{database.ISOTOPE_DICT} (values from
\emph{Neutron Scattering Fundamentals}, vol.~44, App.~2). Attributes:
\code{symbol}, \code{iso}, \code{name}, \code{A}, \code{Z}, \code{weight}
(u), \code{abundance}, \code{spin}, \code{re_bcoh}, \code{im_bcoh},
\code{re_binc}, \code{im_binc} (fm), \code{sig_coh}, \code{sig_inc},
\code{sig_sca}, \code{sig_abs} (b), complex \code{bcoh}, \code{binc},
\code{scoh_bound} $=4\pi|b_\mathrm{coh}|^2/100$, \code{sinc_bound} (b),
\code{neut} $=A-Z$. Keyword arguments override table values (with
plausibility warnings). \code{table=1} prints the table of the isotope,
\code{table=2} of all isotopes of the element. Unknown names raise
\code{ValueError}. \code{Element.getIsotope(name)} lists the isotopes of an
element name, or all names with \code{"All"}. \code{Elemento},
\code{elemento}: deprecated aliases.

\fnx{getFree(\allowbreak{}isot,\allowbreak\ key)}
Free value $[A/(A+1)]^2\times$\code{key} (complex) for an \code{Element}.

\noindent Validators \code{Positive}, \code{Negative}, \code{SemiWhole},
\code{Top}, \code{Int}, \code{String}, \code{Complex}: print a warning when
an override is implausible. \code{database.NOMBRES} and \code{UNITS}: names
and units of the table columns.

% ---------------------------------------------------------------------
\tssec{Packages \texttt{data}, \texttt{plotting}, \texttt{utils}, \texttt{chemistry}}

\fnx{class DataXYE(\allowbreak{}filename)} \fixed{C22}
Reads a 2- or 3-column file (header lines start with \texttt{\#} or \code{!};
a missing error is stored as $-1$); attributes \code{x}, \code{y},
\code{e}, \code{head}, averages, minima and maxima, \code{peaks_x},
\code{peaks_y} and the minimum \code{xminr}, \code{yminr} in
$x\in[5,40]$ (empty / \code{NaN} if no point there). Methods
\code{plot(file_format=0, xmin, xmax, ymin, ymax)} (axis label from the
extension; saves to \code{../regdata/<basename>.<format>} if a format is
given), \code{show()}, \code{header(lines=1000)}. \code{DataXYE_old} is an
alias.

\fnx{class DataXYE0(\allowbreak{}filename,\allowbreak\ symbol,\allowbreak\ save,\allowbreak\ scale,\allowbreak\ superplot,\allowbreak\ ...)}, \code{class DataXYE2(file)}
Simpler containers read with \code{np.loadtxt}, with \code{plot} methods
(scale \code{"a"}, \code{"q"} or \code{"r"} for the axis label;
\code{superplot} to overlay); \code{DataXYE0.save(filename)}. \legacy

\fnx{plotQ(\allowbreak{}plt,\allowbreak\ ...)}, \code{plotA(plt, ...)}, \code{plotR(plt, ...)}
Set title, labels, limits, grid and legend of the current matplotlib plot
for $Q$, angle or $r$ scales.

\fnx{defColors(\allowbreak{})}, \code{makePlotsA(runInfo, diffA, head, numor)}, \code{makePlotsQ(runInfo, diffQbin, head, numor)}
Colour list and the diagnostic plots of d4creg (individual numors,
individual banks, total in angle; total in $Q$).

\fnx{getDate(\allowbreak{})}
Current date as \code{" Dayname DD/MM/YYYY HH:MM:SS"}.

\fnx{extractDict(\allowbreak{}dictionary,\allowbreak\ key)}, \code{extractAttr(Dict, attr)}
Sub-dictionary of a dictionary of dictionaries; dictionary of one
attribute of a dictionary of objects (e.g.\ \code{extractAttr(elements, "sig_sca")}).

\fnx{execute\_command\_and\_write\_to\_file(\allowbreak{}command,\allowbreak\ output\_file)}
Runs a command, writing its output to a file; errors are printed, not raised.

\fnx{stop(\allowbreak{}message)}
Prints a message and asks whether to continue (interactive sessions only;
raises \code{RuntimeError} otherwise).

\noindent \code{chemistry.composition} re-exports \code{AtomicAvg},
\code{getConcentrations}, \code{getNofAtoms} (from \code{physics.properties})
and \code{extractAttr}, \code{extractDict} (from \code{utils.dicts}) for
compatibility.

% ---------------------------------------------------------------------
\tssec{Top level: \texttt{\_\_init\_\_} and \texttt{\_aliases}}

\code{toscana/__init__.py} imports the subpackages \code{data},
\code{dataio}, \code{experiment}, \code{isotopes}, \code{models},
\code{numerics}, \code{physics}, \code{plotting}, \code{utils} (not
\code{calibration}, \code{chemistry}, \code{workflows}, which must be
imported explicitly) and defines \code{__version__ = "0.1.0"}.
\code{_aliases.py} installs, at the front of Python's import machinery, a
finder that maps \code{toscana.math...} to \code{toscana.numerics...} and
\code{toscana.io...} to \code{toscana.dataio...}, returning the existing
module objects and emitting a \code{DeprecationWarning}; attribute access
\code{toscana.math} is handled the same way. Remove both once all
notebooks use the new names.

% When \input into the manual, stop here; stand-alone, close the document.
\ifdefined\ToScaNAManual\expandafter\endinput\fi
\end{document}

%  after \appendix. The stand-alone preamble is then skipped. In a report
%  or book class the three appendices become chapters, otherwise sections.
%  Packages needed: inputenc(utf8) fontenc(T1) lmodern textcomp amsmath amssymb
%  booktabs longtable array tabularx enumitem xcolor(dvipsnames) listings hyperref
%  (with \PassOptionsToPackage{obeyspaces,spaces}{url} before hyperref)
% =====================================================================
\ifdefined\ToScaNAManual\else
\documentclass[11pt,a4paper]{article}
\usepackage[utf8]{inputenc}
\usepackage[T1]{fontenc}
\usepackage{lmodern}
\usepackage{textcomp}
\usepackage[margin=2.3cm]{geometry}
\usepackage{amsmath,amssymb}
\usepackage{booktabs,longtable,array,tabularx}
\usepackage{enumitem}
\usepackage[dvipsnames]{xcolor}
\usepackage{listings}
% keep spaces inside \code{...} and allow line breaks there
\PassOptionsToPackage{obeyspaces,spaces}{url}
\usepackage[hidelinks]{hyperref}
\urlstyle{tt}
\begin{document}
\fi

% ------------------------------------------------------------- macros
% Heading levels: chapters in report/book manuals, sections in articles.
\ifdefined\chapter
  \providecommand{\tsapp}{\chapter}\providecommand{\tssec}{\section}\providecommand{\tssub}{\subsection}
\else
  \providecommand{\tsapp}{\section}\providecommand{\tssec}{\subsection}\providecommand{\tssub}{\subsubsection}
\fi
% \code{...}: verbatim identifiers, breakable after . _ ( , = / (url package)
\ifdefined\code\else\DeclareRobustCommand{\code}[1]{\path{#1}}\fi
\makeatletter\g@addto@macro\UrlBreaks{\do\_\do\(\do\,\do\=\do\[}\makeatother
\providecommand{\fnx}[1]{\par\medskip\noindent\begingroup\raggedright\hangindent=2em{\color{NavyBlue}\bfseries\ttfamily #1}\par\endgroup\nopagebreak\smallskip}
\providecommand{\tn}[1]{{\raggedright\ttfamily #1}}
\providecommand{\fn}[1]{\par\medskip\noindent{\color{NavyBlue}\textbf{\texttt{\detokenize{#1}}}}\par\nopagebreak\smallskip}
\providecommand{\st}[1]{\textsf{\footnotesize\colorbox{black!8}{#1}}}
\providecommand{\fixed}[1]{\textsf{\footnotesize\colorbox{green!15}{fixed #1}}}
\providecommand{\newbeh}{\textsf{\footnotesize\colorbox{orange!20}{results change}}}
\providecommand{\legacy}{\textsf{\footnotesize\colorbox{black!8}{legacy}}}
\lstset{basicstyle=\ttfamily\small,columns=fullflexible,keepspaces=true,
        frame=single,framerule=0.3pt,backgroundcolor=\color{black!3},
        breaklines=true,showstringspaces=false,aboveskip=6pt,belowskip=6pt}
\setlist{itemsep=2pt,topsep=3pt}
\setlength{\emergencystretch}{3em}

\ifdefined\ToScaNAManual\else
\begin{center}
  {\LARGE\bfseries ToScaNA}\\[4pt]
  {\large Total Scattering Neutron Analysis}\\[10pt]
  {\Large Appendices: code audit, corrections, tests and module reference}\\[10pt]
  {\normalsize Version 0.1.0 (branch \code{fixes/group-a}) --- November 2026}
\end{center}
\bigskip
\noindent\textbf{How to read these appendices.}
Appendix~\ref{app:audit} explains what was examined in the code, what was
found and what was changed, with the physical reasoning and the numbers
that justify each change. Appendix~\ref{app:tests} describes the automatic
tests that now protect the code, and how to run them.
Appendix~\ref{app:reference} is a module-by-module reference of the whole
package: what each function does, its arguments and outputs, the formulas
it implements, which files it reads or writes, and whether it was modified.
Throughout, labels such as \st{C2} or \st{P9} refer to the identifiers of
the audit report (\code{ToScaNA_technical_audit.md}): \st{C} for confirmed
errors, \st{P} for possible issues. The marker \newbeh{} flags changes that
modify numerical results with respect to the original code.

\tableofcontents
\clearpage
\appendix
\fi

% =====================================================================
\tsapp{Audit of the code and corrections}\label{app:audit}
% =====================================================================

\tssec{Context and method}

ToScaNA (\emph{Total Scattering Neutron Analysis}) is a Python package for
the reduction and analysis of neutron total-scattering data from the D4
diffractometer at the ILL. Its present form is a port of a single legacy
script, \code{legacy/lib/toscana.py}, into a package. The port was
written ``parity first'': the docstrings repeatedly state that legacy
behaviour, including known defects, was preserved on purpose so that new
and old results could be compared number by number. The legacy script and
its list of known divergences (\code{KNOWN__DIVERGENCES.md}) were not part
of the project tree that was audited.

The work was done in four stages:
\begin{enumerate}
  \item \textbf{Audit (inspection only).} All 52 Python files were read in
    full (about 6\,000 lines plus the 8\,857-line isotope table). Suspected
    problems were reproduced with small, read-only numerical experiments on
    synthetic data, run from a scratch directory so that nothing in the
    project or in the experimental data was modified. The findings were
    classified as \emph{confirmed errors} (C1--C30), \emph{possible issues}
    (P1--P21) and \emph{recommendations}.
  \item \textbf{Version control and tests.} A git repository was created.
    The first commit is the code exactly as received, so that every later
    change can be inspected with \code{git diff}. A test suite (Appendix~\ref{app:tests})
    was written \emph{before} most of the corrections: each test first
    failed on the original code and passes after the correction.
  \item \textbf{Corrections.} Mechanical bugs were fixed directly; changes
    that modify physical results were made only after a decision by the
    author of the code (density and scattering length of vanadium,
    reference wavelength, lattice parameter of Ni, window function).
  \item \textbf{New functionality.} Exact statistical error bars for
    $G(r)$, $g(r)$, RDF$(r)$, $T(r)$ and $N(r)$.
\end{enumerate}

\tssec{Version history}

All changes are on the git branch \code{fixes/group-a}; the branch
\code{main} contains the first two commits only. Each commit message
explains the change in detail (\code{git log}, \code{git show <id>}).

\begin{longtable}{@{}l p{0.78\textwidth}@{}}
\toprule
Commit & Content \\
\midrule\endhead
\tn{1958d36} & Baseline: the code as received, before any change. \\
\tn{67b3663} & Confirmed errors of the D4 reduction chain (Section~\ref{sec:reduction}). \\
\tn{f98cf26} & Test suite (standard library only). \\
\tn{1a9da55} & Mechanical bugs outside the reduction chain (``group A''). \\
\tn{43354d9} & Vanadium density 6.11\,g/cm$^3$. \\
\tn{6785d7d} & Thermal reference wavelength 1.798\,\AA; vanadium $b_\mathrm{coh}=-0.443$\,fm. \\
\tn{337bbaa} & Isotope class renamed \code{Elemento} $\to$ \code{Element}. \\
\tn{67d188a} & Ni lattice parameter 3.52\,\AA{} by default, user-configurable. \\
\tn{d76b9cd} & Fourier window: none, or the standard Lorch function. \\
\tn{0e2a14f} & Subpackages renamed \code{math}$\to$\code{numerics}, \code{io}$\to$\code{dataio}. \\
\tn{255c508} & Physical choices: annulus volume, free cross sections, sigmoid, $Q_0$, weighted Ni fit. \\
\tn{369e286} & Error bars for $G$, $g$, RDF, $T$ and $N$. \\
\bottomrule
\end{longtable}

\noindent To compare any file with the original: \code{git diff 1958d36 -- numerics/fourier.py}
(note that after commit \code{0e2a14f} the old path is \code{math/fourier.py}).
To return to the original code entirely: \code{git checkout 1958d36}.

% ---------------------------------------------------------------------
\tssec{The data-processing chain}\label{sec:chain}

The figure below summarises what the code does from raw data to
real-space functions, and where each step lives.

\begin{lstlisting}
Raw numors (NeXus .nxs, or ILL ASCII)             dataio.acquisition
  -> counts[9 banks x 64 cells], metadata
Angle check |2theta_read - 2theta_requested|       getOneNumor
Poisson errors on RAW counts  sigma = sqrt(N)      getErrorsNumor
Dead time (non-paralysable, per cell)              dataio.corrections
Relative efficiency  N/eps, sigma/eps              getEff
Normalisation to monitor (or time)                 normalise
Angles: 2theta = 2theta_lue + shift[det]
        + 15 deg*det + 0.125 deg*cell              getAngles
Stack numors, subtract zero angle 2theta0          getNumors
Regroup cells onto a 2theta grid                   dataio.d4creg_outputs.getDiffA
  -> .adat (CORRECT format), .reg (D4 format)      saveDiffAngle
Convert to Q and rebin onto a Q grid               numerics.binning.rebin
  -> .qdat (CORRECT format)                        saveDiffQ
--- external program CORRECT: container, environment, absorption,
    multiple scattering, vanadium normalisation (dataio.correct) ---
S(Q)  ->  sine Fourier transform                   numerics.fourier.sineFT
  -> G(r), g(r), RDF(r), T(r), N(r)  [+ errors]    dataio.saving.saveRSCF
Optional: back-transform with low-r cut            numerics.fourier.backFT
Ni powder calibration of lambda and 2theta0        calibration.d4nifit
\end{lstlisting}

\noindent The whole regrouping chain is driven by a parameter file read by
\code{dataio.parameters.readParam}; the line \code{<run>} in that file
launches \code{workflows.d4creg.d4creg}. Corrections for attenuation,
multiple scattering and container are \emph{not} computed in ToScaNA: they
are delegated to CORRECT, for which ToScaNA writes the input file.

% ---------------------------------------------------------------------
\tssec{Conventions, units and formulas}\label{sec:conventions}

\begin{longtable}{@{}l l p{0.52\textwidth}@{}}
\toprule
Quantity & Unit & Notes \\
\midrule\endhead
Scattering angle $2\theta$ & degree & zero angle $2\theta_0$ subtracted in \code{getNumors} \\
Momentum transfer $Q$ & \AA$^{-1}$ & $Q = \frac{4\pi}{\lambda}\sin\theta$ \\
Distance $r$ & \AA & \\
Wavelength $\lambda$ & \AA & \\
Atomic number density $\rho$ & atoms/\AA$^3$ & $\rho = d\,N_A/\bar M \times 10^{-24} = 0.602214076\,d/\bar M$ \\
Macroscopic density $d$ & g/cm$^3$ & \\
Scattering length $b$ & fm & \\
Cross sections $\sigma$ & barn & $1\,\mathrm{b} = 100\,\mathrm{fm}^2$; $\sigma = 4\pi|b|^2/100$ with $b$ in fm \\
Sample dimensions & mm & volumes converted to cm$^3$ (division by 1000) \\
CORRECT input & cm & radius $=$ diameter$/20$ \\
Counts & per monitor & normalised to \code{norm_mon} monitor counts (default $10^6$) or \code{norm_time} seconds \\
Thermal reference & & $v=2200$\,m/s, $\lambda_\mathrm{th}=1.798$\,\AA, $E_\mathrm{th}=25.30$\,meV \\
\bottomrule
\end{longtable}

\paragraph{Real-space functions.} With $F(Q) = Q[S(Q)-1]$ and a window
(modification) function $M(Q)$ that vanishes for $Q \ge Q_\mathrm{max}$,
\code{sineFT} computes (columns of its output)
\begin{align}
  G(r) &= \frac{2}{\pi}\int_0^{Q_\mathrm{max}} Q\,[S(Q)-1]\,M(Q)\,\sin(Qr)\,\mathrm{d}Q
        \;=\; 4\pi\rho r\,[g(r)-1], \label{eq:G}\\
  g(r) &= 1 + \frac{G(r)}{4\pi\rho r} \quad (g(0) \equiv 0),\\
  \mathrm{RDF}(r) &= 4\pi\rho r^2 g(r) = r\,G(r) + 4\pi\rho r^2,\\
  T(r) &= 4\pi\rho r\, g(r) = G(r) + 4\pi\rho r,\\
  N(r) &= \int_{r_0}^{r} \mathrm{RDF}(r')\,\mathrm{d}r' \qquad (r_0 = \text{first point of the $r$ grid}).
\end{align}
$T(r)$ is defined here without the factor $\langle b\rangle^2$ used by some
authors. The integrals are evaluated with Simpson's rule
(\code{scipy.integrate.simpson}) on the given grids; no FFT is used, so
arbitrary, non-uniform grids are allowed. The back-transform \code{backFT}
uses
\begin{equation}
  S(Q) = 1 + \frac{4\pi\rho}{Q}\int r\,[g(r)-1]\sin(Qr)\,\mathrm{d}r,
  \qquad S(0) = 1 + 4\pi\rho\int r^2[g(r)-1]\,\mathrm{d}r .
\end{equation}
These formulas were checked analytically and numerically during the audit
and are correct.

% ---------------------------------------------------------------------
\tssec{Corrections to the reduction chain}\label{sec:reduction}

All items in this section are in commit \code{67b3663}. Before these
corrections, the \emph{intensities} produced by the chain were correct;
the main defects concerned the \emph{error bars}, and the robustness of
the regrouping. The intensities of regrouped and rebinned data were
verified to be identical to the original code (synthetic data, all
in-range cells).

\paragraph{\fixed{C2} Error of the monitor normalisation.} \newbeh{}
\code{normalise} replaced the counts $y$ by $y_n = f\,y$, $f = N_\mathrm{norm}/N_\mathrm{mon}$,
and only afterwards formed the relative error $e/y_n$ with the
\emph{un-normalised} error $e$. The relative error was therefore multiplied
by $N_\mathrm{mon}/N_\mathrm{norm}$. For a cell with $10^4$ counts, the
relative error should be $\sqrt{10^{-4}+1/N_\mathrm{mon}}\approx 0.010$;
the code gave 0.0033 for $N_\mathrm{mon}=10^5$ (3 times too small) and
0.100 for $N_\mathrm{mon}=10^7$ (10 times too large). It was correct only
when the monitor happened to equal $10^6$. Now
\begin{equation}
  y_n = f y, \qquad \sigma_n^2 = (f\,\sigma)^2 + \Bigl(y_n\,\frac{\sigma_{N_\mathrm{mon}}}{N_\mathrm{mon}}\Bigr)^2 .
\end{equation}
Cells with zero counts are flagged \code{NaN} and are therefore excluded by
the regrouping, which is what the original code did in practice (it
produced $0/0$).

\paragraph{\fixed{C5} Poisson errors.} \newbeh{} The error was computed as
$\sqrt{N}$ \emph{after} dead-time and efficiency corrections, i.e.\
$\sqrt{N_\mathrm{raw}/\varepsilon}$ instead of $\sqrt{N_\mathrm{raw}}/\varepsilon$.
Errors are now taken on raw counts and scaled by exactly the same factors
as the counts (dead-time factor, then $1/\varepsilon$). For
$N_\mathrm{raw}=400$ and $\varepsilon=0.8$ the error per cell is now
$25.0$ (it was $22.4$).

\paragraph{\fixed{C4} Dead-time check.} The test
\code{if (correction.all()) > 0} only checks that no factor is exactly
zero: when the measured rate times the dead time exceeded 1, negative
factors were accepted and produced \emph{negative counts} (reproduced:
$\nu\tau = 1.4$ gave $-5\times10^5$ counts). All factors must now be
positive, otherwise a \code{ValueError} is raised. The corrected total
detector counts were computed with the factor of cell [0,0] only; they now
use $\sum N_\mathrm{corr}/\sum N_\mathrm{raw}$. The log message said
``Monitor DT'' for the detector; corrected.

\paragraph{\fixed{C29} Combination of errors when regrouping and rebinning.} \newbeh{}
Both \code{getDiffA} (cells $\to$ angular grid) and \code{rebin} (angle
$\to$ $Q$ grid) form weighted means $\bar y = \sum w_i y_i/\sum w_i$ (weights
= overlap fraction $\times$ monitor share). The errors were averaged in the
same way, $\sum w_i\sigma_i/\sum w_i$, so that they did not decrease when
points were combined. The error of a weighted mean is
\begin{equation}
  \sigma_{\bar y} = \frac{\sqrt{\sum_i w_i^2\sigma_i^2}}{\sum_i w_i}.
\end{equation}
Example: ten points with $\sigma=10$ gave 10.0; now 3.16. On a synthetic
two-numor regrouping the median error fell from 3.0 to 1.58.

\paragraph{\fixed{C6} Size of the angular grid.} \code{getDiffA} sized its
arrays with \code{int((amax-amin)/step)} but filled them with
\code{np.arange(amin, amax, step)}, which has $\lceil(\cdot)\rceil$
elements. Any step that does not divide the range exactly crashed
(reproduced: \code{(1, 140, 0.3)} gave ``could not broadcast (464,) into
(463,)''). The grid size is now taken from the grid itself.

\paragraph{\fixed{C7} Cells outside the angular range.} Bin indices were
obtained with \code{int()}, which truncates towards zero, and were not
checked. A cell just below \code{amin} was counted in bin 0; a cell further
below gave a \emph{negative} index, which Python interprets as counting
from the end of the array: counts measured at $-5^\circ$ appeared near
$135^\circ$. Cells above \code{amax} raised \code{IndexError}. Indices now
use \code{np.floor} and out-of-range parts of a cell are discarded.

\paragraph{\fixed{C8} Last $Q$ point erased.} After rebinning, points with
$y\le0$ are masked together with the previous point; for the first point
the ``previous'' index $-1$ is the \emph{last} point of the array, which was
erased. Corrected.

\paragraph{\fixed{C9} Width used for the $Q$ rebinning.} The angular width
of each regrouped point was hard-coded to $0.125^\circ$ in \code{d4creg}.
It is now the actual step of the regrouped grid (which \code{getDiffA}
enlarges to the cell width if the requested step is smaller). The width
in $Q$ is $\Delta Q = \frac{2\pi}{\lambda}\cos\theta\,\Delta(2\theta)\,\frac{\pi}{180}$.

\paragraph{\fixed{C19} Number of bins.} \code{get_xlim} computed the number
of bins with \code{int()} of a floating-point ratio:
\code{get_xlim(0, 0.3, 0.1)} gave 3 bins instead of 4 because
$0.4/0.1 = 3.9999\ldots$. A tolerance of $10^{-9}$ is now added before
taking the floor; genuinely non-commensurate ranges keep floor semantics.
The same floor/truncation fix was applied to \code{get_bins}.

\paragraph{\fixed{C25} Parsing of the d4creg parameter file.} Lines were
split with \code{split(" ")}, which kept the newline on the last token and
produced empty tokens for repeated spaces. When a value was the last word
of its line it silently became, e.g., \verb|"Nexus\n"|: the data format
then fell back to the ASCII reader, \verb|"True\n"| disabled writing and
plotting, and numor ranges failed their length check. Lines are now cut
at \texttt{\#} (inline comments allowed) and split on any white space.

\paragraph{\fixed{C27} Last line of input files.} Several readers removed
the last character of each line assuming it was a newline; a final line
without newline lost a digit (an efficiency 1.0234 became 1.023). Lines
are now stripped. Affected: \code{getEff}, \code{getDec}, \code{read_xye}.

\paragraph{\fixed{C3} Weighted sum of two data sets.} \newbeh{}
\code{binary_sum} combined $\sqrt{w_1^2 e_1 + w_2^2 e_2}$ instead of
$\sqrt{w_1^2 e_1^2 + w_2^2 e_2^2}$ (dimensionally wrong; $e_1=e_2=2$ gave
2.0 instead of 2.83).

% ---------------------------------------------------------------------
\tssec{Fourier transform and real-space functions}\label{sec:fourier}

\paragraph{\fixed{C1} Window function.} \newbeh{}
The original code offered a step (truncation) and a ``Lorch'' window that
was rescaled so that $\int_0^{Q_\mathrm{max}} M\,\mathrm{d}Q = Q_\mathrm{max}$.
This rescaling gives $M(0) = \pi/\mathrm{Si}(\pi) = 1.696$ instead of 1
($\mathrm{Si}$ is the sine integral), and multiplies $G(r)$ by 1.70 where
the window is smooth. At low $r$, where $g(r)=0$ and $G(r)=-4\pi\rho r$,
the result was $g(r) = 1 - 1.70 = -0.70$ (reproduced). After discussion, the
rescaled version was removed. \code{sineFT(window=...)} and
\code{getSineFT(w=...)} now accept only
\begin{itemize}
  \item \code{0}: no window (truncation at $Q_\mathrm{max}$);
  \item \code{1}: the standard Lorch function
    $M(Q) = \dfrac{\sin(\pi Q/Q_\mathrm{max})}{\pi Q/Q_\mathrm{max}}$, $M(0)=1$;
\end{itemize}
any other value raises \code{ValueError} (the original code treated any
non-zero value as Lorch). The array function \code{LorchN} was replaced by
\code{LorchArray}.

\paragraph{Why the standard Lorch window preserves coordination numbers.}
$M(Q)$ above is the Fourier transform of a box of half-width
$\Delta = \pi/Q_\mathrm{max}$. Multiplying $F(Q)$ by $M(Q)$ is therefore
exactly equivalent to replacing the (truncated) $G(r)$ by its running
average over $[r-\Delta, r+\Delta]$; this is Lorch's original construction
(1969). A running average with a kernel of unit integral, i.e.\ $M(0)=1$:
\begin{enumerate}
  \item conserves the total area of $G(r)$ (the rescaled version inflated
    all areas by 1.70);
  \item changes an integral between finite limits $[r_1, r_2]$ only through
    edge terms of order $\frac{\Delta^2}{6}[G'(r_2) - G'(r_1)]$, which vanish
    when the limits lie where the curve is flat (the region $g=0$ and the
    first minimum);
  \item has no negative lobes, unlike the kernel $\sin(Q_\mathrm{max} r)/r$
    of plain truncation, so it does not move area by oscillation.
\end{enumerate}
Since $N(r)$ integrates $r\,G(r)$ rather than $G(r)$, an additional term
$\frac{\Delta^2}{3}[G(r_2)-G(r_1)]$ appears; with $Q_\mathrm{max}=23.5$\,\AA$^{-1}$
it is about $0.006\,\Delta G$, negligible. Exact conservation for
\emph{arbitrary} limits is impossible with a finite $Q_\mathrm{max}$: it
would require $M(Q)=1$ up to infinite $Q$. Limits that cut through a peak
always lose part of its area. Numerically, on a model liquid whose exact
coordination number is known (Appendix~\ref{sec:synthetic}):

\begin{center}
\begin{tabular}{@{}c c c c c@{}}
\toprule
$Q_\mathrm{max}$ (\AA$^{-1}$) & Limits (\AA) & Exact $N$ & No window & Standard Lorch \\
\midrule
15.0 & 0.0--3.3 & 10.498 & 10.487 & 10.498 \\
15.0 & 2.3--3.3 & 10.498 & 10.497 & 10.490 \\
23.5 & 0.0--3.3 & 10.498 & 10.499 & 10.498 \\
23.5 & 2.3--3.3 & 10.498 & 10.481 & 10.502 \\
\bottomrule
\end{tabular}
\end{center}
With the rescaled window the same coordination numbers came out 4\,\% too
high from the edge of the first shell, and 30\,\% too low from $r=0$.
\textbf{Practical rule:} integrate from a point where $g(r)=0$ to a
minimum. A related pitfall: on a grid produced by \code{np.arange}, the
value \code{3.3} may be stored as 3.3000000000000003 and a mask
\code{r <= 3.3} then drops the end point; since RDF$(3.3\,\text{\AA})\approx 14$
atoms/\AA, this alone changes $N$ by about 1\,\%. Rounding the grid
(\code{np.round(r, 10)}) avoids it.

\paragraph{\fixed{C18} Running coordination number.} \newbeh{} The slice
end of the running integral was computed as \code{1 + int(r/dr)}, which
assumes that the grid starts at 0 and that the division is exact. Even on
a grid starting at 0 in steps of 0.02\,\AA, rounding (e.g.\
$0.58/0.02 = 28.999\ldots$) stopped the integral one point short at 16 of
400 points, producing small steps in $N(r)$; on a grid not starting at 0,
all values were wrong. The loop index is used now; $N(r_0)=0$.

\paragraph{\fixed{C17} Back-transform.} \newbeh{} \code{backFT} divided by
$Q$ and $r$ without protection, producing \code{NaN}/\code{inf} at
$Q=0$ and $r=0$; the analytic limits are now used ($S(0)$ above;
$g(0)\equiv0$ as in \code{sineFT}; $T = 4\pi\rho r g$ computed directly). Its
running integral excluded the current point, unlike \code{sineFT}; made
consistent. The \code{cut} argument: \code{[-1]} zeroes $g(r)$ up to its
last negative value (previously, when $g$ was never negative, the first two
points were zeroed anyway); otherwise any odd-length list
$[r_0, a_1, b_1, a_2, b_2, \ldots]$ zeroes $r<r_0$ and every interval
$(a_k, b_k)$; even lengths raise (previously lengths above 13 were silently
ignored).

\paragraph{\fixed{C20} Sine integral.} \code{getSineIntegral} integrated
numerically on a grid that stopped at $x(1-10^{-4})$ (relative error
$\sim10^{-4}$: $\mathrm{Si}(1)$ gave 0.945999 instead of 0.946083) and
returned \code{NaN} for $x=0$. It now uses the exact \code{scipy.special.sici}.

\paragraph{\fixed{P4} Truncation in \code{getSineFT}.} With \code{w=0} the
argument \code{vmax} was ignored; the integral now stops at \code{vmax}
(results unchanged when \code{vmax} is the default, the end of the data).

\paragraph{\fixed{P5} Error values used as data.} \code{Lorch} and
\code{step} returned $-999$ for $Q_\mathrm{max}\le 0$ after printing a
message; that value was then silently used as a window. They now raise
\code{ValueError}.

% ---------------------------------------------------------------------
\tssec{Error bars of the real-space functions}\label{sec:errors}

\paragraph{Method.} Commit \code{369e286}. Because the transform is linear
in $S(Q)$, statistical errors can be propagated exactly, without Monte
Carlo noise. Writing the Simpson quadrature of Eq.~\eqref{eq:G} as a sum
over the $Q$ points with weights $w_k$,
\begin{equation}
  G(r_j) = \sum_k W_{jk}\,[S(Q_k)-1], \qquad
  W_{jk} = \frac{2}{\pi}\,c\;Q_k\,M(Q_k)\,\sin(Q_k r_j)\,w_k ,
\end{equation}
($c$ is the argument \code{constant}), and assuming uncorrelated $S(Q)$ points
with errors $\sigma_k$,
\begin{equation}
  \sigma_G^2(r_j) = \sum_k W_{jk}^2\,\sigma_k^2,\qquad
  \sigma_T = \sigma_G,\qquad
  \sigma_\mathrm{RDF} = r\,\sigma_G,\qquad
  \sigma_g = \frac{\sigma_G}{4\pi\rho r}\ (r>0).
\end{equation}
The weights $w_k$ are obtained from SciPy itself (by integrating unit
vectors), so they are exactly those used for the values.

\paragraph{The running integral needs the correlations.} Neighbouring
$G(r)$ points are built from the same $S(Q)$ data and are strongly
correlated; adding their variances as if they were independent would be
badly wrong. Instead $N$ is written directly as a linear function of $S$:
$N(r_i) = \sum_j C_{ij}\, r_j\, G(r_j) + \text{const}$, where $C$ is the
matrix of the running Simpson rule, so that
$\sigma_N^2(r_i) = \sum_k \bigl[(C\,\mathrm{diag}(r)\,W)_{ik}\bigr]^2 \sigma_k^2$.
A first version used trapezoidal weights for $C$ and the Monte Carlo test
showed errors up to 11\,\% wrong without window; the final version uses
exactly the running Simpson rule used for $N(r)$ itself. That matrix is
built in $O(N_r^2)$ operations (0.1\,s for 2000 points): odd-length
prefixes are composite Simpson, extended by one three-point panel at a
time; for even lengths SciPy adds a correction on the last interval that
depends only on the last four points. The construction is compared with
SciPy at run time and rebuilt directly (slowly) if a future SciPy uses a
different rule.

\paragraph{Validation.} $S(Q)$ was perturbed 1500 times with Gaussian noise
of the stated errors, each copy was transformed with the ordinary
\code{sineFT}, and the standard deviation of each output was compared with
the analytic error. Ratio analytic/Monte Carlo, median over $r$ [range]:

\begin{center}\small\setlength{\tabcolsep}{4pt}
\begin{tabular}{@{}l c c c c c@{}}
\toprule
Window & $G$ & $g$ & RDF & $T$ & $N$ \\
\midrule
none    & 1.005 [0.95--1.05] & 1.005 [0.95--1.05] & 1.005 [0.95--1.05] & 1.005 [0.95--1.05] & 1.003 [0.96--1.05] \\
Lorch   & 0.998 [0.96--1.06] & 0.997 [0.96--1.06] & 0.998 [0.96--1.06] & 0.998 [0.96--1.06] & 0.998 [0.95--1.06] \\
\bottomrule
\end{tabular}
\end{center}
The spread of 5--6\,\% is the expected statistical spread of a standard
deviation estimated from 1500 samples.

\paragraph{Usage.}
\begin{lstlisting}
fou = sineFT(SQ, r, qmax=23.5, density=rho, window=1, with_errors=True)
saveRSCF("sample.rscf", fou, ["Sample: ...", "Qmax = 23.5, Lorch"])
\end{lstlisting}
\code{SQ} must have three columns ($Q$, $S$, $\sigma_S$). The output has 11
columns, each error beside its function:
$r, G, \sigma_G, g, \sigma_g, \mathrm{RDF}, \sigma_\mathrm{RDF}, T, \sigma_T, N, \sigma_N$.
The file header describes every column, its formula and units. Without
\code{with_errors} the output keeps its 6 columns.

\paragraph{Limits of these error bars.}
\begin{itemize}
  \item They are \emph{statistical} only. Systematic errors (normalisation,
    density, self-scattering subtraction, truncation) are not included, and
    usually dominate.
  \item They assume uncorrelated $S(Q)$ points. The fractional rebinning
    correlates neighbouring points slightly, so the errors may be somewhat
    underestimated.
  \item Errors at different $r$ are strongly correlated. They are correct
    point by point, but a least-squares fit of $G(r)$ that treats them as
    independent overstates the information contained in the data.
\end{itemize}

% ---------------------------------------------------------------------
\tssec{Line shapes, fitting and the Ni calibration}\label{sec:fitting}

\paragraph{\fixed{C11} \code{Lorentzian_error}.} The $x_0$ term contained
\code{e_x0 * 22} (typing error for \code{e_x0**2}) and lacked one factor of
the normalised peak $p = (\gamma/2)^2/[(x-x_0)^2+(\gamma/2)^2]$; the correct
derivative is $\partial L/\partial x_0 = 2A(x-x_0)\,p^2/(\gamma/2)^2$.
Reproduced error: 2.37 instead of 0.032. Now checked against numerical
derivatives.

\paragraph{\fixed{C12} \code{LorGau_error}.} It evaluated an area-normalised
Gaussian while \code{LorGau} uses a unit-height Gaussian, so the returned
value differed from \code{LorGau} (0.600 instead of 1.000 at $x=0$); the
$\eta$, $\sigma$ and $\gamma$ error terms lacked the factor $f_0$. Rewritten
with explicit derivatives.

\paragraph{\fixed{C13} \code{wsum2}.} With $w_1=0$ the error was $w_2 e_2$,
negative for $w_2<0$; now $|w_2|e_2$. With $w_2=0$ and $w_1\notin(0,1]$ it
printed a message and then crashed; it now raises \code{ValueError}.

\paragraph{\fixed{C10} Initial values of the Ni fit.} The amplitudes A2, A3,
A4 started from \code{initial[6..8]} instead of \code{initial[7..9]}
(shifted by one, so A1 and A2 started equal and \code{initial[9]} was never
used). Only starting values were affected, but a non-linear fit can
converge differently.

\paragraph{\fixed{P11} Weighted Ni fit.} \newbeh{} The fit ignored the data
errors. It now uses weights $1/\sigma$; if no valid errors exist
(two-column files store $\sigma=-1$) it fits without weights and prints a
warning. lmfit keeps \code{scale_covar=True}: the reported uncertainties
are multiplied by $\sqrt{\chi^2_\nu}$ and so also account for model misfit.
\code{setting_model_d4nifit} divides the intensities by the value in the
valley between the second and third peaks; the errors are now divided by
the same value.

\paragraph{\fixed{P7} Ni lattice parameter.} \newbeh{} The value
3.52024\,\AA{} was hard-coded in five places. There is now one default,
\code{physics.crystallography.NI_LATTICE_PARAMETER = 3.52}, and every
function that uses it accepts an argument \code{lattice}:
\begin{lstlisting}
setting_model_d4nifit("ni.adat", lattice=3.5238)
fitting_model_d4nifit(nickel, limits, initial, lattice=3.5238)
\end{lstlisting}
In the fit the lattice parameter is a fixed lmfit parameter: the peak
positions determine only $\lambda/a$, so $\lambda$ and $a$ cannot both be
fitted, and the fitted $\lambda$ scales linearly with $a$ (going from
3.52024 to 3.52 lowers $\lambda$ by 0.007\,\%). \code{showing_results_d4nifit}
reads the value back from the fit result.

\paragraph{\fixed{P20} Fitting ranges.} \code{fittingRange0},
\code{fittingRange} and \code{fit_range} were three copies of the same code;
they now share one helper with identical outputs.

% ---------------------------------------------------------------------
\tssec{Physical constants and sample physics}\label{sec:physics}

\paragraph{\fixed{P6} Vanadium density.} \newbeh{} The defaults used
6.51\,g/cm$^3$ (\code{setVanaInfo}) and the corresponding atomic density
0.0769591\,atoms/\AA$^3$ (\code{setSampleInfo}), typed separately. Both now
derive from one constant, \code{VANADIUM_DENSITY_GCC = 6.11}, giving
0.0722305\,atoms/\AA$^3$. The default sample/vanadium density ratio
\code{DenRatioSV} increases by 6.5\,\%.

\paragraph{\fixed{C16} Vanadium scattering length.} \newbeh{} The isotope
table had $b_\mathrm{coh}=-0.443$\,fm but $\sigma_\mathrm{coh}=0.01838$\,b,
which corresponds to $-0.3824$\,fm; \code{setVanaInfo} used $-0.3824$\,fm.
The value $-0.443$\,fm was adopted everywhere and the cross sections now
follow from it: $\sigma_\mathrm{coh} = 4\pi(0.443)^2/100 = 0.02466$\,b and
$\sigma_\mathrm{sca} = \sigma_\mathrm{coh}+\sigma_\mathrm{inc} = 5.10466$\,b
(were 0.01838 and 5.09838\,b). $\sigma_\mathrm{inc}=5.08$\,b and
$\sigma_\mathrm{abs}=5.08$\,b are unchanged, so the vanadium normalisation
changes very little ($\sigma_\mathrm{sca}$ $+0.12$\,\%). A test checks that
the defaults of \code{setVanaInfo} equal the table values.

\paragraph{\fixed{C30} Thermal reference wavelength.} \newbeh{} Tabulated
absorption cross sections refer to neutrons of 2200\,m/s, i.e.\
$\lambda = h/(m_n v) = 1.798$\,\AA{} and $E = 81.8042/\lambda^2 = 25.30$\,meV.
Three different values were in use: 1.8089\,\AA{} (\code{getAbsXS},
$\sigma_\mathrm{abs}$ 0.6\,\% too small), 1.8\,\AA{} (\code{setSampleInfo})
and 25.3\,meV (\code{scattering_probability}). All now derive from
\code{physics.properties.THERMAL_WAVELENGTH = 1.798} and
\code{THERMAL_ENERGY_MEV}; absorption follows the $1/v$ law
$\sigma_\mathrm{abs}(\lambda) = \sigma_\mathrm{abs}^\mathrm{th}\,\lambda/1.798$.

\paragraph{\fixed{C14} Volume of annular samples.} \newbeh{}
\code{Measurement} computed $\pi h (r_o-r_i)^2$ instead of
$\pi h (r_o^2-r_i^2)$. Solid cylinders ($r_i=0$) are unaffected; for a
4/6\,mm annulus the volume was 5 times too small, and the packing fraction
5 times too large.

\paragraph{\fixed{P9} Free-atom cross sections.} \newbeh{} A bound nucleus
of mass ratio $A = m/m_n$ has a free cross section
$\sigma_\mathrm{free} = \sigma_\mathrm{bound}\,[A/(A+1)]^2$. For a mixture
this factor must be applied element by element,
\begin{equation}
  \sigma_\mathrm{free} = \sum_i c_i\,\sigma_i\left(\frac{A_i}{A_i+1}\right)^2 ,
\end{equation}
whereas the code applied it once with the mean mass. The mean mass hides
how light hydrogen is:

\begin{center}
\begin{tabular}{@{}l r r r@{}}
\toprule
Sample & Per element (b) & Mean mass (old, b) & Change \\
\midrule
H$_2$O & 14.908 & 41.116 & $-63.7$\,\% \\
D$_2$O & 3.508 & 4.907 & $-28.5$\,\% \\
SiO$_2$ & 3.170 & 3.212 & $-1.3$\,\% \\
Ni & 17.880 & 17.880 & 0 \\
\bottomrule
\end{tabular}
\end{center}
Applies to \code{Measurement.free} and to \code{FreeCohXS},
\code{FreeIncXS}, \code{FreeScaXS} of \code{setSampleInfo}. Both use
\code{getMassNumber} (neutron mass 1.0086649\,u); previously the isotope
table value 1.00866\,u was also used.

\paragraph{\fixed{C26} \code{sigmoidal}.} \newbeh{} The function returned
$\delta\,\mathrm{lowQ} + (1-\delta)\,\mathrm{highQ}$ with
$\delta = e^{(Q-Q_0)/\mathrm{d}Q}$: unbounded, with negative weights above
$Q_0$, and running from highQ to lowQ, the opposite of what the parameter
names state. It is now the logistic step
\begin{equation}
  f(Q) = \frac{\mathrm{lowQ} + \mathrm{highQ}\,\delta}{1+\delta},\qquad
  f(0)\approx\mathrm{lowQ},\ f(\infty) = \mathrm{highQ},\ f(Q_0) = \tfrac12(\mathrm{lowQ}+\mathrm{highQ}),
\end{equation}
computed with \code{scipy.special.expit} so that large $|Q-Q_0|/\mathrm{d}Q$
cannot overflow. This is the same shape as \code{models.scattering.inelastic}
(with $\mathrm{highQ} = \mathrm{lowQ}\,[A/(A+1)]^2$). \code{siglin}
ignored its argument \code{a1}; it is now $\texttt{sigmoidal}(Q) + a_1 Q$.
\emph{This interpretation of \code{a1} should be confirmed by the author.}

\paragraph{\fixed{P21} One default transition.} \code{inelastic} used
$Q_0=7.0$\,\AA$^{-1}$ and \code{vanaQdep} 7.4\,\AA$^{-1}$. One pair of
defaults, \code{SIGMOID_Q0 = 7.0} and \code{SIGMOID_DQ = 2.4}
(\code{numerics.line_shapes}), now serves \code{sigmoidal}, \code{siglin},
\code{inelastic} and \code{vanaQdep} (effect on \code{vanaQdep}
$\le 0.16$\,\%).

\paragraph{\fixed{C15} Bound cross sections of \code{Element}.} The
attributes \code{scoh_bound} and \code{sinc_bound} were $4\pi|b|^2$ in fm$^2$
but labelled barns (100 times too large: Ni gave 1333 instead of 13.3).
Divided by 100. They were not used elsewhere in the package.

% ---------------------------------------------------------------------
\tssec{Structure and robustness of the code}\label{sec:structure}

\paragraph{\fixed{C21} Names that hid the standard library.} The
subpackages \code{math/} and \code{io/} have the names of standard Python
modules. When Python was started in the project directory (for example a
notebook opened there), \code{import math} found the project folder
instead; \code{import lmfit} failed with ``partially initialized module
`math' has no attribute `floor'''. They were renamed (with \code{git mv},
keeping their history):
\begin{center}
\code{toscana.math} $\to$ \code{toscana.numerics} \qquad
\code{toscana.io} $\to$ \code{toscana.dataio}
\end{center}
The old dotted names remain importable through a small import hook
(\code{_aliases.py}) and emit a \code{DeprecationWarning}; they return the
\emph{same} module objects (an earlier version of the hook loaded a second
copy of each module, which a test detected). No file or folder of the
project now has the name of a standard-library module, and a test checks
this against \code{sys.stdlib_module_names}.

\paragraph{English naming.} The isotope class \code{Elemento} is now
\code{Element}. \code{Elemento} and \code{elemento} remain available as
deprecated aliases. Other Spanish names in the same module
(\code{NOMBRES}, local variables) were not changed.

\paragraph{Other fixes.}
\begin{itemize}
  \item \fixed{C22} \code{DataXYE} crashed on files without any point in
    $[5,40]$ (the range of its automatic peak search), e.g.\ $r$-space
    files; it saved figures \emph{after} \code{plt.show()}, which gives
    empty images with interactive back-ends.
  \item \fixed{P15} An unknown isotope name now raises \code{ValueError}
    instead of returning a half-built object.
  \item \fixed{P13} \code{Measurement.showTable} replaced \code{sys.stdout}
    with an unclosed file and then restored \code{sys.__stdout__}, which
    broke Jupyter output; it now uses \code{contextlib.redirect_stdout}.
  \item \fixed{C23} Characters such as \AA, $^\circ$ and $\pm$ were stored
    doubly encoded (\AA{} appeared as two wrong characters); repaired. \fixed{C24} The \code{.reg}
    header gave the wavelength in degrees; now \AA.
  \item \st{P19} The audit claimed that \code{smooth_curve} had an
    asymmetric window; this was wrong (even sizes become the next odd
    size). Only its documentation was completed.
\end{itemize}

% ---------------------------------------------------------------------
\tssec{What changes for the user}\label{sec:migration}

\begin{longtable}{@{}p{0.36\textwidth} p{0.58\textwidth}@{}}
\toprule
Before & Now \\
\midrule\endhead
\code{from toscana.math...}, \code{toscana.io...} & \code{toscana.numerics...}, \code{toscana.dataio...} (old names work, with a warning) \\
\code{Elemento("Ni")} & \code{Element("Ni")} (old name works, with a warning) \\
\code{LorchN(Q, Qmax)} & removed; \code{LorchArray(Q, Qmax)} is the standard Lorch \\
\code{sineFT(..., window=1)} rescaled Lorch & standard Lorch, $M(0)=1$; values other than 0, 1 raise \\
\code{sineFT} 6 columns & unchanged; \code{with_errors=True} gives 11 columns \\
\code{saveRSCF} 6 columns, no header & 6 or 11 columns, header describing each column \\
Error bars of \code{.adat}, \code{.reg}, \code{.qdat} & smaller and statistically correct (C2, C5, C29) \\
$N(r)$ & correct at every point (C18) \\
\code{Element.scoh_bound} & in barns (100 times smaller) \\
Vanadium defaults & 6.11\,g/cm$^3$, $b=-0.443$\,fm, $\sigma_\mathrm{coh}=0.02466$\,b \\
$\sigma_\mathrm{abs}(\lambda)$ & $\propto\lambda/1.798$ \\
Free cross sections & per element (large change for H, D) \\
\tn{sigmoidal} & a true logistic step \\
Ni calibration & weighted fit, lattice 3.52\,\AA{} by default, \code{lattice=} argument \\
Silent prints on bad input & \code{ValueError} (unknown isotope, invalid \code{cut}, $Q_\mathrm{max}\le0$, invalid window or weights, negative dead-time factor) \\
\bottomrule
\end{longtable}

% ---------------------------------------------------------------------
\tssec{Open items}\label{sec:open}

\paragraph{Need data or instrument knowledge.}
\begin{itemize}
  \item \st{P1} Angle labels of \code{getDiffA}: the angle written for a bin
    is its lower edge, while a cell centred in the bin contributes fully to
    it, i.e.\ a shift of half a step ($0.0625^\circ$ for $0.125^\circ$). It
    is absorbed by the zero angle if the Ni calibration was regrouped in the
    same way; to be confirmed with a real Ni run.
  \item \st{P2} The dead time is applied to the rate of each \emph{cell};
    for position-sensitive banks it is usually governed by the rate of the
    whole bank or of its electronics.
  \item \st{P3} Data are binned twice (cells $\to 2\theta$ $\to Q$), which
    adds broadening and correlations; a direct cells $\to Q$ rebinning
    would avoid one step.
  \item \st{P8} The reactor power is fixed to 1\,MW for NeXus files (only
    the logged rates per MW are affected); the NeXus field is needed.
  \item \st{P16} The ASCII numor reader assumes 10 values per line and a
    block of 11 lines per detector; a sample file is needed to check the
    case \code{ncell} multiple of 10.
\end{itemize}

\paragraph{Found while writing this appendix (not yet changed).}
\begin{itemize}
  \item In \code{getOneNumor}, the $Q$ column of each numor (used in the
    per-numor \code{.cdat} files) is computed from the angles \emph{before}
    the zero angle is subtracted; the main \code{.qdat} output is computed
    from corrected angles and is not affected.
  \item With the default settings the log file is renamed to
    \code{d4creg.log.log} (the extension is added twice).
  \item \code{generate_com} still runs the placeholder command
    \code{executable_name arg1} and writes \code{output.txt} in the current
    directory.
  \item \code{saveCORRECT} writes the keyword \code{wavelenght}; to be
    checked against what CORRECT accepts.
  \item \code{getRunningParams} (sample-description file) still splits lines
    on single spaces and keeps the newline on a text value that ends its
    line (the defect corrected in \code{readParam}, C25): \code{<Cmater> V}
    gives \code{"V\textbackslash n"} and a false ``not available'' warning.
    Work-around: end every line with a comment.
  \item \sloppy \code{saveCORRECT} needs attributes (\code{sampleTitle},
    \code{sampleOuterRadius}, \code{sampleTemperature}, \code{sampleDensity},
    \code{samplePackingFraction}, \code{sampleFullness}, \code{containerFile},
    \code{environmentFile}, \code{vanadium}, \code{vanadiumFile}) that
    \code{Measurement} does not provide.
\end{itemize}

\paragraph{Found while writing the docstrings (not yet changed).}
\begin{itemize}
  \item \code{getDiffA}: for an angular step larger than the cell width, the
    split fractions of a cell straddling a bin edge are computed in bin
    units (0.25/0.75 instead of 0.5/0.5 for a $0.25^\circ$ step); exact at
    the default step $0.125^\circ$.
  \item \code{get_bins}/\code{rebin}: different treatment of the two ends of
    the grid (see Appendix~\ref{app:reference}); \code{rebin} still prints a
    debug line; $Q\lambda/4\pi>1$ or a \code{NaN} abscissa raises a
    \code{ValueError} from \code{int(nan)}.
  \item \code{readParam}: \code{writeNumor} forced to 1 by a single-numor run
    stays 1 for the following \code{<run>} lines of the same file; unknown
    tags and extra values on \code{<num>}/\code{<add>} lines are ignored
    without a message.
  \item \code{getEff}/\code{getDec}: bank or cell numbers $\le 0$ are written
    to elements counted from the end of the array.
  \item \code{setSampleInfo} treats the sample as a solid cylinder (no inner
    diameter); the annulus correction C14 exists only in \code{Measurement}.
  \item Two meanings of ``A'': \code{setSampleInfo()["A"]} averages the
    integer nucleon numbers, \code{setVanaInfo()["A"]} is the mass ratio
    $M/m_n$; \code{isotopes.core.getFree} uses the integer value.
  \item \code{Measurement}: an element given in two \code{<SatoNN>} lines is
    counted twice in \code{natoms} but once in \code{atoms}, so the
    concentrations do not sum to 1.
  \item \code{Element}: overriding \code{bcoh}/\code{binc} does not update
    \code{re_bcoh} etc.; misspelt override keys are ignored silently;
    \code{sinc_bound} is 0 when the table has no incoherent length (e.g.\
    natural V, although $\sigma_\mathrm{inc}=5.08$\,b).
  \item \code{DataXYE.plot} saves to \code{../regdata/} followed by the full
    path of the file without extension, which fails for files in
    sub-directories.
  \item Ni calibration: \code{setting_model_d4nifit} fails (\code{IndexError},
    \code{UnboundLocalError}) if fewer than three peaks are found or no point
    lies above the valley; the fit drops all weights if a single error is
    not positive; the fit-statistics box is cut from lmfit's report at
    fixed character positions; \code{reflections_fcc} gives fewer than ten
    reflections for $\lambda\ge1.355$\,\AA{} (not reachable on D4).
  \item \code{stop()} raises in Jupyter, where standard input is not a
    terminal.
\end{itemize}

\paragraph{Engineering.} A \code{pyproject.toml} so that the package
installs as \code{toscana}; logging instead of \code{print}; output and log
paths as explicit arguments instead of fixed relative paths
(\code{../logfiles}, \code{../regdata}); vectorised transforms and binning;
Monte Carlo propagation of \emph{systematic} parameters ($\lambda$,
$2\theta_0$, density, normalisation); \st{P10} keeping the sign of
$\langle b\rangle$ in \code{Measurement.bcoh}.
% =====================================================================
\tsapp{The test suite}\label{app:tests}
% =====================================================================

\tssec{Purpose and principles}

The tests are small programs that call ToScaNA functions on inputs whose
correct answer is known, and compare. They serve three purposes: they
\emph{document} what each function is expected to do; they \emph{prove}
that the corrections of Appendix~\ref{app:audit} work; and they
\emph{protect} the code: after any future change, running them shows
immediately whether something that used to work is broken.

The suite follows a few rules.
\begin{itemize}
  \item \textbf{Known answers.} Each test compares with an analytic result
    (Bragg's law, an exact integral, a derivative), with an independent
    calculation (finite differences, a Monte Carlo simulation, a direct
    SciPy call), or with an invariant (constant input gives constant
    output; values must not change when only errors are requested).
  \item \textbf{Synthetic data only.} No experimental file is read or
    written. Functions that write files (logs, \code{.adat}, figures) are
    run inside a temporary directory that is deleted afterwards. Functions
    that read raw numors are given a fake reader.
  \item \textbf{Written before the fix.} For every correction, the test was
    written first and seen to fail on the original code (19 tests failed
    for the ``group A'' bugs, for example), then the code was corrected.
  \item \textbf{Standard library only.} The suite uses Python's
    \code{unittest}, so nothing needs to be installed. It can also be run
    by \code{pytest} if available. One test uses lmfit and is skipped
    automatically when lmfit is not installed.
\end{itemize}

\tssec{Running the tests}

From the project directory (or any other directory):
\begin{lstlisting}
python tests/run_tests.py                  # all tests (about 35 s)
python tests/run_tests.py -v               # one line per test
python tests/run_tests.py test_fourier     # one file
python tests/run_tests.py test_fourier.TestBackFT.test_q_zero_limit
\end{lstlisting}
The last line of the output is \code{OK} when every test passes. Otherwise
each failure is listed with the name of the test, the line of the
comparison that failed and both values, for example
\begin{lstlisting}
FAIL: test_bin_count (test_reduction.TestRebin.test_bin_count)
AssertionError: 3 != 4
\end{lstlisting}
which reads: the test expected 4 bins and the function returned 3. About
30 of the 35 seconds are spent in the Monte Carlo validation of the error
bars.

\tssec{Organisation of the \texttt{tests} folder}

\begin{longtable}{@{}l p{0.64\textwidth}@{}}
\toprule
File & Role \\
\midrule\endhead
\code{run_tests.py} & Runs the suite (discovers every \code{test_*.py}). \\
\code{_pkg.py} & Makes the project importable as \code{toscana} without installing it: the directory is called \code{ToScaNA-pilot}, so the package is registered under the name \code{toscana} with \code{importlib}. It also keeps the project directory off \code{sys.path} and disables the writing of \code{.pyc} files into the project. \\
\code{helpers.py} & Shared tools: \code{quiet(f, ...)} calls a function with its printed output suppressed; \code{in_tempdir()} runs a block of code in a fresh temporary directory \code{<tmp>/work} (so that the relative paths \code{../logfiles} and \code{../regdata} used by the code land in \code{<tmp>}); \code{synthetic_liquid()} builds the model liquid described below. \\
\code{test_physics.py} & Unit conversions, crystallography, densities, cross sections, thermal reference, vanadium defaults, scattering models. \\
\code{test_reduction.py} & The D4 reduction chain, including an end-to-end run. \\
\code{test_fourier.py} & Fourier transforms, windows, coordination numbers, back-transform. \\
\code{test_fourier_errors.py} & Error bars of the real-space functions and the output file. \\
\code{test_line_shapes_ops.py} & Line shapes with errors, weighted sums, fitting ranges. \\
\code{test_physical_choices.py} & Annulus volume, free cross sections, sigmoid, $Q_0$, weighted Ni fit. \\
\code{test_misc.py} & Isotopes, data files, Ni calibration, \code{showTable}, project layout, text encoding. \\
\bottomrule
\end{longtable}

\tssec{The synthetic liquid}\label{sec:synthetic}

Several tests need an $S(Q)$ whose real-space functions are known exactly.
\code{helpers.synthetic_liquid(rho=0.1)} builds one from a model pair
correlation function on $r\in[0,30)$\,\AA{} in steps of 0.005\,\AA:
\begin{equation}
  g(r) = \begin{cases} 0 & r < 2.5\ \text{\AA} \\[2pt]
   1 + 0.8\,\exp\!\bigl[-(r-2.8)^2/0.02\bigr] & r \ge 2.5\ \text{\AA} \end{cases}
\end{equation}
a hard core of 2.5\,\AA{} and one coordination shell centred at 2.8\,\AA{}
(Gaussian of standard deviation 0.1\,\AA). Its structure factor is
computed on $Q\in[0,30)$\,\AA$^{-1}$ in steps of 0.01\,\AA$^{-1}$ by direct
integration, $S(Q) = 1 + \frac{4\pi\rho}{Q}\int r[g(r)-1]\sin(Qr)\,\mathrm{d}r$,
with $S(0)$ set equal to the first point. The function returns the
two-column array $(Q, S)$ and the exact $r$ and $g$, from which exact
coordination numbers are obtained. For the error-bar tests an error column
$\sigma(Q) = 0.01\,(1+Q/10)$ is added, growing with $Q$ as in real data.

\tssec{What each test checks}

In the tables below, the test name is given without its \code{test_}
prefix, grouped by test class.

\tssub*{\texttt{test\_physics.py}}
\begin{longtable}{@{}p{0.36\textwidth} p{0.58\textwidth}@{}}
\toprule Test & What is checked \\ \midrule\endhead
\multicolumn{2}{@{}l}{\textit{TestConversions}} \\
\tn{round\_\allowbreak{}trip} & \code{q2ang(ang2q(x))} returns $x$ for $2\theta = 1\ldots140^\circ$. \\
\tn{known\_\allowbreak{}value} & $2\theta=60^\circ$, $\lambda=0.5$\,\AA{} gives $Q = 4\pi$. \\
\tn{bad\_\allowbreak{}wavelength} & $\lambda=0$ raises \code{ValueError}. \\
\multicolumn{2}{@{}l}{\textit{TestCrystallography}} \\
\tn{ni\_\allowbreak{}reflections\_\allowbreak{}match\_\allowbreak{}bragg} & The ten first fcc reflections equal $2\arcsin(\lambda\sqrt{h^2+k^2+l^2}/2a)$ for (111)\ldots(333). \\
\tn{default\_\allowbreak{}lattice\_\allowbreak{}is\_\allowbreak{}ni\_\allowbreak{}3\_\allowbreak{}52} & Default lattice parameter is 3.52\,\AA. \\
\tn{zero\_\allowbreak{}angle\_\allowbreak{}shift} & A zero angle shifts all reflections by the same amount. \\
\multicolumn{2}{@{}l}{\textit{TestProperties}} \\
\tn{water\_\allowbreak{}atomic\_\allowbreak{}density} & 1\,g/cm$^3$ and 6\,g/mol/atom give $0.602214076/6$ atoms/\AA$^3$. \\
\tn{cylinder\_\allowbreak{}volume} & $\pi r^2 h$. \\
\tn{free\_\allowbreak{}cross\_\allowbreak{}section} & $\sigma[A/(A+1)]^2$. \\
\tn{absorption\_\allowbreak{}reference\_\allowbreak{}wavelength} & $\sigma_\mathrm{abs}(1.798\,\text{\AA})$ equals the tabulated value; $1/v$ scaling. \\
\tn{thermal\_\allowbreak{}references\_\allowbreak{}consistent} & 1.798\,\AA{} $\leftrightarrow$ 25.30\,meV $\leftrightarrow$ 2200\,m/s; the energy route (\code{scattering_probability}) and the wavelength route (\code{getAbsXS}) give the same absorption. \\
\multicolumn{2}{@{}l}{\textit{TestVanadiumStandard}} \\
\tn{default\_\allowbreak{}density\_\allowbreak{}and\_\allowbreak{}consistency} & Vanadium density 6.11\,g/cm$^3$; the atomic density default of \code{setSampleInfo} equals that of \code{setVanaInfo}. \\
\multicolumn{2}{@{}l}{\textit{TestModels}} \\
\tn{inelastic\_\allowbreak{}limits} & \code{inelastic} tends to lowQ at low $Q$ and to lowQ$\,[A/(A+1)]^2$ at high $Q$. \\
\tn{xs\_\allowbreak{}model\_\allowbreak{}high\_\allowbreak{}energy\_\allowbreak{}limit} & Epithermal cross section tends to the free value at high energy. \\
\bottomrule
\end{longtable}

\tssub*{\texttt{test\_reduction.py}}
\begin{longtable}{@{}p{0.36\textwidth} p{0.58\textwidth}@{}}
\toprule Test & What is checked \\ \midrule\endhead
\multicolumn{2}{@{}l}{\textit{TestCorrections}} \\
\tn{binary\_\allowbreak{}sum\_\allowbreak{}variances} & $10\pm2$ minus $5\pm2$ gives $5\pm\sqrt8$ (C3). \\
\tn{normalise\_\allowbreak{}relative\_\allowbreak{}error\_\allowbreak{}independent\_\allowbreak{}of\_\allowbreak{}monitor} & For monitors $10^5$, $10^6$, $10^7$ the relative error is $\sqrt{1/N+1/N_\mathrm{mon}}$ (C2). \\
\tn{normalise\_\allowbreak{}zero\_\allowbreak{}counts\_\allowbreak{}flagged} & A cell with zero counts becomes \code{NaN}. \\
\tn{dead\_\allowbreak{}time} & $\nu\tau>1$ raises; corrected counts $=N/(1-N\tau/t)$; total = sum of corrected cells (C4). \\
\tn{efficiency\_\allowbreak{}file\_\allowbreak{}without\_\allowbreak{}final\_\allowbreak{}newline} & The last value of an efficiency file is read intact (C27). \\
\tn{numor\_\allowbreak{}list\_\allowbreak{}expansion} & \code{"000020-000022"} expands to three numors; a reversed range raises. \\
\multicolumn{2}{@{}l}{\textit{TestOneNumor}} \\
\tn{poisson\_\allowbreak{}errors\_\allowbreak{}from\_\allowbreak{}raw\_\allowbreak{}counts} & With a fake NeXus reader (400 counts per cell, efficiency 0.8): counts 500, error $\sqrt{(20/0.8)^2 + (500\times10^{-3})^2}$, i.e.\ Poisson on raw counts plus monitor statistics (C5). \\
\multicolumn{2}{@{}l}{\textit{TestGetDiffA}} (two synthetic numors of 576 cells) \\
\tn{constant\_\allowbreak{}input\_\allowbreak{}gives\_\allowbreak{}constant\_\allowbreak{}output} & Constant counts give the same constant on the regrouped grid. \\
\tn{errors\_\allowbreak{}are\_\allowbreak{}combined\_\allowbreak{}not\_\allowbreak{}averaged} & Combined errors never exceed the input error and their median is clearly smaller (C29). \\
\tn{non\_\allowbreak{}commensurate\_\allowbreak{}step} & Step $0.3^\circ$ on $[1,140)$ works; the grid has the length of \code{np.arange} (C6). \\
\tn{out\_\allowbreak{}of\_\allowbreak{}range\_\allowbreak{}cells\_\allowbreak{}are\_\allowbreak{}dropped} & Cells at $-5^\circ$ and $200^\circ$ do not appear anywhere (C7). \\
\multicolumn{2}{@{}l}{\textit{TestRebin}} \\
\tn{bin\_\allowbreak{}count} & $[0,0.3]$ step 0.1 gives 4 bins; $[0,0.25]$ gives 3 (C19). \\
\tn{error\_\allowbreak{}of\_\allowbreak{}mean} & Ten points of $100\pm10$ per bin give $100\pm10/\sqrt{10}$ (C29). \\
\multicolumn{2}{@{}l}{\textit{TestWorkflow}} \\
\tn{read\_\allowbreak{}param\_\allowbreak{}and\_\allowbreak{}d4creg\_\allowbreak{}end\_\allowbreak{}to\_\allowbreak{}end} & A parameter file (with an inline comment and values at the end of lines) is read by \code{readParam}; its \code{<run>} line executes the whole d4creg chain on four fake numors; \code{.adat} and \code{.qdat} are written, start at the requested $Q$ and contain only finite values (C25, C8). \\
\bottomrule
\end{longtable}

\tssub*{\texttt{test\_fourier.py}}
\begin{longtable}{@{}p{0.36\textwidth} p{0.58\textwidth}@{}}
\toprule Test & What is checked \\ \midrule\endhead
\multicolumn{2}{@{}l}{\textit{TestSineFT}} (synthetic liquid, $Q_\mathrm{max}=25$, $r$ step 0.02\,\AA) \\
\tn{column\_\allowbreak{}relations} & $g = 1+G/4\pi\rho r$, RDF $=4\pi\rho r^2 g$, $T = 4\pi\rho r g$. \\
\tn{low\_\allowbreak{}r\_\allowbreak{}limit} & No window: $|g(r)| < 0.02$ for $0.5<r<1.5$\,\AA. \\
\tn{lorch\_\allowbreak{}low\_\allowbreak{}r\_\allowbreak{}limit} & Same with the Lorch window (C1; it was $-0.70$). \\
\tn{coordination\_\allowbreak{}number} & First-shell $N$ within 5\,\% of the exact value. \\
\tn{running\_\allowbreak{}integral\_\allowbreak{}from\_\allowbreak{}zero} & $N(r_i)$ equals a direct Simpson integral up to $r_i$ (C18). \\
\tn{running\_\allowbreak{}integral\_\allowbreak{}grid\_\allowbreak{}not\_\allowbreak{}starting\_\allowbreak{}at\_\allowbreak{}zero} & Same on a grid starting at 1\,\AA; $N(r_0)=0$. \\
\tn{invalid\_\allowbreak{}window\_\allowbreak{}raises} & \code{window=2} raises. \\
\multicolumn{2}{@{}l}{\textit{TestCoordinationNumbers}} \\
\tn{lorch\_\allowbreak{}conserves\_\allowbreak{}coordination\_\allowbreak{}number} & For $Q_\mathrm{max}=15$ and 23.5\,\AA$^{-1}$, and limits 0--3.3 and 2.3--3.3\,\AA, $N$ with the Lorch window agrees within 0.5\,\% with the exact value and with the no-window value (the property discussed in Section~\ref{sec:fourier}). \\
\multicolumn{2}{@{}l}{\textit{TestBackFT}} \\
\tn{all\_\allowbreak{}outputs\_\allowbreak{}finite} & No \code{NaN}/\code{inf} although the grids contain $Q=0$ and $r=0$ (C17). \\
\tn{q\_\allowbreak{}zero\_\allowbreak{}limit} & $S(0) = 1+4\pi\rho\int r^2(g-1)\,\mathrm{d}r$, continuous with the next point. \\
\tn{r\_\allowbreak{}relations} & $g(0)=0$; $T$, RDF consistent with $g$; running integral includes the current point. \\
\tn{auto\_\allowbreak{}cut\_\allowbreak{}without\_\allowbreak{}negative\_\allowbreak{}values\_\allowbreak{}cuts\_\allowbreak{}nothing} & \code{cut=[-1]} on a $g(r)$ that is never negative cuts nothing. \\
\tn{cut\_\allowbreak{}intervals} & A 17-value cut list works; an even-length list raises. \\
\multicolumn{2}{@{}l}{\textit{TestSineHelpers}} \\
\tn{sine\_\allowbreak{}integral} & \code{getSineIntegral} equals \code{scipy.special.sici} to $10^{-10}$; $\mathrm{Si}(0)=0$ (C20). \\
\tn{get\_\allowbreak{}sine\_\allowbreak{}ft\_\allowbreak{}vmax\_\allowbreak{}truncates} & \code{vmax} truncates the integral; default unchanged (P4). \\
\tn{windows\_\allowbreak{}reject\_\allowbreak{}bad\_\allowbreak{}qmax} & $Q_\mathrm{max}\le 0$ raises (P5). \\
\tn{lorch\_\allowbreak{}shape} & $M(0)=1$, $M(Q_\mathrm{max}/2)=2/\pi$, $M(Q_\mathrm{max})=0$. \\
\tn{lorch\_\allowbreak{}array\_\allowbreak{}is\_\allowbreak{}standard\_\allowbreak{}lorch} & \code{LorchArray} equals the scalar \code{Lorch} everywhere. \\
\tn{get\_\allowbreak{}sine\_\allowbreak{}ft\_\allowbreak{}lorch} & \code{getSineFT(w=1)} equals a direct Simpson integral with the Lorch window; \code{w=2} raises. \\
\bottomrule
\end{longtable}

\tssub*{\texttt{test\_fourier\_errors.py}}
\begin{longtable}{@{}p{0.36\textwidth} p{0.58\textwidth}@{}}
\toprule Test & What is checked \\ \midrule\endhead
\multicolumn{2}{@{}l}{\textit{TestErrorColumns}} \\
\tn{values\_\allowbreak{}identical\_\allowbreak{}to\_\allowbreak{}plain\_\allowbreak{}transform} & Requesting errors does not change the values (bit-identical), both windows. \\
\tn{error\_\allowbreak{}relations} & $\sigma_g=\sigma_G/4\pi\rho r$, $\sigma_\mathrm{RDF}=r\sigma_G$, $\sigma_T=\sigma_G$; $\sigma_N(r_0)=0$ and $\sigma_N$ increases with $r$. \\
\tn{missing\_\allowbreak{}error\_\allowbreak{}column\_\allowbreak{}raises} & Two-column input with \code{with_errors=True} raises. \\
\tn{errors\_\allowbreak{}scale\_\allowbreak{}linearly} & Doubling $\sigma_S$ doubles every error. \\
\multicolumn{2}{@{}l}{\textit{TestRunningSimpsonMatrix}} \\
\tn{matches\_\allowbreak{}scipy\_\allowbreak{}on\_\allowbreak{}any\_\allowbreak{}grid} & The fast running-Simpson matrix reproduces SciPy on uniform, offset, random non-uniform and two-point grids to $10^{-12}$. \\
\multicolumn{2}{@{}l}{\textit{TestAgainstMonteCarlo}} \\
\tn{analytic\_\allowbreak{}errors\_\allowbreak{}match\_\allowbreak{}monte\_\allowbreak{}carlo} & 1500 noisy copies of $S(Q)$ are transformed; for $G$, $g$, RDF, $T$, $N$ and both windows the analytic error agrees with the Monte Carlo standard deviation within 8\,\% at every $r$ and within 2\,\% on average (the statistical precision of the Monte Carlo is about 1.8\,\%). \\
\multicolumn{2}{@{}l}{\textit{TestSaveRSCF}} \\
\tn{file\_\allowbreak{}with\_\allowbreak{}errors} & The 11-column file reloads with \code{np.loadtxt} to the precision written; the header contains the user heading, all 11 column names in order, and the description of the errors. \\
\tn{file\_\allowbreak{}without\_\allowbreak{}errors\_\allowbreak{}keeps\_\allowbreak{}six\_\allowbreak{}columns} & Without errors the file has 6 columns. \\
\tn{wrong\_\allowbreak{}shape\_\allowbreak{}raises} & An array with 5 columns raises. \\
\bottomrule
\end{longtable}

\tssub*{\texttt{test\_line\_shapes\_ops.py}}
\begin{longtable}{@{}p{0.36\textwidth} p{0.58\textwidth}@{}}
\toprule Test & What is checked \\ \midrule\endhead
\multicolumn{2}{@{}l}{\textit{TestLineShapeErrors}} (errors compared with central finite differences, $10^{-5}$) \\
\tn{lorentzian\_\allowbreak{}error} & Value equals \code{Lorentzian}; error equals the propagated error (C11). \\
\tn{gaussian\_\allowbreak{}error} & Same for \code{Gaussian_error} (correct in the original code). \\
\tn{lorgau\_\allowbreak{}error} & Same for \code{LorGau_error} versus \code{LorGau} (C12). \\
\multicolumn{2}{@{}l}{\textit{TestWsum2}} \\
\tn{general\_\allowbreak{}weights} & $w_1 y_1 + w_2 y_2$ with $\sqrt{w_1^2e_1^2+w_2^2e_2^2}$. \\
\tn{negative\_\allowbreak{}single\_\allowbreak{}weight\_\allowbreak{}gives\_\allowbreak{}positive\_\allowbreak{}error} & $w_1=0$, $w_2=-1$: error $e_2$ (C13). \\
\tn{inverse\_\allowbreak{}variance\_\allowbreak{}weighting} & $w_1=w_2=0$: weights $e_2^2/(e_1^2+e_2^2)$, error $e_1e_2/\sqrt{e_1^2+e_2^2}$. \\
\tn{invalid\_\allowbreak{}mixing\_\allowbreak{}weight\_\allowbreak{}raises} & $w_2=0$, $w_1=2$ raises (C13). \\
\multicolumn{2}{@{}l}{\textit{TestFittingRanges}} \\
\tn{three\_\allowbreak{}variants\_\allowbreak{}agree} & \code{fittingRange0}, \code{fittingRange}, \code{fit_range} select the same points as a direct mask (P20). \\
\bottomrule
\end{longtable}

\tssub*{\texttt{test\_physical\_choices.py}}
\begin{longtable}{@{}p{0.36\textwidth} p{0.58\textwidth}@{}}
\toprule Test & What is checked \\ \midrule\endhead
\multicolumn{2}{@{}l}{\textit{TestMeasurement}} (a complete \code{Measurement} of water built from a minimal parameter set) \\
\tn{annulus\_\allowbreak{}volume} & 4/6\,mm annulus, 50\,mm: $\pi h(r_o^2-r_i^2)$ (C14). \\
\tn{solid\_\allowbreak{}cylinder\_\allowbreak{}volume\_\allowbreak{}unchanged} & Solid cylinder: $\pi h r_o^2$. \\
\tn{free\_\allowbreak{}cross\_\allowbreak{}section\_\allowbreak{}per\_\allowbreak{}element} & $\sigma_\mathrm{free}=\sum c_i\sigma_i[A_i/(A_i+1)]^2$; for water it is less than half the old mean-mass value (P9). \\
\multicolumn{2}{@{}l}{\textit{TestSetSampleInfo}} \\
\tn{free\_\allowbreak{}cross\_\allowbreak{}sections\_\allowbreak{}per\_\allowbreak{}element} & Same formula for the coherent, incoherent and total free cross sections of \code{setSampleInfo}. \\
\multicolumn{2}{@{}l}{\textit{TestSigmoid}} \\
\tn{limits\_\allowbreak{}midpoint\_\allowbreak{}monotonic\_\allowbreak{}no\_\allowbreak{}overflow} & Limits lowQ and highQ, value $(\mathrm{lowQ}+\mathrm{highQ})/2$ at $Q_0$, monotonic, bounded, no overflow up to $Q=10^4$ (C26). \\
\tn{siglin\_\allowbreak{}adds\_\allowbreak{}linear\_\allowbreak{}term} & \code{siglin} $=$ \code{sigmoidal} $+\,a_1Q$. \\
\tn{inelastic\_\allowbreak{}is\_\allowbreak{}the\_\allowbreak{}same\_\allowbreak{}logistic} & \code{inelastic} equals \code{sigmoidal} with highQ $=$ lowQ$\,[A/(A+1)]^2$. \\
\tn{single\_\allowbreak{}q0\_\allowbreak{}default} & The four functions share $Q_0=7.0$ (P21). \\
\multicolumn{2}{@{}l}{\textit{TestWeightedNiFit}} (fake lmfit recording the call) \\
\tn{weights\_\allowbreak{}are\_\allowbreak{}inverse\_\allowbreak{}errors} & The fit receives weights $1/\sigma$ (P11). \\
\tn{no\_\allowbreak{}valid\_\allowbreak{}errors\_\allowbreak{}means\_\allowbreak{}unweighted} & With $\sigma=-1$ no weights are passed. \\
\bottomrule
\end{longtable}

\tssub*{\texttt{test\_misc.py}}
\begin{longtable}{@{}p{0.36\textwidth} p{0.58\textwidth}@{}}
\toprule Test & What is checked \\ \midrule\endhead
\multicolumn{2}{@{}l}{\textit{TestIsotopes}} \\
\tn{bound\_\allowbreak{}cross\_\allowbreak{}sections\_\allowbreak{}in\_\allowbreak{}barns} & Ni: \code{scoh_bound} $=4\pi(10.3)^2/100$ b, equal to the tabulated $\sigma_\mathrm{coh}$ (C15). \\
\tn{deprecated\_\allowbreak{}spanish\_\allowbreak{}names} & \code{Elemento}, \code{elemento} still return \code{Element}, with a warning. \\
\tn{unknown\_\allowbreak{}isotope\_\allowbreak{}raises} & An unknown name raises (P15). \\
\tn{vanadium\_\allowbreak{}b\_\allowbreak{}consistent\_\allowbreak{}with\_\allowbreak{}sigma} & V: $b=-0.443$\,fm and $4\pi b^2/100=\sigma_\mathrm{coh}$; $\sigma_\mathrm{sca}=\sigma_\mathrm{coh}+\sigma_\mathrm{inc}$ (C16). \\
\tn{vanadium\_\allowbreak{}standard\_\allowbreak{}defaults\_\allowbreak{}match\_\allowbreak{}database} & \code{setVanaInfo} defaults equal the table values. \\
\multicolumn{2}{@{}l}{\textit{TestDataXYE}} \\
\tn{file\_\allowbreak{}without\_\allowbreak{}points\_\allowbreak{}in\_\allowbreak{}peak\_\allowbreak{}range} & A file with $x<4$ loads; no peaks found (C22). \\
\tn{plot\_\allowbreak{}saved\_\allowbreak{}before\_\allowbreak{}show} & The saved figure contains the curve even when \code{show()} closes the figure (C22). \\
\multicolumn{2}{@{}l}{\textit{TestNiCalibration}} \\
\tn{initial\_\allowbreak{}amplitudes\_\allowbreak{}mapping} & Each A$_i$, G$_i$, S$_i$ starts from its own entry of \code{initial}; lattice 3.52 by default, or the value given (C10, P7). \\
\tn{real\_\allowbreak{}fit\_\allowbreak{}lattice\_\allowbreak{}fixed\_\allowbreak{}and\_\allowbreak{}wavelength\_\allowbreak{}scales} & A real lmfit fit of a synthetic Ni pattern ($\lambda=0.4986$\,\AA) recovers $\lambda$ to $10^{-6}$ with the lattice fixed at 3.52\,\AA; with 3.524\,\AA{} the fitted $\lambda$ is multiplied by $3.524/3.52$. \\
\multicolumn{2}{@{}l}{\textit{TestMeasurementShowTable}} \\
\tn{stdout\_\allowbreak{}restored\_\allowbreak{}and\_\allowbreak{}file\_\allowbreak{}closed} & \code{showTable} writes its file and restores the previous output stream (P13). \\
\multicolumn{2}{@{}l}{\textit{TestProjectLayout}} \\
\tn{no\_\allowbreak{}module\_\allowbreak{}shadows\_\allowbreak{}the\_\allowbreak{}standard\_\allowbreak{}library} & No file or folder name is in \code{sys.stdlib_module_names} (C21). \\
\tn{stdlib\_\allowbreak{}importable\_\allowbreak{}from\_\allowbreak{}project\_\allowbreak{}root} & \code{import math, io, lmfit, numpy} works with the project directory as current directory. \\
\tn{old\_\allowbreak{}package\_\allowbreak{}names\_\allowbreak{}still\_\allowbreak{}work\_\allowbreak{}with\_\allowbreak{}warning} & \code{toscana.math}/\code{toscana.io} return the same modules as \code{numerics}/\code{dataio}, with warnings. \\
\multicolumn{2}{@{}l}{\textit{TestSourceHygiene}} \\
\tn{no\_\allowbreak{}mojibake} & No doubly-encoded characters in any source file (C23). \\
\bottomrule
\end{longtable}

\tssec{Adding a test}

To protect a new function or a new correction, add a method whose name
starts with \code{test_} to a class deriving from \code{unittest.TestCase}
in one of the \code{test_*.py} files (or in a new file of that name). A
minimal example:
\begin{lstlisting}
import unittest
import helpers                      # makes `toscana` importable
import numpy as np
from toscana.physics.conversions import ang2q

class TestMyCase(unittest.TestCase):
    def test_backscattering(self):
        # 2theta = 180 deg -> Q = 4 pi / lambda
        self.assertAlmostEqual(float(ang2q(180.0, 0.5)), 8 * np.pi)
\end{lstlisting}
Useful comparisons: \code{assertAlmostEqual(a, b, places=6)} or
\code{delta=...}; \code{np.testing.assert_allclose(x, y, rtol=...)} for
arrays; \code{with self.assertRaises(ValueError): ...} to require an error.
Wrap calls that print with \code{helpers.quiet}, and calls that write files
in \code{with helpers.in_tempdir():}.
% =====================================================================
\tsapp{Module reference}\label{app:reference}
% =====================================================================

\tssec{Overview of the package}

\begin{lstlisting}
toscana/                       (the project directory, imported as `toscana`)
  __init__.py      imports the subpackages; installs the deprecated aliases
  _aliases.py      toscana.math -> numerics, toscana.io -> dataio (warnings)
  calibration/     Ni powder calibration of wavelength and zero angle
  chemistry/       compatibility re-exports (composition helpers)
  data/            ASCII x-y-e data containers with plotting
  dataio/          raw data, corrections, regrouping, parameter files, output
  experiment/      experiment, beam, container, vanadium, sample bookkeeping
  isotopes/        nuclear data table and the Element class
  models/          model functions: vanadium, self-scattering, cross sections, Ni
  numerics/        binning, Fourier transforms, windows, fitting, line shapes
  physics/         unit conversions, crystallography, densities, cross sections
  plotting/        matplotlib helpers
  utils/           dates, dictionaries, subprocess, interactive stop
  workflows/       the d4creg regrouping workflow
  tests/           the test suite (Appendix B)
  docs/            this document
\end{lstlisting}

\paragraph{Dependencies between modules.} Low-level packages
(\code{numerics}, \code{physics}, \code{utils}, \code{isotopes}) do not import
higher-level ones. \code{dataio} uses \code{physics}; \code{workflows} uses
\code{dataio}, \code{numerics}, \code{plotting}; \code{calibration} uses
\code{data}, \code{numerics}, \code{models}, \code{physics}; \code{experiment}
uses \code{dataio}, \code{isotopes}, \code{physics}. One cycle exists:
\code{dataio.parameters} imports \code{workflows.d4creg} (so that the
\code{<run>} line can launch it), which imports \code{dataio}; it works
because \code{dataio/__init__.py} imports \code{parameters} last.

\paragraph{External libraries.} NumPy and SciPy are required.
matplotlib (plots), lmfit (Ni fit) and h5py (NeXus files) are imported only
inside the functions that need them.

\paragraph{Reading the entries.} Each function is given with its full
signature and default values. \fixed{C$n$} marks functions corrected during
the audit (Appendix~\ref{app:audit}), \legacy{} functions kept as in the
original code. Arrays called ``3-column'' have columns $(x, y, e)$.
``Logs'' means that the function appends to \code{../logfiles/d4creg.log},
a path relative to the current working directory.

% ---------------------------------------------------------------------
\tssec{Package \texttt{dataio}: raw data, corrections and files}

\tssub{\texttt{dataio.acquisition} --- reading numors}

\fnx{printRates(\allowbreak{}detector,\allowbreak\ counts,\allowbreak\ time,\allowbreak\ reactor\_power=52)}
Logs the counts and the counting rate of \code{detector} (a label such as
\code{"monitor"}), and the rate per MW of reactor power. A non-positive
time is replaced by $10^{-6}$\,s, a non-positive power by 52\,MW. Prints
nothing. \legacy

\fnx{getRefLines(\allowbreak{}Lines)}
Locates the marker lines of an ILL ASCII numor (lines starting with ten
\code{R}, \code{A}, \code{F} or \code{S}) and returns the tuple
\code{(lineR, lineAL, lineAI, lineFD, lineFS, lineFM, lineSI)}: run block,
first two \code{A} blocks, three \code{F} (float parameter) blocks, first
\code{S} (spectrum) block. Raises \code{ValueError} if a marker is missing.

\fnx{readD4numorASCII(\allowbreak{}filename)}
Reads one numor in ILL ASCII format. Returns \code{(data, counts)}:
\code{data} is a dictionary of metadata (numor, label, user, local contact,
title, subtitle, proposal, start and end times, and every parameter of the
three float blocks under its label, e.g.\ \code{"CntTime (sec)"},
\code{"MonitorCnts"}, \code{"TotalCnts"}, \code{"Theta_lue(deg)"},
\code{"Theta_des(deg)"}, plus \code{ndet}, \code{ncell}); \code{counts} is
an array (ndet, ncell). The counts are read 10 per line, in blocks of 11
lines per detector. Prints a summary and logs the rates. \legacy

\fnx{load\_nxs(\allowbreak{}nxs\_path,\allowbreak\ ncell=64)}
Reads one numor in NeXus format with h5py (group \code{entry0}). Returns
\code{(metadata, counts)} with the same keys as the ASCII reader where they
exist: monitor (\code{monitor/monsum}), time, angles
(\code{instrument/masked2theta_s} value and target), scanned variables
(\code{TotalAcquisitionSpy}/\code{TotalCount} renamed \code{TotalCnts}),
and counts reshaped to (9, ncell). The reactor power is set to 1\,MW
(open item P8). \legacy

\fnx{getOneNumor(\allowbreak{}numorfile,\allowbreak\ runInfo)} \fixed{C5}
Reads and corrects one numor; this is the core of the reduction. Steps:
\begin{enumerate}
  \item read with \code{load_nxs} (if \code{runInfo["dataFormat"]=="Nexus"},
    adding the extension \code{runInfo["ext"][4]}) or \code{readD4numorASCII};
  \item compare the read and requested angles: more than
    \code{cellTolerance[0]} cell width gives a warning, more than
    \code{cellTolerance[1]} raises \code{ValueError} (cell width =
    \code{angular_range_bank}/\code{ncell});
  \item Poisson errors $\sqrt{N}$ on the raw counts;
  \item dead-time correction of monitor and detector
    (\code{runInfo["dead_time"]}); errors scaled by the same factor;
  \item division of counts and errors by the efficiencies (\code{efffile});
  \item normalisation to monitor or time (\code{runInfo["normalisation"]});
  \item angles of every cell from \code{getAngles}, with the bank shifts of
    \code{decfile}.
\end{enumerate}
Returns lists \code{x} (angles), \code{y} (counts), \code{e} (errors),
array \code{q}, list \code{c} (cell identifiers) and the metadata. If
\code{runInfo["writeNumor"]==1} the numor is also written by
\code{saveOneNumor}. Note: \code{q} is computed before the zero angle is
subtracted (open item, Section~\ref{sec:open}). Logs.

\fnx{getNumors(\allowbreak{}runInfo)}
Expands \code{runInfo["numorLst"]} with \code{getNumorFiles}, reads every
numor from \code{runInfo["path_raw"]} with \code{getOneNumor}, and stacks
them in an array of shape (n\_numors, 5, n\_points), the five rows being:
angle $-$ \code{twotheta0}, counts, errors, $Q$, cell identifier. Returns
\code{(numor, head)} with \code{head} the list of metadata dictionaries.
If there is only one numor, \code{writeNumor} is forced to 1. Logs.

\tssub{\texttt{dataio.corrections} --- detector corrections}

\fnx{getDTC\_mcounts(\allowbreak{}monitor,\allowbreak\ time,\allowbreak\ dead\_time=2.4e-06)}
Dead-time correction of the monitor (non-paralysable):
$\nu = \nu_m/(1-\nu_m\tau)$, $\nu_m = $ monitor$/$time. Returns the
corrected counts; raises \code{ValueError} if $\nu_m\tau\ge1$. Logs the factor.

\fnx{getDTC\_dcounts(\allowbreak{}counts,\allowbreak\ total\_counts,\allowbreak\ time,\allowbreak\ dead\_time=7e-06)} \fixed{C4}
Same correction cell by cell for the detector array. Returns the corrected
array and the total counts scaled by $\sum N_\mathrm{corr}/\sum N_\mathrm{raw}$.
Raises \code{ValueError} if any factor is not positive. Logs the overall
and maximum factor. (Whether the dead time should apply per cell or per
bank is open item P2.)

\fnx{getEff(\allowbreak{}filename=\textquotesingle{}effd4c.eff\textquotesingle{},\allowbreak\ ndet=9,\allowbreak\ ncells=64)} \fixed{C27}
Reads relative efficiencies, one line \code{det cell efficiency} per cell
(1-based); lines containing \texttt{\#} are ignored. Values $\le0$ become
\code{NaN}, which removes the cell from the regrouping. Cells not listed
keep efficiency 1. \code{filename="ones"} returns all ones.

\fnx{getDec(\allowbreak{}filename=\textquotesingle{}dec.dec\textquotesingle{},\allowbreak\ ndet=9)} \fixed{C27}
Reads the angular shift of each bank, lines \code{det shift} (degrees).
\code{"zeros"} returns zeros; \code{"manual"} returns a built-in list.

\fnx{getErrorsNumor(\allowbreak{}counts)}
Poisson errors $\sqrt{N}$. Must be called on raw counts.

\fnx{getNormalFactors(\allowbreak{}monitor,\allowbreak\ time)}
Returns the 2$\times$2 array $[[M, \sqrt M],[t, 0.01]]$: normaliser value
and its absolute error for monitor and time.

\fnx{normalise(\allowbreak{}counts,\allowbreak\ errors,\allowbreak\ normal,\allowbreak\ norm\_type=\textquotesingle{}monitor\textquotesingle{},\allowbreak\ norm\_time=120,\allowbreak\ norm\_mon=1000000)} \fixed{C2}
Scales counts to \code{norm_mon} monitor counts (or \code{norm_time}
seconds if \code{norm_type="time"}): $y_n = f y$, $f=$ norm$/M$, and
$\sigma_n^2 = (f\sigma)^2 + (y_n\sigma_M/M)^2$. Zero-count cells become
\code{NaN}. Returns \code{(counts, errors)}.

\fnx{getAngles(\allowbreak{}counts,\allowbreak\ ttheta\_ref,\allowbreak\ zeros,\allowbreak\ cell\_step=0.125,\allowbreak\ det\_step=15.0)}
Angle of every cell: $2\theta = 2\theta_\mathrm{ref} + \mathrm{zeros}[d] + 15^\circ d + 0.125^\circ c$
($d$, $c$ 0-based); cell identifier $1000(d+1) + c + 1$. Returns two
arrays of the shape of \code{counts}.

\tssub{\texttt{dataio.d4creg\_outputs} --- regrouping and D4 files}

\fnx{getDiffA(\allowbreak{}numor,\allowbreak\ head,\allowbreak\ runInfo)} \fixed{C6, C7, C29}
Regroups all cells of all numors onto the regular angular grid
\code{np.arange(amin, amax, step)} of \code{runInfo["angular_scale"]}. If
the requested step is smaller than a cell
(\code{angular_range_bank}/\code{ncell}), the cell width is used. Each cell,
of the cell width and centred at its angle, is split between the (at
most two) bins it overlaps with fractions $f$ (exact only when the step
equals the cell width, the default; for a larger step the fractions are
computed in units of the bin, not of the cell: with $0.25^\circ$ a cell
overlapping two bins equally is split 0.25/0.75; open item); its weight is $w = f\times
M_i/\sum M$ (monitor share of its numor). For each bin
$y = \sum w y_i/\sum w$ and $\sigma = \sqrt{\sum w^2\sigma_i^2}/\sum w$.
Cells with \code{NaN} counts or errors, and parts of cells outside the grid,
are ignored; empty bins are \code{NaN}. Returns an array (10, 5, n\_angles):
index 0 is the sum of all banks, 1--9 the individual banks; rows are angle,
counts, error, sum of weights, $Q$. The angle written is the lower edge of
the bin (open item P1).

\fnx{saveDiffAngle(\allowbreak{}diffA,\allowbreak\ head,\allowbreak\ runInfo)}
Writes two files: \code{<file>.adat} (CORRECT format: header block, then
angle, counts, error of the total diffractogram, positive finite points
only) and \code{<file>.reg} (D4 format: header with the equivalent d4creg
command line, efficiency and shift files, normalisation, zero angle,
wavelength, dead times; then one block per bank with angle, counts, error,
$Q$). Extensions from \code{runInfo["ext"]}.

\fnx{saveDiffQ(\allowbreak{}diffQ,\allowbreak\ head,\allowbreak\ runInfo)}
Writes \code{<file>.qdat} (CORRECT format) from a 3-column $(Q, y, e)$
array, positive finite points only.

\tssub{\texttt{dataio.parameters} --- the d4creg parameter file}

\fnx{readParam(\allowbreak{}parametersFile)} \fixed{C25}
Reads a d4creg parameter file and returns the dictionary \code{runInfo}
used by the whole chain. Every non-empty line must start with \texttt{\#} or
\code{!} (comments) or with a five-character tag \code{<xxx>} followed by
values; text after \texttt{\#} on a tag line is ignored. Recognised tags and
defaults:

\begin{longtable}{@{}l l p{0.46\textwidth}@{}}
\toprule Tag & Values & Meaning (default) \\ \midrule\endhead
\code{<tol>} & w e & angle tolerance in cells for warning / error (0.1, 1.0) \\
\code{<tau>} & $\tau_M$ $\tau_D$ & dead times in s ($2.4\times10^{-6}$, $7\times10^{-6}$) \\
\code{<ext>} & 6 strings & extensions (\code{.reg .adat .qdat .cdat .nxs .log}) \\
\code{<eff>} & file & efficiency file (\code{effd4c.eff}; \code{ones}) \\
\code{<dec>} & file & bank shifts (\code{dec.dec}; \code{zeros}, \code{manual}) \\
\code{<log>} & name & log file name (\code{d4creg.log}) \\
\code{<fmt>} & \code{Nexus}/other & raw data format (\code{Nexus}) \\
\code{<asc>} & min max step & angular grid, degrees (0, 140, 0.125) \\
\code{<qsc>} & min max step & $Q$ grid, \AA$^{-1}$ (0, 23.8, 0.02); prints a recommended step \\
\code{<wri>} & \code{True}/\code{False} & write one file per numor (False) \\
\code{<plo>} & \code{True}/\code{False} & plot diffractograms (False) \\
\code{<wle>} & $\lambda$ & wavelength, \AA{} (0.49857) \\
\code{<zac>} & $2\theta_0$ & zero angle, degrees ($-0.1013$) \\
\code{<rdp>} & path & raw data directory (\code{rawdata/}) \\
\code{<nor>} & type t M & \code{monitor} or \code{time}, normalisation time and monitor (monitor, 80, $10^6$) \\
\code{<out>} & name & base name of output files (\code{vanadium}) \\
\code{<num>} & list & numors: \code{123456} or \code{123456-123460} (replaces the list) \\
\code{<add>} & list & appends to the list of numors \\
\code{<run>} & & runs \code{d4creg(runInfo)} with the values read so far \\
\bottomrule
\end{longtable}
\noindent Side effect: (re)creates \code{../logfiles/d4creg.log}. The file
is read sequentially, so several \code{<run>} lines with changes in between
perform several reductions.

\tssub{\texttt{dataio.running\_params} --- the analysis parameter file}

\fnx{getRunningParams(\allowbreak{}parfile)}
Reads the parameter file of the data-treatment notebooks and returns a
dictionary \code{inputPar} keyed by tag; each value is a tuple whose first
element is a label and the following ones the values. Tags: experiment
(\code{<ins> <pro> <mpr> <lco> <dte> <cyc> <log> <env>}; \code{<dte>} gives
\code{<sta>}/\code{<end>}, \code{<log>} gives \code{<nbk>}/\code{<pag>}), beam
(\code{<zac> <wle> <sli> <fla> <apr>}; \code{<sli>} gives \code{<lft>}/\code{<rgt>},
\code{<fla>} gives \code{<top>}/\code{<bot>}), grids (\code{<asc> <qsc> <rsc>}),
container, absorber, normaliser, environment (\code{<Cmater> <Cshape>},
\code{<Amater> <Ashape>}, \code{<Nmater> <Nshape>}, \code{<Emater> <Eshape>};
shapes: type, inner diameter, outer diameter, height), sample
(\code{<Sdescr> <StempK> <Sato01>... <Sdegcc> <Smassc> <Sshape> <Sfulln> <Stitle>};
each \code{<SatoNN>} is \code{symbol number}, counted into \code{<SNelem>}),
and data files (\code{<CoFile> <AbFile> <NoFile> <EnFile> <SaFile>}). Unknown
materials or shapes print a warning. Its output is the input of
\code{experiment.Measurement}. \legacy

\tssub{\texttt{dataio.loading}, \texttt{saving}, \texttt{numors}, \texttt{correct}, \texttt{\_logfiles}}

\fnx{read\_xye(\allowbreak{}filename)} \fixed{C27}
Reads 2 or 3 columns, ignoring lines containing \texttt{\#}; a missing or
non-numeric (\code{i}, \code{o}) error becomes 0. Returns three 1-D arrays.

\fnx{read\_3col(\allowbreak{}filename)}
Same, returned as one (n, 3) array.

\fnx{saveFile\_xye(\allowbreak{}filename,\allowbreak\ x,\allowbreak\ y,\allowbreak\ e,\allowbreak\ heading)}, \code{saveFile_3col(filename, data, heading)}
Write a 3-column file with \texttt{\#} header lines (formats \code{9.3f},
\code{18.6f}, \code{18.6f}).

\fnx{saveCorrelations(\allowbreak{}filename,\allowbreak\ ar,\allowbreak\ pcf,\allowbreak\ pdf,\allowbreak\ rdf,\allowbreak\ tor,\allowbreak\ run,\allowbreak\ heading)}
Writes six columns $r, G, g$, RDF, $T, N$ from separate arrays.

\fnx{saveRSCF(\allowbreak{}filename,\allowbreak\ fou,\allowbreak\ heading)} \textsf{\footnotesize\colorbox{green!15}{new header, 11 columns}}
Writes the output of \code{sineFT}: 6 columns, or 11 with errors. After
the file name and the user \code{heading}, a block describes each column
(name, formula, units), the meaning and limits of the errors, and a line of
column names \code{r G dG g dg RDF dRDF T dT N dN}. Other shapes raise.

\fnx{saveOneNumor(\allowbreak{}x,\allowbreak\ y,\allowbreak\ e,\allowbreak\ q,\allowbreak\ c,\allowbreak\ head,\allowbreak\ runInfo)}
Writes \code{<numor><ext[3]>} (angle, counts, error, $Q$, cell) for one numor.

\fnx{getNumorFiles(\allowbreak{}numorLst)}
Expands entries \code{"123456"} and \code{"123456-123460"} into a list of
six-digit numors; malformed or reversed ranges raise.

\fnx{generate\_com(\allowbreak{}CORRECT\_template=\textquotesingle{}correct\_template.com\textquotesingle{},\allowbreak\ name=\textquotesingle{}prueba\textquotesingle{},\allowbreak\ rsample=0.5,\allowbreak\ density=0.11,\allowbreak\ ...)}
Fills the placeholders \texttt{\{pyfname\}}, \texttt{\{pyradius\_cm\}}, \ldots{} of a
CORRECT template and writes \code{<name>.com}. It then runs a placeholder
command (open item). \legacy

\fnx{saveCORRECT(\allowbreak{}expt)}
Writes a CORRECT command file \code{<sampleTitle>.com} from an experiment
object (sample, container, background, vanadium, wavelength, beam size;
zero angle written as 0 since already subtracted). \legacy

\fnx{d4creg\_log\_path(\allowbreak{})}
Returns \code{../logfiles/d4creg.log}, creating the folder and file if needed.

% ---------------------------------------------------------------------
\tssec{Package \texttt{workflows}}

\fnx{d4creg(\allowbreak{}runInfo)} \fixed{C8, C9}
The complete regrouping: \code{getNumors} $\to$ \code{getDiffA} $\to$
\code{saveDiffAngle} (+ plots if \code{plotDiff}) $\to$ conversion of the
regrouped angles to $Q$ $\to$ \code{rebin} onto \code{q_scale} with the
actual angular step as width $\to$ masking of points with $y\le0$ and of
their predecessor $\to$ \code{saveDiffQ} (+ plot) $\to$ renaming of the log
file to \code{../logfiles/<logfile>.log}. Prints the equivalent d4creg
command line. Adds \code{runDate}, \code{nbrNumors}, \code{d4creg_cl} to
\code{runInfo}.

% ---------------------------------------------------------------------
\tssec{Package \texttt{numerics}: numerical methods}

\tssub{\texttt{numerics.fourier} --- sine Fourier transforms}

\fnx{sineFT(\allowbreak{}StrFactor,\allowbreak\ nr,\allowbreak\ qmax=23.5,\allowbreak\ density=0.1,\allowbreak\ constant=1.0,\allowbreak\ selfsca=1.0,\allowbreak\ window=0,\allowbreak\ with\_errors=False)} \fixed{C1, C18} \newbeh
The main transform, Section~\ref{sec:conventions}. \code{StrFactor} is an
array with columns $Q$, $S$ (and $\sigma_S$); \code{nr} the $r$ grid;
\code{density} $\rho$ in atoms/\AA$^3$; \code{constant} a multiplicative
factor; \code{selfsca} the value subtracted from $S$ (1 for a normalised
$S(Q)$); \code{window} 0 (none) or 1 (standard Lorch). Returns an
$(N_r, 6)$ array $r, G, g$, RDF, $T, N$, or with \code{with_errors=True}
an $(N_r, 11)$ array with each error beside its function
(Section~\ref{sec:errors}). The $Q$ integral runs over the whole $Q$ grid
with $M(Q)=0$ for $Q\ge$\code{qmax}, so the zeroed points still take part
in the Simpson panels next to $Q_\mathrm{max}$ (\code{getSineFT} instead
cuts the grid at \code{vmax}); the running integral starts at \code{nr[0]}.

\fnx{getSineFT(\allowbreak{}vvec,\allowbreak\ yvec,\allowbreak\ uvec,\allowbreak\ vmax=0,\allowbreak\ c=1.0,\allowbreak\ s=1.0,\allowbreak\ w=0)} \fixed{P4}
Generic transform
$\mathrm{stf}(u) = \frac2\pi c\int_0^{v_\mathrm{max}} [y(v)-s]\,v\,\sin(vu)\,W(v)\,\mathrm{d}v$,
$W$ = 1 (\code{w=0}) or Lorch (\code{w=1}). \code{vmax}$\le0$ means the
end of the data. Returns \code{None} if the lengths differ. Used by
\code{backFT} in both directions.

\fnx{backFT(\allowbreak{}qu,\allowbreak\ ar,\allowbreak\ pdf,\allowbreak\ density,\allowbreak\ cut)} \fixed{C17}
Cuts $g(r)$ (\code{pdf} on grid \code{ar}) according to \code{cut}
(Section~\ref{sec:fourier}), transforms to $S(Q)$ on \code{qu} and back.
Returns the six arrays $(S, G, g, \mathrm{RDF}, T, N)$. Typical use: remove
the unphysical oscillations of $g(r)$ below the first peak and see their
effect on $S(Q)$.

\fnx{getSineIntegral(\allowbreak{}x)} \fixed{C20}
$\mathrm{Si}(x)=\int_0^x \sin t/t\,\mathrm{d}t$ (\code{scipy.special.sici}).

\noindent Internal helpers: \code{_simpson_weights(x)} (the weights $w_k$ of
SciPy's Simpson rule), \code{_sine_ft_errors(...)} (the error
propagation), \code{_running_simpson_matrix(x)} (matrix $C$ of the running
integral, with self-check), \code{_apply_cut(ar, pdf, cut)}.

\tssub{\texttt{numerics.windows} --- modification functions}

\fnx{Lorch(\allowbreak{}q,\allowbreak\ qmax=23.5)} \fixed{P5}
Standard Lorch value at one point: $\sin(\pi q/q_\mathrm{max})/(\pi q/q_\mathrm{max})$,
1 at 0, 0 outside $[0, q_\mathrm{max})$. Raises if $q_\mathrm{max}\le0$.

\fnx{LorchArray(\allowbreak{}q,\allowbreak\ qmax)} \textsf{\footnotesize\colorbox{green!15}{new}}
The same for an array (replaces the removed, rescaled \code{LorchN}).

\tssub{\texttt{numerics.binning} --- rebinning}

\fnx{get\_xlim(\allowbreak{}xmin,\allowbreak\ xmax,\allowbreak\ dbin)} \fixed{C19}
Bin edges and centres: centres $x_\mathrm{min} + i\,\delta$, edges at
$\pm\delta/2$; returns \code{(nb, x_lim, x_bin)}.

\fnx{get\_bins(\allowbreak{}x,\allowbreak\ xdel,\allowbreak\ x\_lim)}
Bins covered by a rectangle of width \code{xdel} centred at \code{x}, and
the fraction of the rectangle in each. A rectangle that starts below the
first edge is rejected entirely (empty lists). At the upper end the
treatment differs: a rectangle is rejected only if it ends one full bin
beyond the last edge; otherwise its excess goes to an overflow index that
\code{rebin} discards, so a point lying entirely in the bin-width beyond
the last edge is dropped silently and a straddling point contributes only
its overlap.

\fnx{rebin(\allowbreak{}xdel,\allowbreak\ wlength,\allowbreak\ data,\allowbreak\ xmin,\allowbreak\ xmax,\allowbreak\ dbin)} \fixed{C29}
Rebins a 3-column array onto the grid of \code{get_xlim}. Each point is a
rectangle of width \code{xdel} (in the units of $x$ if \code{wlength}$\le0$;
if \code{wlength}$>0$, $x$ is $Q$ and \code{xdel} an angular width in degrees,
converted to $\Delta Q = \frac{2\pi}{\lambda}\cos\theta\,\Delta(2\theta)\frac{\pi}{180}$).
Weighted mean per bin, error $\sqrt{\sum f^2\sigma^2}/\sum f$. Points with
\code{NaN} are skipped; empty bins are 0. Prints a debug line.

\tssub{\texttt{numerics.operations} --- combining data sets}

\fnx{binary\_sum(\allowbreak{}w1,\allowbreak\ data1,\allowbreak\ w2,\allowbreak\ data2)} \fixed{C3}
$w_1y_1+w_2y_2$, error $\sqrt{w_1^2e_1^2+w_2^2e_2^2}$, for two 3-column
arrays of the same shape (otherwise prints and returns \code{None}).

\fnx{wsum2(\allowbreak{}w1,\allowbreak\ data1,\allowbreak\ w2,\allowbreak\ data2)} \fixed{C13}
Four modes: $w_1=w_2=0$: inverse-variance weighted mean; $w_1=0$: $w_2y_2$
with error $|w_2|e_2$; $w_2=0$: mixture $w_1y_1+(1-w_1)y_2$, requires
$0<w_1\le1$ (raises otherwise); else $w_1y_1+w_2y_2$.

\fnx{ratio(\allowbreak{}y1,\allowbreak\ e1,\allowbreak\ y2,\allowbreak\ e2)}
$y_1/y_2$ with error $\sqrt{(y_2e_1)^2+(y_1e_2)^2}/y_2^2$; returns $0\pm0$
where $y_2=0$. Lists, truncated to the shortest input. \legacy

\tssub{\texttt{numerics.line\_shapes} --- peak and step functions}

Constants \code{SIGMOID_Q0 = 7.0} and \code{SIGMOID_DQ = 2.4}
(\AA$^{-1}$): default centre and width of the low-$Q$/high-$Q$ transition.

\fnx{sigmoidal(\allowbreak{}Q,\allowbreak\ lowQ=0.4,\allowbreak\ highQ=0.2,\allowbreak\ Q0=SIGMOID\_Q0,\allowbreak\ dQ=SIGMOID\_DQ)} \fixed{C26} \newbeh
Logistic step $(\mathrm{lowQ}+\mathrm{highQ}\,\delta)/(1+\delta)$, $\delta = e^{(Q-Q_0)/\mathrm{d}Q}$.

\fnx{siglin(\allowbreak{}Q,\allowbreak\ lowQ=0.4,\allowbreak\ highQ=0.2,\allowbreak\ Q0=SIGMOID\_Q0,\allowbreak\ dQ=SIGMOID\_DQ,\allowbreak\ a1=0.001)} \newbeh
\code{sigmoidal(Q)} $+\;a_1Q$.

\fnx{step(\allowbreak{}q,\allowbreak\ qmax=23.5)} \fixed{P5}
1 for $0\le q<q_\mathrm{max}$, else 0; raises if $q_\mathrm{max}\le0$.

\fnx{Lorentzian(\allowbreak{}x,\allowbreak\ A,\allowbreak\ x0,\allowbreak\ gamma,\allowbreak\ bckg)}
$A(\gamma/2)^2/[(x-x_0)^2+(\gamma/2)^2] + b$: height $A$, FWHM $\gamma$.

\fnx{Lorentzian\_error(\allowbreak{}x,\allowbreak\ A=1,\allowbreak\ x0=0,\allowbreak\ gamma=1,\allowbreak\ bckg=0,\allowbreak\ e\_A=0,\allowbreak\ e\_x0=0,\allowbreak\ e\_gamma=0,\allowbreak\ e\_bckg=0)} \fixed{C11}
Value and propagated error of \code{Lorentzian}.

\fnx{Gaussian(\allowbreak{}x,\allowbreak\ A,\allowbreak\ x0,\allowbreak\ sigma,\allowbreak\ bckg)}
Area-normalised Gaussian $\frac{A}{\sigma\sqrt{2\pi}}e^{-(x-x_0)^2/2\sigma^2}+b$ (area $A$).

\fnx{Gaussian\_error(\allowbreak{}x,\allowbreak\ A=1,\allowbreak\ x0=0,\allowbreak\ sigma=1,\allowbreak\ bckg=0,\allowbreak\ e\_A=0,\allowbreak\ e\_x0=0,\allowbreak\ e\_sigma=0,\allowbreak\ e\_bckg=0)}
Value and propagated error of \code{Gaussian}.

\fnx{GaussianA(\allowbreak{}x,\allowbreak\ A,\allowbreak\ x0,\allowbreak\ sigma,\allowbreak\ asym=0)}
Area-normalised Gaussian (area $A$, no background) multiplied by
$e^{-\mathrm{asym}(x-x_0)}$ (asymmetric peak, used for the Ni reflections).
For $\mathrm{asym}\neq0$ the area becomes $A\,e^{\mathrm{asym}^2\sigma^2/2}$
and the maximum moves to $x_0-\mathrm{asym}\,\sigma^2$.

\fnx{LorGau(\allowbreak{}x,\allowbreak\ f0=1.0,\allowbreak\ eta=0.5,\allowbreak\ sigma=2.0,\allowbreak\ gamma=2.0,\allowbreak\ bckg=0.0)}
Mixture centred at 0: $f_0[\eta e^{-x^2/2\sigma^2} + (1-\eta)L(x)] + b$,
both components of unit height.

\fnx{LorGau\_error(\allowbreak{}x,\allowbreak\ f0=(\allowbreak{}1.0,\allowbreak\ 0.0),\allowbreak\ eta=(\allowbreak{}0.5,\allowbreak\ 0.0),\allowbreak\ sigma=(\allowbreak{}2.0,\allowbreak\ 0.0),\allowbreak\ gamma=(\allowbreak{}2.0,\allowbreak\ 0.0),\allowbreak\ bckg=(\allowbreak{}0.0,\allowbreak\ 0.0))} \fixed{C12}
Value and error of \code{LorGau}; each argument is a pair (value, error).

\tssub{\texttt{numerics.fitting}, \texttt{polynomials}, \texttt{peak\_finding}, \texttt{signal\_processing}}

\fnx{get\_chi(\allowbreak{}y1,\allowbreak\ y2)}
Root-mean-square difference $\sqrt{\langle(y_1-y_2)^2\rangle}$.

\fnx{fit\_and\_find\_extremum(\allowbreak{}x,\allowbreak\ y)}
Fits a parabola and returns the extremum $x=-b/2a$, its value and the
fitted curve.

\fnx{fittingRange0(\allowbreak{}xmin,\allowbreak\ xmax,\allowbreak\ ymin,\allowbreak\ ymax,\allowbreak\ abcissas,\allowbreak\ ordinates,\allowbreak\ errors)}, \code{fittingRange(...)}, \code{fit_range(xmin, xmax, ymin, ymax, data)} \fixed{P20}
Select the points inside the rectangle $[x_\mathrm{min},x_\mathrm{max}]\times[y_\mathrm{min},y_\mathrm{max}]$.
\code{fittingRange} returns three $(n,1)$ arrays, the other two one
$(n,3)$ array. Warn below 10 points.

\fnx{polyQ1(\allowbreak{}x,\allowbreak\ a0,\allowbreak\ a1)}, \code{polyQ2(x, a0, a1, a2)}, \code{polyQ4(x, a0, ..., a4)}
Polynomials $\sum a_i x^i$ (Horner scheme).

\fnx{find\_minimum\_within\_range(\allowbreak{}x,\allowbreak\ y,\allowbreak\ x\_min,\allowbreak\ x\_max)}
$(x, y)$ of the minimum of $y$ in $[x_\mathrm{min}, x_\mathrm{max}]$ (raises
on an empty range).

\fnx{find\_peaks\_in\_range(\allowbreak{}x,\allowbreak\ y,\allowbreak\ x\_min,\allowbreak\ x\_max)}
Local maxima (\code{scipy.signal.find_peaks}) whose height divided by the
minimum of the range exceeds 3.5; returns positions and these ratios.
\legacy

\fnx{smooth\_curve(\allowbreak{}x,\allowbreak\ y,\allowbreak\ smoothing\_factor)}
Centred moving average over $2\lfloor w/2\rfloor+1$ points,
$w = \lfloor\text{factor}\times n\rfloor$, shrinking at the edges; \code{x}
is not used.

% ---------------------------------------------------------------------
\tssec{Package \texttt{physics}}

\tssub{\texttt{physics.conversions}, \texttt{geometry}}

\fnx{ang2q(\allowbreak{}x,\allowbreak\ wlength=0.5)}, \code{q2ang(x, wlength=0.5)}
$Q = \frac{4\pi}{\lambda}\sin(\pi\,2\theta/360)$ and its inverse; raise if
$\lambda\le0$.

\fnx{getCylVolume(\allowbreak{}diameter=5.0,\allowbreak\ height=50.0)}
$\pi(d/2)^2h$ (mm$^3$ for mm input).

\tssub{\texttt{physics.crystallography}}

Constant \code{NI_LATTICE_PARAMETER = 3.52} (\AA).

\fnx{sf\_fcc(\allowbreak{}h,\allowbreak\ k,\allowbreak\ l)}, \code{sf_bcc(h, k, l)}, \code{sf_sc(h, k, l)}
Geometric structure factors $\sum_j e^{-i\pi(\ldots)}$ of the fcc (4 or 0),
bcc (2 or 0) and simple cubic (1) lattices.

\fnx{reflections\_fcc(\allowbreak{}wavelength=0.5,\allowbreak\ twotheta0=0.0,\allowbreak\ lattice=NI\_LATTICE\_PARAMETER)}
Angles $2\theta+2\theta_0$ (degrees) of the allowed fcc reflections in
increasing order, followed by the mirrored set $180^\circ-2\theta-2\cdot2\theta_0$.

\tssub{\texttt{physics.properties}}

Constants \code{THERMAL_WAVELENGTH = 1.798} (\AA) and
\code{THERMAL_ENERGY_MEV} $=81.8042/1.798^2 = 25.30$ (meV).

\fnx{getAtomicDensity(\allowbreak{}density=1.0,\allowbreak\ molarMass=6.0)}
$\rho = 0.602214076\,d/\bar M$ atoms/\AA$^3$ ($\bar M$ per atom). A
non-positive density gives the value of water (0.1004) with a message.

\fnx{getDensity(\allowbreak{}atomic\_density=0.1,\allowbreak\ molarM=6.0)}
Inverse of the above.

\fnx{getMassNumber(\allowbreak{}molarMass=6.0)}
$A = M/m_n$, $m_n = 1.0086649$\,u.

\fnx{getFreeXS(\allowbreak{}BoundXS=5.08,\allowbreak\ A=50.9415)}
$\sigma_\mathrm{bound}[A/(A+1)]^2$ for one nucleus, $A=m/m_n$. Note: the
default \code{A=50.9415} is the molar mass of vanadium, not its mass ratio
50.50 (0.03\,\% effect on the factor).

\fnx{getAbsXS(\allowbreak{}AbsXS=5.08,\allowbreak\ wavelength=THERMAL\_WAVELENGTH)} \fixed{C30}
$\sigma_\mathrm{abs}\lambda/1.798$.

\fnx{getNofAtoms(\allowbreak{}atoms)}, \code{getConcentrations(atoms)}, \code{AtomicAvg(concentration, magnitude)}
Number of atoms of a formula unit given as \texttt{\{"Si": 1, "O": 2\}};
concentrations $c_i = n_i/\sum n$; average $\sum c_i m_i$ of a property
given as a dictionary.

% ---------------------------------------------------------------------
\tssec{Package \texttt{models}}

\fnx{inelastic(\allowbreak{}x,\allowbreak\ A=51.0,\allowbreak\ lowQ=0.4,\allowbreak\ Q0=SIGMOID\_Q0,\allowbreak\ dQ=SIGMOID\_DQ)}
Inelasticity (Placzek-like) decrease of the self scattering:
$\mathrm{lowQ}\,[1+(A/(1+A))^2\delta]/(1+\delta)$, from the bound value at
low $Q$ to the free value at high $Q$.

\fnx{self024(\allowbreak{}x,\allowbreak\ self0,\allowbreak\ q2,\allowbreak\ q4)}
$s_0 + q_2Q^2 + q_4Q^4$, an empirical self-scattering curve to fit.

\fnx{vanaQdep(\allowbreak{}x,\allowbreak\ a0,\allowbreak\ a1,\allowbreak\ a2,\allowbreak\ A=51.0,\allowbreak\ lowQ=0.4,\allowbreak\ Q0=SIGMOID\_Q0,\allowbreak\ dQ=SIGMOID\_DQ)}
Vanadium $Q$ dependence: \code{polyQ2} (resolution, attenuation trends)
$\times$ \code{inelastic}.

\fnx{XS\_model(\allowbreak{}Eval,\allowbreak\ Eeff,\allowbreak\ composition\_vec,\allowbreak\ bound\_xs\_vec,\allowbreak\ A\_vec)}
Epithermal total scattering cross section
$\sum c_i\sigma_i[A_i/(A_i+1)]^2[1+E_\mathrm{eff}/(2A_iE)]$.

\fnx{scattering\_probability(\allowbreak{}E0,\allowbreak\ scat\_xs,\allowbreak\ composition\_vec,\allowbreak\ abs\_xs\_vec)}
$1-\sigma_a/(\sigma_a+\sigma_s)$ with $\sigma_a = \sum c_i\sigma_{a,i}/\sqrt{E_0/E_\mathrm{th}}$ ($E_0$ in meV).

\fnx{niPeaks10(\allowbreak{}x,\allowbreak\ I0,\allowbreak\ slope,\allowbreak\ quad,\allowbreak\ wavelength,\allowbreak\ twotheta0,\allowbreak\ A0..A9,\allowbreak\ G0..G9,\allowbreak\ S0..S9,\allowbreak\ lattice=NI\_LATTICE\_PARAMETER)}
Ni powder model: background $I_0 + \mathrm{slope}\,x + \mathrm{quad}\,x^2$
plus ten \code{GaussianA} peaks (areas A$_i$, widths S$_i$, asymmetries
G$_i$) at the first ten fcc reflections for the given wavelength, zero
angle and lattice parameter. $x$ is $2\theta$ in degrees.

% ---------------------------------------------------------------------
\tssec{Package \texttt{calibration}}

\fnx{setting\_model\_d4nifit(\allowbreak{}nifile,\allowbreak\ zac\_ini=0.0,\allowbreak\ lattice=NI\_LATTICE\_PARAMETER)}
Prepares the Ni calibration. Reads the regrouped Ni diffractogram
(\code{DataXYE}), estimates $\lambda$ from the first peak, assumed to be
(111): $\lambda = 2(a/\sqrt3)\sin\theta_{111}$; normalises intensities and
errors to the first data point above the midpoint between the second and
third peaks, truncated to whole degrees; builds the
35 initial parameters (background, $\lambda$, $2\theta_0$, ten areas
proportional to typical Ni intensities, asymmetries 0, widths 0.2--0.4);
plots data, initial model and residuals. Returns
\code{(axes, limits, xy_res, xy_ref, xy_fit, planes, initial, nickel)}.

\fnx{fitting\_model\_d4nifit(\allowbreak{}nickel,\allowbreak\ limits,\allowbreak\ initial,\allowbreak\ lattice=NI\_LATTICE\_PARAMETER)} \fixed{C10, P11}
Fits \code{niPeaks10} with lmfit to the points inside \code{limits}
(rectangle in $2\theta$ and intensity), weighted by $1/\sigma$. Free: I0,
slope, wavelength, twotheta0, A$_i$, S$_i$; fixed: quad, G$_i$, lattice.
Returns the lmfit \code{ModelResult}.

\fnx{showing\_results\_d4nifit(\allowbreak{}result,\allowbreak\ nickel,\allowbreak\ axes,\allowbreak\ limits,\allowbreak\ xy\_res,\allowbreak\ xy\_ref,\allowbreak\ xy\_fit,\allowbreak\ planes)}
Prints $\lambda$, $2\theta_0$ and background with errors, a table of the
ten reflections (position, area, width), plots the fit with the fit
statistics, and saves \code{<basename>.png}.

% ---------------------------------------------------------------------
\tssec{Package \texttt{experiment}}

\tssub{\texttt{experiment.measurement}}

\fnx{class Measurement(\allowbreak{}inputPar)} \fixed{C14, P9, P13}
Collects everything about a measurement from the dictionary returned by
\code{getRunningParams}, and reads the five data files (sample, environment,
container, normaliser, absorber) with \code{read_3col}. Main attributes:
experiment and beam information; \code{vslits}, \code{hslits},
\code{beamHeight}, \code{beamWidth}; \code{aScale}, \code{qScale},
\code{rScale}; \code{container}, \code{normaliser}, \code{absorber},
\code{environ} (material, shape, dimensions); sample \code{Shape},
\code{InnerDiam}, \code{OuterDiam}, \code{Height}, \code{Volume}
$=\pi h(r_o^2-r_i^2)/1000$ cm$^3$, \code{Mass}, \code{Density},
\code{EffDensity} $=$ mass/volume, \code{PackingFraction}; composition
\code{atoms}, \code{natoms}, \code{conc}, \code{nuclei}; averages
\code{molarMass}, \code{scoh}, \code{sinc}, \code{ssca}, \code{sabs}
(barns/atom), \code{bcoh}, \code{binc} (moduli of the averaged lengths,
P10); \code{free} (per-element free cross section); \code{b2lowQ}
$=\sigma_\mathrm{sca}/4\pi$, \code{b2highQ} $=\sigma_\mathrm{free}/4\pi$;
\code{sabswl} (absorption at the wavelength); \code{AtomicDensity},
\code{EffAtomicDensity}; data arrays \code{Data}, \code{envData},
\code{conData}, \code{norData}, \code{absData}. A warning is printed if
the container inner diameter differs from the sample outer diameter.
\code{showTable(outfile=None)} prints the identification, beam and grid
parameters (not the sample, container or composition), to a file if given.
\code{<Nmater>} is not read (the normaliser is always vanadium) and
\code{Fullness} is stored but not used.

\tssub{\texttt{experiment.info}}

Constants \code{VANADIUM_DENSITY_GCC = 6.11}, \code{VANADIUM_MOLAR_MASS = 50.9415},
\code{VANADIUM_ATOMIC_DENSITY} $=0.0722305$ atoms/\AA$^3$.

\fnx{setExpInfo(\allowbreak{}Proposal,\allowbreak\ mainProposer,\allowbreak\ experimenters,\allowbreak\ LC,\allowbreak\ otherUsers,\allowbreak\ startDate,\allowbreak\ endDate,\allowbreak\ environment,\allowbreak\ logBook,\allowbreak\ logPage,\allowbreak\ instrument)}
\fnx{setBeamInfo(\allowbreak{}zeroAngle=0,\allowbreak\ wavelength=0.5,\allowbreak\ LohenSlit=2.5,\allowbreak\ GamsSlit=-2.5,\allowbreak\ topFlag=25.0,\allowbreak\ bottomFlag=-25.0)}
\fnx{setCanInfo(\allowbreak{}material=\textquotesingle{}Vanadium\textquotesingle{},\allowbreak\ shape=\textquotesingle{}Cylinder\textquotesingle{},\allowbreak\ outerDiam=5,\allowbreak\ innerDiam=4.8,\allowbreak\ height=60.0)}
\fnx{setBinInfo(\allowbreak{}AngularResolution=0.125,\allowbreak\ AMin=0.0,\allowbreak\ AMax=140.0,\allowbreak\ AStep=0.125,\allowbreak\ QMin=0.0,\allowbreak\ QMax=23.5,\allowbreak\ QStep=0.02,\allowbreak\ RMin=0.0,\allowbreak\ RMax=20.0,\allowbreak\ RStep=0.01)}
\fnx{setNumorInfo(\allowbreak{}totalNumors,\allowbreak\ containerNumors,\allowbreak\ environmentNumors,\allowbreak\ nickelNumors,\allowbreak\ vanadiumNumors,\allowbreak\ absorberNumors,\allowbreak\ sampleNumors)}
Each builds and prints a dictionary describing one part of the
experiment; every argument has a default (e.g.\ \code{LC="Cuello"},
numor ranges \code{("Exp", 0, 1)}), see the docstrings. The dictionaries describe one part of the
experiment (beam width = left $-$ right slit, height = top $-$ bottom
flag; can wall thickness; number of points of each grid; numor ranges
given as (label, first, last)). \legacy

\fnx{setVanaInfo(\allowbreak{}IncXS=5.08,\allowbreak\ CohXS=0.02466,\allowbreak\ ScaXS=5.10466,\allowbreak\ AbsXS=5.08,\allowbreak\ CohSL=-0.443,\allowbreak\ molarM=50.9415,\allowbreak\ NAtoms=1.0,\allowbreak\ Diam=6.08,\allowbreak\ Height=50.0,\allowbreak\ density=VANADIUM\_DENSITY\_GCC)} \fixed{C16, P6}
Dictionary of the vanadium standard: cross sections, $b$, self scattering
at $Q=0$ ($\sigma_\mathrm{inc}/4\pi$), atomic density, mass number, free
incoherent cross section, volume in cm$^3$.

\fnx{setSampleInfo(\allowbreak{}comp\_sample,\allowbreak\ atoms\_sample,\allowbreak\ natoms\_sample,\allowbreak\ elements,\allowbreak\ c\_sample,\allowbreak\ wavelength=0.5,\allowbreak\ beamH=50.0,\allowbreak\ vanavol=1.0,\allowbreak\ vanadens=VANADIUM\_ATOMIC\_DENSITY,\allowbreak\ height=60.0,\allowbreak\ diameter=5.0,\allowbreak\ mass=1.0,\allowbreak\ density=1.0,\allowbreak\ title=\textquotesingle{}sample\textquotesingle{})} \fixed{C30, P9}
Dictionary of the sample: height in the beam, volumes, effective density,
packing and filling fractions, average molar mass, atomic and effective
atomic densities, averaged cross sections, absorption at $\lambda$,
average $b_\mathrm{coh}$, per-element free cross sections, self scattering
at $Q=0$, and the sample/vanadium volume and density ratios used for
normalisation. \code{elements} maps symbols to \code{Element} objects,
\code{c_sample} symbols to concentrations.

% ---------------------------------------------------------------------
\tssec{Package \texttt{isotopes}}

\fnx{class Element(\allowbreak{}iso\_name,\allowbreak\ table=0,\allowbreak\ *{}*kwargs)} \fixed{C15, P15}
Nuclear data of an element (\code{"Ni"}) or isotope (\code{"Li6"},
\code{"H2"} for deuterium) from \code{database.ISOTOPE_DICT} (values from
\emph{Neutron Scattering Fundamentals}, vol.~44, App.~2). Attributes:
\code{symbol}, \code{iso}, \code{name}, \code{A}, \code{Z}, \code{weight}
(u), \code{abundance}, \code{spin}, \code{re_bcoh}, \code{im_bcoh},
\code{re_binc}, \code{im_binc} (fm), \code{sig_coh}, \code{sig_inc},
\code{sig_sca}, \code{sig_abs} (b), complex \code{bcoh}, \code{binc},
\code{scoh_bound} $=4\pi|b_\mathrm{coh}|^2/100$, \code{sinc_bound} (b),
\code{neut} $=A-Z$. Keyword arguments override table values (with
plausibility warnings). \code{table=1} prints the table of the isotope,
\code{table=2} of all isotopes of the element. Unknown names raise
\code{ValueError}. \code{Element.getIsotope(name)} lists the isotopes of an
element name, or all names with \code{"All"}. \code{Elemento},
\code{elemento}: deprecated aliases.

\fnx{getFree(\allowbreak{}isot,\allowbreak\ key)}
Free value $[A/(A+1)]^2\times$\code{key} (complex) for an \code{Element}.

\noindent Validators \code{Positive}, \code{Negative}, \code{SemiWhole},
\code{Top}, \code{Int}, \code{String}, \code{Complex}: print a warning when
an override is implausible. \code{database.NOMBRES} and \code{UNITS}: names
and units of the table columns.

% ---------------------------------------------------------------------
\tssec{Packages \texttt{data}, \texttt{plotting}, \texttt{utils}, \texttt{chemistry}}

\fnx{class DataXYE(\allowbreak{}filename)} \fixed{C22}
Reads a 2- or 3-column file (header lines start with \texttt{\#} or \code{!};
a missing error is stored as $-1$); attributes \code{x}, \code{y},
\code{e}, \code{head}, averages, minima and maxima, \code{peaks_x},
\code{peaks_y} and the minimum \code{xminr}, \code{yminr} in
$x\in[5,40]$ (empty / \code{NaN} if no point there). Methods
\code{plot(file_format=0, xmin, xmax, ymin, ymax)} (axis label from the
extension; saves to \code{../regdata/<basename>.<format>} if a format is
given), \code{show()}, \code{header(lines=1000)}. \code{DataXYE_old} is an
alias.

\fnx{class DataXYE0(\allowbreak{}filename,\allowbreak\ symbol,\allowbreak\ save,\allowbreak\ scale,\allowbreak\ superplot,\allowbreak\ ...)}, \code{class DataXYE2(file)}
Simpler containers read with \code{np.loadtxt}, with \code{plot} methods
(scale \code{"a"}, \code{"q"} or \code{"r"} for the axis label;
\code{superplot} to overlay); \code{DataXYE0.save(filename)}. \legacy

\fnx{plotQ(\allowbreak{}plt,\allowbreak\ ...)}, \code{plotA(plt, ...)}, \code{plotR(plt, ...)}
Set title, labels, limits, grid and legend of the current matplotlib plot
for $Q$, angle or $r$ scales.

\fnx{defColors(\allowbreak{})}, \code{makePlotsA(runInfo, diffA, head, numor)}, \code{makePlotsQ(runInfo, diffQbin, head, numor)}
Colour list and the diagnostic plots of d4creg (individual numors,
individual banks, total in angle; total in $Q$).

\fnx{getDate(\allowbreak{})}
Current date as \code{" Dayname DD/MM/YYYY HH:MM:SS"}.

\fnx{extractDict(\allowbreak{}dictionary,\allowbreak\ key)}, \code{extractAttr(Dict, attr)}
Sub-dictionary of a dictionary of dictionaries; dictionary of one
attribute of a dictionary of objects (e.g.\ \code{extractAttr(elements, "sig_sca")}).

\fnx{execute\_command\_and\_write\_to\_file(\allowbreak{}command,\allowbreak\ output\_file)}
Runs a command, writing its output to a file; errors are printed, not raised.

\fnx{stop(\allowbreak{}message)}
Prints a message and asks whether to continue (interactive sessions only;
raises \code{RuntimeError} otherwise).

\noindent \code{chemistry.composition} re-exports \code{AtomicAvg},
\code{getConcentrations}, \code{getNofAtoms} (from \code{physics.properties})
and \code{extractAttr}, \code{extractDict} (from \code{utils.dicts}) for
compatibility.

% ---------------------------------------------------------------------
\tssec{Top level: \texttt{\_\_init\_\_} and \texttt{\_aliases}}

\code{toscana/__init__.py} imports the subpackages \code{data},
\code{dataio}, \code{experiment}, \code{isotopes}, \code{models},
\code{numerics}, \code{physics}, \code{plotting}, \code{utils} (not
\code{calibration}, \code{chemistry}, \code{workflows}, which must be
imported explicitly) and defines \code{__version__ = "0.1.0"}.
\code{_aliases.py} installs, at the front of Python's import machinery, a
finder that maps \code{toscana.math...} to \code{toscana.numerics...} and
\code{toscana.io...} to \code{toscana.dataio...}, returning the existing
module objects and emitting a \code{DeprecationWarning}; attribute access
\code{toscana.math} is handled the same way. Remove both once all
notebooks use the new names.

% When \input into the manual, stop here; stand-alone, close the document.
\ifdefined\ToScaNAManual\expandafter\endinput\fi
\end{document}


\end{document}
